\documentclass[10pt]{article}
% Byte-stable output: no creation date, no trailer id, no TeX banner in the PDF.
\ifdefined\pdfinfoomitdate\pdfinfoomitdate=1\fi
\ifdefined\pdftrailerid\pdftrailerid{}\fi
\ifdefined\pdfsuppressptexinfo\pdfsuppressptexinfo=-1\fi
\usepackage[margin=0.78in,headheight=14pt]{geometry}
\usepackage[T1]{fontenc}
\usepackage[utf8]{inputenc}
\usepackage{lmodern}
\usepackage{xcolor}
\usepackage{amsmath,amssymb}
\usepackage{booktabs,longtable,array}
\usepackage{enumitem}
\usepackage{fancyvrb,fvextra}
\usepackage{needspace}
\usepackage{newunicodechar}
\usepackage[hidelinks]{hyperref}
\usepackage{bookmark}
\usepackage{fancyhdr}
\definecolor{navy}{HTML}{17324D}
\hypersetup{colorlinks=true,linkcolor=navy,urlcolor=navy,pdftitle={The benchmark inventory},pdfauthor={Benchmark/report/build\_report.py}}
\pagestyle{fancy}\fancyhf{}\lhead{benchmark inventory}\rhead{\thepage}
\setlength{\parindent}{0pt}\setlength{\parskip}{5pt}\emergencystretch=3em\Urlmuskip=0mu plus 1mu\sloppy
\setcounter{tocdepth}{1}
\newunicodechar{π}{\ensuremath{\pi}}
\newunicodechar{σ}{\ensuremath{\sigma}}
\newunicodechar{∪}{\ensuremath{\cup}}
\newunicodechar{∩}{\ensuremath{\cap}}
\newunicodechar{⊆}{\ensuremath{\subseteq}}
\newunicodechar{⊇}{\ensuremath{\supseteq}}
\newunicodechar{⊂}{\ensuremath{\subset}}
\newunicodechar{⊊}{\ensuremath{\subsetneq}}
\newunicodechar{∅}{\ensuremath{\emptyset}}
\newunicodechar{≠}{\ensuremath{\neq}}
\newunicodechar{∧}{\ensuremath{\wedge}}
\newunicodechar{∨}{\ensuremath{\vee}}
\newunicodechar{∖}{\ensuremath{\setminus}}
\newunicodechar{¬}{\ensuremath{\neg}}
\newunicodechar{↔}{\ensuremath{\leftrightarrow}}
\newunicodechar{≡}{\ensuremath{\equiv}}
\newunicodechar{∈}{\ensuremath{\in}}
\newunicodechar{∉}{\ensuremath{\notin}}
\newunicodechar{∀}{\ensuremath{\forall}}
\newunicodechar{∃}{\ensuremath{\exists}}
\newunicodechar{≤}{\ensuremath{\leq}}
\newunicodechar{≥}{\ensuremath{\geq}}
\newunicodechar{∘}{\ensuremath{\circ}}
\newunicodechar{×}{\ensuremath{\times}}
\newunicodechar{↑}{\ensuremath{\uparrow}}
\newunicodechar{→}{\ensuremath{\rightarrow}}
\newunicodechar{←}{\ensuremath{\leftarrow}}
\newunicodechar{⇒}{\ensuremath{\Rightarrow}}
\newunicodechar{⇔}{\ensuremath{\Leftrightarrow}}
\newunicodechar{↦}{\ensuremath{\mapsto}}
\newunicodechar{⋈}{\ensuremath{\bowtie}}
\newunicodechar{∞}{\ensuremath{\infty}}
\newunicodechar{≈}{\ensuremath{\approx}}
\newunicodechar{≅}{\ensuremath{\cong}}
\newunicodechar{⊢}{\ensuremath{\vdash}}
\newunicodechar{⊨}{\ensuremath{\models}}
\newunicodechar{⊤}{\ensuremath{\top}}
\newunicodechar{⊥}{\ensuremath{\bot}}
\newunicodechar{Δ}{\ensuremath{\Delta}}
\newunicodechar{Σ}{\ensuremath{\Sigma}}
\newunicodechar{Π}{\ensuremath{\Pi}}
\newunicodechar{Ω}{\ensuremath{\Omega}}
\newunicodechar{Γ}{\ensuremath{\Gamma}}
\newunicodechar{Θ}{\ensuremath{\Theta}}
\newunicodechar{Λ}{\ensuremath{\Lambda}}
\newunicodechar{Φ}{\ensuremath{\Phi}}
\newunicodechar{Ψ}{\ensuremath{\Psi}}
\newunicodechar{α}{\ensuremath{\alpha}}
\newunicodechar{β}{\ensuremath{\beta}}
\newunicodechar{γ}{\ensuremath{\gamma}}
\newunicodechar{δ}{\ensuremath{\delta}}
\newunicodechar{ε}{\ensuremath{\varepsilon}}
\newunicodechar{ζ}{\ensuremath{\zeta}}
\newunicodechar{η}{\ensuremath{\eta}}
\newunicodechar{θ}{\ensuremath{\theta}}
\newunicodechar{ι}{\ensuremath{\iota}}
\newunicodechar{κ}{\ensuremath{\kappa}}
\newunicodechar{λ}{\ensuremath{\lambda}}
\newunicodechar{μ}{\ensuremath{\mu}}
\newunicodechar{ν}{\ensuremath{\nu}}
\newunicodechar{ξ}{\ensuremath{\xi}}
\newunicodechar{ρ}{\ensuremath{\rho}}
\newunicodechar{τ}{\ensuremath{\tau}}
\newunicodechar{υ}{\ensuremath{\upsilon}}
\newunicodechar{φ}{\ensuremath{\varphi}}
\newunicodechar{χ}{\ensuremath{\chi}}
\newunicodechar{ψ}{\ensuremath{\psi}}
\newunicodechar{ω}{\ensuremath{\omega}}
\newunicodechar{′}{\ensuremath{\prime}}
\newunicodechar{″}{\ensuremath{\prime\prime}}
\newunicodechar{−}{\ensuremath{-}}
\newunicodechar{±}{\ensuremath{\pm}}
\newunicodechar{∗}{\ensuremath{\ast}}
\newunicodechar{∙}{\ensuremath{\cdot}}
\newunicodechar{·}{\ensuremath{\cdot}}
\newunicodechar{•}{\ensuremath{\bullet}}
\newunicodechar{…}{\ensuremath{\ldots}}
\newunicodechar{⟨}{\ensuremath{\langle}}
\newunicodechar{⟩}{\ensuremath{\rangle}}
\newunicodechar{√}{\ensuremath{\surd}}
\newunicodechar{∇}{\ensuremath{\nabla}}
\newunicodechar{⊕}{\ensuremath{\oplus}}
\newunicodechar{⊗}{\ensuremath{\otimes}}
\newunicodechar{⊑}{\ensuremath{\sqsubseteq}}
\newunicodechar{⊓}{\ensuremath{\sqcap}}
\newunicodechar{⊔}{\ensuremath{\sqcup}}
\newunicodechar{∑}{\ensuremath{\sum}}
\newunicodechar{∏}{\ensuremath{\prod}}
\newunicodechar{≔}{\ensuremath{:=}}
\newunicodechar{≃}{\ensuremath{\simeq}}
\newunicodechar{∝}{\ensuremath{\propto}}
\newunicodechar{∂}{\ensuremath{\partial}}
\newunicodechar{∄}{\ensuremath{\nexists}}
\newunicodechar{⊄}{\ensuremath{\not\subseteq}}
\newunicodechar{⊈}{\ensuremath{\nsubseteq}}
\newunicodechar{≺}{\ensuremath{\prec}}
\newunicodechar{⁻}{\ensuremath{^{-}}}
\newunicodechar{⁺}{\ensuremath{^{+}}}
\newunicodechar{⁰}{\ensuremath{^{0}}}
\newunicodechar{ⁿ}{\ensuremath{^{n}}}
\newunicodechar{ᵏ}{\ensuremath{^{k}}}
\newunicodechar{≻}{\ensuremath{\succ}}
\newunicodechar{⌈}{\ensuremath{\lceil}}
\newunicodechar{⌉}{\ensuremath{\rceil}}
\newunicodechar{⌊}{\ensuremath{\lfloor}}
\newunicodechar{⌋}{\ensuremath{\rfloor}}
\title{The benchmark inventory\\\large 86 cases: titles, categories, kinds, status, sources, tags, Datalog and Whiel encodings}
\author{Generated by Benchmark/report/build\_report.py}
\date{}

\begin{document}
\maketitle

\tableofcontents

\clearpage

\section{Overview}\label{sec:overview}

The folder \path{Benchmark} holds 86 cases. Each case is one section below with its title, category, kind, status, sources, tags, the Datalog program(s) it encodes and the Whiel Hoare triple as written in \texttt{Input.lean}. The tag vocabulary is in Section~\ref{sec:tags}.

\paragraph{Status.} The status of a case is read off its folder and off nothing else: a case is certified only when the certificate emitter wrote the certificate that stands beside its input, and solved when the frozen record a search writes is there but the certificate built from it is not. A solved case has an answer and a record of it, not a kernel-checked proof. What the metadata expects of a case is reported next to the status, never in place of it, so an expectation can never be mistaken for a checked result.
\begin{tabular}{@{}lrp{0.52\textwidth}@{}}
\toprule
Status & Cases & Read off \\
\midrule
certified valid & 69 & Certificate/Valid.lean, written by the certificate emitter \\
certified invalid & 17 & Certificate/Invalid.lean, written by the certificate emitter \\
solved valid, not yet certified & 0 & Core.json, the frozen invariant record; no certificate built from it yet \\
solved invalid, not yet certified & 0 & Counterexample.json, the frozen counterexample record; no certificate built from it yet \\
open & 0 & no certificate and no frozen record in the case folder \\
\bottomrule
\end{tabular}
\par of which outside the library build: 10 (counted as certified above; the kernel check runs on its own, under a raised memory limit, and each such case says so in its own section).
\paragraph{Status against expectation.} certified invalid / expected invalid: 17; certified valid / expected valid: 68; certified valid / expected valid (expected; strengthened claim): 1.

\paragraph{Categories.} llm: 11; production: 7; production-pair: 22; souffle-benchmark: 4; specification: 11; synthetic: 21; verification-task: 10.

\paragraph{Kinds.} conditional claim: 1; containment: 18; correctness: 2; equivalence: 64; implication: 1.

\begin{longtable}{@{}lp{0.34\textwidth}p{0.12\textwidth}p{0.10\textwidth}p{0.11\textwidth}p{0.20\textwidth}@{}}
\toprule
Case & Title & Category & Kind & Status & Tags \\
\midrule
\endhead
\hyperref[case:Example0001]{Example0001} & Transitive closure by naive iteration under the weak transitivity specification & synthetic & correctness & certified valid & {\scriptsize \hyperref[tag:prophecy]{\texttt{prophecy}}} \\
\hyperref[case:Example0013]{Example0013} & Transitive closure by naive iteration under a false specification (T ⊆ E) & synthetic & correctness & certified invalid & {\scriptsize } \\
\hyperref[case:Example0106]{Example0106} & Prefix-equivalence: naive tc\_left computes lfp of tc\_nonlinear & synthetic & equivalence & certified valid & {\scriptsize \hyperref[tag:fixpoint-witness]{\texttt{fixpoint-witness}} \hyperref[tag:linear-vs-nonlinear]{\texttt{linear-vs-nonlinear}} \hyperref[tag:nonlinear-recursion]{\texttt{nonlinear-recursion}} \hyperref[tag:prophecy]{\texttt{prophecy}}} \\
\hyperref[case:Example0107]{Example0107} & Prefix-equivalence: naive tc\_nonlinear computes lfp of tc\_left & synthetic & equivalence & certified valid & {\scriptsize \hyperref[tag:fixpoint-witness]{\texttt{fixpoint-witness}} \hyperref[tag:linear-vs-nonlinear]{\texttt{linear-vs-nonlinear}} \hyperref[tag:nonlinear-recursion]{\texttt{nonlinear-recursion}}} \\
\hyperref[case:Example0108]{Example0108} & Prefix-equivalence: naive tc\_right computes lfp of tc\_left & synthetic & equivalence & certified valid & {\scriptsize \hyperref[tag:fixpoint-witness]{\texttt{fixpoint-witness}} \hyperref[tag:linear-vs-nonlinear]{\texttt{linear-vs-nonlinear}}} \\
\hyperref[case:Example0109]{Example0109} & Containment: red-blue alternating paths ⊆ TC(R ∪ B) & synthetic & containment & certified valid & {\scriptsize \hyperref[tag:fixpoint-witness]{\texttt{fixpoint-witness}} \hyperref[tag:mutual-recursion]{\texttt{mutual-recursion}}} \\
\hyperref[case:Example0115]{Example0115} & Prefix-equivalence: naive all\_paths\_via\_even\_odd (linear) computes lfp of tc\_left & synthetic & equivalence & certified valid & {\scriptsize \hyperref[tag:fixpoint-witness]{\texttt{fixpoint-witness}} \hyperref[tag:mutual-recursion]{\texttt{mutual-recursion}}} \\
\hyperref[case:Example0117]{Example0117} & Prefix-equivalence: naive all\_paths\_via\_even\_odd (non-linear) computes lfp of tc\_left & synthetic & equivalence & certified valid & {\scriptsize \hyperref[tag:fixpoint-witness]{\texttt{fixpoint-witness}} \hyperref[tag:linear-vs-nonlinear]{\texttt{linear-vs-nonlinear}} \hyperref[tag:mutual-recursion]{\texttt{mutual-recursion}} \hyperref[tag:nonlinear-recursion]{\texttt{nonlinear-recursion}}} \\
\hyperref[case:Example0119]{Example0119} & Prefix-equivalence: naive all\_paths\_via\_even\_odd (linear) computes lfp of tc\_nonlinear & synthetic & equivalence & certified valid & {\scriptsize \hyperref[tag:fixpoint-witness]{\texttt{fixpoint-witness}} \hyperref[tag:linear-vs-nonlinear]{\texttt{linear-vs-nonlinear}} \hyperref[tag:mutual-recursion]{\texttt{mutual-recursion}} \hyperref[tag:nonlinear-recursion]{\texttt{nonlinear-recursion}} \hyperref[tag:prophecy]{\texttt{prophecy}}} \\
\hyperref[case:Example0123]{Example0123} & Prefix-equivalence: naive even\_odd\_nonlinear computes lfp of even\_odd\_linear & synthetic & equivalence & certified valid & {\scriptsize \hyperref[tag:fixpoint-witness]{\texttt{fixpoint-witness}} \hyperref[tag:linear-vs-nonlinear]{\texttt{linear-vs-nonlinear}} \hyperref[tag:mutual-recursion]{\texttt{mutual-recursion}} \hyperref[tag:nonlinear-recursion]{\texttt{nonlinear-recursion}}} \\
\hyperref[case:Example0125]{Example0125} & Prefix-equivalence: naive even\_odd\_linear computes lfp of even\_odd\_nonlinear & synthetic & equivalence & certified valid & {\scriptsize \hyperref[tag:fixpoint-witness]{\texttt{fixpoint-witness}} \hyperref[tag:linear-vs-nonlinear]{\texttt{linear-vs-nonlinear}} \hyperref[tag:mutual-recursion]{\texttt{mutual-recursion}} \hyperref[tag:nonlinear-recursion]{\texttt{nonlinear-recursion}} \hyperref[tag:prophecy]{\texttt{prophecy}}} \\
\hyperref[case:Example0128]{Example0128} & Containment: naive red\_blue\_nonlinear outputs ⊆ pre-fixpoint of TC(R ∪ B) & synthetic & containment & certified valid & {\scriptsize \hyperref[tag:fixpoint-witness]{\texttt{fixpoint-witness}} \hyperref[tag:mutual-recursion]{\texttt{mutual-recursion}} \hyperref[tag:nonlinear-recursion]{\texttt{nonlinear-recursion}}} \\
\hyperref[case:Example0130]{Example0130} & Containment: even/odd/ambiguous-depth path-system program ⊆ every set closed under the basic path-system rules & synthetic & containment & certified valid & {\scriptsize \hyperref[tag:fixpoint-witness]{\texttt{fixpoint-witness}} \hyperref[tag:mutual-recursion]{\texttt{mutual-recursion}} \hyperref[tag:nonlinear-recursion]{\texttt{nonlinear-recursion}}} \\
\hyperref[case:Example0132]{Example0132} & Odd and even paths: linear versus non-linear program (two compiled loops) & synthetic & equivalence & certified valid & {\scriptsize \hyperref[tag:linear-vs-nonlinear]{\texttt{linear-vs-nonlinear}} \hyperref[tag:mutual-recursion]{\texttt{mutual-recursion}} \hyperref[tag:nonlinear-recursion]{\texttt{nonlinear-recursion}} \hyperref[tag:prophecy]{\texttt{prophecy}}} \\
\hyperref[case:Example0133]{Example0133} & All paths via odd/even paths: linear versus non-linear program (two compiled loops) & synthetic & equivalence & certified valid & {\scriptsize \hyperref[tag:linear-vs-nonlinear]{\texttt{linear-vs-nonlinear}} \hyperref[tag:mutual-recursion]{\texttt{mutual-recursion}} \hyperref[tag:nonlinear-recursion]{\texttt{nonlinear-recursion}} \hyperref[tag:prophecy]{\texttt{prophecy}}} \\
\hyperref[case:Example0134]{Example0134} & Red/blue alternating paths: linear versus non-linear program (two compiled loops) & synthetic & equivalence & certified valid & {\scriptsize \hyperref[tag:linear-vs-nonlinear]{\texttt{linear-vs-nonlinear}} \hyperref[tag:mutual-recursion]{\texttt{mutual-recursion}} \hyperref[tag:nonlinear-recursion]{\texttt{nonlinear-recursion}} \hyperref[tag:prophecy]{\texttt{prophecy}}} \\
\hyperref[case:Example0136]{Example0136} & Right-linear transitive closure versus the non-linear odd/even path program (two compiled loops) & synthetic & equivalence & certified valid & {\scriptsize \hyperref[tag:linear-vs-nonlinear]{\texttt{linear-vs-nonlinear}} \hyperref[tag:mutual-recursion]{\texttt{mutual-recursion}} \hyperref[tag:nonlinear-recursion]{\texttt{nonlinear-recursion}} \hyperref[tag:prophecy]{\texttt{prophecy}}} \\
\hyperref[case:Example0137]{Example0137} & Non-linear transitive closure versus the linear odd/even path program (two compiled loops) & synthetic & equivalence & certified valid & {\scriptsize \hyperref[tag:linear-vs-nonlinear]{\texttt{linear-vs-nonlinear}} \hyperref[tag:mutual-recursion]{\texttt{mutual-recursion}} \hyperref[tag:nonlinear-recursion]{\texttt{nonlinear-recursion}} \hyperref[tag:prophecy]{\texttt{prophecy}}} \\
\hyperref[case:Example0138]{Example0138} & Joint-loop equivalence: tc\_nonlinear vs all\_paths\_via\_even\_odd (non-linear, naive) & synthetic & equivalence & certified valid & {\scriptsize \hyperref[tag:mutual-recursion]{\texttt{mutual-recursion}} \hyperref[tag:nonlinear-recursion]{\texttt{nonlinear-recursion}} \hyperref[tag:prophecy]{\texttt{prophecy}}} \\
\hyperref[case:Example0162]{Example0162} & Two-player game reachability mixing existential and universal successor steps & synthetic & equivalence & certified valid & {\scriptsize \hyperref[tag:fixpoint-witness]{\texttt{fixpoint-witness}} \hyperref[tag:negation]{\texttt{negation}}} \\
\hyperref[case:Example0163]{Example0163} & Greatest-fixpoint safe set: shrinking iteration bounded below by Rb & synthetic & equivalence & certified valid & {\scriptsize \hyperref[tag:fixpoint-witness]{\texttt{fixpoint-witness}} \hyperref[tag:greatest-fixpoint]{\texttt{greatest-fixpoint}} \hyperref[tag:negation]{\texttt{negation}}} \\
\hyperref[case:Example1030]{Example1030} & Foreign-key support core for sharded materialized rows & llm & equivalence & certified valid & {\scriptsize \hyperref[tag:fixpoint-witness]{\texttt{fixpoint-witness}} \hyperref[tag:greatest-fixpoint]{\texttt{greatest-fixpoint}} \hyperref[tag:integrity-constraints]{\texttt{integrity-constraints}} \hyperref[tag:negation]{\texttt{negation}}} \\
\hyperref[case:Example1061]{Example1061} & Equivalence of computed same-generation with lfp of the same-generation rules (pre-fixpoint) & llm & equivalence & certified valid & {\scriptsize \hyperref[tag:fixpoint-witness]{\texttt{fixpoint-witness}}} \\
\hyperref[case:Example1081]{Example1081} & Audited delegation contained in broad authorization closure & llm & containment & certified valid & {\scriptsize \hyperref[tag:dependencies]{\texttt{dependencies}} \hyperref[tag:mutual-recursion]{\texttt{mutual-recursion}} \hyperref[tag:nonlinear-recursion]{\texttt{nonlinear-recursion}} \hyperref[tag:prophecy]{\texttt{prophecy}}} \\
\hyperref[case:Example1088]{Example1088} & Equivalent forward dependency closure and transposed reverse resolution & llm & equivalence & certified valid & {\scriptsize \hyperref[tag:dependencies]{\texttt{dependencies}} \hyperref[tag:mutual-recursion]{\texttt{mutual-recursion}} \hyperref[tag:nonlinear-recursion]{\texttt{nonlinear-recursion}}} \\
\hyperref[case:Example1090]{Example1090} & Equivalent taint reachability with ternary modes and projected frontiers & llm & equivalence & certified valid & {\scriptsize \hyperref[tag:dependencies]{\texttt{dependencies}} \hyperref[tag:mutual-recursion]{\texttt{mutual-recursion}} \hyperref[tag:nonlinear-recursion]{\texttt{nonlinear-recursion}}} \\
\hyperref[case:Example1091]{Example1091} & Delegation reachability bounded by a transitively closed policy view & llm & containment & certified valid & {\scriptsize \hyperref[tag:dependencies]{\texttt{dependencies}} \hyperref[tag:mutual-recursion]{\texttt{mutual-recursion}} \hyperref[tag:nonlinear-recursion]{\texttt{nonlinear-recursion}}} \\
\hyperref[case:Example1093]{Example1093} & Lineage triples contained in a closed ternary cover & llm & containment & certified valid & {\scriptsize \hyperref[tag:dependencies]{\texttt{dependencies}} \hyperref[tag:mutual-recursion]{\texttt{mutual-recursion}} \hyperref[tag:nonlinear-recursion]{\texttt{nonlinear-recursion}}} \\
\hyperref[case:Example1094]{Example1094} & Routing table contained in an administratively closed distance view & llm & containment & certified valid & {\scriptsize \hyperref[tag:dependencies]{\texttt{dependencies}} \hyperref[tag:mutual-recursion]{\texttt{mutual-recursion}} \hyperref[tag:nonlinear-recursion]{\texttt{nonlinear-recursion}}} \\
\hyperref[case:Example2004]{Example2004} & Magic-set transformation: demanded answers equal the original program's answers (Soufflé) & production & equivalence & certified valid & {\scriptsize \hyperref[tag:database-engine]{\texttt{database-engine}} \hyperref[tag:magic-sets]{\texttt{magic-sets}} \hyperref[tag:prophecy]{\texttt{prophecy}}} \\
\hyperref[case:Example2005]{Example2005} & Restricted-edge reachability is contained in full-edge reachability (edge-monotonicity) & production & containment & certified valid & {\scriptsize \hyperref[tag:fixpoint-witness]{\texttt{fixpoint-witness}} \hyperref[tag:negation]{\texttt{negation}}} \\
\hyperref[case:Example2006]{Example2006} & Policy-filtered connectivity is contained in raw topological connectivity & production & containment & certified valid & {\scriptsize \hyperref[tag:fixpoint-witness]{\texttt{fixpoint-witness}} \hyperref[tag:negation]{\texttt{negation}}} \\
\hyperref[case:Example2007]{Example2007} & E[f U g] is contained in EF g & production & containment & certified valid & {\scriptsize \hyperref[tag:fixpoint-witness]{\texttt{fixpoint-witness}} \hyperref[tag:semi-naive]{\texttt{semi-naive}}} \\
\hyperref[case:Example2010]{Example2010} & Semi-naive evaluation ≡ naive/declarative transitive closure & production & equivalence & certified valid & {\scriptsize \hyperref[tag:database-engine]{\texttt{database-engine}} \hyperref[tag:evaluation-strategies]{\texttt{evaluation-strategies}} \hyperref[tag:fixpoint-witness]{\texttt{fixpoint-witness}} \hyperref[tag:semi-naive]{\texttt{semi-naive}}} \\
\hyperref[case:Example2011]{Example2011} & Incremental view maintenance ≡ from-scratch recomputation & production & equivalence & certified valid & {\scriptsize \hyperref[tag:evaluation-strategies]{\texttt{evaluation-strategies}} \hyperref[tag:fixpoint-witness]{\texttt{fixpoint-witness}} \hyperref[tag:view-maintenance]{\texttt{view-maintenance}}} \\
\hyperref[case:Example2012]{Example2012} & Frontier (grey-set) marking ≡ naive mark propagation & production & equivalence & certified valid & {\scriptsize \hyperref[tag:evaluation-strategies]{\texttt{evaluation-strategies}} \hyperref[tag:fixpoint-witness]{\texttt{fixpoint-witness}} \hyperref[tag:semi-naive]{\texttt{semi-naive}}} \\
\hyperref[case:Example3001]{Example3001} & Audited delegation alternation loses an uncertified reviewed hop & llm & containment & certified invalid & {\scriptsize \hyperref[tag:dependencies]{\texttt{dependencies}} \hyperref[tag:mutual-recursion]{\texttt{mutual-recursion}} \hyperref[tag:nonlinear-recursion]{\texttt{nonlinear-recursion}}} \\
\hyperref[case:Example3009]{Example3009} & Certified delegation closure exceeds audit-bound authorization & llm & containment & certified invalid & {\scriptsize \hyperref[tag:dependencies]{\texttt{dependencies}} \hyperref[tag:mutual-recursion]{\texttt{mutual-recursion}} \hyperref[tag:nonlinear-recursion]{\texttt{nonlinear-recursion}}} \\
\hyperref[case:Example3014]{Example3014} & Auditor trust diverges from resource trust under recursive elevation & llm & containment & certified invalid & {\scriptsize \hyperref[tag:dependencies]{\texttt{dependencies}} \hyperref[tag:mutual-recursion]{\texttt{mutual-recursion}} \hyperref[tag:nonlinear-recursion]{\texttt{nonlinear-recursion}}} \\
\hyperref[case:Example4001]{Example4001} & Directed triangle query versus its path-view reformulation & verification-task & equivalence & certified valid & {\scriptsize \hyperref[tag:views]{\texttt{views}}} \\
\hyperref[case:Example4002]{Example4002} & Transitive closure versus its length-two path-view reformulation & verification-task & equivalence & certified invalid & {\scriptsize \hyperref[tag:views]{\texttt{views}}} \\
\hyperref[case:Example4003]{Example4003} & Containment of the path-view reformulation in transitive closure & verification-task & containment & certified valid & {\scriptsize \hyperref[tag:prophecy]{\texttt{prophecy}} \hyperref[tag:views]{\texttt{views}}} \\
\hyperref[case:Example4004]{Example4004} & Path-view reformulation versus the even-length path query & verification-task & equivalence & certified valid & {\scriptsize \hyperref[tag:mutual-recursion]{\texttt{mutual-recursion}} \hyperref[tag:prophecy]{\texttt{prophecy}} \hyperref[tag:views]{\texttt{views}}} \\
\hyperref[case:Example4034]{Example4034} & Containment of a reformulation over sound odd/even views in transitive closure & verification-task & conditional claim & certified valid & {\scriptsize \hyperref[tag:mutual-recursion]{\texttt{mutual-recursion}} \hyperref[tag:prophecy]{\texttt{prophecy}} \hyperref[tag:views]{\texttt{views}}} \\
\hyperref[case:Example4035]{Example4035} & Good-interior path query versus its begin-good and end-good edge-view reformulation & verification-task & equivalence & certified valid & {\scriptsize \hyperref[tag:prophecy]{\texttt{prophecy}} \hyperref[tag:views]{\texttt{views}}} \\
\hyperref[case:Example4036]{Example4036} & Airline reachability query versus its three-view reformulation, forward containment & verification-task & containment & certified valid & {\scriptsize \hyperref[tag:prophecy]{\texttt{prophecy}} \hyperref[tag:views]{\texttt{views}}} \\
\hyperref[case:Example4037]{Example4037} & Airline reachability query versus its three-view reformulation, reverse containment & verification-task & containment & certified invalid & {\scriptsize \hyperref[tag:views]{\texttt{views}}} \\
\hyperref[case:Example4040]{Example4040} & Closure of red and blue edges versus its reformulation over Datalog path views & verification-task & equivalence & certified valid & {\scriptsize \hyperref[tag:linear-vs-nonlinear]{\texttt{linear-vs-nonlinear}} \hyperref[tag:nonlinear-recursion]{\texttt{nonlinear-recursion}} \hyperref[tag:prophecy]{\texttt{prophecy}} \hyperref[tag:views]{\texttt{views}}} \\
\hyperref[case:Example4041]{Example4041} & Transitive closure versus its reformulation over exact odd/even views & verification-task & equivalence & certified valid & {\scriptsize \hyperref[tag:mutual-recursion]{\texttt{mutual-recursion}} \hyperref[tag:prophecy]{\texttt{prophecy}} \hyperref[tag:views]{\texttt{views}}} \\
\hyperref[case:Example5001]{Example5001} & PostgreSQL recursive UNION versus naive seeded closure & production-pair & equivalence & certified valid & {\scriptsize \hyperref[tag:database-engine]{\texttt{database-engine}} \hyperref[tag:evaluation-strategies]{\texttt{evaluation-strategies}} \hyperref[tag:prophecy]{\texttt{prophecy}} \hyperref[tag:semi-naive]{\texttt{semi-naive}}} \\
\hyperref[case:Example5002]{Example5002} & Datafrog semi-naive delta evaluation versus naive evaluation & production-pair & equivalence & certified valid & {\scriptsize \hyperref[tag:database-engine]{\texttt{database-engine}} \hyperref[tag:evaluation-strategies]{\texttt{evaluation-strategies}} \hyperref[tag:prophecy]{\texttt{prophecy}} \hyperref[tag:semi-naive]{\texttt{semi-naive}}} \\
\hyperref[case:Example5003]{Example5003} & Worklist breadth-first reachability versus naive evaluation & production-pair & equivalence & certified valid & {\scriptsize \hyperref[tag:evaluation-strategies]{\texttt{evaluation-strategies}} \hyperref[tag:prophecy]{\texttt{prophecy}} \hyperref[tag:semi-naive]{\texttt{semi-naive}}} \\
\hyperref[case:Example5004]{Example5004} & Materialized role-hierarchy closure with a join versus a recursive has-role query & production-pair & equivalence & certified valid & {\scriptsize \hyperref[tag:access-control]{\texttt{access-control}} \hyperref[tag:evaluation-strategies]{\texttt{evaluation-strategies}} \hyperref[tag:prophecy]{\texttt{prophecy}}} \\
\hyperref[case:Example5005]{Example5005} & Incremental maintenance of a reachability view versus recomputation from scratch & production-pair & equivalence & certified valid & {\scriptsize \hyperref[tag:evaluation-strategies]{\texttt{evaluation-strategies}} \hyperref[tag:prophecy]{\texttt{prophecy}} \hyperref[tag:view-maintenance]{\texttt{view-maintenance}}} \\
\hyperref[case:Example5006]{Example5006} & Semi-naive closure with an appended delta versus a prepended delta & production-pair & equivalence & certified valid & {\scriptsize \hyperref[tag:evaluation-strategies]{\texttt{evaluation-strategies}} \hyperref[tag:linear-vs-nonlinear]{\texttt{linear-vs-nonlinear}} \hyperref[tag:prophecy]{\texttt{prophecy}} \hyperref[tag:semi-naive]{\texttt{semi-naive}}} \\
\hyperref[case:Example5007]{Example5007} & Transitive closure by repeated squaring versus linear closure & production-pair & equivalence & certified valid & {\scriptsize \hyperref[tag:evaluation-strategies]{\texttt{evaluation-strategies}} \hyperref[tag:linear-vs-nonlinear]{\texttt{linear-vs-nonlinear}} \hyperref[tag:nonlinear-recursion]{\texttt{nonlinear-recursion}} \hyperref[tag:prophecy]{\texttt{prophecy}}} \\
\hyperref[case:Example5008]{Example5008} & Soufflé tc benchmark: naive versus semi-naive evaluation & souffle-benchmark & equivalence & certified valid & {\scriptsize \hyperref[tag:database-engine]{\texttt{database-engine}} \hyperref[tag:evaluation-strategies]{\texttt{evaluation-strategies}} \hyperref[tag:nonlinear-recursion]{\texttt{nonlinear-recursion}} \hyperref[tag:prophecy]{\texttt{prophecy}} \hyperref[tag:semi-naive]{\texttt{semi-naive}}} \\
\hyperref[case:Example5009]{Example5009} & Soufflé family benchmark: naive versus semi-naive evaluation & souffle-benchmark & equivalence & certified valid & {\scriptsize \hyperref[tag:database-engine]{\texttt{database-engine}} \hyperref[tag:evaluation-strategies]{\texttt{evaluation-strategies}} \hyperref[tag:negation]{\texttt{negation}} \hyperref[tag:prophecy]{\texttt{prophecy}} \hyperref[tag:semi-naive]{\texttt{semi-naive}}} \\
\hyperref[case:Example5010]{Example5010} & Soufflé orbits benchmark: naive versus semi-naive evaluation & souffle-benchmark & equivalence & certified valid & {\scriptsize \hyperref[tag:database-engine]{\texttt{database-engine}} \hyperref[tag:evaluation-strategies]{\texttt{evaluation-strategies}} \hyperref[tag:negation]{\texttt{negation}} \hyperref[tag:nonlinear-recursion]{\texttt{nonlinear-recursion}} \hyperref[tag:prophecy]{\texttt{prophecy}} \hyperref[tag:semi-naive]{\texttt{semi-naive}}} \\
\hyperref[case:Example5011]{Example5011} & Soufflé topological\_ordering benchmark: naive versus semi-naive evaluation & souffle-benchmark & equivalence & certified invalid & {\scriptsize \hyperref[tag:database-engine]{\texttt{database-engine}} \hyperref[tag:evaluation-strategies]{\texttt{evaluation-strategies}} \hyperref[tag:negation]{\texttt{negation}} \hyperref[tag:nonlinear-recursion]{\texttt{nonlinear-recursion}} \hyperref[tag:semi-naive]{\texttt{semi-naive}}} \\
\hyperref[case:Example5012]{Example5012} & Casbin depth-bounded role resolution contained in the unbounded role closure & production-pair & containment & certified valid & {\scriptsize \hyperref[tag:access-control]{\texttt{access-control}} \hyperref[tag:prophecy]{\texttt{prophecy}} \hyperref[tag:recursion-limit]{\texttt{recursion-limit}}} \\
\hyperref[case:Example5013]{Example5013} & Casbin depth-bounded role resolution versus the unbounded role closure & production-pair & equivalence & certified invalid & {\scriptsize \hyperref[tag:access-control]{\texttt{access-control}} \hyperref[tag:recursion-limit]{\texttt{recursion-limit}}} \\
\hyperref[case:Example5014]{Example5014} & W3C SPARQL 1.1 ALP evaluation of p* and p+ versus the Datalog closure with reflexive pairs over nodes(G) & specification & equivalence & certified valid & {\scriptsize \hyperref[tag:evaluation-strategies]{\texttt{evaluation-strategies}} \hyperref[tag:semi-naive]{\texttt{semi-naive}}} \\
\hyperref[case:Example5015]{Example5015} & W3C SPARQL 1.1 ALP evaluation of p* versus closure plus identity over the p endpoints (invalid twin of Example5014) & specification & equivalence & certified invalid & {\scriptsize \hyperref[tag:evaluation-strategies]{\texttt{evaluation-strategies}} \hyperref[tag:semi-naive]{\texttt{semi-naive}}} \\
\hyperref[case:Example5016]{Example5016} & DRed deletion maintenance of a materialised transitive closure versus recomputation & specification & equivalence & certified valid & {\scriptsize \hyperref[tag:evaluation-strategies]{\texttt{evaluation-strategies}} \hyperref[tag:negation]{\texttt{negation}} \hyperref[tag:prophecy]{\texttt{prophecy}} \hyperref[tag:semi-naive]{\texttt{semi-naive}} \hyperref[tag:view-maintenance]{\texttt{view-maintenance}}} \\
\hyperref[case:Example5017]{Example5017} & Overdeletion without rederivation versus recomputation (invalid twin of Example5016) & specification & equivalence & certified invalid & {\scriptsize \hyperref[tag:evaluation-strategies]{\texttt{evaluation-strategies}} \hyperref[tag:negation]{\texttt{negation}} \hyperref[tag:semi-naive]{\texttt{semi-naive}} \hyperref[tag:view-maintenance]{\texttt{view-maintenance}}} \\
\hyperref[case:Example5018]{Example5018} & Chase order-independence on the full-TGD subset of ChaseBench correctness/tgds: dependency-at-a-time versus parallel chase & specification & equivalence & certified valid & {\scriptsize \hyperref[tag:chase]{\texttt{chase}} \hyperref[tag:evaluation-strategies]{\texttt{evaluation-strategies}} \hyperref[tag:mutual-recursion]{\texttt{mutual-recursion}} \hyperref[tag:prophecy]{\texttt{prophecy}}} \\
\hyperref[case:Example5019]{Example5019} & One pass over the dependencies versus the parallel chase on ChaseBench correctness/tgds (invalid twin of Example5018) & specification & equivalence & certified invalid & {\scriptsize \hyperref[tag:chase]{\texttt{chase}} \hyperref[tag:evaluation-strategies]{\texttt{evaluation-strategies}} \hyperref[tag:mutual-recursion]{\texttt{mutual-recursion}}} \\
\hyperref[case:Example5020]{Example5020} & RDFS closure by the rho-df rules versus the full W3C RDFS rule set: equal up to the reflexive subPropertyOf/subClassOf triples & specification & equivalence & certified valid & {\scriptsize \hyperref[tag:mutual-recursion]{\texttt{mutual-recursion}} \hyperref[tag:nonlinear-recursion]{\texttt{nonlinear-recursion}}} \\
\hyperref[case:Example5021]{Example5021} & rho-df versus full RDFS with the reflexive subClassOf triples forgotten (invalid twin of Example5020) & specification & equivalence & certified invalid & {\scriptsize \hyperref[tag:mutual-recursion]{\texttt{mutual-recursion}} \hyperref[tag:nonlinear-recursion]{\texttt{nonlinear-recursion}}} \\
\hyperref[case:Example5022]{Example5022} & SQL recursion limits (MySQL cte\_max\_recursion\_depth, SQL Server MAXRECURSION): completion within the limit implies the closure & specification & implication & certified valid & {\scriptsize \hyperref[tag:database-engine]{\texttt{database-engine}} \hyperref[tag:prophecy]{\texttt{prophecy}} \hyperref[tag:recursion-limit]{\texttt{recursion-limit}}} \\
\hyperref[case:Example5023]{Example5023} & Win-move: the alternating fixpoint of the well-founded semantics (nested loops) versus the attractor computation & specification & equivalence & certified valid & {\scriptsize \hyperref[tag:evaluation-strategies]{\texttt{evaluation-strategies}} \hyperref[tag:greatest-fixpoint]{\texttt{greatest-fixpoint}} \hyperref[tag:negation]{\texttt{negation}} \hyperref[tag:prophecy]{\texttt{prophecy}}} \\
\hyperref[case:Example5024]{Example5024} & Win-move: alternating fixpoint versus attractor, claiming the odd limit is the win set (invalid twin of Example5023) & specification & equivalence & certified invalid & {\scriptsize \hyperref[tag:evaluation-strategies]{\texttt{evaluation-strategies}} \hyperref[tag:greatest-fixpoint]{\texttt{greatest-fixpoint}} \hyperref[tag:negation]{\texttt{negation}}} \\
\hyperref[case:Example5030]{Example5030} & NetworkX per-source transitive closure versus PostgreSQL recursive-union rounds & production-pair & equivalence & certified valid & {\scriptsize \hyperref[tag:database-engine]{\texttt{database-engine}} \hyperref[tag:evaluation-strategies]{\texttt{evaluation-strategies}} \hyperref[tag:prophecy]{\texttt{prophecy}} \hyperref[tag:semi-naive]{\texttt{semi-naive}}} \\
\hyperref[case:Example5031]{Example5031} & NetworkX per-source search with cutoff 2 versus PostgreSQL recursive-union rounds (invalid twin of Example5030) & production-pair & equivalence & certified invalid & {\scriptsize \hyperref[tag:database-engine]{\texttt{database-engine}} \hyperref[tag:evaluation-strategies]{\texttt{evaluation-strategies}} \hyperref[tag:recursion-limit]{\texttt{recursion-limit}} \hyperref[tag:semi-naive]{\texttt{semi-naive}}} \\
\hyperref[case:Example5032]{Example5032} & Delta-driven incremental maintenance of a transitive closure under insertions versus PostgreSQL recomputation on the updated edges & production-pair & equivalence & certified valid & {\scriptsize \hyperref[tag:database-engine]{\texttt{database-engine}} \hyperref[tag:evaluation-strategies]{\texttt{evaluation-strategies}} \hyperref[tag:prophecy]{\texttt{prophecy}} \hyperref[tag:semi-naive]{\texttt{semi-naive}} \hyperref[tag:view-maintenance]{\texttt{view-maintenance}}} \\
\hyperref[case:Example5033]{Example5033} & Soufflé magic-set demand-driven evaluation versus PostgreSQL bulk evaluation restricted to the demanded sources & production-pair & equivalence & certified valid & {\scriptsize \hyperref[tag:database-engine]{\texttt{database-engine}} \hyperref[tag:evaluation-strategies]{\texttt{evaluation-strategies}} \hyperref[tag:magic-sets]{\texttt{magic-sets}} \hyperref[tag:prophecy]{\texttt{prophecy}} \hyperref[tag:semi-naive]{\texttt{semi-naive}}} \\
\hyperref[case:Example5034]{Example5034} & Warshall's intermediate-vertex sweep (boolean Floyd-Warshall) versus PostgreSQL recursive-union rounds & production-pair & equivalence & certified valid & {\scriptsize \hyperref[tag:database-engine]{\texttt{database-engine}} \hyperref[tag:evaluation-strategies]{\texttt{evaluation-strategies}} \hyperref[tag:prophecy]{\texttt{prophecy}} \hyperref[tag:semi-naive]{\texttt{semi-naive}}} \\
\hyperref[case:Example5035]{Example5035} & Warshall's sweep truncated to the first two vertices versus PostgreSQL recursive-union rounds (invalid twin of Example5034) & production-pair & equivalence & certified invalid & {\scriptsize \hyperref[tag:database-engine]{\texttt{database-engine}} \hyperref[tag:evaluation-strategies]{\texttt{evaluation-strategies}} \hyperref[tag:recursion-limit]{\texttt{recursion-limit}} \hyperref[tag:semi-naive]{\texttt{semi-naive}}} \\
\hyperref[case:Example5036]{Example5036} & NetworkX DAG transitive closure by a reverse topological sweep versus PostgreSQL recursive-union rounds & production-pair & equivalence & certified valid & {\scriptsize \hyperref[tag:database-engine]{\texttt{database-engine}} \hyperref[tag:evaluation-strategies]{\texttt{evaluation-strategies}} \hyperref[tag:prophecy]{\texttt{prophecy}} \hyperref[tag:semi-naive]{\texttt{semi-naive}}} \\
\hyperref[case:Example5037]{Example5037} & Go runtime frontier marking versus the materialised Soufflé closure of the pointer graph joined with the roots & production-pair & equivalence & certified valid & {\scriptsize \hyperref[tag:evaluation-strategies]{\texttt{evaluation-strategies}} \hyperref[tag:nonlinear-recursion]{\texttt{nonlinear-recursion}} \hyperref[tag:prophecy]{\texttt{prophecy}} \hyperref[tag:semi-naive]{\texttt{semi-naive}}} \\
\hyperref[case:Example5038]{Example5038} & OpenFGA Check by top-down userset dispatch versus OpenFGA ListObjects by reverse expansion from the user & production-pair & equivalence & certified valid & {\scriptsize \hyperref[tag:access-control]{\texttt{access-control}} \hyperref[tag:evaluation-strategies]{\texttt{evaluation-strategies}} \hyperref[tag:magic-sets]{\texttt{magic-sets}} \hyperref[tag:prophecy]{\texttt{prophecy}}} \\
\hyperref[case:Example5039]{Example5039} & Permission exclusion applied at every object (SpiceDB/OpenFGA) is contained in the exclusion applied once to the finished closure & production-pair & containment & certified valid & {\scriptsize \hyperref[tag:access-control]{\texttt{access-control}} \hyperref[tag:negation]{\texttt{negation}}} \\
\hyperref[case:Example5040]{Example5040} & Permission exclusion applied at every object versus once to the finished closure, claimed equal (invalid twin of Example5039) & production-pair & equivalence & certified invalid & {\scriptsize \hyperref[tag:access-control]{\texttt{access-control}} \hyperref[tag:negation]{\texttt{negation}}} \\
\hyperref[case:Example5041]{Example5041} & PostgreSQL foreign-key cascade as a trigger worklist with the surviving-reference RESTRICT check versus the declarative delete closure & production-pair & equivalence & certified valid & {\scriptsize \hyperref[tag:database-engine]{\texttt{database-engine}} \hyperref[tag:evaluation-strategies]{\texttt{evaluation-strategies}} \hyperref[tag:integrity-constraints]{\texttt{integrity-constraints}} \hyperref[tag:negation]{\texttt{negation}} \hyperref[tag:prophecy]{\texttt{prophecy}} \hyperref[tag:semi-naive]{\texttt{semi-naive}}} \\
\hyperref[case:Example5042]{Example5042} & The over-approximating RESTRICT check of legacy Example2024 versus PostgreSQL's surviving-reference check, claimed equal (invalid twin of Example5041) & production-pair & equivalence & certified invalid & {\scriptsize \hyperref[tag:database-engine]{\texttt{database-engine}} \hyperref[tag:evaluation-strategies]{\texttt{evaluation-strategies}} \hyperref[tag:integrity-constraints]{\texttt{integrity-constraints}} \hyperref[tag:negation]{\texttt{negation}} \hyperref[tag:semi-naive]{\texttt{semi-naive}}} \\
\bottomrule
\end{longtable}

\section{Tag vocabulary and index}\label{sec:tags}
A short vocabulary of what makes a program interesting to database people. The three recursion tags are computed from the Datalog rules and the Whiel text by the build script; the others are assigned by hand in \path{Benchmark/report/curated.json}; the vocabulary lives in \path{Benchmark/report/tags.json}. Selecting cases by several tags at once is supported by the HTML page and the spreadsheet next to this report.
\subsection*{rules}
\begin{description}[leftmargin=1.2em,itemsep=2pt]
\item[\texttt{nonlinear-recursion}\label{tag:nonlinear-recursion}] (31) A recursive rule joins two recursive relations (T(x,y) :- T(x,z), T(z,y)), or a hand-written loop composes a relation with itself (squaring): recursion beyond the linear recursion that SQL recursive CTEs support.\\ {\small \hyperref[case:Example0106]{0106}, \hyperref[case:Example0107]{0107}, \hyperref[case:Example0117]{0117}, \hyperref[case:Example0119]{0119}, \hyperref[case:Example0123]{0123}, \hyperref[case:Example0125]{0125}, \hyperref[case:Example0128]{0128}, \hyperref[case:Example0130]{0130}, \hyperref[case:Example0132]{0132}, \hyperref[case:Example0133]{0133}, \hyperref[case:Example0134]{0134}, \hyperref[case:Example0136]{0136}, \hyperref[case:Example0137]{0137}, \hyperref[case:Example0138]{0138}, \hyperref[case:Example1081]{1081}, \hyperref[case:Example1088]{1088}, \hyperref[case:Example1090]{1090}, \hyperref[case:Example1091]{1091}, \hyperref[case:Example1093]{1093}, \hyperref[case:Example1094]{1094}, \hyperref[case:Example3001]{3001}, \hyperref[case:Example3009]{3009}, \hyperref[case:Example3014]{3014}, \hyperref[case:Example4040]{4040}, \hyperref[case:Example5007]{5007}, \hyperref[case:Example5008]{5008}, \hyperref[case:Example5010]{5010}, \hyperref[case:Example5011]{5011}, \hyperref[case:Example5020]{5020}, \hyperref[case:Example5021]{5021}, \hyperref[case:Example5037]{5037}}
\item[\texttt{mutual-recursion}\label{tag:mutual-recursion}] (30) Two or more relations defined in terms of each other (odd/even paths, RDFS classes and properties), which SQL recursive CTEs do not express directly.\\ {\small \hyperref[case:Example0109]{0109}, \hyperref[case:Example0115]{0115}, \hyperref[case:Example0117]{0117}, \hyperref[case:Example0119]{0119}, \hyperref[case:Example0123]{0123}, \hyperref[case:Example0125]{0125}, \hyperref[case:Example0128]{0128}, \hyperref[case:Example0130]{0130}, \hyperref[case:Example0132]{0132}, \hyperref[case:Example0133]{0133}, \hyperref[case:Example0134]{0134}, \hyperref[case:Example0136]{0136}, \hyperref[case:Example0137]{0137}, \hyperref[case:Example0138]{0138}, \hyperref[case:Example1081]{1081}, \hyperref[case:Example1088]{1088}, \hyperref[case:Example1090]{1090}, \hyperref[case:Example1091]{1091}, \hyperref[case:Example1093]{1093}, \hyperref[case:Example1094]{1094}, \hyperref[case:Example3001]{3001}, \hyperref[case:Example3009]{3009}, \hyperref[case:Example3014]{3014}, \hyperref[case:Example4004]{4004}, \hyperref[case:Example4034]{4034}, \hyperref[case:Example4041]{4041}, \hyperref[case:Example5018]{5018}, \hyperref[case:Example5019]{5019}, \hyperref[case:Example5020]{5020}, \hyperref[case:Example5021]{5021}}
\item[\texttt{negation}\label{tag:negation}] (16) Negation with semantic content: negated body atoms (stratified, or unstratified as in win-move), set difference in the claim, or a program whose meaning is an exclusion or a deletion (policy exclusion, edge filtering, overdeletion, a surviving-reference check). The set differences of frontier bookkeeping (Delta ∖ Acc, a cursor over an order) do not count.\\ {\small \hyperref[case:Example0162]{0162}, \hyperref[case:Example0163]{0163}, \hyperref[case:Example1030]{1030}, \hyperref[case:Example2005]{2005}, \hyperref[case:Example2006]{2006}, \hyperref[case:Example5009]{5009}, \hyperref[case:Example5010]{5010}, \hyperref[case:Example5011]{5011}, \hyperref[case:Example5016]{5016}, \hyperref[case:Example5017]{5017}, \hyperref[case:Example5023]{5023}, \hyperref[case:Example5024]{5024}, \hyperref[case:Example5039]{5039}, \hyperref[case:Example5040]{5040}, \hyperref[case:Example5041]{5041}, \hyperref[case:Example5042]{5042}}
\end{description}
\subsection*{query}
\begin{description}[leftmargin=1.2em,itemsep=2pt]
\item[\texttt{linear-vs-nonlinear}\label{tag:linear-vs-nonlinear}] (15) Two formulations of one recursive query compared: linear versus nonlinear (T∘T), or left- versus right-linear, the classic rewriting question for recursive SQL.\\ {\small \hyperref[case:Example0106]{0106}, \hyperref[case:Example0107]{0107}, \hyperref[case:Example0108]{0108}, \hyperref[case:Example0117]{0117}, \hyperref[case:Example0119]{0119}, \hyperref[case:Example0123]{0123}, \hyperref[case:Example0125]{0125}, \hyperref[case:Example0132]{0132}, \hyperref[case:Example0133]{0133}, \hyperref[case:Example0134]{0134}, \hyperref[case:Example0136]{0136}, \hyperref[case:Example0137]{0137}, \hyperref[case:Example4040]{4040}, \hyperref[case:Example5006]{5006}, \hyperref[case:Example5007]{5007}}
\item[\texttt{views}\label{tag:views}] (10) Query reformulation over views (answering queries using views): LAV views, exact views stated as precondition equations, sound views stated as an implication in the claim, Datalog-defined views. A view that only links two schemas inside a claim does not count.\\ {\small \hyperref[case:Example4001]{4001}, \hyperref[case:Example4002]{4002}, \hyperref[case:Example4003]{4003}, \hyperref[case:Example4004]{4004}, \hyperref[case:Example4034]{4034}, \hyperref[case:Example4035]{4035}, \hyperref[case:Example4036]{4036}, \hyperref[case:Example4037]{4037}, \hyperref[case:Example4040]{4040}, \hyperref[case:Example4041]{4041}}
\item[\texttt{magic-sets}\label{tag:magic-sets}] (3) Magic-set / demand transformation of a recursive query.\\ {\small \hyperref[case:Example2004]{2004}, \hyperref[case:Example5033]{5033}, \hyperref[case:Example5038]{5038}}
\item[\texttt{dependencies}\label{tag:dependencies}] (9) Two programs compared under tuple-generating dependencies between their input relations, stated as the precondition, including existential ones (a witness must exist): containment or equivalence under constraints, the setting of the chase.\\ {\small \hyperref[case:Example1081]{1081}, \hyperref[case:Example1088]{1088}, \hyperref[case:Example1090]{1090}, \hyperref[case:Example1091]{1091}, \hyperref[case:Example1093]{1093}, \hyperref[case:Example1094]{1094}, \hyperref[case:Example3001]{3001}, \hyperref[case:Example3009]{3009}, \hyperref[case:Example3014]{3014}}
\item[\texttt{fixpoint-witness}\label{tag:fixpoint-witness}] (21) The claim is stated through a free witness relation, never assigned by the command and occurring in the postcondition, that ranges over all pre-fixpoints (or post-fixpoints) of a rule set (RBound, Ra, Closed, ...): the executed program is compared with the least (or greatest) fixpoint of the rules, the equ.pdf method, instead of with a second executed program.\\ {\small \hyperref[case:Example0106]{0106}, \hyperref[case:Example0107]{0107}, \hyperref[case:Example0108]{0108}, \hyperref[case:Example0109]{0109}, \hyperref[case:Example0115]{0115}, \hyperref[case:Example0117]{0117}, \hyperref[case:Example0119]{0119}, \hyperref[case:Example0123]{0123}, \hyperref[case:Example0125]{0125}, \hyperref[case:Example0128]{0128}, \hyperref[case:Example0130]{0130}, \hyperref[case:Example0162]{0162}, \hyperref[case:Example0163]{0163}, \hyperref[case:Example1030]{1030}, \hyperref[case:Example1061]{1061}, \hyperref[case:Example2005]{2005}, \hyperref[case:Example2006]{2006}, \hyperref[case:Example2007]{2007}, \hyperref[case:Example2010]{2010}, \hyperref[case:Example2011]{2011}, \hyperref[case:Example2012]{2012}}
\end{description}
\subsection*{evaluation}
\begin{description}[leftmargin=1.2em,itemsep=2pt]
\item[\texttt{evaluation-strategies}\label{tag:evaluation-strategies}] (33) Two evaluation algorithms of one recursive query are compared, or a non-naive evaluation algorithm is verified against the declarative meaning of its rules: naive versus semi-naive, worklist/BFS, per-source search, squaring, Warshall or topological sweeps, delta-driven or recomputing maintenance, chase orders, top-down versus bottom-up, materialize-then-join versus recursive plan. A naive iteration verified against the least fixpoint of its own rules does not count.\\ {\small \hyperref[case:Example2010]{2010}, \hyperref[case:Example2011]{2011}, \hyperref[case:Example2012]{2012}, \hyperref[case:Example5001]{5001}, \hyperref[case:Example5002]{5002}, \hyperref[case:Example5003]{5003}, \hyperref[case:Example5004]{5004}, \hyperref[case:Example5005]{5005}, \hyperref[case:Example5006]{5006}, \hyperref[case:Example5007]{5007}, \hyperref[case:Example5008]{5008}, \hyperref[case:Example5009]{5009}, \hyperref[case:Example5010]{5010}, \hyperref[case:Example5011]{5011}, \hyperref[case:Example5014]{5014}, \hyperref[case:Example5015]{5015}, \hyperref[case:Example5016]{5016}, \hyperref[case:Example5017]{5017}, \hyperref[case:Example5018]{5018}, \hyperref[case:Example5019]{5019}, \hyperref[case:Example5023]{5023}, \hyperref[case:Example5024]{5024}, \hyperref[case:Example5030]{5030}, \hyperref[case:Example5031]{5031}, \hyperref[case:Example5032]{5032}, \hyperref[case:Example5033]{5033}, \hyperref[case:Example5034]{5034}, \hyperref[case:Example5035]{5035}, \hyperref[case:Example5036]{5036}, \hyperref[case:Example5037]{5037}, \hyperref[case:Example5038]{5038}, \hyperref[case:Example5041]{5041}, \hyperref[case:Example5042]{5042}}
\item[\texttt{semi-naive}\label{tag:semi-naive}] (25) One side evaluates the recursion by the semi-naive (delta/frontier) technique of Datalog engines and SQL recursive CTE executors: only the facts new in the last round are joined and the loop stops on an empty frontier, whatever the source calls it (PostgreSQL working table, Datafrog delta, GC grey set, BFS worklist, trigger queue, SPARQL visited set, DRed rounds). Computed from the command.\\ {\small \hyperref[case:Example2007]{2007}, \hyperref[case:Example2010]{2010}, \hyperref[case:Example2012]{2012}, \hyperref[case:Example5001]{5001}, \hyperref[case:Example5002]{5002}, \hyperref[case:Example5003]{5003}, \hyperref[case:Example5006]{5006}, \hyperref[case:Example5008]{5008}, \hyperref[case:Example5009]{5009}, \hyperref[case:Example5010]{5010}, \hyperref[case:Example5011]{5011}, \hyperref[case:Example5014]{5014}, \hyperref[case:Example5015]{5015}, \hyperref[case:Example5016]{5016}, \hyperref[case:Example5017]{5017}, \hyperref[case:Example5030]{5030}, \hyperref[case:Example5031]{5031}, \hyperref[case:Example5032]{5032}, \hyperref[case:Example5033]{5033}, \hyperref[case:Example5034]{5034}, \hyperref[case:Example5035]{5035}, \hyperref[case:Example5036]{5036}, \hyperref[case:Example5037]{5037}, \hyperref[case:Example5041]{5041}, \hyperref[case:Example5042]{5042}}
\item[\texttt{view-maintenance}\label{tag:view-maintenance}] (5) Incremental maintenance of a materialized recursive view (insertions, deletions with DRed, delta-driven) against recomputation from scratch.\\ {\small \hyperref[case:Example2011]{2011}, \hyperref[case:Example5005]{5005}, \hyperref[case:Example5016]{5016}, \hyperref[case:Example5017]{5017}, \hyperref[case:Example5032]{5032}}
\item[\texttt{recursion-limit}\label{tag:recursion-limit}] (5) One side truncates the recursion: SQL recursion limits (MySQL cte\_max\_recursion\_depth, SQL Server MAXRECURSION), a depth bound (Casbin maxHierarchyLevel), a search cutoff.\\ {\small \hyperref[case:Example5012]{5012}, \hyperref[case:Example5013]{5013}, \hyperref[case:Example5022]{5022}, \hyperref[case:Example5031]{5031}, \hyperref[case:Example5035]{5035}}
\end{description}
\subsection*{semantics}
\begin{description}[leftmargin=1.2em,itemsep=2pt]
\item[\texttt{chase}\label{tag:chase}] (2) The chase with tuple-generating dependencies (data exchange, ChaseBench scenarios).\\ {\small \hyperref[case:Example5018]{5018}, \hyperref[case:Example5019]{5019}}
\item[\texttt{greatest-fixpoint}\label{tag:greatest-fixpoint}] (4) The meaning is a greatest fixpoint or an alternating fixpoint: safe sets, integrity cores, the well-founded semantics of win-move.\\ {\small \hyperref[case:Example0163]{0163}, \hyperref[case:Example1030]{1030}, \hyperref[case:Example5023]{5023}, \hyperref[case:Example5024]{5024}}
\item[\texttt{integrity-constraints}\label{tag:integrity-constraints}] (3) Foreign keys, ON DELETE CASCADE / RESTRICT, referential-integrity cores.\\ {\small \hyperref[case:Example1030]{1030}, \hyperref[case:Example5041]{5041}, \hyperref[case:Example5042]{5042}}
\end{description}
\subsection*{systems}
\begin{description}[leftmargin=1.2em,itemsep=2pt]
\item[\texttt{database-engine}\label{tag:database-engine}] (18) The program is a production database or Datalog engine's own mechanism, read from its source or documentation: PostgreSQL recursive union and RI triggers, Soufflé's magic-set transformation and semi-naive translation, Datafrog, MySQL / SQL Server recursion limits. A system that is only the setting of the claim (the streaming engines behind an insert-only restart) or a paper algorithm (DRed) does not count.\\ {\small \hyperref[case:Example2004]{2004}, \hyperref[case:Example2010]{2010}, \hyperref[case:Example5001]{5001}, \hyperref[case:Example5002]{5002}, \hyperref[case:Example5008]{5008}, \hyperref[case:Example5009]{5009}, \hyperref[case:Example5010]{5010}, \hyperref[case:Example5011]{5011}, \hyperref[case:Example5022]{5022}, \hyperref[case:Example5030]{5030}, \hyperref[case:Example5031]{5031}, \hyperref[case:Example5032]{5032}, \hyperref[case:Example5033]{5033}, \hyperref[case:Example5034]{5034}, \hyperref[case:Example5035]{5035}, \hyperref[case:Example5036]{5036}, \hyperref[case:Example5041]{5041}, \hyperref[case:Example5042]{5042}}
\item[\texttt{access-control}\label{tag:access-control}] (6) Relationship-based access control and role hierarchies as production systems compute them (Zanzibar-style OpenFGA/SpiceDB, Casbin RBAC): Datalog-like reachability in production authorization databases. The LLM-generated delegation stories carry dependencies instead.\\ {\small \hyperref[case:Example5004]{5004}, \hyperref[case:Example5012]{5012}, \hyperref[case:Example5013]{5013}, \hyperref[case:Example5038]{5038}, \hyperref[case:Example5039]{5039}, \hyperref[case:Example5040]{5040}}
\end{description}
\subsection*{verification}
\begin{description}[leftmargin=1.2em,itemsep=2pt]
\item[\texttt{prophecy}\label{tag:prophecy}] (42) The one verification tag: known or expected to need a prophecy-level invariant, formerly 'obstructed'; the case page says whether this is established or only expected.\\ {\small \hyperref[case:Example0001]{0001}, \hyperref[case:Example0106]{0106}, \hyperref[case:Example0119]{0119}, \hyperref[case:Example0125]{0125}, \hyperref[case:Example0132]{0132}, \hyperref[case:Example0133]{0133}, \hyperref[case:Example0134]{0134}, \hyperref[case:Example0136]{0136}, \hyperref[case:Example0137]{0137}, \hyperref[case:Example0138]{0138}, \hyperref[case:Example1081]{1081}, \hyperref[case:Example2004]{2004}, \hyperref[case:Example4003]{4003}, \hyperref[case:Example4004]{4004}, \hyperref[case:Example4034]{4034}, \hyperref[case:Example4035]{4035}, \hyperref[case:Example4036]{4036}, \hyperref[case:Example4040]{4040}, \hyperref[case:Example4041]{4041}, \hyperref[case:Example5001]{5001}, \hyperref[case:Example5002]{5002}, \hyperref[case:Example5003]{5003}, \hyperref[case:Example5004]{5004}, \hyperref[case:Example5005]{5005}, \hyperref[case:Example5006]{5006}, \hyperref[case:Example5007]{5007}, \hyperref[case:Example5008]{5008}, \hyperref[case:Example5009]{5009}, \hyperref[case:Example5010]{5010}, \hyperref[case:Example5012]{5012}, \hyperref[case:Example5016]{5016}, \hyperref[case:Example5018]{5018}, \hyperref[case:Example5022]{5022}, \hyperref[case:Example5023]{5023}, \hyperref[case:Example5030]{5030}, \hyperref[case:Example5032]{5032}, \hyperref[case:Example5033]{5033}, \hyperref[case:Example5034]{5034}, \hyperref[case:Example5036]{5036}, \hyperref[case:Example5037]{5037}, \hyperref[case:Example5038]{5038}, \hyperref[case:Example5041]{5041}}
\end{description}

\clearpage

\Needspace{0.45\textheight}

\section{Example0001: Transitive closure by naive iteration under the weak transitivity specification}\label{case:Example0001}

\begin{description}[leftmargin=2.6cm,style=nextline,itemsep=1pt]

\item[Category] synthetic (syntactic generation from the classic Datalog programs); family: Core.

\item[Kind] correctness (closure/safety conditions on one program).

\item[Contributors] Jesse Comer, Fangzhu Shen.

\item[Status] certified valid (Certificate/Valid.lean, written by the certificate emitter); expected valid.

\item[Tags] \hyperref[tag:prophecy]{\texttt{prophecy}}.

\item[Input] \path{Benchmark/Example0001/Input.lean} (canonical notation).

\item[Summary] Transitive closure by naive iteration, under the weak specification that its result is transitively closed.

\item[Sources] {\raggedright legacy source record hard-lowerbound-spec of the retired classic-source collection\par}

\item[Prophecy] kernel-checked certificate with a level-1 prophecy clause (Certificate/Valid.lean, Core.json).

\item[Note] Carries the kernel-checked certificate for this triple. The classic Example0160 stated the same triple and was removed as a duplicate of this case.

\end{description}

\paragraph{Datalog program(s).}

{\raggedright\textit{Datalog reading of the executed loop (written for this report; right-linear transitive closure, S = E ∪ E∘T, T := S)} --- hand-written reading (this report).\par}

\begin{Verbatim}[fontsize=\scriptsize,breaklines=true,breakanywhere=true,frame=single,framesep=1.5mm]
T(x, y) :- E(x, y).
T(x, y) :- E(x, z), T(z, y).
\end{Verbatim}

\paragraph{Reading.} The loop performs naive evaluation of the right-linear transitive-closure program T(x,y):-E(x,y), T(x,y):-E(x,z),T(z,y), iterating S:=E∪(E∘T) until S=T, the classical way to compute reachability (transitive closure) in a directed graph E. The Hoare triple states only a weak postcondition: that the computed T is transitively closed (T∘T⊆T), with no assertion pinning T from above, so it does not claim T is the least or exact transitive closure, only that it is closed under composition with itself.

\paragraph{First-order reading of PRE/POST, translated mechanically from the relational-algebra assertion in Input.lean.} The relational-algebra assertion in \texttt{Input.lean} remains the authority.

\textit{PRE:} ⊤\par The precondition is simply true, imposing no constraint on the input relations.

\textit{POST:} ∀x y. ((∃z. T(x,z) ∧ T(z,y)) → T(x,y))\par The postcondition asserts that the computed relation T is transitively closed, i.e., composing T with itself never leaves T: (T∘T) ⊆ T.

\paragraph{Whiel encoding (Hoare triple).}

\begin{Verbatim}[fontsize=\scriptsize,breaklines=true,breakanywhere=true,frame=single,framesep=1.5mm]
SCHEMA
  E/2, T/2, S/2

PRE
  true 

CMD
  T := ∅;
  S := (E ∪ (π[0, 3] (σ[#1 = #2] (E × T))));
  WHILE (S ≠ T) DO
    T := S;
    S := (E ∪ (π[0, 3] (σ[#1 = #2] (E × T))))
  END

POST
  ((π[0, 3] (σ[#1 = #2] (T × T))) ⊆ T)
\end{Verbatim}

\Needspace{0.45\textheight}

\section{Example0013: Transitive closure by naive iteration under a false specification (T ⊆ E)}\label{case:Example0013}

\begin{description}[leftmargin=2.6cm,style=nextline,itemsep=1pt]

\item[Category] synthetic (syntactic generation from the classic Datalog programs); family: Core.

\item[Kind] correctness (closure/safety conditions on one program).

\item[Contributors] Jesse Comer, Fangzhu Shen.

\item[Status] certified invalid (Certificate/Invalid.lean, written by the certificate emitter); expected invalid.

\item[Tags] .

\item[Input] \path{Benchmark/Example0013/Input.lean} (canonical notation).

\item[Summary] Transitive closure by naive iteration, under a false specification that a two-edge instance refutes.

\item[Sources] none recorded.

\item[Note] The refutable twin of Example0001: the same schema, precondition and command under the false postcondition T ⊆ E, refuted by the two-edge instance E = \{(0, 1), (1, 2)\}.

\end{description}

\paragraph{Datalog program(s).}

{\raggedright\textit{Datalog reading of the executed loop (written for this report; right-linear transitive closure, S = E ∪ E∘T, T := S)} --- hand-written reading (this report).\par}

\begin{Verbatim}[fontsize=\scriptsize,breaklines=true,breakanywhere=true,frame=single,framesep=1.5mm]
T(x, y) :- E(x, y).
T(x, y) :- E(x, z), T(z, y).
\end{Verbatim}

\paragraph{Reading.} The program and schema are identical to Example0001 --- naive right-linear evaluation of transitive closure T(x,y):-E(x,y), T(x,y):-E(x,z),T(z,y) --- but the postcondition is replaced by the false claim that the computed closure adds nothing beyond the base edges, T⊆E. The Hoare triple is invalid: on the two-edge instance E=\{(0,1),(1,2)\} the loop computes T=\{(0,1),(1,2),(0,2)\}, and the derived pair (0,2) is not itself an edge, refuting T⊆E.

\paragraph{First-order reading of PRE/POST, translated mechanically from the relational-algebra assertion in Input.lean.} The relational-algebra assertion in \texttt{Input.lean} remains the authority.

\textit{PRE:} ⊤\par The precondition is simply true, imposing no constraint on the input relations.

\textit{POST:} ∀x y. (T(x,y) → E(x,y))\par The postcondition claims the computed transitive closure T is contained in the base edge relation E, i.e., that transitive closure adds no new pairs.

\paragraph{Whiel encoding (Hoare triple).}

\begin{Verbatim}[fontsize=\scriptsize,breaklines=true,breakanywhere=true,frame=single,framesep=1.5mm]
SCHEMA
  E/2, T/2, S/2

PRE
  true 

CMD
  T := ∅;
  S := (E ∪ (π[0, 3] (σ[#1 = #2] (E × T))));
  WHILE (S ≠ T) DO
    T := S;
    S := (E ∪ (π[0, 3] (σ[#1 = #2] (E × T))))
  END

POST
  (T ⊆ E)
\end{Verbatim}

\Needspace{0.45\textheight}

\section{Example0106: Prefix-equivalence: naive tc\_left computes lfp of tc\_nonlinear}\label{case:Example0106}

\begin{description}[leftmargin=2.6cm,style=nextline,itemsep=1pt]

\item[Category] synthetic (syntactic generation from the classic Datalog programs); family: Synthetic.

\item[Kind] equivalence (with a least fixpoint, via a pre-fixpoint witness).

\item[Contributors] Leo Zhang.

\item[Status] certified valid (Certificate/Valid.lean, written by the certificate emitter); expected valid.

\item[Tags] \hyperref[tag:fixpoint-witness]{\texttt{fixpoint-witness}} \hyperref[tag:linear-vs-nonlinear]{\texttt{linear-vs-nonlinear}} \hyperref[tag:nonlinear-recursion]{\texttt{nonlinear-recursion}} \hyperref[tag:prophecy]{\texttt{prophecy}}.

\item[Input] \path{Benchmark/Example0106/Input.lean} (canonical notation).

\item[Summary] A left-linear closure loop, specified as the least fixpoint of the non-linear closure program.

\item[Sources] {\raggedright retired-corpus source record synthetic\_6 of the Synthetic family \par retired-corpus input record synthetic\_6 of the Synthetic family \par classic Datalog program tc-left of the retired generator suite \par classic Datalog program tc-nonlinear of the retired generator suite\par}

\item[Prophecy] kernel-checked certificate with a level-1 prophecy clause (Certificate/Valid.lean, Core.json).

\end{description}

\paragraph{Datalog program(s).}

{\raggedright\textit{tc\_left} --- legacy inventory: source programs with one executed WHIEL program and one formal PRE/POST reference (not executed); sources: tc-left.dl, tc-nonlinear.dl.\par}

\begin{Verbatim}[fontsize=\scriptsize,breaklines=true,breakanywhere=true,frame=single,framesep=1.5mm]
-- tc_left
-- Executed WHIEL program
tcl(x, y) :- base(x, y).
tcl(x, y) :- tcl(x, z), base(z, y).
\end{Verbatim}

{\raggedright\textit{tc\_nonlinear} --- legacy inventory: source programs with one executed WHIEL program and one formal PRE/POST reference (not executed); sources: tc-left.dl, tc-nonlinear.dl.\par}

\begin{Verbatim}[fontsize=\scriptsize,breaklines=true,breakanywhere=true,frame=single,framesep=1.5mm]
-- tc_nonlinear
-- Formal PRE/POST reference only; not executed by the WHIEL command
tc(x, y) :- base(x, y).
tc(x, y) :- tc(x, z), tc(z, y).
\end{Verbatim}

\paragraph{Reading.} The executed loop naively evaluates the left-linear transitive-closure program T(x,y):-Base(x,y), T(x,y):-T(x,z),Base(z,y), while the reference program against which it is measured is the ordinary non-linear transitive closure Tc(x,y):-Base(x,y), Tc(x,y):-Tc(x,z),Tc(z,y). The Hoare triple is the standard least-fixpoint introduction obligation: given any relation RBound that is already a pre-fixpoint of the non-linear closure program, the naively computed T must itself be a pre-fixpoint of that same program and lie below RBound, establishing that the left-linear loop computes the least pre-fixpoint (i.e. the transitive closure) of the non-linear program.

\paragraph{First-order reading of PRE/POST, translated mechanically from the relational-algebra assertion in Input.lean.} The relational-algebra assertion in \texttt{Input.lean} remains the authority.

\textit{PRE:} ∀x y. (Base(x,y) ∨ (∃z. RBound(x,z) ∧ RBound(z,y)) → RBound(x,y))\par The precondition states that RBound is a pre-fixpoint of the non-linear closure program: Base ∪ (RBound∘RBound) ⊆ RBound.

\textit{POST:} (∀x y. (Base(x,y) ∨ (∃z. T(x,z) ∧ T(z,y)) → T(x,y))) ∧ ∀u v. (T(u,v) → RBound(u,v))\par The postcondition states that the computed T is itself a pre-fixpoint of the non-linear closure program, Base ∪ (T∘T) ⊆ T, and that T ⊆ RBound.

\paragraph{Whiel encoding (Hoare triple).}

\begin{Verbatim}[fontsize=\scriptsize,breaklines=true,breakanywhere=true,frame=single,framesep=1.5mm]
SCHEMA
  Base/2, T/2, T_aux/2, RBound/2

PRE
  ((Base ∪ (π[0, 3] (σ[#1 = #2] (RBound × RBound))))
    ⊆ RBound)

CMD
  T_aux := ∅;
  T := Base;
  WHILE T ≠ T_aux DO
    T_aux := T;
    T :=
      (Base ∪
        (π[0, 3] (σ[#1 = #2] (T_aux × Base))))
  END

POST
  (((Base ∪ (π[0, 3] (σ[#1 = #2] (T × T)))) ⊆ T)
    ∧ (T ⊆ RBound))
\end{Verbatim}

\Needspace{0.45\textheight}

\section{Example0107: Prefix-equivalence: naive tc\_nonlinear computes lfp of tc\_left}\label{case:Example0107}

\begin{description}[leftmargin=2.6cm,style=nextline,itemsep=1pt]

\item[Category] synthetic (syntactic generation from the classic Datalog programs); family: Synthetic.

\item[Kind] equivalence (with a least fixpoint, via a pre-fixpoint witness).

\item[Contributors] Leo Zhang.

\item[Status] certified valid (Certificate/Valid.lean, written by the certificate emitter); expected valid.

\item[Tags] \hyperref[tag:fixpoint-witness]{\texttt{fixpoint-witness}} \hyperref[tag:linear-vs-nonlinear]{\texttt{linear-vs-nonlinear}} \hyperref[tag:nonlinear-recursion]{\texttt{nonlinear-recursion}}.

\item[Input] \path{Benchmark/Example0107/Input.lean} (canonical notation).

\item[Summary] Canonical form of the earlier single-level encoding of this case: the same schema, precondition, command and postcondition; the snapshot relation T\_2 is now T\_aux.

\item[Sources] {\raggedright retired-corpus source record synthetic\_7 of the Synthetic family \par retired-corpus input record synthetic\_7 of the Synthetic family \par classic Datalog program tc-left of the retired generator suite \par classic Datalog program tc-nonlinear of the retired generator suite\par}

\end{description}

\paragraph{Datalog program(s).}

{\raggedright\textit{tc\_left} --- legacy inventory: source programs with one executed WHIEL program and one formal PRE/POST reference (not executed); sources: tc-left.dl, tc-nonlinear.dl.\par}

\begin{Verbatim}[fontsize=\scriptsize,breaklines=true,breakanywhere=true,frame=single,framesep=1.5mm]
-- tc_left
-- Formal PRE/POST reference only; not executed by the WHIEL command
tcl(x, y) :- base(x, y).
tcl(x, y) :- tcl(x, z), base(z, y).
\end{Verbatim}

{\raggedright\textit{tc\_nonlinear} --- legacy inventory: source programs with one executed WHIEL program and one formal PRE/POST reference (not executed); sources: tc-left.dl, tc-nonlinear.dl.\par}

\begin{Verbatim}[fontsize=\scriptsize,breaklines=true,breakanywhere=true,frame=single,framesep=1.5mm]
-- tc_nonlinear
-- Executed WHIEL program
tc(x, y) :- base(x, y).
tc(x, y) :- tc(x, z), tc(z, y).
\end{Verbatim}

\paragraph{Reading.} The executed loop naively evaluates the non-linear transitive-closure program by squaring, T:=Base∪(T∘T), while the reference program is the left-linear closure T(x,y):-Base(x,y), T(x,y):-T(x,z),Base(z,y). The Hoare triple is the least-fixpoint introduction obligation for that reference: given any RBound already a pre-fixpoint of the left-linear program, the naively computed non-linear T must itself be a pre-fixpoint of the left-linear program and lie below RBound, i.e. the non-linear loop computes the least pre-fixpoint of the left-linear program, which is again the transitive closure of Base.

\paragraph{First-order reading of PRE/POST, translated mechanically from the relational-algebra assertion in Input.lean.} The relational-algebra assertion in \texttt{Input.lean} remains the authority.

\textit{PRE:} ∀x y. (Base(x,y) ∨ (∃z. RBound(x,z) ∧ Base(z,y)) → RBound(x,y))\par The precondition states that RBound is a pre-fixpoint of the left-linear closure program: Base ∪ (RBound∘Base) ⊆ RBound.

\textit{POST:} (∀x y. (Base(x,y) ∨ (∃z. T(x,z) ∧ Base(z,y)) → T(x,y))) ∧ ∀u v. (T(u,v) → RBound(u,v))\par The postcondition states that the computed T is itself a pre-fixpoint of the left-linear closure program, Base ∪ (T∘Base) ⊆ T, and that T ⊆ RBound.

\paragraph{Whiel encoding (Hoare triple).}

\begin{Verbatim}[fontsize=\scriptsize,breaklines=true,breakanywhere=true,frame=single,framesep=1.5mm]
SCHEMA
  Base/2, T/2, T_aux/2, RBound/2

PRE
  ((Base ∪ (π[0, 3] (σ[#1 = #2] (RBound × Base)))) ⊆ RBound)

CMD
  T_aux := ∅;
  T := Base;
  WHILE T ≠ T_aux DO
    T_aux := T;
    T := (Base ∪ (π[0, 3] (σ[#1 = #2] (T_aux × T_aux))))
  END

POST
  (((Base ∪ (π[0, 3] (σ[#1 = #2] (T × Base)))) ⊆ T) ∧ (T ⊆ RBound))
\end{Verbatim}

\Needspace{0.45\textheight}

\section{Example0108: Prefix-equivalence: naive tc\_right computes lfp of tc\_left}\label{case:Example0108}

\begin{description}[leftmargin=2.6cm,style=nextline,itemsep=1pt]

\item[Category] synthetic (syntactic generation from the classic Datalog programs); family: Synthetic.

\item[Kind] equivalence (with a least fixpoint, via a pre-fixpoint witness).

\item[Contributors] Leo Zhang.

\item[Status] certified valid (Certificate/Valid.lean, written by the certificate emitter); expected valid.

\item[Tags] \hyperref[tag:fixpoint-witness]{\texttt{fixpoint-witness}} \hyperref[tag:linear-vs-nonlinear]{\texttt{linear-vs-nonlinear}}.

\item[Input] \path{Benchmark/Example0108/Input.lean} (canonical notation).

\item[Summary] Canonical form of the earlier single-level encoding of this case: the same schema, precondition, command and postcondition; the snapshot relation T\_2 is now T\_aux.

\item[Sources] {\raggedright retired-corpus source record synthetic\_8 of the Synthetic family \par retired-corpus input record synthetic\_8 of the Synthetic family \par classic Datalog program tc-left of the retired generator suite \par classic Datalog program tc-right of the retired generator suite\par}

\end{description}

\paragraph{Datalog program(s).}

{\raggedright\textit{tc\_left} --- legacy inventory: source programs with one executed WHIEL program and one formal PRE/POST reference (not executed); sources: tc-left.dl, tc-right.dl.\par}

\begin{Verbatim}[fontsize=\scriptsize,breaklines=true,breakanywhere=true,frame=single,framesep=1.5mm]
-- tc_left
-- Formal PRE/POST reference only; not executed by the WHIEL command
tcl(x, y) :- base(x, y).
tcl(x, y) :- tcl(x, z), base(z, y).
\end{Verbatim}

{\raggedright\textit{tc\_right} --- legacy inventory: source programs with one executed WHIEL program and one formal PRE/POST reference (not executed); sources: tc-left.dl, tc-right.dl.\par}

\begin{Verbatim}[fontsize=\scriptsize,breaklines=true,breakanywhere=true,frame=single,framesep=1.5mm]
-- tc_right
-- Executed WHIEL program
tcr(x, y) :- base(x, y).
tcr(x, y) :- base(x, z), tcr(z, y).
\end{Verbatim}

\paragraph{Reading.} The executed loop naively evaluates the right-linear transitive-closure program T(x,y):-Base(x,y), T(x,y):-Base(x,z),T(z,y), while the reference program is the left-linear closure T(x,y):-Base(x,y), T(x,y):-T(x,z),Base(z,y). The Hoare triple is the least-fixpoint introduction obligation for the left-linear reference: given any RBound already a pre-fixpoint of the left-linear program, the naively computed right-linear T must itself be a pre-fixpoint of the left-linear program and lie below RBound, i.e. right-linear naive evaluation computes the least pre-fixpoint of the left-linear program (both compute the transitive closure of Base).

\paragraph{First-order reading of PRE/POST, translated mechanically from the relational-algebra assertion in Input.lean.} The relational-algebra assertion in \texttt{Input.lean} remains the authority.

\textit{PRE:} ∀x y. (Base(x,y) ∨ (∃z. RBound(x,z) ∧ Base(z,y)) → RBound(x,y))\par The precondition states that RBound is a pre-fixpoint of the left-linear closure program: Base ∪ (RBound∘Base) ⊆ RBound.

\textit{POST:} (∀x y. (Base(x,y) ∨ (∃z. T(x,z) ∧ Base(z,y)) → T(x,y))) ∧ ∀u v. (T(u,v) → RBound(u,v))\par The postcondition states that the computed T is itself a pre-fixpoint of the left-linear closure program, Base ∪ (T∘Base) ⊆ T, and that T ⊆ RBound.

\paragraph{Whiel encoding (Hoare triple).}

\begin{Verbatim}[fontsize=\scriptsize,breaklines=true,breakanywhere=true,frame=single,framesep=1.5mm]
SCHEMA
  Base/2, T/2, T_aux/2, RBound/2

PRE
  ((Base ∪ (π[0, 3] (σ[#1 = #2] (RBound × Base)))) ⊆ RBound)

CMD
  T_aux := ∅;
  T := Base;
  WHILE T ≠ T_aux DO
    T_aux := T;
    T := (Base ∪ (π[0, 3] (σ[#1 = #2] (Base × T_aux))))
  END

POST
  (((Base ∪ (π[0, 3] (σ[#1 = #2] (T × Base)))) ⊆ T) ∧ (T ⊆ RBound))
\end{Verbatim}

\Needspace{0.45\textheight}

\section{Example0109: Containment: red-blue alternating paths ⊆ TC(R ∪ B)}\label{case:Example0109}

\begin{description}[leftmargin=2.6cm,style=nextline,itemsep=1pt]

\item[Category] synthetic (syntactic generation from the classic Datalog programs); family: Synthetic.

\item[Kind] containment (via a pre-fixpoint witness).

\item[Contributors] Leo Zhang.

\item[Status] certified valid (Certificate/Valid.lean, written by the certificate emitter); expected valid.

\item[Tags] \hyperref[tag:fixpoint-witness]{\texttt{fixpoint-witness}} \hyperref[tag:mutual-recursion]{\texttt{mutual-recursion}}.

\item[Input] \path{Benchmark/Example0109/Input.lean} (canonical notation).

\item[Summary] Canonical form of the earlier single-level encoding of this case: the same schema, precondition, command and postcondition; the snapshot relations RR\_2 is now RR\_aux, BB\_2 is now BB\_aux, RB\_2 is now RB\_aux, BR\_2 is now BR\_aux.

\item[Sources] {\raggedright retired-corpus source record synthetic\_9 of the Synthetic family \par retired-corpus input record synthetic\_9 of the Synthetic family\par}

\end{description}

\paragraph{Datalog program(s).}

{\raggedright\textit{red\_blue\_linear} --- legacy inventory: \path{catalog/source} intent plus actual left-recursive formal PRE/POST reference-operator reconstruction; sources: red\_blue\_linear.dl, tc\_of\_union.dl, \path{Benchmark/Example0109/Input.lean}.\par}

\begin{Verbatim}[fontsize=\scriptsize,breaklines=true,breakanywhere=true,frame=single,framesep=1.5mm]
-- red_blue_linear
-- Executed WHIEL program
RR(x,y) :- R(x,y).
BB(x,y) :- B(x,y).
RB(x,y) :- RR(x,z), B(z,y).
BR(x,y) :- BB(x,z), R(z,y).
RR(x,y) :- R(x,z), BR(z,y).
BB(x,y) :- B(x,z), RB(z,y).
\end{Verbatim}

{\raggedright\textit{tc\_of\_union} --- legacy inventory: \path{catalog/source} intent plus actual left-recursive formal PRE/POST reference-operator reconstruction; sources: red\_blue\_linear.dl, tc\_of\_union.dl, \path{Benchmark/Example0109/Input.lean}.\par}

\begin{Verbatim}[fontsize=\scriptsize,breaklines=true,breakanywhere=true,frame=single,framesep=1.5mm]
-- tc_of_union
-- Catalog/source target intent only; not executed by the WHIEL command
E(x,y) :- R(x,y).
E(x,y) :- B(x,y).
T(x,y) :- E(x,y).
T(x,y) :- E(x,z), T(z,y).
-- Actual formal PRE/POST reference only; not executed (left-recursive T composed with R union B)
E_union(x,y) :- R(x,y).
E_union(x,y) :- B(x,y).
T_left(x,y) :- E_union(x,y).
T_left(x,z) :- T_left(x,y), E_union(y,z).
\end{Verbatim}

\paragraph{Reading.} The loop naively evaluates a red/blue alternating-path program: RR and BB collect same-colour-start walks and RB, BR collect colour-switching walks, computed by mutual recursion matching red\_blue\_linear (e.g. RR(x,y):-R(x,z),BR(z,y), RB(x,y):-RR(x,z),B(z,y), etc.), so together the four relations characterize walks in a two-coloured graph that strictly alternate colour. The Hoare triple is a containment claim rather than an equivalence: it asserts that the union of these four alternating-path relations is contained in any pre-fixpoint bound RBound of the plain transitive closure of the underlying uncoloured graph R∪B, since every colour-alternating walk is in particular an ordinary walk in R∪B.

\paragraph{First-order reading of PRE/POST, translated mechanically from the relational-algebra assertion in Input.lean.} The relational-algebra assertion in \texttt{Input.lean} remains the authority.

\textit{PRE:} ∀x y. (R(x,y) ∨ B(x,y) ∨ (∃z. RBound(x,z) ∧ (R(z,y) ∨ B(z,y))) → RBound(x,y))\par The precondition states that RBound is a pre-fixpoint of the transitive-closure program on the union graph R∪B: (R∪B) ∪ (RBound∘(R∪B)) ⊆ RBound.

\textit{POST:} ∀x y. (RR(x,y) ∨ BB(x,y) ∨ RB(x,y) ∨ BR(x,y) → RBound(x,y))\par The postcondition states that the union of the four computed red/blue relations, RR∪BB∪RB∪BR, is contained in RBound.

\paragraph{Whiel encoding (Hoare triple).}

\begin{Verbatim}[fontsize=\scriptsize,breaklines=true,breakanywhere=true,frame=single,framesep=1.5mm]
SCHEMA
  R/2, B/2, RR/2, BB/2, RB/2, BR/2, RR_aux/2, BB_aux/2, RB_aux/2, BR_aux/2, RBound/2

PRE
  (((R ∪ B) ∪ (π[0, 3] (σ[#1 = #2] (RBound × (R ∪ B))))) ⊆ RBound)

CMD
  RR_aux := ∅;
  BB_aux := ∅;
  RB_aux := ∅;
  BR_aux := ∅;
  RR := R;
  BB := B;
  RB := (π[0, 3] (σ[#1 = #2] (R × B)));
  BR := (π[0, 3] (σ[#1 = #2] (B × R)));
  WHILE ((RR ≠ RR_aux) ∨ (BB ≠ BB_aux)) ∨ ((RB ≠ RB_aux) ∨ (BR ≠ BR_aux)) DO
    RR_aux := RR;
    BB_aux := BB;
    RB_aux := RB;
    BR_aux := BR;
    RB := (RB ∪ (π[0, 3] (σ[#1 = #2] (RR_aux × B))));
    BR := (BR ∪ (π[0, 3] (σ[#1 = #2] (BB_aux × R))));
    RR := (RR ∪ (π[0, 3] (σ[#1 = #2] (R × BR_aux))));
    BB := (BB ∪ (π[0, 3] (σ[#1 = #2] (B × RB_aux))))
  END

POST
  ((((RR ∪ BB) ∪ RB) ∪ BR) ⊆ RBound)
\end{Verbatim}

\Needspace{0.45\textheight}

\section{Example0115: Prefix-equivalence: naive all\_paths\_via\_even\_odd (linear) computes lfp of tc\_left}\label{case:Example0115}

\begin{description}[leftmargin=2.6cm,style=nextline,itemsep=1pt]

\item[Category] synthetic (syntactic generation from the classic Datalog programs); family: Synthetic.

\item[Kind] equivalence (with a least fixpoint, via a pre-fixpoint witness).

\item[Contributors] Fangzhu Shen, Leo Zhang.

\item[Status] certified valid (Certificate/Valid.lean, written by the certificate emitter); expected valid.

\item[Tags] \hyperref[tag:fixpoint-witness]{\texttt{fixpoint-witness}} \hyperref[tag:mutual-recursion]{\texttt{mutual-recursion}}.

\item[Input] \path{Benchmark/Example0115/Input.lean} (canonical notation).

\item[Summary] Canonical form of the earlier single-level encoding of this case: the same schema, precondition, command and postcondition.

\item[Sources] {\raggedright legacy source record v2\_\_prefix\_equivalence\_\_tc\_left\_\_all\_paths\_via\_even\_odd\_linear\_\_naive of the retired classic-source collection\par}

\end{description}

\paragraph{Datalog program(s).}

{\raggedright\textit{all\_paths\_via\_even\_odd\_linear} --- legacy inventory: \path{catalog/source} intent plus actual right-recursive formal PRE/POST reference-operator reconstruction; sources: all\_paths\_via\_even\_odd.dl, tc-left.dl, \path{Benchmark/Example0115/Input.lean}.\par}

\begin{Verbatim}[fontsize=\scriptsize,breaklines=true,breakanywhere=true,frame=single,framesep=1.5mm]
-- all_paths_via_even_odd_linear
-- Executed WHIEL program
OL(x,y) :- E(x,y).
OL(x,y) :- E(x,z), EL(z,y).
EL(x,y) :- E(x,z), OL(z,y).
T(x,y) :- EL(x,y).
T(x,y) :- OL(x,y).
\end{Verbatim}

{\raggedright\textit{tc\_left} --- legacy inventory: \path{catalog/source} intent plus actual right-recursive formal PRE/POST reference-operator reconstruction; sources: all\_paths\_via\_even\_odd.dl, tc-left.dl, \path{Benchmark/Example0115/Input.lean}.\par}

\begin{Verbatim}[fontsize=\scriptsize,breaklines=true,breakanywhere=true,frame=single,framesep=1.5mm]
-- tc_left
-- Catalog/source target intent only; not executed by the WHIEL command
tcl(x, y) :- base(x, y).
tcl(x, y) :- tcl(x, z), base(z, y).
-- Actual formal PRE/POST reference only; not executed (right-recursive E composed with T)
T_right(x,y) :- E(x,y).
T_right(x,y) :- E(x,z), T_right(z,y).
\end{Verbatim}

\paragraph{Reading.} The loop naively evaluates the all\_paths\_via\_even\_odd program: Ta and Tb are computed by mutual recursion matching the even/odd path relations (Ta(x,y):-E(x,z),Tb(z,y), Tb(x,y):-E(x,y) or E(x,z),Ta(z,y)), and Tc:=Ta∪Tb collects all paths of length one or more. Although the title names tc\_left as the reference, the precondition and postcondition actually encode the right-linear closure T(x,y):-E(x,y), T(x,y):-E(x,z),T(z,y) (as the recorded Datalog reference for this case acknowledges), so the Hoare triple is the least-fixpoint introduction obligation stating that the naively computed union Tc of even and odd paths is itself a pre-fixpoint of that right-linear closure program and lies below any given pre-fixpoint R --- i.e. the even/odd decomposition computes the transitive closure, since every nonempty path has either even or odd length.

\paragraph{First-order reading of PRE/POST, translated mechanically from the relational-algebra assertion in Input.lean.} The relational-algebra assertion in \texttt{Input.lean} remains the authority.

\textit{PRE:} ∀x y. (E(x,y) ∨ (∃z. E(x,z) ∧ R(z,y)) → R(x,y))\par The precondition states that R is a pre-fixpoint of the right-linear closure program: E ∪ (E∘R) ⊆ R.

\textit{POST:} (∀x y. (E(x,y) ∨ (∃z. E(x,z) ∧ Tc(z,y)) → Tc(x,y))) ∧ ∀u v. (Tc(u,v) → R(u,v))\par The postcondition states that the computed union Tc is itself a pre-fixpoint of the right-linear closure program, E ∪ (E∘Tc) ⊆ Tc, and that Tc ⊆ R.

\paragraph{Whiel encoding (Hoare triple).}

\begin{Verbatim}[fontsize=\scriptsize,breaklines=true,breakanywhere=true,frame=single,framesep=1.5mm]
SCHEMA
  E/2, R/2, Ta/2, Sa/2, Tb/2, Sb/2, Tc/2, Sc/2

PRE
  ((E ∪ (π[0, 3] (σ[#1 = #2] (E × R)))) ⊆ R)

CMD
  Ta := ∅;
  Tb := ∅;
  Tc := ∅;
  Sa := (π[0, 3] (σ[#1 = #2] (E × Tb)));
  Sb := (E ∪ (π[0, 3] (σ[#1 = #2] (E × Ta))));
  Sc := (Ta ∪ Tb);
  WHILE ¬(((Sa = Ta) ∧ ((Sb = Tb) ∧ (Sc = Tc)))) DO
    Ta := Sa;
    Tb := Sb;
    Tc := Sc;
    Sa := (π[0, 3] (σ[#1 = #2] (E × Tb)));
    Sb := (E ∪ (π[0, 3] (σ[#1 = #2] (E × Ta))));
    Sc := (Ta ∪ Tb)
  END

POST
  (((E ∪ (π[0, 3] (σ[#1 = #2] (E × Tc)))) ⊆ Tc) ∧ (Tc ⊆ R))
\end{Verbatim}

\Needspace{0.45\textheight}

\section{Example0117: Prefix-equivalence: naive all\_paths\_via\_even\_odd (non-linear) computes lfp of tc\_left}\label{case:Example0117}

\begin{description}[leftmargin=2.6cm,style=nextline,itemsep=1pt]

\item[Category] synthetic (syntactic generation from the classic Datalog programs); family: Synthetic.

\item[Kind] equivalence (with a least fixpoint, via a pre-fixpoint witness).

\item[Contributors] Fangzhu Shen, Leo Zhang.

\item[Status] certified valid (Certificate/Valid.lean, written by the certificate emitter); expected valid.

\item[Tags] \hyperref[tag:fixpoint-witness]{\texttt{fixpoint-witness}} \hyperref[tag:linear-vs-nonlinear]{\texttt{linear-vs-nonlinear}} \hyperref[tag:mutual-recursion]{\texttt{mutual-recursion}} \hyperref[tag:nonlinear-recursion]{\texttt{nonlinear-recursion}}.

\item[Input] \path{Benchmark/Example0117/Input.lean} (canonical notation).

\item[Summary] Canonical form of the earlier single-level encoding of this case: the same schema, precondition, command and postcondition.

\item[Sources] {\raggedright legacy source record v2\_\_prefix\_equivalence\_\_tc\_left\_\_all\_paths\_via\_even\_odd\_nonlinear\_\_naive of the retired classic-source collection\par}

\end{description}

\paragraph{Datalog program(s).}

{\raggedright\textit{all\_paths\_via\_even\_odd\_nonlinear} --- legacy inventory: \path{catalog/source} intent plus actual right-recursive formal PRE/POST reference-operator reconstruction; sources: all\_paths\_via\_even\_odd\_nonlinear.dl, tc-left.dl, \path{Benchmark/Example0117/Input.lean}.\par}

\begin{Verbatim}[fontsize=\scriptsize,breaklines=true,breakanywhere=true,frame=single,framesep=1.5mm]
-- all_paths_via_even_odd_nonlinear
-- Executed WHIEL program
OL(x,y) :- E(x,y).
EL(x,y) :- E(x,z), E(z,y).
OL(x,y) :- OL(x,z), EL(z,y).
OL(x,y) :- EL(x,z), OL(z,y).
EL(x,y) :- EL(x,z), EL(z,y).
EL(x,y) :- OL(x,z), OL(z,y).
T(x,y) :- EL(x,y).
T(x,y) :- OL(x,y).
\end{Verbatim}

{\raggedright\textit{tc\_left} --- legacy inventory: \path{catalog/source} intent plus actual right-recursive formal PRE/POST reference-operator reconstruction; sources: all\_paths\_via\_even\_odd\_nonlinear.dl, tc-left.dl, \path{Benchmark/Example0117/Input.lean}.\par}

\begin{Verbatim}[fontsize=\scriptsize,breaklines=true,breakanywhere=true,frame=single,framesep=1.5mm]
-- tc_left
-- Catalog/source target intent only; not executed by the WHIEL command
tcl(x, y) :- base(x, y).
tcl(x, y) :- tcl(x, z), base(z, y).
-- Actual formal PRE/POST reference only; not executed (right-recursive E composed with T)
T_right(x,y) :- E(x,y).
T_right(x,y) :- E(x,z), T_right(z,y).
\end{Verbatim}

\paragraph{Reading.} The loop naively evaluates the non-linear all\_paths\_via\_even\_odd program, where Ta (even-length paths) and Tb (odd-length paths) are each computed by squaring/composing Ta and Tb against each other and against E (Ta gets contributions from E∘E, Ta∘Ta and Tb∘Tb), and Tc:=Ta∪Tb collects all paths of length one or more. As in the related Example0115, although the title names tc\_left, the precondition and postcondition actually encode the right-linear closure T(x,y):-E(x,y), T(x,y):-E(x,z),T(z,y); the Hoare triple is the least-fixpoint introduction obligation stating that the naively computed union Tc is itself a pre-fixpoint of that right-linear program and lies below any given pre-fixpoint R, again because every nonempty path has even or odd length.

\paragraph{First-order reading of PRE/POST, translated mechanically from the relational-algebra assertion in Input.lean.} The relational-algebra assertion in \texttt{Input.lean} remains the authority.

\textit{PRE:} ∀x y. (E(x,y) ∨ (∃z. E(x,z) ∧ R(z,y)) → R(x,y))\par The precondition states that R is a pre-fixpoint of the right-linear closure program: E ∪ (E∘R) ⊆ R.

\textit{POST:} (∀x y. (E(x,y) ∨ (∃z. E(x,z) ∧ Tc(z,y)) → Tc(x,y))) ∧ ∀u v. (Tc(u,v) → R(u,v))\par The postcondition states that the computed union Tc is itself a pre-fixpoint of the right-linear closure program, E ∪ (E∘Tc) ⊆ Tc, and that Tc ⊆ R.

\paragraph{Whiel encoding (Hoare triple).}

\begin{Verbatim}[fontsize=\scriptsize,breaklines=true,breakanywhere=true,frame=single,framesep=1.5mm]
SCHEMA
  E/2, R/2, Ta/2, Sa/2, Tb/2, Sb/2, Tc/2, Sc/2

PRE
  ((E ∪ (π[0, 3] (σ[#1 = #2] (E × R)))) ⊆ R)

CMD
  Ta := ∅;
  Tb := ∅;
  Tc := ∅;
  Sa := (((π[0, 3] (σ[#1 = #2] (E × E))) ∪ (π[0, 3] (σ[#1 = #2] (Ta × Ta)))) ∪ (π[0, 3] (σ[#1 = #2] (Tb × Tb))));
  Sb := ((E ∪ (π[0, 3] (σ[#1 = #2] (Tb × Ta)))) ∪ (π[0, 3] (σ[#1 = #2] (Ta × Tb))));
  Sc := (Ta ∪ Tb);
  WHILE ¬(((Sa = Ta) ∧ ((Sb = Tb) ∧ (Sc = Tc)))) DO
    Ta := Sa;
    Tb := Sb;
    Tc := Sc;
    Sa := (((π[0, 3] (σ[#1 = #2] (E × E))) ∪ (π[0, 3] (σ[#1 = #2] (Ta × Ta)))) ∪ (π[0, 3] (σ[#1 = #2] (Tb × Tb))));
    Sb := ((E ∪ (π[0, 3] (σ[#1 = #2] (Tb × Ta)))) ∪ (π[0, 3] (σ[#1 = #2] (Ta × Tb))));
    Sc := (Ta ∪ Tb)
  END

POST
  (((E ∪ (π[0, 3] (σ[#1 = #2] (E × Tc)))) ⊆ Tc) ∧ (Tc ⊆ R))
\end{Verbatim}

\Needspace{0.45\textheight}

\section{Example0119: Prefix-equivalence: naive all\_paths\_via\_even\_odd (linear) computes lfp of tc\_nonlinear}\label{case:Example0119}

\begin{description}[leftmargin=2.6cm,style=nextline,itemsep=1pt]

\item[Category] synthetic (syntactic generation from the classic Datalog programs); family: Synthetic.

\item[Kind] equivalence (with a least fixpoint, via a pre-fixpoint witness).

\item[Contributors] Fangzhu Shen, Leo Zhang.

\item[Status] certified valid (Certificate/Valid.lean, written by the certificate emitter); expected valid.

\item[Tags] \hyperref[tag:fixpoint-witness]{\texttt{fixpoint-witness}} \hyperref[tag:linear-vs-nonlinear]{\texttt{linear-vs-nonlinear}} \hyperref[tag:mutual-recursion]{\texttt{mutual-recursion}} \hyperref[tag:nonlinear-recursion]{\texttt{nonlinear-recursion}} \hyperref[tag:prophecy]{\texttt{prophecy}}.

\item[Input] \path{Benchmark/Example0119/Input.lean} (canonical notation).

\item[Summary] All paths, computed through mutually recursive odd and even halves, specified as the least fixpoint of the non-linear closure program.

\item[Sources] {\raggedright legacy source record v2\_\_prefix\_equivalence\_\_tc\_nonlin\_\_all\_paths\_via\_even\_odd\_linear\_\_naive of the retired classic-source collection \par classic Datalog program all\_paths\_via\_even\_odd of the retired generator suite \par classic Datalog program tc-nonlinear of the retired generator suite\par}

\item[Prophecy] kernel-checked certificate with a level-1 prophecy clause (Certificate/Valid.lean, Core.json).

\end{description}

\paragraph{Datalog program(s).}

{\raggedright\textit{all\_paths\_via\_even\_odd\_linear} --- legacy inventory: source programs with one executed WHIEL program and one formal PRE/POST reference (not executed); sources: all\_paths\_via\_even\_odd.dl, tc-nonlinear.dl.\par}

\begin{Verbatim}[fontsize=\scriptsize,breaklines=true,breakanywhere=true,frame=single,framesep=1.5mm]
-- all_paths_via_even_odd_linear
-- Executed WHIEL program
OL(x,y) :- E(x,y).
OL(x,y) :- E(x,z), EL(z,y).
EL(x,y) :- E(x,z), OL(z,y).
T(x,y) :- EL(x,y).
T(x,y) :- OL(x,y).
\end{Verbatim}

{\raggedright\textit{tc\_nonlin} --- legacy inventory: source programs with one executed WHIEL program and one formal PRE/POST reference (not executed); sources: all\_paths\_via\_even\_odd.dl, tc-nonlinear.dl.\par}

\begin{Verbatim}[fontsize=\scriptsize,breaklines=true,breakanywhere=true,frame=single,framesep=1.5mm]
-- tc_nonlin
-- Formal PRE/POST reference only; not executed by the WHIEL command
tc(x, y) :- base(x, y).
tc(x, y) :- tc(x, z), tc(z, y).
\end{Verbatim}

\paragraph{Reading.} The loop naively evaluates the all\_paths\_via\_even\_odd program (Tb the odd-length paths via Tb:=E∪(E∘Ta), Ta the even-length paths via Ta:=E∘Tb, and Tc:=Ta∪Tb their union), matching the linear all\_paths\_via\_even\_odd.dl program. The Hoare triple is the least-fixpoint introduction obligation for the non-linear transitive-closure program T(x,y):-E(x,y), T(x,y):-T(x,z),T(z,y): given any R that is already a pre-fixpoint of that non-linear program, the computed Tc must itself be a pre-fixpoint of it and lie below R, so the even/odd decomposition computes the (non-linear) transitive closure, since every path has odd or even length.

\paragraph{First-order reading of PRE/POST, translated mechanically from the relational-algebra assertion in Input.lean.} The relational-algebra assertion in \texttt{Input.lean} remains the authority.

\textit{PRE:} ∀x y. (E(x,y) ∨ (∃z. R(x,z) ∧ R(z,y)) → R(x,y))\par The precondition states that R is a pre-fixpoint of the non-linear transitive-closure program: E ∪ (R∘R) ⊆ R.

\textit{POST:} (∀x y. (E(x,y) ∨ (∃z. Tc(x,z) ∧ Tc(z,y)) → Tc(x,y))) ∧ ∀u v. (Tc(u,v) → R(u,v))\par The postcondition states that the computed Tc is itself a pre-fixpoint of the non-linear transitive-closure program, E ∪ (Tc∘Tc) ⊆ Tc, and that Tc ⊆ R.

\paragraph{Whiel encoding (Hoare triple).}

\begin{Verbatim}[fontsize=\scriptsize,breaklines=true,breakanywhere=true,frame=single,framesep=1.5mm]
SCHEMA
  E/2, R/2, Ta/2, Sa/2, Tb/2, Sb/2, Tc/2, Sc/2

PRE
  ((E ∪ (π[0, 3] (σ[#1 = #2] (R × R)))) ⊆ R)

CMD
  Ta := ∅;
  Tb := ∅;
  Tc := ∅;
  Sa := (π[0, 3] (σ[#1 = #2] (E × Tb)));
  Sb := (E ∪ (π[0, 3] (σ[#1 = #2] (E × Ta))));
  Sc := (Ta ∪ Tb);
  WHILE
      ¬(((Sa = Ta) ∧
        ((Sb = Tb) ∧ (Sc = Tc)))) DO
    Ta := Sa;
    Tb := Sb;
    Tc := Sc;
    Sa := (π[0, 3] (σ[#1 = #2] (E × Tb)));
    Sb :=
      (E ∪ (π[0, 3] (σ[#1 = #2] (E × Ta))));
    Sc := (Ta ∪ Tb)
  END

POST
  (((E ∪ (π[0, 3] (σ[#1 = #2] (Tc × Tc)))) ⊆ Tc)
    ∧ (Tc ⊆ R))
\end{Verbatim}

\Needspace{0.45\textheight}

\section{Example0123: Prefix-equivalence: naive even\_odd\_nonlinear computes lfp of even\_odd\_linear}\label{case:Example0123}

\begin{description}[leftmargin=2.6cm,style=nextline,itemsep=1pt]

\item[Category] synthetic (syntactic generation from the classic Datalog programs); family: Synthetic.

\item[Kind] equivalence (with a least fixpoint, via a pre-fixpoint witness).

\item[Contributors] Fangzhu Shen, Leo Zhang.

\item[Status] certified valid (Certificate/Valid.lean, written by the certificate emitter); expected valid.

\item[Tags] \hyperref[tag:fixpoint-witness]{\texttt{fixpoint-witness}} \hyperref[tag:linear-vs-nonlinear]{\texttt{linear-vs-nonlinear}} \hyperref[tag:mutual-recursion]{\texttt{mutual-recursion}} \hyperref[tag:nonlinear-recursion]{\texttt{nonlinear-recursion}}.

\item[Input] \path{Benchmark/Example0123/Input.lean} (canonical notation).

\item[Summary] Canonical form of the earlier single-level encoding of this case: the same schema, precondition, command and postcondition.

\item[Sources] {\raggedright legacy source record v2\_\_prefix\_equivalence\_\_even\_odd\_linear\_\_even\_odd\_nonlinear\_\_naive of the retired classic-source collection \par classic Datalog program even\_odd\_linear of the retired generator suite \par classic Datalog program even\_odd\_nonlinear of the retired generator suite\par}

\end{description}

\paragraph{Datalog program(s).}

{\raggedright\textit{even\_odd\_linear} --- legacy inventory: source programs with one executed WHIEL program and one formal PRE/POST reference (not executed); sources: even\_odd\_linear.dl, even\_odd\_nonlinear.dl.\par}

\begin{Verbatim}[fontsize=\scriptsize,breaklines=true,breakanywhere=true,frame=single,framesep=1.5mm]
-- even_odd_linear
-- Formal PRE/POST reference only; not executed by the WHIEL command
OL(x,y) :- E(x,y).
OL(x,y) :- E(x,z), EL(z,y).
EL(x,y) :- E(x,z), OL(z,y).
\end{Verbatim}

{\raggedright\textit{even\_odd\_nonlinear} --- legacy inventory: source programs with one executed WHIEL program and one formal PRE/POST reference (not executed); sources: even\_odd\_linear.dl, even\_odd\_nonlinear.dl.\par}

\begin{Verbatim}[fontsize=\scriptsize,breaklines=true,breakanywhere=true,frame=single,framesep=1.5mm]
-- even_odd_nonlinear
-- Executed WHIEL program
OL(x,y) :- E(x,y).
EL(x,y) :- E(x,z), E(z,y).
OL(x,y) :- OL(x,z), EL(z,y).
OL(x,y) :- EL(x,z), OL(z,y).
EL(x,y) :- EL(x,z), EL(z,y).
EL(x,y) :- OL(x,z), OL(z,y).
\end{Verbatim}

\paragraph{Reading.} The loop naively evaluates the non-linear even\_odd\_nonlinear program, computing even-length paths Ta and odd-length paths Tb by mutual squaring recursion (Ta from E∘E, Ta∘Ta and Tb∘Tb; Tb from E, Tb∘Ta and Ta∘Tb), matching even\_odd\_nonlinear.dl. The Hoare triple is the least-fixpoint introduction obligation for the linear even\_odd\_linear recursion (Ta(x,y):-E(x,z),Tb(z,y); Tb(x,y):-E(x,y) or E(x,z),Ta(z,y)): given any pair Ra,Rb already a pre-fixpoint of that linear recursion, the computed Ta,Tb must themselves form a pre-fixpoint of it and lie below Ra,Rb respectively, i.e. the non-linear program computes the least pre-fixpoint pair of the linear even/odd program.

\paragraph{First-order reading of PRE/POST, translated mechanically from the relational-algebra assertion in Input.lean.} The relational-algebra assertion in \texttt{Input.lean} remains the authority.

\textit{PRE:} (∀x y. ((∃z. E(x,z) ∧ Rb(z,y)) → Ra(x,y))) ∧ ∀u v. (E(u,v) ∨ (∃w. E(u,w) ∧ Ra(w,v)) → Rb(u,v))\par The precondition states that the pair Ra, Rb is a pre-fixpoint of the linear even/odd recursion: E∘Rb ⊆ Ra and E∪(E∘Ra) ⊆ Rb.

\textit{POST:} (∀x y. ((∃z. E(x,z) ∧ Tb(z,y)) → Ta(x,y))) ∧ (∀u v. (E(u,v) ∨ (∃w. E(u,w) ∧ Ta(w,v)) → Tb(u,v))) ∧ (∀x1 x2. (Ta(x1,x2) → Ra(x1,x2))) ∧ ∀x3 x4. (Tb(x3,x4) → Rb(x3,x4))\par The postcondition states that the computed pair Ta, Tb is itself a pre-fixpoint of the linear even/odd recursion, E∘Tb ⊆ Ta and E∪(E∘Ta) ⊆ Tb, and that Ta ⊆ Ra and Tb ⊆ Rb.

\paragraph{Whiel encoding (Hoare triple).}

\begin{Verbatim}[fontsize=\scriptsize,breaklines=true,breakanywhere=true,frame=single,framesep=1.5mm]
SCHEMA
  E/2, Ra/2, Rb/2, Ta/2, Sa/2, Tb/2, Sb/2

PRE
  (((π[0, 3] (σ[#1 = #2] (E × Rb))) ⊆ Ra) ∧ ((E ∪ (π[0, 3] (σ[#1 = #2] (E × Ra)))) ⊆ Rb))

CMD
  Ta := ∅;
  Tb := ∅;
  Sa := (((π[0, 3] (σ[#1 = #2] (E × E))) ∪ (π[0, 3] (σ[#1 = #2] (Ta × Ta)))) ∪ (π[0, 3] (σ[#1 = #2] (Tb × Tb))));
  Sb := ((E ∪ (π[0, 3] (σ[#1 = #2] (Tb × Ta)))) ∪ (π[0, 3] (σ[#1 = #2] (Ta × Tb))));
  WHILE ¬(((Sa = Ta) ∧ (Sb = Tb))) DO
    Ta := Sa;
    Tb := Sb;
    Sa := (((π[0, 3] (σ[#1 = #2] (E × E))) ∪ (π[0, 3] (σ[#1 = #2] (Ta × Ta)))) ∪ (π[0, 3] (σ[#1 = #2] (Tb × Tb))));
    Sb := ((E ∪ (π[0, 3] (σ[#1 = #2] (Tb × Ta)))) ∪ (π[0, 3] (σ[#1 = #2] (Ta × Tb))))
  END

POST
  (((π[0, 3] (σ[#1 = #2] (E × Tb))) ⊆ Ta) ∧ (((E ∪ (π[0, 3] (σ[#1 = #2] (E × Ta)))) ⊆ Tb) ∧ ((Ta ⊆ Ra) ∧ (Tb ⊆ Rb))))
\end{Verbatim}

\Needspace{0.45\textheight}

\section{Example0125: Prefix-equivalence: naive even\_odd\_linear computes lfp of even\_odd\_nonlinear}\label{case:Example0125}

\begin{description}[leftmargin=2.6cm,style=nextline,itemsep=1pt]

\item[Category] synthetic (syntactic generation from the classic Datalog programs); family: Synthetic.

\item[Kind] equivalence (with a least fixpoint, via a pre-fixpoint witness).

\item[Contributors] Fangzhu Shen, Leo Zhang.

\item[Status] certified valid (Certificate/Valid.lean, written by the certificate emitter); expected valid.

\item[Tags] \hyperref[tag:fixpoint-witness]{\texttt{fixpoint-witness}} \hyperref[tag:linear-vs-nonlinear]{\texttt{linear-vs-nonlinear}} \hyperref[tag:mutual-recursion]{\texttt{mutual-recursion}} \hyperref[tag:nonlinear-recursion]{\texttt{nonlinear-recursion}} \hyperref[tag:prophecy]{\texttt{prophecy}}.

\item[Input] \path{Benchmark/Example0125/Input.lean} (canonical notation).

\item[Summary] Odd and even paths computed by a linear program, specified as the least fixpoint of the non-linear odd/even program.

\item[Sources] {\raggedright legacy source record v2\_\_prefix\_equivalence\_\_even\_odd\_nonlinear\_\_even\_odd\_linear\_\_naive of the retired classic-source collection \par classic Datalog program even\_odd\_linear of the retired generator suite \par classic Datalog program even\_odd\_nonlinear of the retired generator suite\par}

\item[Prophecy] kernel-checked certificate with a level-1 prophecy clause (Certificate/Valid.lean, Core.json).

\end{description}

\paragraph{Datalog program(s).}

{\raggedright\textit{even\_odd\_linear} --- legacy inventory: source programs with one executed WHIEL program and one formal PRE/POST reference (not executed); sources: even\_odd\_linear.dl, even\_odd\_nonlinear.dl.\par}

\begin{Verbatim}[fontsize=\scriptsize,breaklines=true,breakanywhere=true,frame=single,framesep=1.5mm]
-- even_odd_linear
-- Executed WHIEL program
OL(x,y) :- E(x,y).
OL(x,y) :- E(x,z), EL(z,y).
EL(x,y) :- E(x,z), OL(z,y).
\end{Verbatim}

{\raggedright\textit{even\_odd\_nonlinear} --- legacy inventory: source programs with one executed WHIEL program and one formal PRE/POST reference (not executed); sources: even\_odd\_linear.dl, even\_odd\_nonlinear.dl.\par}

\begin{Verbatim}[fontsize=\scriptsize,breaklines=true,breakanywhere=true,frame=single,framesep=1.5mm]
-- even_odd_nonlinear
-- Formal PRE/POST reference only; not executed by the WHIEL command
OL(x,y) :- E(x,y).
EL(x,y) :- E(x,z), E(z,y).
OL(x,y) :- OL(x,z), EL(z,y).
OL(x,y) :- EL(x,z), OL(z,y).
EL(x,y) :- EL(x,z), EL(z,y).
EL(x,y) :- OL(x,z), OL(z,y).
\end{Verbatim}

\paragraph{Reading.} The loop naively evaluates the linear even\_odd\_linear program, computing even-length paths Ta and odd-length paths Tb by the mutual recursion Ta(x,y):-E(x,z),Tb(z,y) and Tb(x,y):-E(x,y) or E(x,z),Ta(z,y). The Hoare triple is the least-fixpoint introduction obligation for the non-linear even\_odd\_nonlinear recursion, which computes the same even/odd relations by squaring (e.g. Ta from E∘E, Ta∘Ta and Tb∘Tb): given any pair Ra,Rb already a pre-fixpoint of that non-linear program, the computed Ta,Tb must themselves form a pre-fixpoint of it and lie below Ra,Rb respectively, so the linear program computes the least pre-fixpoint pair of the non-linear one.

\paragraph{First-order reading of PRE/POST, translated mechanically from the relational-algebra assertion in Input.lean.} The relational-algebra assertion in \texttt{Input.lean} remains the authority.

\textit{PRE:} (∀x y. ((∃z. E(x,z) ∧ E(z,y)) ∨ (∃u. Ra(x,u) ∧ Ra(u,y)) ∨ (∃v. Rb(x,v) ∧ Rb(v,y)) → Ra(x,y))) ∧ ∀w x1. (E(w,x1) ∨ (∃x2. Rb(w,x2) ∧ Ra(x2,x1)) ∨ (∃x3. Ra(w,x3) ∧ Rb(x3,x1)) → Rb(w,x1))\par The precondition states that the pair Ra, Rb is a pre-fixpoint of the non-linear even/odd recursion: (E∘E) ∪ (Ra∘Ra) ∪ (Rb∘Rb) ⊆ Ra, and E ∪ (Rb∘Ra) ∪ (Ra∘Rb) ⊆ Rb.

\textit{POST:} (∀x y. ((∃z. E(x,z) ∧ E(z,y)) ∨ (∃u. Ta(x,u) ∧ Ta(u,y)) ∨ (∃v. Tb(x,v) ∧ Tb(v,y)) → Ta(x,y))) ∧ (∀w x1. (E(w,x1) ∨ (∃x2. Tb(w,x2) ∧ Ta(x2,x1)) ∨ (∃x3. Ta(w,x3) ∧ Tb(x3,x1)) → Tb(w,x1))) ∧ (∀x4 x5. (Ta(x4,x5) → Ra(x4,x5))) ∧ ∀x6 x7. (Tb(x6,x7) → Rb(x6,x7))\par The postcondition states that the computed pair Ta, Tb is itself a pre-fixpoint of the non-linear even/odd recursion, (E∘E)∪(Ta∘Ta)∪(Tb∘Tb) ⊆ Ta and E∪(Tb∘Ta)∪(Ta∘Tb) ⊆ Tb, and that Ta ⊆ Ra and Tb ⊆ Rb.

\paragraph{Whiel encoding (Hoare triple).}

\begin{Verbatim}[fontsize=\scriptsize,breaklines=true,breakanywhere=true,frame=single,framesep=1.5mm]
SCHEMA
  E/2, Ra/2, Rb/2, Ta/2, Sa/2, Tb/2, Sb/2

PRE
  (((((π[0, 3] (σ[#1 = #2] (E × E))) ∪
        (π[0, 3] (σ[#1 = #2] (Ra × Ra)))) ∪
      (π[0, 3] (σ[#1 = #2] (Rb × Rb)))) ⊆ Ra)
    ∧ (((E ∪
          (π[0, 3] (σ[#1 = #2] (Rb × Ra)))) ∪
        (π[0, 3] (σ[#1 = #2] (Ra × Rb)))) ⊆ Rb))

CMD
  Ta := ∅;
  Tb := ∅;
  Sa := (π[0, 3] (σ[#1 = #2] (E × Tb)));
  Sb := (E ∪ (π[0, 3] (σ[#1 = #2] (E × Ta))));
  WHILE ¬(((Sa = Ta) ∧ (Sb = Tb))) DO
    Ta := Sa;
    Tb := Sb;
    Sa := (π[0, 3] (σ[#1 = #2] (E × Tb)));
    Sb :=
      (E ∪ (π[0, 3] (σ[#1 = #2] (E × Ta))))
  END

POST
  (((((π[0, 3] (σ[#1 = #2] (E × E))) ∪
        (π[0, 3] (σ[#1 = #2] (Ta × Ta)))) ∪
      (π[0, 3] (σ[#1 = #2] (Tb × Tb)))) ⊆ Ta)
    ∧ ((((E ∪
            (π[0, 3] (σ[#1 = #2] (Tb × Ta)))) ∪
          (π[0, 3] (σ[#1 = #2] (Ta × Tb)))) ⊆ Tb)
      ∧ ((Ta ⊆ Ra) ∧ (Tb ⊆ Rb))))
\end{Verbatim}

\Needspace{0.45\textheight}

\section{Example0128: Containment: naive red\_blue\_nonlinear outputs ⊆ pre-fixpoint of TC(R ∪ B)}\label{case:Example0128}

\begin{description}[leftmargin=2.6cm,style=nextline,itemsep=1pt]

\item[Category] synthetic (syntactic generation from the classic Datalog programs); family: Synthetic.

\item[Kind] containment (via a pre-fixpoint witness).

\item[Contributors] Fangzhu Shen, Leo Zhang.

\item[Status] certified valid (Certificate/Valid.lean, written by the certificate emitter); expected valid.

\item[Tags] \hyperref[tag:fixpoint-witness]{\texttt{fixpoint-witness}} \hyperref[tag:mutual-recursion]{\texttt{mutual-recursion}} \hyperref[tag:nonlinear-recursion]{\texttt{nonlinear-recursion}}.

\item[Input] \path{Benchmark/Example0128/Input.lean} (canonical notation).

\item[Summary] Canonical form of the earlier single-level encoding of this case: the same schema, precondition, command and postcondition.

\item[Sources] {\raggedright legacy source record v2\_\_containment\_\_tc\_of\_union\_\_red\_blue\_nonlinear\_\_naive of the retired classic-source collection \par classic Datalog program red\_blue\_nonlinear of the retired generator suite \par classic Datalog program tc\_of\_union of the retired generator suite\par}

\end{description}

\paragraph{Datalog program(s).}

{\raggedright\textit{red\_blue\_nonlinear} --- legacy inventory: source programs with one executed WHIEL program and one formal PRE/POST reference (not executed); sources: red\_blue\_nonlinear.dl, tc\_of\_union.dl.\par}

\begin{Verbatim}[fontsize=\scriptsize,breaklines=true,breakanywhere=true,frame=single,framesep=1.5mm]
-- red_blue_nonlinear
-- Executed WHIEL program
RR(x,y) :- R(x,y).
BB(x,y) :- B(x,y).
RB(x,y) :- R(x,z), B(z,y).
BR(x,y) :- B(x,z), R(z,y).
RR(x,y) :- RR(x,z), BR(z,y).
BB(x,y) :- BB(x,z), RB(z,y).
RB(x,y) :- RR(x,z), BB(z,y).
BR(x,y) :- BB(x,z), RR(z,y).
\end{Verbatim}

{\raggedright\textit{tc\_of\_union} --- legacy inventory: source programs with one executed WHIEL program and one formal PRE/POST reference (not executed); sources: red\_blue\_nonlinear.dl, tc\_of\_union.dl.\par}

\begin{Verbatim}[fontsize=\scriptsize,breaklines=true,breakanywhere=true,frame=single,framesep=1.5mm]
-- tc_of_union
-- Formal PRE/POST reference only; not executed by the WHIEL command
E(x,y) :- R(x,y).
E(x,y) :- B(x,y).
T(x,y) :- E(x,y).
T(x,y) :- E(x,z), T(z,y).
\end{Verbatim}

\paragraph{Reading.} The loop naively evaluates the non-linear red\_blue\_nonlinear program, computing four colour-alternating path relations (same-colour-start relations analogous to RR/BB and colour-switching relations analogous to RB/BR) by mutual recursion over a red edge relation RX and a blue edge relation B, matching red\_blue\_nonlinear.dl. The Hoare triple is a containment claim: it asserts that each of the four computed relations is contained in any relation R that is already a pre-fixpoint of the (uncoloured) transitive-closure program on RX∪B, since every colour-alternating walk is in particular an ordinary walk in the union graph.

\paragraph{First-order reading of PRE/POST, translated mechanically from the relational-algebra assertion in Input.lean.} The relational-algebra assertion in \texttt{Input.lean} remains the authority.

\textit{PRE:} ∀x y. (RX(x,y) ∨ B(x,y) ∨ (∃z. RX(x,z) ∧ R(z,y)) ∨ (∃u. B(x,u) ∧ R(u,y)) → R(x,y))\par The precondition states that R is a pre-fixpoint of the transitive-closure program on the union graph RX∪B: (RX∪B) ∪ (RX∘R) ∪ (B∘R) ⊆ R.

\textit{POST:} (∀x y. (Ta(x,y) → R(x,y))) ∧ (∀z u. (Tb(z,u) → R(z,u))) ∧ (∀v w. (Tc(v,w) → R(v,w))) ∧ ∀x1 x2. (Td(x1,x2) → R(x1,x2))\par The postcondition states that all four computed relations Ta, Tb, Tc, Td are each contained in R.

\paragraph{Whiel encoding (Hoare triple).}

\begin{Verbatim}[fontsize=\scriptsize,breaklines=true,breakanywhere=true,frame=single,framesep=1.5mm]
SCHEMA
  B/2, RX/2, R/2, Ta/2, Sa/2, Tb/2, Sb/2, Tc/2, Sc/2, Td/2, Sd/2

PRE
  ((((RX ∪ B) ∪ (π[0, 3] (σ[#1 = #2] (RX × R)))) ∪ (π[0, 3] (σ[#1 = #2] (B × R)))) ⊆ R)

CMD
  Ta := ∅;
  Tb := ∅;
  Tc := ∅;
  Td := ∅;
  Sa := (B ∪ (π[0, 3] (σ[#1 = #2] (Ta × Tc))));
  Sb := ((π[0, 3] (σ[#1 = #2] (B × RX))) ∪ (π[0, 3] (σ[#1 = #2] (Ta × Td))));
  Sc := ((π[0, 3] (σ[#1 = #2] (RX × B))) ∪ (π[0, 3] (σ[#1 = #2] (Td × Ta))));
  Sd := (RX ∪ (π[0, 3] (σ[#1 = #2] (Td × Tb))));
  WHILE ¬(((Sa = Ta) ∧ ((Sb = Tb) ∧ ((Sc = Tc) ∧ (Sd = Td))))) DO
    Ta := Sa;
    Tb := Sb;
    Tc := Sc;
    Td := Sd;
    Sa := (B ∪ (π[0, 3] (σ[#1 = #2] (Ta × Tc))));
    Sb := ((π[0, 3] (σ[#1 = #2] (B × RX))) ∪ (π[0, 3] (σ[#1 = #2] (Ta × Td))));
    Sc := ((π[0, 3] (σ[#1 = #2] (RX × B))) ∪ (π[0, 3] (σ[#1 = #2] (Td × Ta))));
    Sd := (RX ∪ (π[0, 3] (σ[#1 = #2] (Td × Tb))))
  END

POST
  ((Ta ⊆ R) ∧ ((Tb ⊆ R) ∧ ((Tc ⊆ R) ∧ (Td ⊆ R))))
\end{Verbatim}

\Needspace{0.45\textheight}

\section{Example0130: Containment: even/odd/ambiguous-depth path-system program ⊆ every set closed under the basic path-system rules}\label{case:Example0130}

\begin{description}[leftmargin=2.6cm,style=nextline,itemsep=1pt]

\item[Category] synthetic (syntactic generation from the classic Datalog programs); family: Synthetic.

\item[Kind] containment.

\item[Contributors] Fangzhu Shen, Leo Zhang.

\item[Status] certified valid (Certificate/Valid.lean, written by the certificate emitter); expected valid.

\item[Tags] \hyperref[tag:fixpoint-witness]{\texttt{fixpoint-witness}} \hyperref[tag:mutual-recursion]{\texttt{mutual-recursion}} \hyperref[tag:nonlinear-recursion]{\texttt{nonlinear-recursion}}.

\item[Input] \path{Benchmark/Example0130/Input.lean} (canonical notation).

\item[Summary] Path systems (Cook): unary axioms A and ternary rules R(x,y,z), read as "if y and z are accessible then x is accessible". The basic program T(x) :- A(x). T(x) :- R(x,y,z), T(y), T(z). computes the accessible nodes T.

\item[Sources] {\raggedright legacy source record v2\_\_containment\_\_path\_systems\_basic\_\_path\_systems\_height\_\_naive of the retired classic-source collection \par classic Datalog program path\_systems\_basic of the retired generator suite \par classic Datalog program path\_systems\_height of the retired generator suite\par}

\end{description}

\paragraph{Datalog program(s).}

{\raggedright\textit{rules as stated in the header comment of Input.lean} --- header comment.\par}

\begin{Verbatim}[fontsize=\scriptsize,breaklines=true,breakanywhere=true,frame=single,framesep=1.5mm]
T(x) :- A(x).
T(x) :- R(x,y,z), T(y), T(z).
EH(x) :- A(x). x has a derivation whose leaves all lie at even depth
OH(x) :- R(x,y,z), EH(y), EH(z). ... all at odd depth
EH(x) :- R(x,y,z), OH(y), OH(z).
AH(x) :- R(x,y,z), Y(y), Z(z) for every other pair (Y,Z) over EH/OH/AH:
\end{Verbatim}

\paragraph{Reading.} The Datalog programs classify the nodes a Cook path system (unary axioms A, ternary rules R(x,y,z) meaning x is accessible if y and z both are) makes accessible, by the parity of their derivation trees' leaf depths: EH, OH and AH hold for a node with a derivation whose leaves are all at even depth, all at odd depth, or mixed, respectively. The precondition says a free relation Closed already contains every axiom and is closed under the path-system rule (a pre-fixpoint); the postcondition says the three depth-parity outputs are each contained in Closed. Since Closed ranges over every such closed set and the accessible-node relation is the least one, this is the standard least-fixpoint-containment argument, applied to three refinements of the accessible set at once.

\paragraph{First-order reading of PRE/POST, translated mechanically from the relational-algebra assertion in Input.lean.} The relational-algebra assertion in \texttt{Input.lean} remains the authority.

\textit{PRE:} ∀x. (A(x) ∨ (∃y z. R(x,y,z) ∧ Closed(y) ∧ Closed(z)) → Closed(x))\par PRE says Closed contains every axiom A(x) and is closed under the path-system rule: whenever R(x,y,z) holds and y and z are both already in Closed, x is too.

\textit{POST:} (∀x. (AH(x) → Closed(x))) ∧ (∀y. (EH(y) → Closed(y))) ∧ ∀z. (OH(z) → Closed(z))\par POST says the even-depth, odd-depth and mixed-depth output relations AH, EH and OH are each contained in Closed.

\paragraph{Whiel encoding (Hoare triple).}

\begin{Verbatim}[fontsize=\tiny,breaklines=true,breakanywhere=true,frame=single,framesep=1.5mm]
SCHEMA
  A/1, Closed/1, AH/1, EH/1, OH/1, AHNext/1, EHNext/1, OHNext/1, R/3

PRE
  ((A ∪ (π[0] (σ[#2 = #4] (σ[#1 = #3] ((R × Closed) × Closed))))) ⊆ Closed)

CMD
  AH := ∅;
  EH := ∅;
  OH := ∅;
  AHNext :=
    (((((((π[0] (σ[#2 = #4] (σ[#1 = #3] ((R × EH) × OH))))   -- AH(x) :- R(x,y,z), EH(y), OH(z)
    ∪ (π[0] (σ[#2 = #4] (σ[#1 = #3] ((R × OH) × EH)))))   -- AH(x) :- R(x,y,z), OH(y), EH(z)
    ∪ (π[0] (σ[#2 = #4] (σ[#1 = #3] ((R × AH) × OH)))))   -- AH(x) :- R(x,y,z), AH(y), OH(z)
    ∪ (π[0] (σ[#2 = #4] (σ[#1 = #3] ((R × OH) × AH)))))   -- AH(x) :- R(x,y,z), OH(y), AH(z)
    ∪ (π[0] (σ[#2 = #4] (σ[#1 = #3] ((R × AH) × EH)))))   -- AH(x) :- R(x,y,z), AH(y), EH(z)
    ∪ (π[0] (σ[#2 = #4] (σ[#1 = #3] ((R × EH) × AH)))))   -- AH(x) :- R(x,y,z), EH(y), AH(z)
    ∪ (π[0] (σ[#2 = #4] (σ[#1 = #3] ((R × AH) × AH)))));   -- AH(x) :- R(x,y,z), AH(y), AH(z)
  EHNext := (A ∪ (π[0] (σ[#2 = #4] (σ[#1 = #3] ((R × OH) × OH)))));   -- EH(x) :- A(x).   EH(x) :- R(x,y,z), OH(y), OH(z)
  OHNext := (π[0] (σ[#2 = #4] (σ[#1 = #3] ((R × EH) × EH))));   -- OH(x) :- R(x,y,z), EH(y), EH(z)
  WHILE ¬(((AHNext = AH) ∧ ((EHNext = EH) ∧ (OHNext = OH)))) DO
    AH := AHNext;
    EH := EHNext;
    OH := OHNext;
    AHNext :=
      (((((((π[0] (σ[#2 = #4] (σ[#1 = #3] ((R × EH) × OH))))   -- AH(x) :- R(x,y,z), EH(y), OH(z)
      ∪ (π[0] (σ[#2 = #4] (σ[#1 = #3] ((R × OH) × EH)))))   -- AH(x) :- R(x,y,z), OH(y), EH(z)
      ∪ (π[0] (σ[#2 = #4] (σ[#1 = #3] ((R × AH) × OH)))))   -- AH(x) :- R(x,y,z), AH(y), OH(z)
      ∪ (π[0] (σ[#2 = #4] (σ[#1 = #3] ((R × OH) × AH)))))   -- AH(x) :- R(x,y,z), OH(y), AH(z)
      ∪ (π[0] (σ[#2 = #4] (σ[#1 = #3] ((R × AH) × EH)))))   -- AH(x) :- R(x,y,z), AH(y), EH(z)
      ∪ (π[0] (σ[#2 = #4] (σ[#1 = #3] ((R × EH) × AH)))))   -- AH(x) :- R(x,y,z), EH(y), AH(z)
      ∪ (π[0] (σ[#2 = #4] (σ[#1 = #3] ((R × AH) × AH)))));   -- AH(x) :- R(x,y,z), AH(y), AH(z)
    EHNext := (A ∪ (π[0] (σ[#2 = #4] (σ[#1 = #3] ((R × OH) × OH)))));   -- EH(x) :- A(x).   EH(x) :- R(x,y,z), OH(y), OH(z)
    OHNext := (π[0] (σ[#2 = #4] (σ[#1 = #3] ((R × EH) × EH))))   -- OH(x) :- R(x,y,z), EH(y), EH(z)
  END

POST
  ((AH ⊆ Closed) ∧ ((EH ⊆ Closed) ∧ (OH ⊆ Closed)))
\end{Verbatim}

\Needspace{0.45\textheight}

\section{Example0132: Odd and even paths: linear versus non-linear program (two compiled loops)}\label{case:Example0132}

\begin{description}[leftmargin=2.6cm,style=nextline,itemsep=1pt]

\item[Category] synthetic (syntactic generation from the classic Datalog programs); family: Obstructed.

\item[Kind] equivalence.

\item[Contributors] Fangzhu Shen.

\item[Status] certified valid (Certificate/Valid.lean, written by the certificate emitter); expected valid.

\item[Tags] \hyperref[tag:linear-vs-nonlinear]{\texttt{linear-vs-nonlinear}} \hyperref[tag:mutual-recursion]{\texttt{mutual-recursion}} \hyperref[tag:nonlinear-recursion]{\texttt{nonlinear-recursion}} \hyperref[tag:prophecy]{\texttt{prophecy}}.

\item[Input] \path{Benchmark/Example0132/Input.lean} (canonical notation).

\item[Summary] Odd and even paths, computed twice: by a linear program and by a non-linear one.

\item[Sources] {\raggedright legacy source record v2\_\_equiv\_pairs\_\_even\_odd\_\_linear\_vs\_nonlinear\_\_naive of the retired classic-source collection \par classic Datalog program even\_odd\_linear of the retired generator suite \par classic Datalog program even\_odd\_nonlinear of the retired generator suite\par}

\item[Prophecy] kernel-checked certificate with a level-1 prophecy clause (Certificate/Valid.lean, Core.json).

\end{description}

\paragraph{Datalog program(s).}

{\raggedright\textit{programX --- Odd and even paths, linearly.} --- fidelity module (kernel-checked: the compiled sides of the command are this program's naive compilation).\par}

\begin{Verbatim}[fontsize=\scriptsize,breaklines=true,breakanywhere=true,frame=single,framesep=1.5mm]
TbX(x1, y1) :- E(x1, y1).
TbX(x1, y1) :- E(x1, z1), TaX(z1, y1).
TaX(x1, y1) :- E(x1, z1), TbX(z1, y1).
\end{Verbatim}

{\raggedright\textit{programY --- Odd and even paths, non-linearly.} --- fidelity module (kernel-checked: the compiled sides of the command are this program's naive compilation).\par}

\begin{Verbatim}[fontsize=\scriptsize,breaklines=true,breakanywhere=true,frame=single,framesep=1.5mm]
TaY(x1, y1) :- E(x1, z1), E(z1, y1).
TaY(x1, y1) :- TaY(x1, z1), TaY(z1, y1).
TaY(x1, y1) :- TbY(x1, z1), TbY(z1, y1).
TbY(x1, y1) :- E(x1, y1).
TbY(x1, y1) :- TbY(x1, z1), TaY(z1, y1).
TbY(x1, y1) :- TaY(x1, z1), TbY(z1, y1).
\end{Verbatim}

\paragraph{Reading.} The command runs two independently compiled Datalog programs in sequence --- a linear even/odd program producing TaX (even-length paths) and TbX (odd-length paths) via naive accumulation, and a non-linear even/odd program producing TaY and TbY via squaring --- and the tool's preprocessor merges the two into a single verified loop. The postcondition states that the two computations agree relation by relation, TaX=TaY and TbX=TbY, i.e. the Hoare triple is an equivalence claim between the linear and non-linear encodings of the even/odd path decomposition, both of which converge to the same least fixpoints since the underlying immediate-consequence operators are monotone.

\paragraph{First-order reading of PRE/POST, translated mechanically from the relational-algebra assertion in Input.lean.} The relational-algebra assertion in \texttt{Input.lean} remains the authority.

\textit{PRE:} ⊤\par The precondition is simply true, imposing no constraint on the input relations.

\textit{POST:} (∀x y. (TaX(x,y) ↔ TaY(x,y))) ∧ ∀z u. (TbX(z,u) ↔ TbY(z,u))\par The postcondition asserts that the linear program's outputs equal the non-linear program's outputs: TaX = TaY and TbX = TbY.

\paragraph{Whiel encoding (Hoare triple).}

\begin{Verbatim}[fontsize=\tiny,breaklines=true,breakanywhere=true,frame=single,framesep=1.5mm]
SCHEMA
  E/2, TaX/2, TbX/2, TaY/2, TbY/2, TaX_aux/2, TbX_aux/2, TaY_aux/2, TbY_aux/2

PRE
  true 

CMD
  TbX_aux := ∅[2];
  TaX_aux := ∅[2];
  TbX :=
    (E ∪ (π[0, 3] (σ[#1 = #2] (E × TaX_aux))));
  TaX := (π[0, 3] (σ[#1 = #2] (E × TbX_aux)));
  WHILE
      (¬ ((TbX = TbX_aux) ∧ (TaX = TaX_aux))) DO
    TbX_aux := TbX;
    TaX_aux := TaX;
    TbX :=
      (TbX ∪
        (E ∪
          (π[0, 3]
            (σ[#1 = #2] (E × TaX_aux)))));
    TaX :=
      (TaX ∪
        (π[0, 3] (σ[#1 = #2] (E × TbX_aux))))
  END;
  TaY_aux := ∅[2];
  TbY_aux := ∅[2];
  TaY :=
    ((π[0, 3] (σ[#1 = #2] (E × E))) ∪
      ((π[0, 3]
          (σ[#1 = #2] (TaY_aux × TaY_aux))) ∪
        (π[0, 3]
          (σ[#1 = #2] (TbY_aux × TbY_aux)))));
  TbY :=
    (E ∪
      ((π[0, 3]
          (σ[#1 = #2] (TbY_aux × TaY_aux))) ∪
        (π[0, 3]
          (σ[#1 = #2] (TaY_aux × TbY_aux)))));
  WHILE
      (¬ ((TaY = TaY_aux) ∧ (TbY = TbY_aux))) DO
    TaY_aux := TaY;
    TbY_aux := TbY;
    TaY :=
      (TaY ∪
        ((π[0, 3] (σ[#1 = #2] (E × E))) ∪
          ((π[0, 3]
              (σ[#1 = #2]
                (TaY_aux × TaY_aux))) ∪
            (π[0, 3]
              (σ[#1 = #2]
                (TbY_aux × TbY_aux))))));
    TbY :=
      (TbY ∪
        (E ∪
          ((π[0, 3]
              (σ[#1 = #2]
                (TbY_aux × TaY_aux))) ∪
            (π[0, 3]
              (σ[#1 = #2]
                (TaY_aux × TbY_aux))))))
  END

POST
  ((TaX = TaY) ∧ (TbX = TbY))
\end{Verbatim}

\Needspace{0.45\textheight}

\section{Example0133: All paths via odd/even paths: linear versus non-linear program (two compiled loops)}\label{case:Example0133}

\begin{description}[leftmargin=2.6cm,style=nextline,itemsep=1pt]

\item[Category] synthetic (syntactic generation from the classic Datalog programs); family: Obstructed.

\item[Kind] equivalence.

\item[Contributors] Fangzhu Shen.

\item[Status] certified valid (Certificate/Valid.lean, written by the certificate emitter); expected valid.

\item[Tags] \hyperref[tag:linear-vs-nonlinear]{\texttt{linear-vs-nonlinear}} \hyperref[tag:mutual-recursion]{\texttt{mutual-recursion}} \hyperref[tag:nonlinear-recursion]{\texttt{nonlinear-recursion}} \hyperref[tag:prophecy]{\texttt{prophecy}}.

\item[Input] \path{Benchmark/Example0133/Input.lean} (canonical notation).

\item[Summary] All paths through odd and even halves, computed twice: by a linear program and by a non-linear one.

\item[Sources] {\raggedright legacy source record v2\_\_equiv\_pairs\_\_all\_paths\_via\_even\_odd\_\_linear\_vs\_nonlinear\_\_naive of the retired classic-source collection \par classic Datalog program all\_paths\_via\_even\_odd of the retired generator suite \par classic Datalog program all\_paths\_via\_even\_odd\_nonlinear of the retired generator suite\par}

\item[Prophecy] kernel-checked certificate with a level-1 prophecy clause (Certificate/Valid.lean, Core.json).

\item[Note] Kept after the duplicate review at Leo's decision: it is 0132 plus the union outputs and an implied conjunct, but Val's Part 1 review kept both and the all-paths program is a classic program in its own right.

\end{description}

\paragraph{Datalog program(s).}

{\raggedright\textit{programX --- Odd and even paths and their union, linearly.} --- fidelity module (kernel-checked: the compiled sides of the command are this program's naive compilation).\par}

\begin{Verbatim}[fontsize=\scriptsize,breaklines=true,breakanywhere=true,frame=single,framesep=1.5mm]
TbX(x1, y1) :- E(x1, y1).
TbX(x1, y1) :- E(x1, z1), TaX(z1, y1).
TaX(x1, y1) :- E(x1, z1), TbX(z1, y1).
TcX(x1, y1) :- TaX(x1, y1).
TcX(x1, y1) :- TbX(x1, y1).
\end{Verbatim}

{\raggedright\textit{programY --- Odd and even paths and their union, non-linearly.} --- fidelity module (kernel-checked: the compiled sides of the command are this program's naive compilation).\par}

\begin{Verbatim}[fontsize=\scriptsize,breaklines=true,breakanywhere=true,frame=single,framesep=1.5mm]
TaY(x1, y1) :- E(x1, z1), E(z1, y1).
TaY(x1, y1) :- TaY(x1, z1), TaY(z1, y1).
TaY(x1, y1) :- TbY(x1, z1), TbY(z1, y1).
TbY(x1, y1) :- E(x1, y1).
TbY(x1, y1) :- TbY(x1, z1), TaY(z1, y1).
TbY(x1, y1) :- TaY(x1, z1), TbY(z1, y1).
TcY(x1, y1) :- TaY(x1, y1).
TcY(x1, y1) :- TbY(x1, y1).
\end{Verbatim}

\paragraph{Reading.} The command runs two independently compiled Datalog programs in sequence --- a linear even/odd program producing TaX, TbX and their union TcX, and a non-linear even/odd program producing TaY, TbY and their union TcY --- merged by the preprocessor into a single verified loop. The postcondition states that all three pairs of relations agree, TaX=TaY, TbX=TbY and TcX=TcY, so the Hoare triple is an equivalence claim between the linear and non-linear implementations of the even/odd path decomposition and its union (all paths of length one or more), both reaching the same least fixpoints.

\paragraph{First-order reading of PRE/POST, translated mechanically from the relational-algebra assertion in Input.lean.} The relational-algebra assertion in \texttt{Input.lean} remains the authority.

\textit{PRE:} ⊤\par The precondition is simply true, imposing no constraint on the input relations.

\textit{POST:} (∀x y. (TaX(x,y) ↔ TaY(x,y))) ∧ (∀z u. (TbX(z,u) ↔ TbY(z,u))) ∧ ∀v w. (TcX(v,w) ↔ TcY(v,w))\par The postcondition asserts that the linear program's three outputs equal the non-linear program's corresponding outputs: TaX = TaY, TbX = TbY, and TcX = TcY.

\paragraph{Whiel encoding (Hoare triple).}

\begin{Verbatim}[fontsize=\tiny,breaklines=true,breakanywhere=true,frame=single,framesep=1.5mm]
SCHEMA
  E/2, TaX/2, TbX/2, TcX/2, TaY/2, TbY/2, TcY/2, TaX_aux/2, TbX_aux/2, TcX_aux/2, TaY_aux/2, TbY_aux/2, TcY_aux/2

PRE
  true 

CMD
  TbX_aux := ∅[2];
  TaX_aux := ∅[2];
  TcX_aux := ∅[2];
  TbX :=
    (E ∪ (π[0, 3] (σ[#1 = #2] (E × TaX_aux))));
  TaX := (π[0, 3] (σ[#1 = #2] (E × TbX_aux)));
  TcX := (TaX_aux ∪ TbX_aux);
  WHILE
      (¬ ((TbX = TbX_aux) ∧
        ((TaX = TaX_aux) ∧ (TcX = TcX_aux)))) DO
    TbX_aux := TbX;
    TaX_aux := TaX;
    TcX_aux := TcX;
    TbX :=
      (TbX ∪
        (E ∪
          (π[0, 3]
            (σ[#1 = #2] (E × TaX_aux)))));
    TaX :=
      (TaX ∪
        (π[0, 3] (σ[#1 = #2] (E × TbX_aux))));
    TcX := (TcX ∪ (TaX_aux ∪ TbX_aux))
  END;
  TaY_aux := ∅[2];
  TbY_aux := ∅[2];
  TcY_aux := ∅[2];
  TaY :=
    ((π[0, 3] (σ[#1 = #2] (E × E))) ∪
      ((π[0, 3]
          (σ[#1 = #2] (TaY_aux × TaY_aux))) ∪
        (π[0, 3]
          (σ[#1 = #2] (TbY_aux × TbY_aux)))));
  TbY :=
    (E ∪
      ((π[0, 3]
          (σ[#1 = #2] (TbY_aux × TaY_aux))) ∪
        (π[0, 3]
          (σ[#1 = #2] (TaY_aux × TbY_aux)))));
  TcY := (TaY_aux ∪ TbY_aux);
  WHILE
      (¬ ((TaY = TaY_aux) ∧
        ((TbY = TbY_aux) ∧ (TcY = TcY_aux)))) DO
    TaY_aux := TaY;
    TbY_aux := TbY;
    TcY_aux := TcY;
    TaY :=
      (TaY ∪
        ((π[0, 3] (σ[#1 = #2] (E × E))) ∪
          ((π[0, 3]
              (σ[#1 = #2]
                (TaY_aux × TaY_aux))) ∪
            (π[0, 3]
              (σ[#1 = #2]
                (TbY_aux × TbY_aux))))));
    TbY :=
      (TbY ∪
        (E ∪
          ((π[0, 3]
              (σ[#1 = #2]
                (TbY_aux × TaY_aux))) ∪
            (π[0, 3]
              (σ[#1 = #2]
                (TaY_aux × TbY_aux))))));
    TcY := (TcY ∪ (TaY_aux ∪ TbY_aux))
  END

POST
  ((TaX = TaY) ∧ ((TbX = TbY) ∧ (TcX = TcY)))
\end{Verbatim}

\Needspace{0.45\textheight}

\section{Example0134: Red/blue alternating paths: linear versus non-linear program (two compiled loops)}\label{case:Example0134}

\begin{description}[leftmargin=2.6cm,style=nextline,itemsep=1pt]

\item[Category] synthetic (syntactic generation from the classic Datalog programs); family: Obstructed.

\item[Kind] equivalence.

\item[Contributors] Fangzhu Shen.

\item[Status] certified valid (certificate checked outside the library build) (Certificate/Valid.lean, written by the certificate emitter); expected valid. Placement note: a proof module imported by \path{Certificate/Valid.lean} does not build under the library build's 4194304 KB limit; it builds under 12582912 KB (peak 1705 MB)

\item[Tags] \hyperref[tag:linear-vs-nonlinear]{\texttt{linear-vs-nonlinear}} \hyperref[tag:mutual-recursion]{\texttt{mutual-recursion}} \hyperref[tag:nonlinear-recursion]{\texttt{nonlinear-recursion}} \hyperref[tag:prophecy]{\texttt{prophecy}}.

\item[Input] \path{Benchmark/Example0134/Input.lean} (canonical notation).

\item[Summary] Alternating red and blue walks, computed twice: by a linear program and by a non-linear one.

\item[Sources] {\raggedright legacy source record v2\_\_equiv\_pairs\_\_red\_blue\_\_linear\_vs\_nonlinear\_\_naive of the retired classic-source collection \par classic Datalog program red\_blue\_linear of the retired generator suite \par classic Datalog program red\_blue\_nonlinear of the retired generator suite\par}

\item[Prophecy] kernel-checked certificate with a level-1 prophecy clause (Certificate/Valid.lean, Core.json).

\end{description}

\paragraph{Datalog program(s).}

{\raggedright\textit{programX --- Colour-alternating walks, linearly.} --- fidelity module (kernel-checked: the compiled sides of the command are this program's naive compilation).\par}

\begin{Verbatim}[fontsize=\scriptsize,breaklines=true,breakanywhere=true,frame=single,framesep=1.5mm]
TaX(x1, y1) :- B(x1, y1).
TaX(x1, y1) :- B(x1, z1), TcX(z1, y1).
TbX(x1, y1) :- TaX(x1, z1), R(z1, y1).
TcX(x1, y1) :- TdX(x1, z1), B(z1, y1).
TdX(x1, y1) :- R(x1, y1).
TdX(x1, y1) :- R(x1, z1), TbX(z1, y1).
\end{Verbatim}

{\raggedright\textit{programY --- Colour-alternating walks, non-linearly.} --- fidelity module (kernel-checked: the compiled sides of the command are this program's naive compilation).\par}

\begin{Verbatim}[fontsize=\scriptsize,breaklines=true,breakanywhere=true,frame=single,framesep=1.5mm]
TaY(x1, y1) :- B(x1, y1).
TaY(x1, y1) :- TaY(x1, z1), TcY(z1, y1).
TbY(x1, y1) :- B(x1, z1), R(z1, y1).
TbY(x1, y1) :- TaY(x1, z1), TdY(z1, y1).
TcY(x1, y1) :- R(x1, z1), B(z1, y1).
TcY(x1, y1) :- TdY(x1, z1), TaY(z1, y1).
TdY(x1, y1) :- R(x1, y1).
TdY(x1, y1) :- TdY(x1, z1), TbY(z1, y1).
\end{Verbatim}

\paragraph{Reading.} The command runs two independently compiled Datalog programs in sequence, each computing the four families of colour-alternating walks in a red/blue-edged graph (walks that start and end blue, start blue end red, start red end blue, and start and end red) --- one program linear, the other non-linear --- merged by the preprocessor into a single verified loop. The postcondition states that all four output relations from the linear program equal the corresponding relations from the non-linear program, so the Hoare triple is an equivalence claim between two implementations of colour-alternating reachability, both converging to the same least fixpoints.

\paragraph{First-order reading of PRE/POST, translated mechanically from the relational-algebra assertion in Input.lean.} The relational-algebra assertion in \texttt{Input.lean} remains the authority.

\textit{PRE:} ⊤\par The precondition is simply true, imposing no constraint on the input relations.

\textit{POST:} (∀x y. (TaX(x,y) ↔ TaY(x,y))) ∧ (∀z u. (TbX(z,u) ↔ TbY(z,u))) ∧ (∀v w. (TcX(v,w) ↔ TcY(v,w))) ∧ ∀x1 x2. (TdX(x1,x2) ↔ TdY(x1,x2))\par The postcondition asserts that all four linear-program outputs equal the corresponding non-linear-program outputs: TaX=TaY, TbX=TbY, TcX=TcY, and TdX=TdY.

\paragraph{Whiel encoding (Hoare triple).}

\begin{Verbatim}[fontsize=\tiny,breaklines=true,breakanywhere=true,frame=single,framesep=1.5mm]
SCHEMA
  B/2, R/2, TaX/2, TbX/2, TcX/2, TdX/2, TaY/2, TbY/2, TcY/2, TdY/2, TaX_aux/2, TbX_aux/2, TcX_aux/2, TdX_aux/2, TaY_aux/2, TbY_aux/2, TcY_aux/2, TdY_aux/2

PRE
  true 

CMD
  TaX_aux := ∅[2];
  TbX_aux := ∅[2];
  TcX_aux := ∅[2];
  TdX_aux := ∅[2];
  TaX :=
    (B ∪ (π[0, 3] (σ[#1 = #2] (B × TcX_aux))));
  TbX := (π[0, 3] (σ[#1 = #2] (TaX_aux × R)));
  TcX := (π[0, 3] (σ[#1 = #2] (TdX_aux × B)));
  TdX :=
    (R ∪ (π[0, 3] (σ[#1 = #2] (R × TbX_aux))));
  WHILE
      (¬ ((TaX = TaX_aux) ∧
        ((TbX = TbX_aux) ∧
          ((TcX = TcX_aux) ∧
            (TdX = TdX_aux))))) DO
    TaX_aux := TaX;
    TbX_aux := TbX;
    TcX_aux := TcX;
    TdX_aux := TdX;
    TaX :=
      (TaX ∪
        (B ∪
          (π[0, 3]
            (σ[#1 = #2] (B × TcX_aux)))));
    TbX :=
      (TbX ∪
        (π[0, 3] (σ[#1 = #2] (TaX_aux × R))));
    TcX :=
      (TcX ∪
        (π[0, 3] (σ[#1 = #2] (TdX_aux × B))));
    TdX :=
      (TdX ∪
        (R ∪
          (π[0, 3]
            (σ[#1 = #2] (R × TbX_aux)))))
  END;
  TaY_aux := ∅[2];
  TbY_aux := ∅[2];
  TcY_aux := ∅[2];
  TdY_aux := ∅[2];
  TaY :=
    (B ∪
      (π[0, 3]
        (σ[#1 = #2] (TaY_aux × TcY_aux))));
  TbY :=
    ((π[0, 3] (σ[#1 = #2] (B × R))) ∪
      (π[0, 3]
        (σ[#1 = #2] (TaY_aux × TdY_aux))));
  TcY :=
    ((π[0, 3] (σ[#1 = #2] (R × B))) ∪
      (π[0, 3]
        (σ[#1 = #2] (TdY_aux × TaY_aux))));
  TdY :=
    (R ∪
      (π[0, 3]
        (σ[#1 = #2] (TdY_aux × TbY_aux))));
  WHILE
      (¬ ((TaY = TaY_aux) ∧
        ((TbY = TbY_aux) ∧
          ((TcY = TcY_aux) ∧
            (TdY = TdY_aux))))) DO
    TaY_aux := TaY;
    TbY_aux := TbY;
    TcY_aux := TcY;
    TdY_aux := TdY;
    TaY :=
      (TaY ∪
        (B ∪
          (π[0, 3]
            (σ[#1 = #2]
              (TaY_aux × TcY_aux)))));
    TbY :=
      (TbY ∪
        ((π[0, 3] (σ[#1 = #2] (B × R))) ∪
          (π[0, 3]
            (σ[#1 = #2]
              (TaY_aux × TdY_aux)))));
    TcY :=
      (TcY ∪
        ((π[0, 3] (σ[#1 = #2] (R × B))) ∪
          (π[0, 3]
            (σ[#1 = #2]
              (TdY_aux × TaY_aux)))));
    TdY :=
      (TdY ∪
        (R ∪
          (π[0, 3]
            (σ[#1 = #2]
              (TdY_aux × TbY_aux)))))
  END

POST
  ((TaX = TaY) ∧
    ((TbX = TbY) ∧ ((TcX = TcY) ∧ (TdX = TdY))))
\end{Verbatim}

\Needspace{0.45\textheight}

\section{Example0136: Right-linear transitive closure versus the non-linear odd/even path program (two compiled loops)}\label{case:Example0136}

\begin{description}[leftmargin=2.6cm,style=nextline,itemsep=1pt]

\item[Category] synthetic (syntactic generation from the classic Datalog programs); family: Obstructed.

\item[Kind] equivalence.

\item[Contributors] Fangzhu Shen.

\item[Status] certified valid (Certificate/Valid.lean, written by the certificate emitter); expected valid.

\item[Tags] \hyperref[tag:linear-vs-nonlinear]{\texttt{linear-vs-nonlinear}} \hyperref[tag:mutual-recursion]{\texttt{mutual-recursion}} \hyperref[tag:nonlinear-recursion]{\texttt{nonlinear-recursion}} \hyperref[tag:prophecy]{\texttt{prophecy}}.

\item[Input] \path{Benchmark/Example0136/Input.lean} (canonical notation).

\item[Summary] Transitive closure computed right-linearly, against all paths computed through non-linear odd and even halves.

\item[Sources] {\raggedright legacy source record v2\_\_equiv\_pairs\_\_tc\_left\_vs\_all\_paths\_via\_eo\_nonlinear\_\_naive of the retired classic-source collection\par}

\item[Prophecy] kernel-checked certificate with a level-1 prophecy clause (Certificate/Valid.lean, Core.json).

\end{description}

\paragraph{Datalog program(s).}

{\raggedright\textit{programX --- Transitive closure, right-linearly.} --- fidelity module (kernel-checked: the compiled sides of the command are this program's naive compilation).\par}

\begin{Verbatim}[fontsize=\scriptsize,breaklines=true,breakanywhere=true,frame=single,framesep=1.5mm]
TX(x1, y1) :- E(x1, y1).
TX(x1, y1) :- E(x1, z1), TX(z1, y1).
\end{Verbatim}

{\raggedright\textit{programY --- Odd and even paths and their union, non-linearly.} --- fidelity module (kernel-checked: the compiled sides of the command are this program's naive compilation).\par}

\begin{Verbatim}[fontsize=\scriptsize,breaklines=true,breakanywhere=true,frame=single,framesep=1.5mm]
TaY(x1, y1) :- E(x1, z1), E(z1, y1).
TaY(x1, y1) :- TaY(x1, z1), TaY(z1, y1).
TaY(x1, y1) :- TbY(x1, z1), TbY(z1, y1).
TbY(x1, y1) :- E(x1, y1).
TbY(x1, y1) :- TbY(x1, z1), TaY(z1, y1).
TbY(x1, y1) :- TaY(x1, z1), TbY(z1, y1).
TcY(x1, y1) :- TaY(x1, y1).
TcY(x1, y1) :- TbY(x1, y1).
\end{Verbatim}

\paragraph{Reading.} The command runs two independently compiled Datalog programs in sequence --- a right-linear transitive-closure program producing TX, and a non-linear even/odd program producing TaY, TbY and their union TcY --- merged by the preprocessor into a single verified loop. The postcondition TX=TcY is an equivalence claim: the ordinary transitive closure of E equals the union of the even-length and odd-length path relations, since every nonempty path has odd or even length.

\paragraph{First-order reading of PRE/POST, translated mechanically from the relational-algebra assertion in Input.lean.} The relational-algebra assertion in \texttt{Input.lean} remains the authority.

\textit{PRE:} ⊤\par The precondition is simply true, imposing no constraint on the input relations.

\textit{POST:} ∀x y. (TX(x,y) ↔ TcY(x,y))\par The postcondition asserts that the transitive closure TX computed by the first program equals TcY, the union of even- and odd-length paths computed by the second.

\paragraph{Whiel encoding (Hoare triple).}

\begin{Verbatim}[fontsize=\tiny,breaklines=true,breakanywhere=true,frame=single,framesep=1.5mm]
SCHEMA
  E/2, TX/2, TaY/2, TbY/2, TcY/2, TX_aux/2, TaY_aux/2, TbY_aux/2, TcY_aux/2

PRE
  true 

CMD
  TX_aux := ∅[2];
  TX :=
    (E ∪ (π[0, 3] (σ[#1 = #2] (E × TX_aux))));
  WHILE (¬ (TX = TX_aux)) DO
    TX_aux := TX;
    TX :=
      (TX ∪
        (E ∪
          (π[0, 3]
            (σ[#1 = #2] (E × TX_aux)))))
  END;
  TaY_aux := ∅[2];
  TbY_aux := ∅[2];
  TcY_aux := ∅[2];
  TaY :=
    ((π[0, 3] (σ[#1 = #2] (E × E))) ∪
      ((π[0, 3]
          (σ[#1 = #2] (TaY_aux × TaY_aux))) ∪
        (π[0, 3]
          (σ[#1 = #2] (TbY_aux × TbY_aux)))));
  TbY :=
    (E ∪
      ((π[0, 3]
          (σ[#1 = #2] (TbY_aux × TaY_aux))) ∪
        (π[0, 3]
          (σ[#1 = #2] (TaY_aux × TbY_aux)))));
  TcY := (TaY_aux ∪ TbY_aux);
  WHILE
      (¬ ((TaY = TaY_aux) ∧
        ((TbY = TbY_aux) ∧ (TcY = TcY_aux)))) DO
    TaY_aux := TaY;
    TbY_aux := TbY;
    TcY_aux := TcY;
    TaY :=
      (TaY ∪
        ((π[0, 3] (σ[#1 = #2] (E × E))) ∪
          ((π[0, 3]
              (σ[#1 = #2]
                (TaY_aux × TaY_aux))) ∪
            (π[0, 3]
              (σ[#1 = #2]
                (TbY_aux × TbY_aux))))));
    TbY :=
      (TbY ∪
        (E ∪
          ((π[0, 3]
              (σ[#1 = #2]
                (TbY_aux × TaY_aux))) ∪
            (π[0, 3]
              (σ[#1 = #2]
                (TaY_aux × TbY_aux))))));
    TcY := (TcY ∪ (TaY_aux ∪ TbY_aux))
  END

POST
  (TX = TcY)
\end{Verbatim}

\Needspace{0.45\textheight}

\section{Example0137: Non-linear transitive closure versus the linear odd/even path program (two compiled loops)}\label{case:Example0137}

\begin{description}[leftmargin=2.6cm,style=nextline,itemsep=1pt]

\item[Category] synthetic (syntactic generation from the classic Datalog programs); family: Obstructed.

\item[Kind] equivalence.

\item[Contributors] Fangzhu Shen.

\item[Status] certified valid (Certificate/Valid.lean, written by the certificate emitter); expected valid.

\item[Tags] \hyperref[tag:linear-vs-nonlinear]{\texttt{linear-vs-nonlinear}} \hyperref[tag:mutual-recursion]{\texttt{mutual-recursion}} \hyperref[tag:nonlinear-recursion]{\texttt{nonlinear-recursion}} \hyperref[tag:prophecy]{\texttt{prophecy}}.

\item[Input] \path{Benchmark/Example0137/Input.lean} (canonical notation).

\item[Summary] Transitive closure computed non-linearly, against all paths computed through linear odd and even halves.

\item[Sources] {\raggedright legacy source record v2\_\_equiv\_pairs\_\_tc\_nonlin\_vs\_all\_paths\_via\_eo\_linear\_\_naive of the retired classic-source collection \par classic Datalog program all\_paths\_via\_even\_odd of the retired generator suite \par classic Datalog program tc-nonlinear of the retired generator suite\par}

\item[Prophecy] kernel-checked certificate with a level-1 prophecy clause (Certificate/Valid.lean, Core.json).

\end{description}

\paragraph{Datalog program(s).}

{\raggedright\textit{programX --- Transitive closure, non-linearly.} --- fidelity module (kernel-checked: the compiled sides of the command are this program's naive compilation).\par}

\begin{Verbatim}[fontsize=\scriptsize,breaklines=true,breakanywhere=true,frame=single,framesep=1.5mm]
TX(x1, y1) :- E(x1, y1).
TX(x1, y1) :- TX(x1, z1), TX(z1, y1).
\end{Verbatim}

{\raggedright\textit{programY --- Odd and even paths and their union, linearly.} --- fidelity module (kernel-checked: the compiled sides of the command are this program's naive compilation).\par}

\begin{Verbatim}[fontsize=\scriptsize,breaklines=true,breakanywhere=true,frame=single,framesep=1.5mm]
TbY(x1, y1) :- E(x1, y1).
TbY(x1, y1) :- E(x1, z1), TaY(z1, y1).
TaY(x1, y1) :- E(x1, z1), TbY(z1, y1).
TcY(x1, y1) :- TaY(x1, y1).
TcY(x1, y1) :- TbY(x1, y1).
\end{Verbatim}

\paragraph{Reading.} The command runs two independently compiled Datalog programs in sequence --- a non-linear transitive-closure program producing TX by squaring, and a linear even/odd program producing TaY, TbY and their union TcY --- merged by the preprocessor into a single verified loop. The postcondition TX=TcY is an equivalence claim: the non-linearly computed transitive closure of E equals the union of the even-length and odd-length path relations computed linearly, since every nonempty path has odd or even length.

\paragraph{First-order reading of PRE/POST, translated mechanically from the relational-algebra assertion in Input.lean.} The relational-algebra assertion in \texttt{Input.lean} remains the authority.

\textit{PRE:} ⊤\par The precondition is simply true, imposing no constraint on the input relations.

\textit{POST:} ∀x y. (TX(x,y) ↔ TcY(x,y))\par The postcondition asserts that the transitive closure TX computed by the first program equals TcY, the union of even- and odd-length paths computed by the second.

\paragraph{Whiel encoding (Hoare triple).}

\begin{Verbatim}[fontsize=\tiny,breaklines=true,breakanywhere=true,frame=single,framesep=1.5mm]
SCHEMA
  E/2, TX/2, TaY/2, TbY/2, TcY/2, TX_aux/2, TaY_aux/2, TbY_aux/2, TcY_aux/2

PRE
  true 

CMD
  TX_aux := ∅[2];
  TX :=
    (E ∪
      (π[0, 3] (σ[#1 = #2] (TX_aux × TX_aux))));
  WHILE (¬ (TX = TX_aux)) DO
    TX_aux := TX;
    TX :=
      (TX ∪
        (E ∪
          (π[0, 3]
            (σ[#1 = #2] (TX_aux × TX_aux)))))
  END;
  TbY_aux := ∅[2];
  TaY_aux := ∅[2];
  TcY_aux := ∅[2];
  TbY :=
    (E ∪ (π[0, 3] (σ[#1 = #2] (E × TaY_aux))));
  TaY := (π[0, 3] (σ[#1 = #2] (E × TbY_aux)));
  TcY := (TaY_aux ∪ TbY_aux);
  WHILE
      (¬ ((TbY = TbY_aux) ∧
        ((TaY = TaY_aux) ∧ (TcY = TcY_aux)))) DO
    TbY_aux := TbY;
    TaY_aux := TaY;
    TcY_aux := TcY;
    TbY :=
      (TbY ∪
        (E ∪
          (π[0, 3]
            (σ[#1 = #2] (E × TaY_aux)))));
    TaY :=
      (TaY ∪
        (π[0, 3] (σ[#1 = #2] (E × TbY_aux))));
    TcY := (TcY ∪ (TaY_aux ∪ TbY_aux))
  END

POST
  (TX = TcY)
\end{Verbatim}

\Needspace{0.45\textheight}

\section{Example0138: Joint-loop equivalence: tc\_nonlinear vs all\_paths\_via\_even\_odd (non-linear, naive)}\label{case:Example0138}

\begin{description}[leftmargin=2.6cm,style=nextline,itemsep=1pt]

\item[Category] synthetic (syntactic generation from the classic Datalog programs); family: Synthetic.

\item[Kind] equivalence.

\item[Contributors] Fangzhu Shen, Leo Zhang.

\item[Status] certified valid (Certificate/Valid.lean, written by the certificate emitter); expected valid.

\item[Tags] \hyperref[tag:mutual-recursion]{\texttt{mutual-recursion}} \hyperref[tag:nonlinear-recursion]{\texttt{nonlinear-recursion}} \hyperref[tag:prophecy]{\texttt{prophecy}}.

\item[Input] \path{Benchmark/Example0138/Input.lean} (canonical notation).

\item[Summary] Transitive closure and the union of odd and even paths, both computed non-linearly, compared.

\item[Sources] {\raggedright legacy source record v2\_\_equiv\_pairs\_\_tc\_nonlin\_vs\_all\_paths\_via\_eo\_nonlinear\_\_naive of the retired classic-source collection \par classic Datalog program all\_paths\_via\_even\_odd\_nonlinear of the retired generator suite \par classic Datalog program tc-nonlinear of the retired generator suite\par}

\item[Prophecy] accepted invariant with a prophecy clause above level zero (Core.json); no certificate built from it yet.

\end{description}

\paragraph{Datalog program(s).}

{\raggedright\textit{programX --- Transitive closure, non-linearly.} --- fidelity module (kernel-checked: the compiled sides of the command are this program's naive compilation).\par}

\begin{Verbatim}[fontsize=\scriptsize,breaklines=true,breakanywhere=true,frame=single,framesep=1.5mm]
TX(x1, y1) :- E(x1, y1).
TX(x1, y1) :- TX(x1, z1), TX(z1, y1).
\end{Verbatim}

{\raggedright\textit{programY --- Odd and even paths and their union, non-linearly.} --- fidelity module (kernel-checked: the compiled sides of the command are this program's naive compilation).\par}

\begin{Verbatim}[fontsize=\scriptsize,breaklines=true,breakanywhere=true,frame=single,framesep=1.5mm]
TaY(x1, y1) :- E(x1, z1), E(z1, y1).
TaY(x1, y1) :- TaY(x1, z1), TaY(z1, y1).
TaY(x1, y1) :- TbY(x1, z1), TbY(z1, y1).
TbY(x1, y1) :- E(x1, y1).
TbY(x1, y1) :- TbY(x1, z1), TaY(z1, y1).
TbY(x1, y1) :- TaY(x1, z1), TbY(z1, y1).
TcY(x1, y1) :- TaY(x1, y1).
TcY(x1, y1) :- TbY(x1, y1).
\end{Verbatim}

\paragraph{Reading.} The command runs two independently compiled Datalog programs in sequence --- a non-linear transitive-closure program producing TX by squaring, and a non-linear even/odd program producing TaY, TbY and their union TcY, also by squaring --- merged by the preprocessor into a single verified loop. The postcondition TX=TcY is an equivalence claim: the transitive closure of E equals the union of the even-length and odd-length path relations, both computed non-linearly, since every nonempty path has odd or even length; the header notes this re-encodes a retired hand-written lockstep-loop version of the same case.

\paragraph{First-order reading of PRE/POST, translated mechanically from the relational-algebra assertion in Input.lean.} The relational-algebra assertion in \texttt{Input.lean} remains the authority.

\textit{PRE:} ⊤\par The precondition is simply true, imposing no constraint on the input relations.

\textit{POST:} ∀x y. (TX(x,y) ↔ TcY(x,y))\par The postcondition asserts that the transitive closure TX computed by the first program equals TcY, the union of even- and odd-length paths computed by the second.

\paragraph{Whiel encoding (Hoare triple).}

\begin{Verbatim}[fontsize=\scriptsize,breaklines=true,breakanywhere=true,frame=single,framesep=1.5mm]
SCHEMA
  E/2, TX/2, TaY/2, TbY/2, TcY/2, TX_aux/2, TaY_aux/2, TbY_aux/2, TcY_aux/2

PRE
  true 

CMD
  TX_aux := ∅[2];
  TX := (E ∪ π[0,3] (σ[#1 = #2] ((TX_aux × TX_aux))));
  WHILE (¬((TX = TX_aux))) DO
    TX_aux := TX;
    TX := (TX ∪ (E ∪ π[0,3] (σ[#1 = #2] ((TX_aux × TX_aux)))))
  END;
  TaY_aux := ∅[2];
  TbY_aux := ∅[2];
  TcY_aux := ∅[2];
  TaY := (π[0,3] (σ[#1 = #2] ((E × E))) ∪ (π[0,3] (σ[#1 = #2] ((TaY_aux × TaY_aux))) ∪ π[0,3] (σ[#1 = #2] ((TbY_aux × TbY_aux)))));
  TbY := (E ∪ (π[0,3] (σ[#1 = #2] ((TbY_aux × TaY_aux))) ∪ π[0,3] (σ[#1 = #2] ((TaY_aux × TbY_aux)))));
  TcY := (TaY_aux ∪ TbY_aux);
  WHILE (¬(((TaY = TaY_aux) ∧ ((TbY = TbY_aux) ∧ (TcY = TcY_aux))))) DO
    TaY_aux := TaY;
    TbY_aux := TbY;
    TcY_aux := TcY;
    TaY := (TaY ∪ (π[0,3] (σ[#1 = #2] ((E × E))) ∪ (π[0,3] (σ[#1 = #2] ((TaY_aux × TaY_aux))) ∪ π[0,3] (σ[#1 = #2] ((TbY_aux × TbY_aux))))));
    TbY := (TbY ∪ (E ∪ (π[0,3] (σ[#1 = #2] ((TbY_aux × TaY_aux))) ∪ π[0,3] (σ[#1 = #2] ((TaY_aux × TbY_aux))))));
    TcY := (TcY ∪ (TaY_aux ∪ TbY_aux))
  END

POST
  (TX = TcY)
\end{Verbatim}

\Needspace{0.45\textheight}

\section{Example0162: Two-player game reachability mixing existential and universal successor steps}\label{case:Example0162}

\begin{description}[leftmargin=2.6cm,style=nextline,itemsep=1pt]

\item[Category] synthetic (syntactic generation from the classic Datalog programs); family: Synthetic.

\item[Kind] equivalence (with a least fixpoint, via a pre-fixpoint witness).

\item[Contributors] Fangzhu Shen, Leo Zhang.

\item[Status] certified valid (Certificate/Valid.lean, written by the certificate emitter); expected valid.

\item[Tags] \hyperref[tag:fixpoint-witness]{\texttt{fixpoint-witness}} \hyperref[tag:negation]{\texttt{negation}}.

\item[Input] \path{Benchmark/Example0162/Input.lean} (canonical notation).

\item[Summary] Canonical form of the earlier single-level encoding of this case: the same schema, precondition, command and postcondition.

\item[Sources] {\raggedright legacy source record hard-game-reach of the retired classic-source collection\par}

\end{description}

\paragraph{Datalog program(s).}

{\raggedright\textit{Datalog reading of the game reachability (written for this report; the operator is monotone, T is its least fixpoint; the universal step is spelled with a negated escape relation)} --- hand-written reading (this report).\par}

\begin{Verbatim}[fontsize=\scriptsize,breaklines=true,breakanywhere=true,frame=single,framesep=1.5mm]
T(x) :- Target(x).
T(x) :- Pa(x), E(x, y), T(y).             % player-A vertex with some successor in T
T(x) :- V(x), not Pa(x), not Escape(x).    % player-B vertex all of whose successors are in T
Escape(x) :- E(x, y), not T(y).
\end{Verbatim}

\paragraph{Reading.} The loop computes the attractor set of a two-player reachability game on a graph E: each vertex belongs to the existential player (the Pa vertices) or the universal player (V minus Pa). A Pa-vertex is winning if it has some successor already known to be winning; a non-Pa vertex in V is winning if all of its successors that lie in V are already winning (so it has no escape from the winning set). T is the least fixpoint of this one-step operator starting from Target: the set of positions from which the existential player can force play to reach Target. The precondition says Rc contains Target and is closed under the same one-step operator (a pre-fixpoint); the postcondition says T is itself a fixpoint of that operator and is contained in Rc --- the standard least-fixpoint-correctness statement, with Rc a universally quantified dummy standing in for every pre-fixpoint.

\paragraph{First-order reading of PRE/POST, translated mechanically from the relational-algebra assertion in Input.lean.} The relational-algebra assertion in \texttt{Input.lean} remains the authority.

\textit{PRE:} (∀x. (Target(x) ∨ Pa(x) ∧ (∃y. E(x,y) ∧ Rc(y)) ∨ V(x) ∧ ¬Pa(x) ∧ (¬∃z. E(x,z) ∧ V(z) ∧ ¬Rc(z)) → Rc(x))) ∧ ∀u. (Target(u) → Rc(u))\par PRE says Rc contains every Target vertex and is closed under the attractor step: any Pa-vertex with an Rc-winning successor, or any non-Pa V-vertex all of whose V-successors are Rc-winning, is itself in Rc.

\textit{POST:} (∀x. (T(x) ↔ Target(x) ∨ Pa(x) ∧ (∃y. E(x,y) ∧ T(y)) ∨ V(x) ∧ ¬Pa(x) ∧ ¬∃z. E(x,z) ∧ V(z) ∧ ¬T(z))) ∧ ∀u. (T(u) → Rc(u))\par POST says T equals the result of one attractor step applied to T itself (T is a fixpoint of the operator), and that T is contained in Rc.

\paragraph{Whiel encoding (Hoare triple).}

\begin{Verbatim}[fontsize=\scriptsize,breaklines=true,breakanywhere=true,frame=single,framesep=1.5mm]
SCHEMA
  Target/1, V/1, Pa/1, Rc/1, T/1, S/1, E/2

PRE
  ((((Target ∪ (π[0] (σ[#0 = #1] (Pa × (π[0] (σ[#1 = #2] (E × Rc))))))) ∪ ((V ∖ Pa) ∖ (π[0] (σ[#1 = #2] (E × (V ∖ Rc)))))) ⊆ Rc) ∧ (Target ⊆ Rc))

CMD
  T := ∅;
  S := ((Target ∪ (π[0] (σ[#0 = #1] (Pa × (π[0] (σ[#1 = #2] (E × T))))))) ∪ ((V ∖ Pa) ∖ (π[0] (σ[#1 = #2] (E × (V ∖ T))))));
  WHILE (S ≠ T) DO
    T := S;
    S := ((Target ∪ (π[0] (σ[#0 = #1] (Pa × (π[0] (σ[#1 = #2] (E × T))))))) ∪ ((V ∖ Pa) ∖ (π[0] (σ[#1 = #2] (E × (V ∖ T))))))
  END

POST
  ((T = ((Target ∪ (π[0] (σ[#0 = #1] (Pa × (π[0] (σ[#1 = #2] (E × T))))))) ∪ ((V ∖ Pa) ∖ (π[0] (σ[#1 = #2] (E × (V ∖ T))))))) ∧ (T ⊆ Rc))
\end{Verbatim}

\Needspace{0.45\textheight}

\section{Example0163: Greatest-fixpoint safe set: shrinking iteration bounded below by Rb}\label{case:Example0163}

\begin{description}[leftmargin=2.6cm,style=nextline,itemsep=1pt]

\item[Category] synthetic (syntactic generation from the classic Datalog programs); family: Synthetic.

\item[Kind] equivalence (with a greatest fixpoint, via a post-fixpoint witness).

\item[Contributors] Fangzhu Shen, Leo Zhang.

\item[Status] certified valid (Certificate/Valid.lean, written by the certificate emitter); expected valid.

\item[Tags] \hyperref[tag:fixpoint-witness]{\texttt{fixpoint-witness}} \hyperref[tag:greatest-fixpoint]{\texttt{greatest-fixpoint}} \hyperref[tag:negation]{\texttt{negation}}.

\item[Input] \path{Benchmark/Example0163/Input.lean} (canonical notation).

\item[Summary] Canonical form of the earlier single-level encoding of this case: the same schema, precondition, command and postcondition.

\item[Sources] {\raggedright legacy source record hard-gfp-safeset of the retired classic-source collection\par}

\end{description}

\paragraph{Datalog program(s).}

{\raggedright\textit{Datalog reading of the greatest-fixpoint safe set (written for this report; complement formulation: T = V ∖ Unsafe)} --- hand-written reading (this report).\par}

\begin{Verbatim}[fontsize=\scriptsize,breaklines=true,breakanywhere=true,frame=single,framesep=1.5mm]
Unsafe(x) :- Bad(x).
Unsafe(x) :- E(x, y), Unsafe(y).
T(x) :- V(x), not Unsafe(x).
\end{Verbatim}

\paragraph{Reading.} The loop computes a greatest-fixpoint safe set over a graph E with vertex set V and designated bad vertices Bad: starting from T:=V and iterating T:=(V∖Bad)∖\{x : some E-successor of x lies outside the current T\}, the decreasing iteration converges to the largest set of non-bad vertices all of whose successors also stay in the set --- the standard safety/invariant computation used in model checking or as the safe region of a reachability game, equivalently the Datalog reading T(x):-V(x), not Unsafe(x) where Unsafe is the least fixpoint of reachability to Bad. The Hoare triple is the coinduction (greatest-fixpoint elimination) obligation: given any Rb that is already a post-fixpoint of the safe-set operator, the computed T must itself be an exact fixpoint of the operator and dominate Rb from above, establishing that the naive decreasing iteration computes the greatest fixpoint.

\paragraph{First-order reading of PRE/POST, translated mechanically from the relational-algebra assertion in Input.lean.} The relational-algebra assertion in \texttt{Input.lean} remains the authority.

\textit{PRE:} ∀x. (Rb(x) → V(x) ∧ ¬Bad(x) ∧ ¬∃y. E(x,y) ∧ V(y) ∧ ¬Rb(y))\par The precondition states that Rb is a post-fixpoint of the safe-set operator: Rb ⊆ (V∖Bad)∖\{x : E(x,y) for some y ∉ Rb\}, written Rb ⊆ ((V∖Bad)∖π[0](σ[\#1=\#2](E×(V∖Rb)))).

\textit{POST:} (∀x. (T(x) ↔ V(x) ∧ ¬Bad(x) ∧ ¬∃y. E(x,y) ∧ V(y) ∧ ¬T(y))) ∧ ∀z. (Rb(z) → T(z))\par The postcondition states that the computed T is an exact fixpoint of the safe-set operator, T = (V∖Bad)∖\{x : E(x,y) for some y∉T\}, and that Rb ⊆ T.

\paragraph{Whiel encoding (Hoare triple).}

\begin{Verbatim}[fontsize=\scriptsize,breaklines=true,breakanywhere=true,frame=single,framesep=1.5mm]
SCHEMA
  Bad/1, V/1, Rb/1, T/1, S/1, E/2

PRE
  (Rb ⊆ ((V ∖ Bad) ∖ (π[0] (σ[#1 = #2] (E × (V ∖ Rb))))))

CMD
  T := V;
  S := (V ∖ Bad);
  WHILE (S ≠ T) DO
    T := S;
    S := ((V ∖ Bad) ∖ (π[0] (σ[#1 = #2] (E × (V ∖ T)))))
  END

POST
  ((T = ((V ∖ Bad) ∖ (π[0] (σ[#1 = #2] (E × (V ∖ T)))))) ∧ (Rb ⊆ T))
\end{Verbatim}

\Needspace{0.45\textheight}

\section{Example1030: Foreign-key support core for sharded materialized rows}\label{case:Example1030}

\begin{description}[leftmargin=2.6cm,style=nextline,itemsep=1pt]

\item[Category] LLM-generated (GPT/Claude benchmark programs); family: LLM.

\item[Kind] equivalence (with a greatest fixpoint, via a post-fixpoint witness).

\item[Contributors] Leo Zhang.

\item[Status] certified valid (Certificate/Valid.lean, written by the certificate emitter); expected valid (expected; strengthened claim).

\item[Tags] \hyperref[tag:fixpoint-witness]{\texttt{fixpoint-witness}} \hyperref[tag:greatest-fixpoint]{\texttt{greatest-fixpoint}} \hyperref[tag:integrity-constraints]{\texttt{integrity-constraints}} \hyperref[tag:negation]{\texttt{negation}}.

\item[Input] \path{Benchmark/Example1030/Input.lean} (canonical notation).

\item[Summary] Rows are pairs (id, shard) in Base; Req(id, ref) says row id needs row ref to exist in the same shard. The loop repeatedly removes rows with an unmet requirement (Bad) until none remain, so Keep is the greatest subset of Base in which every requirement is satisfied within the subset: the greatest fixpoint of G(X) = X ∖ Bad(X).

\item[Sources] {\raggedright retired-corpus source record llm\_30 of the LLM family \par retired-corpus input record llm\_30 of the LLM family\par}

\end{description}

\paragraph{Datalog program(s).}

{\raggedright\textit{Datalog reading of the greatest-fixpoint integrity core (written for this report; complement formulation: Keep = Base ∖ Removed, a row is removed when a requirement points outside the kept rows of its shard)} --- hand-written reading (this report).\par}

\begin{Verbatim}[fontsize=\scriptsize,breaklines=true,breakanywhere=true,frame=single,framesep=1.5mm]
Removed(id, s) :- Base(id, s), Req(id, ref), not Base(ref, s).
Removed(id, s) :- Base(id, s), Req(id, ref), Removed(ref, s).
Keep(id, s) :- Base(id, s), not Removed(id, s).
\end{Verbatim}

\paragraph{Reading.} The Datalog program computes a referential-integrity closure over sharded materialized rows: Removed accumulates ids whose foreign-key requirement (Req) points outside Base, or that depend transitively on another already-removed row, and Keep = Base ∖ Removed is dually the greatest fixpoint of the operator that deletes a row whenever its requirement is unmet within the same shard. The Hoare triple states the greatest-fixpoint correctness result in its maximality form: Rb is a free relation the program never touches, ranging over every subset of Base that is internally self-consistent (a post-fixpoint of the delete operator), and the postcondition asserts that Keep itself is self-consistent (Bad = ∅) and contains every such Rb, i.e. Keep is the greatest self-consistent subset of Base.

\paragraph{First-order reading of PRE/POST, translated mechanically from the relational-algebra assertion in Input.lean.} The relational-algebra assertion in \texttt{Input.lean} remains the authority.

\textit{PRE:} (∀x y. (Rb(x,y) → Base(x,y))) ∧ ∀z u. (Rb(z,u) → Rb(z,u) ∧ ¬∃v. Rb(z,u) ∧ Req(z,v) ∧ ¬(Rb(z,u) ∧ Req(z,v) ∧ Rb(v,u)))\par Rb is a subset of Base such that no row of Rb has an unmet foreign-key requirement within Rb itself.

\textit{POST:} (∀x y. (Keep(x,y) → Base(x,y))) ∧ (¬∃z u. Bad(z,u)) ∧ ∀v w. (Rb(v,w) → Keep(v,w))\par Keep is a subset of Base in which every row's requirement is satisfied (Bad = ∅), and Keep contains every self-consistent subset Rb of Base, i.e. Keep is the greatest such subset.

\paragraph{Whiel encoding (Hoare triple).}

\begin{Verbatim}[fontsize=\scriptsize,breaklines=true,breakanywhere=true,frame=single,framesep=1.5mm]
SCHEMA
  Base/2, Req/2, Keep/2, Keep_aux/2, Bad/2, Rb/2, Need/3, Good/3, BadNeed/3

PRE
  (Rb ⊆ Base) ∧ (Rb ⊆ (Rb ∖ (π[0, 1] ((π[0, 1, 3] (σ[#0 = #2] (Rb × Req))) ∖ (π[0, 1, 2] (σ[(#2 = #3) ∧ (#1 = #4)] ((π[0, 1, 3] (σ[#0 = #2] (Rb × Req))) × Rb)))))))

CMD
  Keep_aux := Base;
  Need := (π[0, 1, 3] (σ[#0 = #2] (Keep_aux × Req)));
  Good := (π[0, 1, 2] (σ[(#2 = #3) ∧ (#1 = #4)] (Need × Keep_aux)));
  BadNeed := (Need ∖ Good);
  Bad := (π[0, 1] BadNeed);
  Keep := (Keep_aux ∖ Bad);
  WHILE Keep ≠ Keep_aux DO
  Keep_aux := Keep;
  Need := (π[0, 1, 3] (σ[#0 = #2] (Keep_aux × Req)));
  Good := (π[0, 1, 2] (σ[(#2 = #3) ∧ (#1 = #4)] (Need × Keep_aux)));
  BadNeed := (Need ∖ Good);
  Bad := (π[0, 1] BadNeed);
  Keep := (Keep_aux ∖ Bad)
  END

POST
  ((Keep ⊆ Base) ∧ (Bad = ∅[2])) ∧ (Rb ⊆ Keep)
\end{Verbatim}

\Needspace{0.45\textheight}

\section{Example1061: Equivalence of computed same-generation with lfp of the same-generation rules (pre-fixpoint)}\label{case:Example1061}

\begin{description}[leftmargin=2.6cm,style=nextline,itemsep=1pt]

\item[Category] LLM-generated (GPT/Claude benchmark programs); family: LLM.

\item[Kind] equivalence (with a least fixpoint, via a pre-fixpoint witness).

\item[Contributors] Leo Zhang.

\item[Status] certified valid (Certificate/Valid.lean, written by the certificate emitter); expected valid.

\item[Tags] \hyperref[tag:fixpoint-witness]{\texttt{fixpoint-witness}}.

\item[Input] \path{Benchmark/Example1061/Input.lean} (canonical notation).

\item[Summary] Canonical form of the earlier single-level encoding of this case: the same schema, precondition, command and postcondition; the snapshot relation SG\_2 is now SG\_aux.

\item[Sources] {\raggedright retired-corpus source record llm\_61 of the LLM family \par retired-corpus input record llm\_61 of the LLM family\par}

\end{description}

\paragraph{Datalog program(s).}

{\raggedright\textit{prog\_p} --- production check (the hand-written loop equals this program's naive compilation on the check instances).\par}

\begin{Verbatim}[fontsize=\scriptsize,breaklines=true,breakanywhere=true,frame=single,framesep=1.5mm]
SGC(x1, y1) :- Flat(x1, y1).
SGC(x1, y1) :- Up(x1, z1), SGC(z1, x2), Down(x2, y1).
\end{Verbatim}

\paragraph{Reading.} The Datalog program computes the classical same-generation query over a hierarchy: SGC holds directly related Flat pairs, plus pairs reached by going Up one level, recursively finding a same-generation pair, then coming Down one level. The Hoare triple is the standard least-fixpoint correctness statement: K is any pre-fixpoint of the same-generation immediate-consequence operator, and the postcondition shows the loop's computed SG is itself a pre-fixpoint of that operator and is contained in every such K, so the loop computes exactly the least fixpoint of the same-generation rules.

\paragraph{First-order reading of PRE/POST, translated mechanically from the relational-algebra assertion in Input.lean.} The relational-algebra assertion in \texttt{Input.lean} remains the authority.

\textit{PRE:} ∀x y. (Flat(x,y) ∨ (∃z. Up(x,z) ∧ ∃u. K(z,u) ∧ Down(u,y)) → K(x,y))\par K contains every Flat pair and is closed under composing an Up-edge, a pair already in K, and a Down-edge.

\textit{POST:} (∀x y. (Flat(x,y) → SG(x,y))) ∧ (∀z u. ((∃v. Up(z,v) ∧ ∃w. SG(v,w) ∧ Down(w,u)) → SG(z,u))) ∧ ∀x1 x2. (SG(x1,x2) → K(x1,x2))\par SG contains every Flat pair and is closed under one more Up/SG/Down composition step, and SG is contained in K, together pinning SG to the least fixpoint of the same-generation rules.

\paragraph{Whiel encoding (Hoare triple).}

\begin{Verbatim}[fontsize=\scriptsize,breaklines=true,breakanywhere=true,frame=single,framesep=1.5mm]
SCHEMA
  Flat/2, Up/2, Down/2, SG/2, SG_aux/2, K/2

PRE
  ((Flat ∪ (π[0, 3] (σ[#1 = #2] (Up × (π[0, 3] (σ[#1 = #2] (K × Down))))))) ⊆ K)

CMD
  SG_aux := ∅;
  SG := Flat;
  WHILE SG ≠ SG_aux DO
    SG_aux := SG;
    SG := (Flat ∪ (π[0, 3] (σ[#1 = #2] (Up × (π[0, 3] (σ[#1 = #2] (SG_aux × Down)))))))
  END

POST
  (((Flat ⊆ SG) ∧ ((π[0, 3] (σ[#1 = #2] (Up × (π[0, 3] (σ[#1 = #2] (SG × Down)))))) ⊆ SG)) ∧ (SG ⊆ K))
\end{Verbatim}

\Needspace{0.45\textheight}

\section{Example1081: Audited delegation contained in broad authorization closure}\label{case:Example1081}

\begin{description}[leftmargin=2.6cm,style=nextline,itemsep=1pt]

\item[Category] LLM-generated (GPT/Claude benchmark programs); family: two programs under TGD constraints (LLM-generated, GPT-5.6; legacy raw\_new pipeline).

\item[Kind] containment.

\item[Contributors] Leo Zhang.

\item[Status] certified valid (Certificate/Valid.lean, written by the certificate emitter); expected valid.

\item[Tags] \hyperref[tag:dependencies]{\texttt{dependencies}} \hyperref[tag:mutual-recursion]{\texttt{mutual-recursion}} \hyperref[tag:nonlinear-recursion]{\texttt{nonlinear-recursion}} \hyperref[tag:prophecy]{\texttt{prophecy}}.

\item[Input] \path{Benchmark/Example1081/Input.lean} (canonical notation).

\item[Summary] Audited delegation contained in broad authorization closure.

\item[Sides] naive compiler: program C1 of the retired corpus case Example1081 (output Grant) \par naive compiler: program C2 of the retired corpus case Example1081 (output Permit)

\item[Sources] {\raggedright LLM-generated (GPT-5.6), legacy raw\_new TGD pipeline this repository the retired corpus case Example1081 (programs C1, C2 and the TGD list tgd0..) TGD pipeline run raw-new-tgd-0002-product: TGDs translated to relational algebra by RelCalc.TGD.listToAssertExpr; programs compiled with the certified naive compiler (read 2026-09-15) \par method this repository the equivalence-under-constraints paper draft (technical appendix: reduction of \path{containment/equivalence} under TGDs to a Hoare triple) (read 2026-09-15) \par legacy source record raw\_tgd\_gpt56\_w\_inv.json \par legacy source record \path{Benchmark/Example1081/Input.lean} \par Datalog source noted in the legacy inventory of Example1081\par}

\item[Prophecy] kernel-checked certificate with a level-1 prophecy clause (Certificate/Valid.lean, Core.json).

\item[Note] Legacy raw\_new TGD case (GPT-5.6); one representative of the specialisation-in-generalisation template (1083, 1085, 1087, 1089 are its siblings, not included).

\end{description}

\paragraph{Datalog program(s).}

{\raggedright\textit{prog\_c1} --- fidelity module (kernel-checked: the compiled sides of the command are this program's naive compilation).\par}

\begin{Verbatim}[fontsize=\scriptsize,breaklines=true,breakanywhere=true,frame=single,framesep=1.5mm]
Grant(x1, y1) :- Trusted(x1, y1).
Delegated(x1, y1) :- Special(x1, y1).
Audited(x1, y1) :- Composite(x1, y1).
Grant(x1, z1) :- Delegated(x1, y1), Audited(y1, z1).
Delegated(x1, z1) :- Grant(x1, y1), Trusted(y1, z1).
Audited(x1, z1) :- Grant(x1, y1), Evidence(y1, z1, x2), Valid(x2).
Grant(x1, z1) :- Grant(x1, y1), Audited(y1, z1).
\end{Verbatim}

{\raggedright\textit{prog\_c2} --- fidelity module (kernel-checked: the compiled sides of the command are this program's naive compilation).\par}

\begin{Verbatim}[fontsize=\scriptsize,breaklines=true,breakanywhere=true,frame=single,framesep=1.5mm]
Permit(x1, y1) :- Broad(x1, y1).
Trace(x1, y1) :- Broad(x1, y1).
Permit(x1, y1) :- Evidence(x1, y1, z1), Valid(z1).
Permit(x1, y1) :- Trace(x1, y1).
Trace(x1, z1) :- Permit(x1, y1), Permit(y1, z1).
Trace(x1, z1) :- Permit(x1, y1), Trace(y1, z1).
Trace(x1, z1) :- Trace(x1, y1), Permit(y1, z1).
Trace(x1, z1) :- Trace(x1, y1), Trace(y1, z1).
Trace(x1, z1) :- Permit(x1, y1), Evidence(y1, z1, x2), Valid(x2).
\end{Verbatim}

\paragraph{Reading.} This is a query-containment-under-dependencies problem, the setting of the relational chase: program C1 computes Grant, the least fixpoint of a delegation program built from Trusted pairs directly and from Delegated/Audited intermediate relations that compose earlier grants (Composite grants backed by an evidence witness); program C2 computes Permit, the least fixpoint of a broader authorization closure built from Broad pairs, evidence-backed pairs, and a transitive Trace relation. The precondition states the four tuple-generating dependencies (one of them existential) relating the two programs' input relations, and the postcondition claims that whenever those dependencies hold, every pair C1's program grants is also permitted by C2's program: Grant is contained in Permit.

\paragraph{First-order reading of PRE/POST, translated mechanically from the relational-algebra assertion in Input.lean.} The relational-algebra assertion in \texttt{Input.lean} remains the authority.

\textit{PRE:} (∀x y. (Trusted(x,y) → Broad(x,y))) ∧ (∀z u. (Special(z,u) → Broad(z,u))) ∧ (∀v w. (Composite(v,w) → ∃x1. Evidence(v,w,x1) ∧ Valid(x1))) ∧ ∀x2 x3. ((∃x4. Evidence(x2,x3,x4) ∧ Valid(x4)) → Broad(x2,x3))\par PRE says every Trusted or Special pair is a Broad pair, every Composite pair has some Evidence witness marked Valid, and every pair with a Valid Evidence witness is itself a Broad pair.

\textit{POST:} ∀x y. (Grant(x,y) → Permit(x,y))\par POST says every pair the delegation program's Grant relation derives is also derived by the authorization program's Permit relation.

\paragraph{Whiel encoding (Hoare triple).}

\begin{Verbatim}[fontsize=\scriptsize,breaklines=true,breakanywhere=true,frame=single,framesep=1.5mm]
SCHEMA
  Valid/1, Trusted/2, Special/2, Composite/2, Broad/2, Grant/2, Delegated/2, Audited/2, Permit/2, Trace/2, Grant_aux/2, Delegated_aux/2, Audited_aux/2, Permit_aux/2, Trace_aux/2, Evidence/3

PRE
  ((Trusted ⊆ Broad) ∧ ((Special ⊆ Broad) ∧ ((Composite ⊆ π[0,1] (σ[#2 = #3] ((Evidence × Valid)))) ∧ (π[0,1] (σ[#2 = #3] ((Evidence × Valid))) ⊆ Broad))))

CMD
  Grant_aux := ∅[2];
  Delegated_aux := ∅[2];
  Audited_aux := ∅[2];
  Grant := (Trusted ∪ (π[0,3] (σ[#1 = #2] ((Delegated_aux × Audited_aux))) ∪ π[0,3] (σ[#1 = #2] ((Grant_aux × Audited_aux)))));
  Delegated := (Special ∪ π[0,3] (σ[#1 = #2] ((Grant_aux × Trusted))));
  Audited := (Composite ∪ π[0,3] (σ[(#1 = #2 ∧ #4 = #5)] ((Grant_aux × (Evidence × Valid)))));
  WHILE (¬(((Grant = Grant_aux) ∧ ((Delegated = Delegated_aux) ∧ (Audited = Audited_aux))))) DO
    Grant_aux := Grant;
    Delegated_aux := Delegated;
    Audited_aux := Audited;
    Grant := (Grant ∪ (Trusted ∪ (π[0,3] (σ[#1 = #2] ((Delegated_aux × Audited_aux))) ∪ π[0,3] (σ[#1 = #2] ((Grant_aux × Audited_aux))))));
    Delegated := (Delegated ∪ (Special ∪ π[0,3] (σ[#1 = #2] ((Grant_aux × Trusted)))));
    Audited := (Audited ∪ (Composite ∪ π[0,3] (σ[(#1 = #2 ∧ #4 = #5)] ((Grant_aux × (Evidence × Valid))))))
  END;
  Permit_aux := ∅[2];
  Trace_aux := ∅[2];
  Permit := (Broad ∪ (π[0,1] (σ[#2 = #3] ((Evidence × Valid))) ∪ Trace_aux));
  Trace := (Broad ∪ (π[0,3] (σ[#1 = #2] ((Permit_aux × Permit_aux))) ∪ (π[0,3] (σ[#1 = #2] ((Permit_aux × Trace_aux))) ∪ (π[0,3] (σ[#1 = #2] ((Trace_aux × Permit_aux))) ∪ (π[0,3] (σ[#1 = #2] ((Trace_aux × Trace_aux))) ∪ π[0,3] (σ[(#1 = #2 ∧ #4 = #5)] ((Permit_aux × (Evidence × Valid)))))))));
  WHILE (¬(((Permit = Permit_aux) ∧ (Trace = Trace_aux)))) DO
    Permit_aux := Permit;
    Trace_aux := Trace;
    Permit := (Permit ∪ (Broad ∪ (π[0,1] (σ[#2 = #3] ((Evidence × Valid))) ∪ Trace_aux)));
    Trace := (Trace ∪ (Broad ∪ (π[0,3] (σ[#1 = #2] ((Permit_aux × Permit_aux))) ∪ (π[0,3] (σ[#1 = #2] ((Permit_aux × Trace_aux))) ∪ (π[0,3] (σ[#1 = #2] ((Trace_aux × Permit_aux))) ∪ (π[0,3] (σ[#1 = #2] ((Trace_aux × Trace_aux))) ∪ π[0,3] (σ[(#1 = #2 ∧ #4 = #5)] ((Permit_aux × (Evidence × Valid))))))))))
  END

POST
  (Grant ⊆ Permit)
\end{Verbatim}

\Needspace{0.45\textheight}

\section{Example1088: Equivalent forward dependency closure and transposed reverse resolution}\label{case:Example1088}

\begin{description}[leftmargin=2.6cm,style=nextline,itemsep=1pt]

\item[Category] LLM-generated (GPT/Claude benchmark programs); family: two programs under TGD constraints (LLM-generated, GPT-5.6; legacy raw\_new pipeline).

\item[Kind] equivalence.

\item[Contributors] Leo Zhang.

\item[Status] certified valid (Certificate/Valid.lean, written by the certificate emitter); expected valid.

\item[Tags] \hyperref[tag:dependencies]{\texttt{dependencies}} \hyperref[tag:mutual-recursion]{\texttt{mutual-recursion}} \hyperref[tag:nonlinear-recursion]{\texttt{nonlinear-recursion}}.

\item[Input] \path{Benchmark/Example1088/Input.lean} (canonical notation).

\item[Summary] Equivalent forward dependency closure and transposed reverse resolution.

\item[Sides] naive compiler: program C1 of the retired corpus case Example1088 (output ForwardReach) \par naive compiler: program C2 of the retired corpus case Example1088 (output ResolvedPair)

\item[Sources] {\raggedright LLM-generated (GPT-5.6), legacy raw\_new TGD pipeline this repository the retired corpus case Example1088 (programs C1, C2 and the TGD list tgd0..) TGD pipeline run raw-new-tgd-0002-product: TGDs translated to relational algebra by RelCalc.TGD.listToAssertExpr; programs compiled with the certified naive compiler (read 2026-09-15) \par method this repository the equivalence-under-constraints paper draft (technical appendix: reduction of \path{containment/equivalence} under TGDs to a Hoare triple) (read 2026-09-15) \par legacy source record raw\_tgd\_gpt56\_w\_inv.json \par legacy source record \path{Benchmark/Example1088/Input.lean} \par Datalog source noted in the legacy inventory of Example1088\par}

\item[Obstruction note] from the metadata: anchored: both sides are naive closures at the same rate and the legacy product loop was certified with a quantifier-free invariant; no prophecy level expected

\item[Note] Legacy raw\_new TGD case (GPT-5.6): forward closure equals the transposed reverse resolution under inverse seed views and bridge-proof witnesses.

\end{description}

\paragraph{Datalog program(s).}

{\raggedright\textit{prog\_c1} --- fidelity module (kernel-checked: the compiled sides of the command are this program's naive compilation).\par}

\begin{Verbatim}[fontsize=\scriptsize,breaklines=true,breakanywhere=true,frame=single,framesep=1.5mm]
ForwardReach(x1, y1) :- ForwardSeed(x1, y1).
ForwardStep(x1, y1) :- ForwardEdge(x1, y1).
ForwardReach(x1, z1) :- ForwardReach(x1, y1), ForwardStep(y1, z1).
ForwardStep(x1, z1) :- ForwardReach(x1, y1), BridgeProof(y1, z1, x2), ValidBridge(x2).
ForwardReach(x1, z1) :- ForwardStep(x1, y1), ForwardReach(y1, z1).
\end{Verbatim}

{\raggedright\textit{prog\_c2} --- fidelity module (kernel-checked: the compiled sides of the command are this program's naive compilation).\par}

\begin{Verbatim}[fontsize=\scriptsize,breaklines=true,breakanywhere=true,frame=single,framesep=1.5mm]
ResolvedPair(x1, y1) :- BackwardSeed(y1, x1).
ReverseStep(y1, x1) :- BackwardEdge(y1, x1).
ResolvedPair(x1, z1) :- ResolvedPair(x1, y1), ReverseStep(z1, y1).
ReverseStep(z1, x1) :- ResolvedPair(x1, y1), BridgeProof(y1, z1, x2), ValidBridge(x2).
ResolvedPair(x1, z1) :- ReverseStep(y1, x1), ResolvedPair(y1, z1).
\end{Verbatim}

\paragraph{Reading.} This is a query-equivalence-under-dependencies problem, the setting of the relational chase: program C1 computes ForwardReach, the transitive closure of a forward step relation (direct ForwardEdge pairs, plus pairs backed by a valid bridge proof) starting from ForwardSeed; program C2 computes ResolvedPair, the same kind of closure computed instead over the transposed backward relations (BackwardSeed, BackwardEdge) with its own reverse step relation. The precondition's six tuple-generating dependencies (two of them existential) say the forward and backward seed relations are transposes of each other, that ForwardEdge equals the pairs backed by a valid bridge proof, and that BackwardEdge equals the transpose of that same relation; under those dependencies the postcondition claims the two closures agree, ForwardReach = ResolvedPair.

\paragraph{First-order reading of PRE/POST, translated mechanically from the relational-algebra assertion in Input.lean.} The relational-algebra assertion in \texttt{Input.lean} remains the authority.

\textit{PRE:} (∀x y. (ForwardSeed(x,y) → BackwardSeed(y,x))) ∧ (∀z u. (BackwardSeed(z,u) → ForwardSeed(u,z))) ∧ (∀v w. (ForwardEdge(v,w) → ∃x1. BridgeProof(v,w,x1) ∧ ValidBridge(x1))) ∧ (∀x2 x3. ((∃x4. BridgeProof(x2,x3,x4) ∧ ValidBridge(x4)) → ForwardEdge(x2,x3))) ∧ (∀x5 x6. ((∃x7. BridgeProof(x5,x6,x7) ∧ ValidBridge(x7)) → BackwardEdge(x6,x5))) ∧ ∀x8 x9. (BackwardEdge(x8,x9) → ∃x10. BridgeProof(x9,x8,x10) ∧ ValidBridge(x10))\par PRE says ForwardSeed and BackwardSeed are transposes of each other, ForwardEdge equals the pairs backed by a valid bridge proof, and BackwardEdge equals the transpose of that same bridge-proof relation.

\textit{POST:} ∀x y. (ForwardReach(x,y) ↔ ResolvedPair(x,y))\par POST says the forward-computed reachability relation ForwardReach equals the reverse-computed relation ResolvedPair.

\paragraph{Whiel encoding (Hoare triple).}

\begin{Verbatim}[fontsize=\scriptsize,breaklines=true,breakanywhere=true,frame=single,framesep=1.5mm]
SCHEMA
  ValidBridge/1, ForwardSeed/2, BackwardSeed/2, ForwardEdge/2, BackwardEdge/2, ForwardReach/2, ForwardStep/2, ResolvedPair/2, ReverseStep/2, ForwardReach_aux/2, ForwardStep_aux/2, ResolvedPair_aux/2, ReverseStep_aux/2, BridgeProof/3

PRE
  ((ForwardSeed ⊆ π[1,0] (BackwardSeed)) ∧ ((BackwardSeed ⊆ π[1,0] (ForwardSeed)) ∧ ((ForwardEdge ⊆ π[0,1] (σ[#2 = #3] ((BridgeProof × ValidBridge)))) ∧ ((π[0,1] (σ[#2 = #3] ((BridgeProof × ValidBridge))) ⊆ ForwardEdge) ∧ ((π[0,1] (σ[#2 = #3] ((BridgeProof × ValidBridge))) ⊆ π[1,0] (BackwardEdge)) ∧ (BackwardEdge ⊆ π[1,0] (σ[#2 = #3] ((BridgeProof × ValidBridge)))))))))

CMD
  ForwardReach_aux := ∅[2];
  ForwardStep_aux := ∅[2];
  ForwardReach := (ForwardSeed ∪ (π[0,3] (σ[#1 = #2] ((ForwardReach_aux × ForwardStep_aux))) ∪ π[0,3] (σ[#1 = #2] ((ForwardStep_aux × ForwardReach_aux)))));
  ForwardStep := (ForwardEdge ∪ π[0,3] (σ[(#1 = #2 ∧ #4 = #5)] ((ForwardReach_aux × (BridgeProof × ValidBridge)))));
  WHILE (¬(((ForwardReach = ForwardReach_aux) ∧ (ForwardStep = ForwardStep_aux)))) DO
    ForwardReach_aux := ForwardReach;
    ForwardStep_aux := ForwardStep;
    ForwardReach := (ForwardReach ∪ (ForwardSeed ∪ (π[0,3] (σ[#1 = #2] ((ForwardReach_aux × ForwardStep_aux))) ∪ π[0,3] (σ[#1 = #2] ((ForwardStep_aux × ForwardReach_aux))))));
    ForwardStep := (ForwardStep ∪ (ForwardEdge ∪ π[0,3] (σ[(#1 = #2 ∧ #4 = #5)] ((ForwardReach_aux × (BridgeProof × ValidBridge))))))
  END;
  ResolvedPair_aux := ∅[2];
  ReverseStep_aux := ∅[2];
  ResolvedPair := (π[1,0] (BackwardSeed) ∪ (π[0,2] (σ[#1 = #3] ((ResolvedPair_aux × ReverseStep_aux))) ∪ π[1,3] (σ[#0 = #2] ((ReverseStep_aux × ResolvedPair_aux)))));
  ReverseStep := (BackwardEdge ∪ π[3,0] (σ[(#1 = #2 ∧ #4 = #5)] ((ResolvedPair_aux × (BridgeProof × ValidBridge)))));
  WHILE (¬(((ResolvedPair = ResolvedPair_aux) ∧ (ReverseStep = ReverseStep_aux)))) DO
    ResolvedPair_aux := ResolvedPair;
    ReverseStep_aux := ReverseStep;
    ResolvedPair := (ResolvedPair ∪ (π[1,0] (BackwardSeed) ∪ (π[0,2] (σ[#1 = #3] ((ResolvedPair_aux × ReverseStep_aux))) ∪ π[1,3] (σ[#0 = #2] ((ReverseStep_aux × ResolvedPair_aux))))));
    ReverseStep := (ReverseStep ∪ (BackwardEdge ∪ π[3,0] (σ[(#1 = #2 ∧ #4 = #5)] ((ResolvedPair_aux × (BridgeProof × ValidBridge))))))
  END

POST
  (ForwardReach = ResolvedPair)
\end{Verbatim}

\Needspace{0.45\textheight}

\section{Example1090: Equivalent taint reachability with ternary modes and projected frontiers}\label{case:Example1090}

\begin{description}[leftmargin=2.6cm,style=nextline,itemsep=1pt]

\item[Category] LLM-generated (GPT/Claude benchmark programs); family: two programs under TGD constraints (LLM-generated, GPT-5.6; legacy raw\_new pipeline).

\item[Kind] equivalence.

\item[Contributors] Leo Zhang.

\item[Status] certified valid (Certificate/Valid.lean, written by the certificate emitter); expected valid.

\item[Tags] \hyperref[tag:dependencies]{\texttt{dependencies}} \hyperref[tag:mutual-recursion]{\texttt{mutual-recursion}} \hyperref[tag:nonlinear-recursion]{\texttt{nonlinear-recursion}}.

\item[Input] \path{Benchmark/Example1090/Input.lean} (canonical notation).

\item[Summary] Equivalent taint reachability with ternary modes and projected frontiers.

\item[Sides] naive compiler: program C1 of the retired corpus case Example1090 (output TaintReach) \par naive compiler: program C2 of the retired corpus case Example1090 (output AnalysisReach)

\item[Sources] {\raggedright LLM-generated (GPT-5.6), legacy raw\_new TGD pipeline this repository the retired corpus case Example1090 (programs C1, C2 and the TGD list tgd0..) TGD pipeline run raw-new-tgd-0002-product: TGDs translated to relational algebra by RelCalc.TGD.listToAssertExpr; programs compiled with the certified naive compiler (read 2026-09-15) \par method this repository the equivalence-under-constraints paper draft (technical appendix: reduction of \path{containment/equivalence} under TGDs to a Hoare triple) (read 2026-09-15) \par legacy source record raw\_tgd\_gpt56\_w\_inv.json \par legacy source record \path{Benchmark/Example1090/Input.lean} \par Datalog source noted in the legacy inventory of Example1090\par}

\item[Obstruction note] from the metadata: anchored: both sides are naive closures at the same rate and the legacy product loop was certified with a quantifier-free invariant; no prophecy level expected

\item[Note] Legacy raw\_new TGD case (GPT-5.6): a ternary mode-carrying state equals a binary projected frontier; exact only because every node has a mode (existential dependency).

\end{description}

\paragraph{Datalog program(s).}

{\raggedright\textit{prog\_c1} --- fidelity module (kernel-checked: the compiled sides of the command are this program's naive compilation).\par}

\begin{Verbatim}[fontsize=\scriptsize,breaklines=true,breakanywhere=true,frame=single,framesep=1.5mm]
TaintReach(x1, y1) :- SeedPair(x1, y1).
TaintState(x1, y1, z1) :- SeedPair(x1, y1), ModeTag(y1, z1).
TaintReach(x1, z1) :- TaintState(x1, y1, x2), FlowEdge(y1, z1).
TaintState(x1, z1, x2) :- TaintReach(x1, y1), FlowEdge(y1, z1), ModeTag(z1, x2).
TaintState(x1, z1, y2) :- TaintState(x1, y1, x2), FlowEdge(y1, z1), ModeTag(z1, y2).
TaintReach(x1, z1) :- TaintReach(x1, y1), TaintReach(y1, z1).
\end{Verbatim}

{\raggedright\textit{prog\_c2} --- fidelity module (kernel-checked: the compiled sides of the command are this program's naive compilation).\par}

\begin{Verbatim}[fontsize=\scriptsize,breaklines=true,breakanywhere=true,frame=single,framesep=1.5mm]
AnalysisReach(x1, y1) :- InputPair(x1, y1).
FrontierState(x1, y1) :- InputPair(x1, y1).
AnalysisReach(x1, z1) :- FrontierState(x1, y1), ViewEdge(y1, z1).
FrontierState(x1, z1) :- AnalysisReach(x1, y1), ViewEdge(y1, z1).
FrontierState(x1, z1) :- FrontierState(x1, y1), StateWitness(y1, z1, x2), EnabledMode(x2).
AnalysisReach(x1, z1) :- AnalysisReach(x1, y1), AnalysisReach(y1, z1).
\end{Verbatim}

\paragraph{Reading.} This is a query-equivalence-under-dependencies problem, the setting of the relational chase: program C1 computes TaintReach, an interprocedural-taint-style reachability relation computed with the help of an auxiliary ternary TaintState that additionally records a mode witness at each frontier vertex; program C2 computes AnalysisReach, the same reachability relation computed with a simpler binary FrontierState that drops the mode tag. The precondition's eight tuple-generating dependencies (three of them existential) say SeedPair equals InputPair, FlowEdge equals ViewEdge, every vertex reached by SeedPair or FlowEdge carries some ModeTag, and FlowEdge coincides exactly with the pairs backed by a valid StateWitness; under these, the postcondition claims the two closures compute the same relation, TaintReach = AnalysisReach --- the mode bookkeeping in C1 turns out not to change what gets reached, once the dependencies guarantee a mode witness always exists.

\paragraph{First-order reading of PRE/POST, translated mechanically from the relational-algebra assertion in Input.lean.} The relational-algebra assertion in \texttt{Input.lean} remains the authority.

\textit{PRE:} (∀x y. (SeedPair(x,y) → InputPair(x,y))) ∧ (∀z u. (InputPair(z,u) → SeedPair(z,u))) ∧ (∀v w. (FlowEdge(v,w) → ViewEdge(v,w))) ∧ (∀x1 x2. (ViewEdge(x1,x2) → FlowEdge(x1,x2))) ∧ (∀x3. ((∃x4. SeedPair(x4,x3)) → ∃x5. ModeTag(x3,x5))) ∧ (∀x6. ((∃x7. FlowEdge(x7,x6)) → ∃x8. ModeTag(x6,x8))) ∧ (∀x9 x10. (FlowEdge(x9,x10) → ∃x11. StateWitness(x9,x10,x11) ∧ EnabledMode(x11))) ∧ ∀x12 x13. ((∃x14. StateWitness(x12,x13,x14) ∧ EnabledMode(x14)) → FlowEdge(x12,x13))\par PRE says SeedPair and InputPair are the same relation, FlowEdge and ViewEdge are the same relation, every vertex reached by SeedPair or FlowEdge carries some ModeTag, and FlowEdge coincides exactly with the pairs backed by a valid StateWitness.

\textit{POST:} ∀x y. (TaintReach(x,y) ↔ AnalysisReach(x,y))\par POST says the ternary mode-carrying reachability relation TaintReach and the binary frontier relation AnalysisReach compute the same pairs.

\paragraph{Whiel encoding (Hoare triple).}

\begin{Verbatim}[fontsize=\scriptsize,breaklines=true,breakanywhere=true,frame=single,framesep=1.5mm]
SCHEMA
  EnabledMode/1, SeedPair/2, InputPair/2, FlowEdge/2, ViewEdge/2, ModeTag/2, TaintReach/2, AnalysisReach/2, FrontierState/2, TaintReach_aux/2, AnalysisReach_aux/2, FrontierState_aux/2, StateWitness/3, TaintState/3, TaintState_aux/3

PRE
  ((SeedPair ⊆ InputPair) ∧ ((InputPair ⊆ SeedPair) ∧ ((FlowEdge ⊆ ViewEdge) ∧ ((ViewEdge ⊆ FlowEdge) ∧ ((π[1] (SeedPair) ⊆ π[0] (ModeTag)) ∧ ((π[1] (FlowEdge) ⊆ π[0] (ModeTag)) ∧ ((FlowEdge ⊆ π[0,1] (σ[#2 = #3] ((StateWitness × EnabledMode)))) ∧ (π[0,1] (σ[#2 = #3] ((StateWitness × EnabledMode))) ⊆ FlowEdge))))))))

CMD
  TaintReach_aux := ∅[2];
  TaintState_aux := ∅[3];
  TaintReach := (SeedPair ∪ (π[0,4] (σ[#1 = #3] ((TaintState_aux × FlowEdge))) ∪ π[0,3] (σ[#1 = #2] ((TaintReach_aux × TaintReach_aux)))));
  TaintState := (π[0,1,3] (σ[#1 = #2] ((SeedPair × ModeTag))) ∪ (π[0,3,5] (σ[(#1 = #2 ∧ #3 = #4)] ((TaintReach_aux × (FlowEdge × ModeTag)))) ∪ π[0,4,6] (σ[(#1 = #3 ∧ #4 = #5)] ((TaintState_aux × (FlowEdge × ModeTag))))));
  WHILE (¬(((TaintReach = TaintReach_aux) ∧ (TaintState = TaintState_aux)))) DO
    TaintReach_aux := TaintReach;
    TaintState_aux := TaintState;
    TaintReach := (TaintReach ∪ (SeedPair ∪ (π[0,4] (σ[#1 = #3] ((TaintState_aux × FlowEdge))) ∪ π[0,3] (σ[#1 = #2] ((TaintReach_aux × TaintReach_aux))))));
    TaintState := (TaintState ∪ (π[0,1,3] (σ[#1 = #2] ((SeedPair × ModeTag))) ∪ (π[0,3,5] (σ[(#1 = #2 ∧ #3 = #4)] ((TaintReach_aux × (FlowEdge × ModeTag)))) ∪ π[0,4,6] (σ[(#1 = #3 ∧ #4 = #5)] ((TaintState_aux × (FlowEdge × ModeTag)))))))
  END;
  AnalysisReach_aux := ∅[2];
  FrontierState_aux := ∅[2];
  AnalysisReach := (InputPair ∪ (π[0,3] (σ[#1 = #2] ((FrontierState_aux × ViewEdge))) ∪ π[0,3] (σ[#1 = #2] ((AnalysisReach_aux × AnalysisReach_aux)))));
  FrontierState := (InputPair ∪ (π[0,3] (σ[#1 = #2] ((AnalysisReach_aux × ViewEdge))) ∪ π[0,3] (σ[(#1 = #2 ∧ #4 = #5)] ((FrontierState_aux × (StateWitness × EnabledMode))))));
  WHILE (¬(((AnalysisReach = AnalysisReach_aux) ∧ (FrontierState = FrontierState_aux)))) DO
    AnalysisReach_aux := AnalysisReach;
    FrontierState_aux := FrontierState;
    AnalysisReach := (AnalysisReach ∪ (InputPair ∪ (π[0,3] (σ[#1 = #2] ((FrontierState_aux × ViewEdge))) ∪ π[0,3] (σ[#1 = #2] ((AnalysisReach_aux × AnalysisReach_aux))))));
    FrontierState := (FrontierState ∪ (InputPair ∪ (π[0,3] (σ[#1 = #2] ((AnalysisReach_aux × ViewEdge))) ∪ π[0,3] (σ[(#1 = #2 ∧ #4 = #5)] ((FrontierState_aux × (StateWitness × EnabledMode)))))))
  END

POST
  (TaintReach = AnalysisReach)
\end{Verbatim}

\Needspace{0.45\textheight}

\section{Example1091: Delegation reachability bounded by a transitively closed policy view}\label{case:Example1091}

\begin{description}[leftmargin=2.6cm,style=nextline,itemsep=1pt]

\item[Category] LLM-generated (GPT/Claude benchmark programs); family: two programs under TGD constraints (LLM-generated, Claude Opus; legacy raw\_new pipeline).

\item[Kind] containment.

\item[Contributors] Leo Zhang.

\item[Status] certified valid (Certificate/Valid.lean, written by the certificate emitter); expected valid.

\item[Tags] \hyperref[tag:dependencies]{\texttt{dependencies}} \hyperref[tag:mutual-recursion]{\texttt{mutual-recursion}} \hyperref[tag:nonlinear-recursion]{\texttt{nonlinear-recursion}}.

\item[Input] \path{Benchmark/Example1091/Input.lean} (canonical notation).

\item[Summary] Delegation reachability bounded by a transitively closed policy view.

\item[Sides] naive compiler: program C1 of the retired corpus case Example1091 (output P) \par naive compiler: program C2 of the retired corpus case Example1091 (output T)

\item[Sources] {\raggedright LLM-generated (Claude Opus), legacy raw\_new TGD pipeline this repository the retired corpus case Example1091 (programs C1, C2 and the TGD list tgd0..) TGD pipeline run raw-new-tgd-0002-product: TGDs translated to relational algebra by RelCalc.TGD.listToAssertExpr; programs compiled with the certified naive compiler (read 2026-09-15) \par method this repository the equivalence-under-constraints paper draft (technical appendix: reduction of \path{containment/equivalence} under TGDs to a Hoare triple) (read 2026-09-15) \par legacy source record raw\_tgd\_opus\_w\_inv.json \par legacy source record \path{Benchmark/Example1091/Input.lean} \par Datalog source noted in the legacy inventory of Example1091\par}

\item[Obstruction note] from the metadata: anchored: both sides are naive closures at the same rate and the legacy product loop was certified with a quantifier-free invariant; no prophecy level expected

\item[Note] Legacy raw\_new TGD case (Claude Opus): every link has a midpoint (existential dependency); containment in a constraint-closed policy view.

\end{description}

\paragraph{Datalog program(s).}

{\raggedright\textit{prog\_c1} --- fidelity module (kernel-checked: the compiled sides of the command are this program's naive compilation).\par}

\begin{Verbatim}[fontsize=\scriptsize,breaklines=true,breakanywhere=true,frame=single,framesep=1.5mm]
P(x1, y1) :- Link(x1, y1).
P(x1, y1) :- Mid(x1, y1).
P(x1, z1) :- P(x1, y1), Q(y1, z1).
P(x1, x2) :- P(x1, y1), Mid(y1, z1), P(z1, x2).
Q(x1, y1) :- Tri(x1, z1, y1).
Q(x1, z1) :- Q(x1, y1), P(y1, z1).
\end{Verbatim}

{\raggedright\textit{prog\_c2} --- fidelity module (kernel-checked: the compiled sides of the command are this program's naive compilation).\par}

\begin{Verbatim}[fontsize=\scriptsize,breaklines=true,breakanywhere=true,frame=single,framesep=1.5mm]
T(x1, y1) :- BaseTwo(x1, y1).
G(x1, y1) :- Jump(x1, y1).
G(x1, z1) :- G(x1, y1), Jump(y1, z1).
T(x1, z1) :- T(x1, y1), G(y1, z1).
T(x1, z1) :- T(x1, y1), T(y1, z1).
\end{Verbatim}

\paragraph{Reading.} Program C1 recursively derives grantable accesses P via direct delegation (Link, Mid) and composition with Q, where Q is built from ternary relay facts Tri and further composed with P, a mutually recursive delegation/role-composition propagation typical of access-control reachability analysis. Program C2 recomputes reachability T as the transitive closure of BaseTwo composed with a separately closed relation G built from Jump edges. The Hoare triple is a containment claim under tuple-generating dependencies (the classical chase setting): the precondition encodes five TGDs forcing C to be transitively closed, to contain Mid and the first/third coordinates of every Tri triple, to guarantee every Link edge a two-hop Mid-Mid witness, and to force BaseTwo to contain all of C; under these constraints the postcondition claims P ⊆ T, expected valid because both programs are naive closures advancing at the same rate.

\paragraph{First-order reading of PRE/POST, translated mechanically from the relational-algebra assertion in Input.lean.} The relational-algebra assertion in \texttt{Input.lean} remains the authority.

\textit{PRE:} (∀x y. ((∃z. C(x,z) ∧ C(z,y)) → C(x,y))) ∧ (∀u v. (Mid(u,v) → C(u,v))) ∧ (∀w x1. (Link(w,x1) → ∃x2. Mid(w,x2) ∧ Mid(x2,x1))) ∧ (∀x3 x4. ((∃x5. Tri(x3,x5,x4)) → C(x3,x4))) ∧ ∀x6 x7. (C(x6,x7) → BaseTwo(x6,x7))\par The precondition is five TGDs: C is transitively closed and contains Mid and the first/third coordinates of every Tri triple; every Link edge has some two-hop Mid-Mid witness; and BaseTwo contains all of C.

\textit{POST:} ∀x y. (P(x,y) → T(x,y))\par The postcondition states that every pair C1 derives into P is also derived by C2 into T, i.e. P ⊆ T.

\paragraph{Whiel encoding (Hoare triple).}

\begin{Verbatim}[fontsize=\scriptsize,breaklines=true,breakanywhere=true,frame=single,framesep=1.5mm]
SCHEMA
  C/2, Mid/2, Link/2, BaseTwo/2, Jump/2, P/2, Q/2, T/2, G/2, P_aux/2, Q_aux/2, T_aux/2, G_aux/2, Tri/3

PRE
  ((π[0,3] (σ[#1 = #2] ((C × C))) ⊆ C) ∧ ((Mid ⊆ C) ∧ ((Link ⊆ π[0,3] (σ[#1 = #2] ((Mid × Mid)))) ∧ ((π[0,2] (Tri) ⊆ C) ∧ (C ⊆ BaseTwo)))))

CMD
  P_aux := ∅[2];
  Q_aux := ∅[2];
  P := (Link ∪ (Mid ∪ (π[0,3] (σ[#1 = #2] ((P_aux × Q_aux))) ∪ π[0,5] (σ[(#1 = #2 ∧ #3 = #4)] ((P_aux × (Mid × P_aux)))))));
  Q := (π[0,2] (Tri) ∪ π[0,3] (σ[#1 = #2] ((Q_aux × P_aux))));
  WHILE (¬(((P = P_aux) ∧ (Q = Q_aux)))) DO
    P_aux := P;
    Q_aux := Q;
    P := (P ∪ (Link ∪ (Mid ∪ (π[0,3] (σ[#1 = #2] ((P_aux × Q_aux))) ∪ π[0,5] (σ[(#1 = #2 ∧ #3 = #4)] ((P_aux × (Mid × P_aux))))))));
    Q := (Q ∪ (π[0,2] (Tri) ∪ π[0,3] (σ[#1 = #2] ((Q_aux × P_aux)))))
  END;
  T_aux := ∅[2];
  G_aux := ∅[2];
  T := (BaseTwo ∪ (π[0,3] (σ[#1 = #2] ((T_aux × G_aux))) ∪ π[0,3] (σ[#1 = #2] ((T_aux × T_aux)))));
  G := (Jump ∪ π[0,3] (σ[#1 = #2] ((G_aux × Jump))));
  WHILE (¬(((T = T_aux) ∧ (G = G_aux)))) DO
    T_aux := T;
    G_aux := G;
    T := (T ∪ (BaseTwo ∪ (π[0,3] (σ[#1 = #2] ((T_aux × G_aux))) ∪ π[0,3] (σ[#1 = #2] ((T_aux × T_aux))))));
    G := (G ∪ (Jump ∪ π[0,3] (σ[#1 = #2] ((G_aux × Jump)))))
  END

POST
  (P ⊆ T)
\end{Verbatim}

\Needspace{0.45\textheight}

\section{Example1093: Lineage triples contained in a closed ternary cover}\label{case:Example1093}

\begin{description}[leftmargin=2.6cm,style=nextline,itemsep=1pt]

\item[Category] LLM-generated (GPT/Claude benchmark programs); family: two programs under TGD constraints (LLM-generated, Claude Opus; legacy raw\_new pipeline).

\item[Kind] containment.

\item[Contributors] Leo Zhang.

\item[Status] certified valid (Certificate/Valid.lean, written by the certificate emitter); expected valid.

\item[Tags] \hyperref[tag:dependencies]{\texttt{dependencies}} \hyperref[tag:mutual-recursion]{\texttt{mutual-recursion}} \hyperref[tag:nonlinear-recursion]{\texttt{nonlinear-recursion}}.

\item[Input] \path{Benchmark/Example1093/Input.lean} (canonical notation).

\item[Summary] Lineage triples contained in a closed ternary cover.

\item[Sides] naive compiler: program C1 of the retired corpus case Example1093 (output Lin) \par naive compiler: program C2 of the retired corpus case Example1093 (output Prov)

\item[Sources] {\raggedright LLM-generated (Claude Opus), legacy raw\_new TGD pipeline this repository the retired corpus case Example1093 (programs C1, C2 and the TGD list tgd0..) TGD pipeline run raw-new-tgd-0002-product: TGDs translated to relational algebra by RelCalc.TGD.listToAssertExpr; programs compiled with the certified naive compiler (read 2026-09-15) \par method this repository the equivalence-under-constraints paper draft (technical appendix: reduction of \path{containment/equivalence} under TGDs to a Hoare triple) (read 2026-09-15) \par legacy source record raw\_tgd\_opus\_w\_inv.json \par legacy source record \path{Benchmark/Example1093/Input.lean} \par Datalog source noted in the legacy inventory of Example1093\par}

\item[Obstruction note] from the metadata: anchored: both sides are naive closures at the same rate and the legacy product loop was certified with a quantifier-free invariant; no prophecy level expected

\item[Note] Legacy raw\_new TGD case (Claude Opus): the provenance cover is closed under propagation up to a re-projected coordinate (head-existential dependencies).

\end{description}

\paragraph{Datalog program(s).}

{\raggedright\textit{prog\_c1} --- fidelity module (kernel-checked: the compiled sides of the command are this program's naive compilation).\par}

\begin{Verbatim}[fontsize=\scriptsize,breaklines=true,breakanywhere=true,frame=single,framesep=1.5mm]
Reach(x1, y1) :- Edge(x1, y1).
Reach(x1, z1) :- Reach(x1, y1), Reach(y1, z1).
Reach(x1, y1) :- Lin(x1, y1, z1).
Lin(x1, y1, z1) :- Fact(x1, y1, z1).
Lin(x1, z1, z2) :- Lin(x1, y1, z2), Reach(y1, z1).
Lin(x2, y1, z1) :- Lin(x1, y1, z1), Gen(x1, x2).
\end{Verbatim}

{\raggedright\textit{prog\_c2} --- fidelity module (kernel-checked: the compiled sides of the command are this program's naive compilation).\par}

\begin{Verbatim}[fontsize=\scriptsize,breaklines=true,breakanywhere=true,frame=single,framesep=1.5mm]
Prov(x1, y1, z1) :- Cover(x1, y1, z1).
Trav(x1, y1) :- Edge(x1, y1).
Trav(x1, z1) :- Trav(x1, y1), Trav(y1, z1).
Prov(x1, z1, z2) :- Prov(x1, y1, z2), Trav(y1, z1).
Prov(x2, y1, z1) :- Prov(x1, y1, z1), Gen(x1, x2).
\end{Verbatim}

\paragraph{Reading.} Program C1 derives ternary lineage/provenance triples Lin(item, derived-item, source) by propagating an initial Fact triple along a computed reachability relation Reach (built from graph edges Edge and fed back by Lin's own second coordinate) and along generation edges Gen. Program C2 instead reads answers from Prov, defined via a ternary Cover relation the TGD constraints force to be closed under the same edge-traversal, generation-transfer, and self-composition steps, so Cover is guaranteed to dominate any lineage C1's rules could derive. The claim is the containment Lin ⊆ Prov under these dependencies, a chase-style containment check with an existential dependency (every Edge pair needs a Fact witness), expected valid.

\paragraph{First-order reading of PRE/POST, translated mechanically from the relational-algebra assertion in Input.lean.} The relational-algebra assertion in \texttt{Input.lean} remains the authority.

\textit{PRE:} (∀x y z. (Fact(x,y,z) → Cover(x,y,z))) ∧ (∀u v w. ((∃x1. Cover(u,x1,v) ∧ Edge(x1,w)) → Cover(u,w,v))) ∧ (∀x2 x3 x4. ((∃x5. Cover(x5,x2,x3) ∧ Gen(x5,x4)) → Cover(x4,x2,x3))) ∧ (∀x6 x7 x8. ((∃x9 x10. Cover(x6,x9,x7) ∧ Cover(x9,x8,x10)) → Cover(x6,x8,x7))) ∧ ∀x11 x12. (Edge(x11,x12) → ∃x13. Fact(x11,x12,x13))\par The precondition is five TGDs: Fact triples are all in Cover; Cover is closed under extending its second coordinate along Edge, under moving its first coordinate along Gen, and under composing two Cover triples through a shared middle value; and every Edge pair has some Fact triple witnessing it.

\textit{POST:} ∀x y z. (Lin(x,y,z) → Prov(x,y,z))\par The postcondition is the containment Lin ⊆ Prov: every lineage triple C1 derives is already present in C2's closed cover relation.

\paragraph{Whiel encoding (Hoare triple).}

\begin{Verbatim}[fontsize=\scriptsize,breaklines=true,breakanywhere=true,frame=single,framesep=1.5mm]
SCHEMA
  Edge/2, Gen/2, Reach/2, Trav/2, Reach_aux/2, Trav_aux/2, Fact/3, Cover/3, Lin/3, Prov/3, Lin_aux/3, Prov_aux/3

PRE
  ((Fact ⊆ Cover) ∧ ((π[0,2,4] (σ[#1 = #3] ((Cover × Edge))) ⊆ π[0,2,1] (Cover)) ∧ ((π[1,2,4] (σ[#0 = #3] ((Cover × Gen))) ⊆ π[1,2,0] (Cover)) ∧ ((π[0,2,4] (σ[#1 = #3] ((Cover × Cover))) ⊆ π[0,2,1] (Cover)) ∧ (Edge ⊆ π[0,1] (Fact))))))

CMD
  Reach_aux := ∅[2];
  Lin_aux := ∅[3];
  Reach := (Edge ∪ (π[0,3] (σ[#1 = #2] ((Reach_aux × Reach_aux))) ∪ π[0,1] (Lin_aux)));
  Lin := (Fact ∪ (π[0,4,2] (σ[#1 = #3] ((Lin_aux × Reach_aux))) ∪ π[4,1,2] (σ[#0 = #3] ((Lin_aux × Gen)))));
  WHILE (¬(((Reach = Reach_aux) ∧ (Lin = Lin_aux)))) DO
    Reach_aux := Reach;
    Lin_aux := Lin;
    Reach := (Reach ∪ (Edge ∪ (π[0,3] (σ[#1 = #2] ((Reach_aux × Reach_aux))) ∪ π[0,1] (Lin_aux))));
    Lin := (Lin ∪ (Fact ∪ (π[0,4,2] (σ[#1 = #3] ((Lin_aux × Reach_aux))) ∪ π[4,1,2] (σ[#0 = #3] ((Lin_aux × Gen))))))
  END;
  Prov_aux := ∅[3];
  Trav_aux := ∅[2];
  Prov := (Cover ∪ (π[0,4,2] (σ[#1 = #3] ((Prov_aux × Trav_aux))) ∪ π[4,1,2] (σ[#0 = #3] ((Prov_aux × Gen)))));
  Trav := (Edge ∪ π[0,3] (σ[#1 = #2] ((Trav_aux × Trav_aux))));
  WHILE (¬(((Prov = Prov_aux) ∧ (Trav = Trav_aux)))) DO
    Prov_aux := Prov;
    Trav_aux := Trav;
    Prov := (Prov ∪ (Cover ∪ (π[0,4,2] (σ[#1 = #3] ((Prov_aux × Trav_aux))) ∪ π[4,1,2] (σ[#0 = #3] ((Prov_aux × Gen))))));
    Trav := (Trav ∪ (Edge ∪ π[0,3] (σ[#1 = #2] ((Trav_aux × Trav_aux)))))
  END

POST
  (Lin ⊆ Prov)
\end{Verbatim}

\Needspace{0.45\textheight}

\section{Example1094: Routing table contained in an administratively closed distance view}\label{case:Example1094}

\begin{description}[leftmargin=2.6cm,style=nextline,itemsep=1pt]

\item[Category] LLM-generated (GPT/Claude benchmark programs); family: two programs under TGD constraints (LLM-generated, Claude Opus; legacy raw\_new pipeline).

\item[Kind] containment.

\item[Contributors] Leo Zhang.

\item[Status] certified valid (Certificate/Valid.lean, written by the certificate emitter); expected valid.

\item[Tags] \hyperref[tag:dependencies]{\texttt{dependencies}} \hyperref[tag:mutual-recursion]{\texttt{mutual-recursion}} \hyperref[tag:nonlinear-recursion]{\texttt{nonlinear-recursion}}.

\item[Input] \path{Benchmark/Example1094/Input.lean} (canonical notation).

\item[Summary] Routing table contained in an administratively closed distance view.

\item[Sides] naive compiler: program C1 of the retired corpus case Example1094 (output Route) \par naive compiler: program C2 of the retired corpus case Example1094 (output Dist)

\item[Sources] {\raggedright LLM-generated (Claude Opus), legacy raw\_new TGD pipeline this repository the retired corpus case Example1094 (programs C1, C2 and the TGD list tgd0..) TGD pipeline run raw-new-tgd-0002-product: TGDs translated to relational algebra by RelCalc.TGD.listToAssertExpr; programs compiled with the certified naive compiler (read 2026-09-15) \par method this repository the equivalence-under-constraints paper draft (technical appendix: reduction of \path{containment/equivalence} under TGDs to a Hoare triple) (read 2026-09-15) \par legacy source record raw\_tgd\_opus\_w\_inv.json \par legacy source record \path{Benchmark/Example1094/Input.lean} \par Datalog source noted in the legacy inventory of Example1094\par}

\item[Obstruction note] from the metadata: anchored: both sides are naive closures at the same rate and the legacy product loop was certified with a quantifier-free invariant; no prophecy level expected

\item[Note] Legacy raw\_new TGD case (Claude Opus): every peer pair has a relay (existential dependency); containment in an administratively closed distance view.

\end{description}

\paragraph{Datalog program(s).}

{\raggedright\textit{prog\_c1} --- fidelity module (kernel-checked: the compiled sides of the command are this program's naive compilation).\par}

\begin{Verbatim}[fontsize=\scriptsize,breaklines=true,breakanywhere=true,frame=single,framesep=1.5mm]
Route(x1, y1) :- Peer(x1, y1).
Hop(x1, y1) :- Adj(x1, y1).
Hop(x1, z1) :- Route(x1, y1), Hop(y1, z1).
Route(x1, z1) :- Route(x1, y1), Hop(y1, z1).
Route(x1, z1) :- Route(x1, y1), Relay(y1, x2, z1).
Route(x1, z1) :- Route(x1, y1), Route(y1, z1).
\end{Verbatim}

{\raggedright\textit{prog\_c2} --- fidelity module (kernel-checked: the compiled sides of the command are this program's naive compilation).\par}

\begin{Verbatim}[fontsize=\scriptsize,breaklines=true,breakanywhere=true,frame=single,framesep=1.5mm]
Dist(x1, y1) :- Admin(x1, y1).
Span(x1, y1) :- Adj(x1, y1).
Span(x1, z1) :- Span(x1, y1), Span(y1, z1).
Dist(x1, z1) :- Dist(x1, y1), Span(y1, z1).
Dist(x1, z1) :- Dist(x1, y1), Dist(y1, z1).
\end{Verbatim}

\paragraph{Reading.} Program C1 builds a routing table Route from direct peer links (Peer), adjacency hops composed through Hop, ternary relay legs (Relay), and transitive composition of Route with itself, a multi-hop route computation typical of network routing. Program C2 recomputes distances Dist by reading an administrative relation Admin that the TGD constraints force to contain Peer, Adj, and the endpoints of every Relay triple, and to be transitively closed, then closes Dist under Adj-based spans and self-composition. The claim is the containment Route ⊆ Dist under these dependencies (an existential dependency guarantees every Peer edge a Relay witness), expected valid because both sides are naive closures advancing at the same rate.

\paragraph{First-order reading of PRE/POST, translated mechanically from the relational-algebra assertion in Input.lean.} The relational-algebra assertion in \texttt{Input.lean} remains the authority.

\textit{PRE:} (∀x y. (Peer(x,y) → Admin(x,y))) ∧ (∀z u. (Adj(z,u) → Admin(z,u))) ∧ (∀v w. ((∃x1. Admin(v,x1) ∧ Admin(x1,w)) → Admin(v,w))) ∧ (∀x2 x3. ((∃x4. Relay(x2,x4,x3)) → Admin(x2,x3))) ∧ ∀x5 x6. (Peer(x5,x6) → ∃x7. Relay(x5,x7,x6))\par The precondition is five TGDs: Peer and Adj are contained in Admin, Admin is transitively closed, the endpoints of every Relay triple are in Admin, and every Peer edge has some Relay triple witnessing it.

\textit{POST:} ∀x y. (Route(x,y) → Dist(x,y))\par The postcondition is the containment Route ⊆ Dist: every route pair C1 computes is also a distance pair C2 computes.

\paragraph{Whiel encoding (Hoare triple).}

\begin{Verbatim}[fontsize=\scriptsize,breaklines=true,breakanywhere=true,frame=single,framesep=1.5mm]
SCHEMA
  Peer/2, Adj/2, Admin/2, Route/2, Hop/2, Dist/2, Span/2, Route_aux/2, Hop_aux/2, Dist_aux/2, Span_aux/2, Relay/3

PRE
  ((Peer ⊆ Admin) ∧ ((Adj ⊆ Admin) ∧ ((π[0,3] (σ[#1 = #2] ((Admin × Admin))) ⊆ Admin) ∧ ((π[0,2] (Relay) ⊆ Admin) ∧ (Peer ⊆ π[0,2] (Relay))))))

CMD
  Route_aux := ∅[2];
  Hop_aux := ∅[2];
  Route := (Peer ∪ (π[0,3] (σ[#1 = #2] ((Route_aux × Hop_aux))) ∪ (π[0,4] (σ[#1 = #2] ((Route_aux × Relay))) ∪ π[0,3] (σ[#1 = #2] ((Route_aux × Route_aux))))));
  Hop := (Adj ∪ π[0,3] (σ[#1 = #2] ((Route_aux × Hop_aux))));
  WHILE (¬(((Route = Route_aux) ∧ (Hop = Hop_aux)))) DO
    Route_aux := Route;
    Hop_aux := Hop;
    Route := (Route ∪ (Peer ∪ (π[0,3] (σ[#1 = #2] ((Route_aux × Hop_aux))) ∪ (π[0,4] (σ[#1 = #2] ((Route_aux × Relay))) ∪ π[0,3] (σ[#1 = #2] ((Route_aux × Route_aux)))))));
    Hop := (Hop ∪ (Adj ∪ π[0,3] (σ[#1 = #2] ((Route_aux × Hop_aux)))))
  END;
  Dist_aux := ∅[2];
  Span_aux := ∅[2];
  Dist := (Admin ∪ (π[0,3] (σ[#1 = #2] ((Dist_aux × Span_aux))) ∪ π[0,3] (σ[#1 = #2] ((Dist_aux × Dist_aux)))));
  Span := (Adj ∪ π[0,3] (σ[#1 = #2] ((Span_aux × Span_aux))));
  WHILE (¬(((Dist = Dist_aux) ∧ (Span = Span_aux)))) DO
    Dist_aux := Dist;
    Span_aux := Span;
    Dist := (Dist ∪ (Admin ∪ (π[0,3] (σ[#1 = #2] ((Dist_aux × Span_aux))) ∪ π[0,3] (σ[#1 = #2] ((Dist_aux × Dist_aux))))));
    Span := (Span ∪ (Adj ∪ π[0,3] (σ[#1 = #2] ((Span_aux × Span_aux)))))
  END

POST
  (Route ⊆ Dist)
\end{Verbatim}

\Needspace{0.45\textheight}

\section{Example2004: Magic-set transformation: demanded answers equal the original program's answers (Soufflé)}\label{case:Example2004}

\begin{description}[leftmargin=2.6cm,style=nextline,itemsep=1pt]

\item[Category] production (derived from a production system's source or documentation); family: Obstructed.

\item[Kind] equivalence (stated as mutual containments).

\item[Contributors] Leo Zhang.

\item[Status] certified valid (Certificate/Valid.lean, written by the certificate emitter); expected valid.

\item[Tags] \hyperref[tag:database-engine]{\texttt{database-engine}} \hyperref[tag:magic-sets]{\texttt{magic-sets}} \hyperref[tag:prophecy]{\texttt{prophecy}}.

\item[Input] \path{Benchmark/Example2004/Input.lean} (canonical notation).

\item[Summary] A magic-set transformation of a right-linear path program, against the program it was derived from.

\item[Sources] {\raggedright \href{https://souffle-lang.github.io/magicset}{Soufflé magic-set transformation (documentation)} \par the equivalence-and-containment paper draft (the pre-fixpoint witness method for equivalence and containment) \par retired-corpus source record production\_4 of the Production family \par retired-corpus input record production\_4 of the Production family \par Datalog source noted in the legacy inventory \path{reports/evidence/gen_examples.py}\par}

\item[Prophecy] kernel-checked certificate with a level-1 prophecy clause (Certificate/Valid.lean, Core.json).

\item[Note] The canonical re-encoding of the legacy obstructed case; the query seed Root is pinned to \{"root"\} by the precondition.

\end{description}

\paragraph{Datalog program(s).}

{\raggedright\textit{programM --- The magic-set transformation: demand and answers.} --- fidelity module (kernel-checked: the compiled sides of the command are this program's naive compilation).\par}

\begin{Verbatim}[fontsize=\scriptsize,breaklines=true,breakanywhere=true,frame=single,framesep=1.5mm]
MPath(x1) :- Root(x1).
MPath(y1) :- MPath(x1), Edge(x1, y1).
PathBF(x1, y1) :- MPath(x1), Edge(x1, y1).
PathBF(x1, x2) :- MPath(x1), Edge(x1, y1), PathBF(y1, x2).
\end{Verbatim}

{\raggedright\textit{programP --- The original right-linear path program.} --- fidelity module (kernel-checked: the compiled sides of the command are this program's naive compilation).\par}

\begin{Verbatim}[fontsize=\scriptsize,breaklines=true,breakanywhere=true,frame=single,framesep=1.5mm]
Path(x1, y1) :- Edge(x1, y1).
Path(x1, x2) :- Edge(x1, y1), Path(y1, x2).
\end{Verbatim}

\paragraph{Reading.} This case checks a magic-set transformation of a right-linear reachability program: the transformed program seeds a demand set MPath from Root and propagates it forward along Edge, then computes PathBF only for pairs whose source is in the demand set, while the original program Path computes unrestricted transitive closure of Edge. This is the standard magic-set/demand-driven query optimization used to avoid computing answers a Datalog engine will not need. The postcondition states both directions of correctness for the transformation restricted to the demand set: every original-program answer whose source is demanded is also produced by the transformed program, and every transformed-program answer is also an original-program answer; each relation starts empty and monotonically accumulates until it stops changing, so both hold the least fixpoints of their respective programs.

\paragraph{First-order reading of PRE/POST, translated mechanically from the relational-algebra assertion in Input.lean.} The relational-algebra assertion in \texttt{Input.lean} remains the authority.

\textit{PRE:} ∀x. (Root(x) ↔ x = "root")\par The precondition pins the query's demand seed to the single source "root", i.e. Root = \{"root"\}.

\textit{POST:} (∀x y. (MPath(x) ∧ Path(x,y) → PathBF(x,y))) ∧ ∀z u. (PathBF(z,u) → Path(z,u))\par The postcondition asserts both completeness and soundness of the transformation on the demanded sources: every Path pair whose source is in MPath is also a PathBF pair, and every PathBF pair is also a Path pair.

\paragraph{Whiel encoding (Hoare triple).}

\begin{Verbatim}[fontsize=\tiny,breaklines=true,breakanywhere=true,frame=single,framesep=1.5mm]
SCHEMA
  Root/1, MPath/1, MPath_aux/1, Edge/2, PathBF/2, PathBF_aux/2, Path/2, Path_aux/2

PRE
  (Root = { "root" }) 

CMD
  MPath_aux := ∅[1];
  PathBF_aux := ∅[2];
  MPath :=
    (Root ∪
      (π[2] (σ[#0 = #1] (MPath_aux × Edge))));
  PathBF :=
    ((π[0, 2]
        (σ[#0 = #1] (MPath_aux × Edge))) ∪
      (π[0, 4]
        (σ[(#0 = #1) ∧ (#2 = #3)]
          (MPath_aux × (Edge × PathBF_aux)))));
  WHILE
      (¬ ((MPath = MPath_aux) ∧
        (PathBF = PathBF_aux))) DO
    MPath_aux := MPath;
    PathBF_aux := PathBF;
    MPath :=
      (MPath ∪
        (Root ∪
          (π[2]
            (σ[#0 = #1] (MPath_aux × Edge)))));
    PathBF :=
      (PathBF ∪
        ((π[0, 2]
            (σ[#0 = #1] (MPath_aux × Edge))) ∪
          (π[0, 4]
            (σ[(#0 = #1) ∧ (#2 = #3)]
              (MPath_aux ×
                (Edge × PathBF_aux))))))
  END;
  Path_aux := ∅[2];
  Path :=
    (Edge ∪
      (π[0, 3] (σ[#1 = #2] (Edge × Path_aux))));
  WHILE (¬ (Path = Path_aux)) DO
    Path_aux := Path;
    Path :=
      (Path ∪
        (Edge ∪
          (π[0, 3]
            (σ[#1 = #2] (Edge × Path_aux)))))
  END

POST
  ((π[1, 2] (σ[#0 = #1] (MPath × Path)))
    ⊆ PathBF) ∧
  (PathBF ⊆ Path)
\end{Verbatim}

\Needspace{0.45\textheight}

\section{Example2005: Restricted-edge reachability is contained in full-edge reachability (edge-monotonicity)}\label{case:Example2005}

\begin{description}[leftmargin=2.6cm,style=nextline,itemsep=1pt]

\item[Category] production (derived from a production system's source or documentation); family: Production.

\item[Kind] containment (via a pre-fixpoint witness).

\item[Contributors] Leo Zhang.

\item[Status] certified valid (Certificate/Valid.lean, written by the certificate emitter); expected valid.

\item[Tags] \hyperref[tag:fixpoint-witness]{\texttt{fixpoint-witness}} \hyperref[tag:negation]{\texttt{negation}}.

\item[Input] \path{Benchmark/Example2005/Input.lean} (canonical notation).

\item[Summary] containment\_03: Restricted-edge reachability is contained in full-edge reachability (edge-monotonicity).

\item[Sources] {\raggedright \href{https://github.com/npm/arborist/blob/main/lib/calc-dep-flags.js}{npm/arborist lib/calc-dep-flags.js} \par the equivalence-and-containment paper draft (the pre-fixpoint witness method for equivalence and containment) \par retired-corpus source record production\_5 of the Production family \par retired-corpus input record production\_5 of the Production family\par}

\end{description}

\paragraph{Datalog program(s).}

{\raggedright\textit{prog\_p} --- production check (the hand-written loop equals this program's naive compilation on the check instances).\par}

\begin{Verbatim}[fontsize=\scriptsize,breaklines=true,breakanywhere=true,frame=single,framesep=1.5mm]
NonOptC(x1) :- RootN(x1).
NonOptC(y1) :- NonOptC(x1), EdgeHardC(x1, y1).
\end{Verbatim}

\paragraph{Reading.} The program computes graph reachability from a root set RootN restricted to "hard" (non-optional) edges only, modeling the analysis a dependency resolver like npm's Arborist uses to decide which packages keep a required (non-optional) dependency path from the root. The Hoare triple is a containment claim rather than full equivalence: the precondition makes RBound a pre-fixpoint of the reachability operator over the unrestricted edge set EdgeAll, and the postcondition asserts the hard-edge-only reachable set NonOpt is contained in RBound, i.e. contained in full-edge reachability, the standard edge-monotonicity fact that removing edges (here, optional ones) can only shrink what is reachable.

\paragraph{First-order reading of PRE/POST, translated mechanically from the relational-algebra assertion in Input.lean.} The relational-algebra assertion in \texttt{Input.lean} remains the authority.

\textit{PRE:} ∀x. (RootN(x) ∨ (∃y. EdgeAll(y,x) ∧ RBound(y)) → RBound(x))\par The precondition asserts RBound contains RootN and is closed under one more step of reachability through the full edge set EdgeAll, i.e. RBound is a pre-fixpoint of full-edge reachability.

\textit{POST:} ∀x. (NonOpt(x) → RBound(x))\par The postcondition asserts the hard-edge-restricted reachable set NonOpt is contained in RBound.

\paragraph{Whiel encoding (Hoare triple).}

\begin{Verbatim}[fontsize=\scriptsize,breaklines=true,breakanywhere=true,frame=single,framesep=1.5mm]
SCHEMA
  RootN/1, NonOpt/1, NonOpt_aux/1, RBound/1, EdgeAll/2, EdgeOpt/2, EdgeHard/2

PRE
  ((RootN ∪ (π[1] (σ[#0 = #2] (EdgeAll × RBound)))) ⊆ RBound)

CMD
  NonOpt_aux := ∅;
  EdgeHard := (EdgeAll ∖ EdgeOpt);
  NonOpt := RootN;
  WHILE NonOpt ≠ NonOpt_aux DO
    NonOpt_aux := NonOpt;
    NonOpt := (RootN ∪
      (π[1] (σ[#0 = #2] (EdgeHard × NonOpt_aux))))
  END

POST
  (NonOpt ⊆ RBound)
\end{Verbatim}

\Needspace{0.45\textheight}

\section{Example2006: Policy-filtered connectivity is contained in raw topological connectivity}\label{case:Example2006}

\begin{description}[leftmargin=2.6cm,style=nextline,itemsep=1pt]

\item[Category] production (derived from a production system's source or documentation); family: Production.

\item[Kind] containment (via a pre-fixpoint witness).

\item[Contributors] Leo Zhang.

\item[Status] certified valid (Certificate/Valid.lean, written by the certificate emitter); expected valid.

\item[Tags] \hyperref[tag:fixpoint-witness]{\texttt{fixpoint-witness}} \hyperref[tag:negation]{\texttt{negation}}.

\item[Input] \path{Benchmark/Example2006/Input.lean} (canonical notation).

\item[Summary] containment\_04: Policy-filtered connectivity is contained in raw topological connectivity.

\item[Sources] {\raggedright \href{https://docs.cilium.io/en/stable/security/policy/}{Cilium network policy (documentation)} \par \href{https://github.com/batfish/batfish}{batfish/batfish (network configuration analysis)} \par the equivalence-and-containment paper draft (the pre-fixpoint witness method for equivalence and containment) \par retired-corpus source record production\_6 of the Production family \par retired-corpus input record production\_6 of the Production family\par}

\end{description}

\paragraph{Datalog program(s).}

{\raggedright\textit{prog\_p} --- production check (the hand-written loop equals this program's naive compilation on the check instances).\par}

\begin{Verbatim}[fontsize=\scriptsize,breaklines=true,breakanywhere=true,frame=single,framesep=1.5mm]
ConnC(x1, y1) :- AllowedC(x1, y1).
ConnC(x1, z1) :- ConnC(x1, y1), AllowedC(y1, z1).
\end{Verbatim}

\paragraph{Reading.} The program computes reachability (Conn) over an edge set filtered by removing denied links (LinkE ∖ Denied), modeling how a network-policy engine like Cilium restricts which nodes remain connected under an access policy, contrasted with the unrestricted topological reachability a tool like Batfish would compute over the full link set LinkE. The Hoare triple is a containment claim: the precondition makes RBound a pre-fixpoint of the reachability operator over the full, unfiltered edge set LinkE, and the postcondition asserts the policy-filtered connectivity Conn is contained in RBound, i.e. contained in unrestricted reachability, since network policies can only remove connectivity, never add it.

\paragraph{First-order reading of PRE/POST, translated mechanically from the relational-algebra assertion in Input.lean.} The relational-algebra assertion in \texttt{Input.lean} remains the authority.

\textit{PRE:} ∀x y. (LinkE(x,y) ∨ (∃z. RBound(x,z) ∧ LinkE(z,y)) → RBound(x,y))\par The precondition asserts RBound contains LinkE and is closed under one more step of composing itself with LinkE, i.e. RBound is a pre-fixpoint of full-topology reachability.

\textit{POST:} ∀x y. (Conn(x,y) → RBound(x,y))\par The postcondition asserts the policy-filtered connectivity relation Conn is contained in RBound.

\paragraph{Whiel encoding (Hoare triple).}

\begin{Verbatim}[fontsize=\scriptsize,breaklines=true,breakanywhere=true,frame=single,framesep=1.5mm]
SCHEMA
  LinkE/2, Denied/2, AllowedE/2, Conn/2, Conn_aux/2, RBound/2

PRE
  ((LinkE ∪ (π[0, 3] (σ[#1 = #2] (RBound × LinkE)))) ⊆ RBound)

CMD
  Conn_aux := ∅;
  AllowedE := (LinkE ∖ Denied);
  Conn := (LinkE ∖ Denied);
  WHILE Conn ≠ Conn_aux DO
    Conn_aux := Conn;
    Conn := (AllowedE ∪
      (π[0, 3] (σ[#1 = #2] (Conn_aux × AllowedE))))
  END

POST
  (Conn ⊆ RBound)
\end{Verbatim}

\Needspace{0.45\textheight}

\section{Example2007: E[f U g] is contained in EF g}\label{case:Example2007}

\begin{description}[leftmargin=2.6cm,style=nextline,itemsep=1pt]

\item[Category] production (derived from a production system's source or documentation); family: Production.

\item[Kind] containment (via a pre-fixpoint witness).

\item[Contributors] Leo Zhang.

\item[Status] certified valid (Certificate/Valid.lean, written by the certificate emitter); expected valid.

\item[Tags] \hyperref[tag:fixpoint-witness]{\texttt{fixpoint-witness}} \hyperref[tag:semi-naive]{\texttt{semi-naive}}.

\item[Input] \path{Benchmark/Example2007/Input.lean} (canonical notation).

\item[Summary] containment\_05: E[f U g] is contained in EF g.

\item[Sources] {\raggedright \href{https://github.com/hklarner/NuSMV-a}{NuSMV 2.6.0 core/mc/mcMc.c, eu() (hklarner/NuSMV-a mirror)} \par \href{https://nusmv.fbk.eu/userman/v21/nusmv_2.html}{NuSMV user manual, CTL model checking} \par the equivalence-and-containment paper draft (the pre-fixpoint witness method for equivalence and containment) \par retired-corpus source record production\_7 of the Production family \par retired-corpus input record production\_7 of the Production family\par}

\end{description}

\paragraph{Datalog program(s).}

{\raggedright\textit{prog\_p} --- production check (the hand-written loop equals this program's naive compilation on the check instances).\par}

\begin{Verbatim}[fontsize=\scriptsize,breaklines=true,breakanywhere=true,frame=single,framesep=1.5mm]
YC(x1) :- SatG(x1).
YC(x1) :- SatF(x1), Trans(x1, y1), YC(y1).
\end{Verbatim}

\paragraph{Reading.} The loop computes a CTL-style backward-reachability fixpoint: starting from states satisfying g (SatG) and repeatedly adding predecessors along Trans that also satisfy f (SatF), it computes the states satisfying the until formula E[f U g], mirroring the frontier-based eu() computation in a model checker such as NuSMV. The Hoare triple is a containment claim: the precondition makes RBound a pre-fixpoint of the unconstrained backward-reachability operator that computes EF g (reachability with no f-restriction), and the postcondition asserts the E[f U g] result Y is contained in RBound, the textbook CTL fact that any until-witness for g is in particular an eventually-witness for g.

\paragraph{First-order reading of PRE/POST, translated mechanically from the relational-algebra assertion in Input.lean.} The relational-algebra assertion in \texttt{Input.lean} remains the authority.

\textit{PRE:} ∀x. (SatG(x) ∨ (∃y. Trans(x,y) ∧ RBound(y)) → RBound(x))\par The precondition asserts RBound contains SatG and is closed under taking one more step of predecessors through Trans, i.e. RBound is a pre-fixpoint of the operator computing EF g.

\textit{POST:} ∀x. (Y(x) → RBound(x))\par The postcondition asserts the computed until-set Y is contained in RBound.

\paragraph{Whiel encoding (Hoare triple).}

\begin{Verbatim}[fontsize=\scriptsize,breaklines=true,breakanywhere=true,frame=single,framesep=1.5mm]
SCHEMA
  SatF/1, SatG/1, Y/1, Yold/1, NewS/1, PreS/1, RBound/1, Trans/2

PRE
  ((SatG ∪ (π[0] (σ[#1 = #2] (Trans × RBound)))) ⊆ RBound)

CMD
  Y := SatG;
  NewS := SatG;
  WHILE NewS ≠ ∅ DO
    Yold := Y;
    PreS := (π[0] (σ[#1 = #2] (Trans × NewS)));
    Y := (Y ∪ (SatF ∖ (SatF ∖ PreS)));
    NewS := (Y ∖ Yold)
  END

POST
  (Y ⊆ RBound)
\end{Verbatim}

\Needspace{0.45\textheight}

\section{Example2010: Semi-naive evaluation ≡ naive/declarative transitive closure}\label{case:Example2010}

\begin{description}[leftmargin=2.6cm,style=nextline,itemsep=1pt]

\item[Category] production (derived from a production system's source or documentation); family: Production.

\item[Kind] equivalence (with a least fixpoint, via a pre-fixpoint witness).

\item[Contributors] Leo Zhang.

\item[Status] certified valid (Certificate/Valid.lean, written by the certificate emitter); expected valid.

\item[Tags] \hyperref[tag:database-engine]{\texttt{database-engine}} \hyperref[tag:evaluation-strategies]{\texttt{evaluation-strategies}} \hyperref[tag:fixpoint-witness]{\texttt{fixpoint-witness}} \hyperref[tag:semi-naive]{\texttt{semi-naive}}.

\item[Input] \path{Benchmark/Example2010/Input.lean} (canonical notation).

\item[Summary] equivalence\_01: Semi-naive evaluation ≡ naive/declarative transitive closure.

\item[Sources] {\raggedright \href{https://github.com/postgres/postgres/blob/master/src/backend/executor/nodeRecursiveunion.c}{postgres/postgres src/backend/executor/nodeRecursiveunion.c (ExecRecursiveUnion)} \par \href{https://souffle-lang.github.io/translate}{Soufflé translation to C++ / semi-naive loop (documentation)} \par the equivalence-and-containment paper draft (the pre-fixpoint witness method for equivalence and containment) \par retired-corpus source record production\_10 of the Production family \par retired-corpus input record production\_10 of the Production family\par}

\end{description}

\paragraph{Datalog program(s).}

{\raggedright\textit{prog\_p} --- production check (the hand-written loop equals this program's naive compilation on the check instances).\par}

\begin{Verbatim}[fontsize=\scriptsize,breaklines=true,breakanywhere=true,frame=single,framesep=1.5mm]
ResC(x1, y1) :- Edge(x1, y1).
ResC(x1, z1) :- ResC(x1, y1), Edge(y1, z1).
\end{Verbatim}

\paragraph{Reading.} The loop computes transitive closure of Edge using a semi-naive evaluation strategy: it maintains an accumulated result Res and a frontier WorkT of newly derived pairs, extends WorkT through Edge, keeps only pairs not already in Res, and terminates when the frontier is empty, the standard semi-naive optimization of naive bottom-up Datalog evaluation, which would otherwise recompute the full join every round. The Hoare triple is the least-fixpoint correctness statement: the precondition makes RBound a pre-fixpoint of the naive transitive-closure operator, and the postcondition shows Res is itself a pre-fixpoint of that operator and is contained in RBound, which together pin Res to the least fixpoint, i.e. semi-naive evaluation computes exactly the same transitive closure as naive evaluation.

\paragraph{First-order reading of PRE/POST, translated mechanically from the relational-algebra assertion in Input.lean.} The relational-algebra assertion in \texttt{Input.lean} remains the authority.

\textit{PRE:} ∀x y. (Edge(x,y) ∨ (∃z. RBound(x,z) ∧ Edge(z,y)) → RBound(x,y))\par The precondition asserts RBound contains Edge and is closed under one more step of composing itself with Edge, i.e. RBound is a pre-fixpoint of the naive transitive-closure operator.

\textit{POST:} (∀x y. (Edge(x,y) ∨ (∃z. Res(x,z) ∧ Edge(z,y)) → Res(x,y))) ∧ ∀u v. (Res(u,v) → RBound(u,v))\par The postcondition asserts the semi-naive result Res is itself a pre-fixpoint of the same operator (contains Edge and is closed under composing with Edge) and is contained in RBound, together forcing Res to equal the least fixpoint.

\paragraph{Whiel encoding (Hoare triple).}

\begin{Verbatim}[fontsize=\scriptsize,breaklines=true,breakanywhere=true,frame=single,framesep=1.5mm]
SCHEMA
  Edge/2, Res/2, WorkT/2, Inter/2, RBound/2

PRE
  ((Edge ∪ (π[0, 3] (σ[#1 = #2] (RBound × Edge)))) ⊆ RBound)

CMD
  Res := Edge;
  WorkT := Edge;
  WHILE WorkT ≠ ∅ DO
    Inter := ((π[0, 3] (σ[#1 = #2] (WorkT × Edge))) ∖ Res);
    Res := (Res ∪ Inter);
    WorkT := Inter
  END

POST
  ((Edge ∪ (π[0, 3] (σ[#1 = #2] (Res × Edge)))) ⊆ Res) ∧
  (Res ⊆ RBound)
\end{Verbatim}

\Needspace{0.45\textheight}

\section{Example2011: Incremental view maintenance ≡ from-scratch recomputation}\label{case:Example2011}

\begin{description}[leftmargin=2.6cm,style=nextline,itemsep=1pt]

\item[Category] production (derived from a production system's source or documentation); family: Production.

\item[Kind] equivalence (with a least fixpoint, via a pre-fixpoint witness).

\item[Contributors] Leo Zhang.

\item[Status] certified valid (Certificate/Valid.lean, written by the certificate emitter); expected valid.

\item[Tags] \hyperref[tag:evaluation-strategies]{\texttt{evaluation-strategies}} \hyperref[tag:fixpoint-witness]{\texttt{fixpoint-witness}} \hyperref[tag:view-maintenance]{\texttt{view-maintenance}}.

\item[Input] \path{Benchmark/Example2011/Input.lean} (canonical notation).

\item[Summary] equivalence\_02: Incremental view maintenance ≡ from-scratch recomputation.

\item[Sources] {\raggedright \href{https://github.com/TimelyDataflow/differential-dataflow}{TimelyDataflow/differential-dataflow} \par \href{https://materialize.com/docs/get-started/arrangements/}{Materialize arrangements / incremental dataflows (documentation)} \par the equivalence-and-containment paper draft (the pre-fixpoint witness method for equivalence and containment) \par retired-corpus source record production\_11 of the Production family \par retired-corpus input record production\_11 of the Production family\par}

\item[Note] Models insert-only re-evaluation restarted from the maintained view; Differential Dataflow/Materialize are the motivation, not the mechanism.

\end{description}

\paragraph{Datalog program(s).}

{\raggedright\textit{prog\_p} --- production check (the hand-written loop equals this program's naive compilation on the check instances).\par}

\begin{Verbatim}[fontsize=\scriptsize,breaklines=true,breakanywhere=true,frame=single,framesep=1.5mm]
ReachC(x1) :- Roots(x1).
ReachC(x1) :- ReachOld(x1).
ReachC(y1) :- ReachC(x1), EBase(x1, y1).
ReachC(y1) :- ReachC(x1), EDelta(x1, y1).
\end{Verbatim}

\paragraph{Reading.} The loop computes reachability incrementally: it restarts from a previously maintained view ReachOld and re-propagates from Roots and ReachOld through the updated edge set EBase ∪ EDelta until stable, modeling insert-only incremental view maintenance in a streaming or materialized-view system (the case notes this is simplified insert-only re-evaluation, not full timestamped differential-dataflow propagation, though the claim being checked is the guarantee that differential dataflow systems provide). The Hoare triple is the least-fixpoint correctness statement: the precondition constrains ReachOld to be a pre-fixpoint of the old edge set's operator and RBound to be a pre-fixpoint of the updated operator with ReachOld ⊆ RBound standing in for ReachOld's minimality, and the postcondition shows the incrementally recomputed Reach is itself a pre-fixpoint of the updated operator and is contained in RBound, together pinning Reach to the least fixpoint over EBase ∪ EDelta, i.e. the incremental update yields exactly what a from-scratch recomputation would.

\paragraph{First-order reading of PRE/POST, translated mechanically from the relational-algebra assertion in Input.lean.} The relational-algebra assertion in \texttt{Input.lean} remains the authority.

\textit{PRE:} (∀x. (Roots(x) ∨ (∃y. EBase(y,x) ∧ ReachOld(y)) → ReachOld(x))) ∧ (∀z. (Roots(z) ∨ (∃u. (EBase(u,z) ∨ EDelta(u,z)) ∧ RBound(u)) → RBound(z))) ∧ ∀v. (ReachOld(v) → RBound(v))\par The precondition asserts ReachOld is a pre-fixpoint of reachability over EBase alone, and that RBound is a pre-fixpoint of reachability over the updated edge set EBase ∪ EDelta and contains ReachOld.

\textit{POST:} (∀x. (Roots(x) ∨ (∃y. (EBase(y,x) ∨ EDelta(y,x)) ∧ Reach(y)) → Reach(x))) ∧ ∀z. (Reach(z) → RBound(z))\par The postcondition asserts the incrementally updated Reach is itself a pre-fixpoint of reachability over EBase ∪ EDelta and is contained in RBound, together forcing Reach to equal the least fixpoint over the updated edges.

\paragraph{Whiel encoding (Hoare triple).}

\begin{Verbatim}[fontsize=\scriptsize,breaklines=true,breakanywhere=true,frame=single,framesep=1.5mm]
SCHEMA
  Roots/1, ReachOld/1, Reach/1, Reach_aux/1, RBound/1, EBase/2, EDelta/2

PRE
  ((Roots ∪ (π[1] (σ[#0 = #2] (EBase × ReachOld)))) ⊆ ReachOld) ∧
  (((Roots ∪ (π[1] (σ[#0 = #2] ((EBase ∪ EDelta) × RBound)))) ⊆ RBound) ∧
   (ReachOld ⊆ RBound))

CMD
  Reach_aux := ∅;
  Reach := (ReachOld ∪ Roots);
  WHILE Reach ≠ Reach_aux DO
    Reach_aux := Reach;
    Reach := ((Roots ∪ ReachOld) ∪
      (π[1] (σ[#0 = #2] ((EBase ∪ EDelta) × Reach_aux))))
  END

POST
  ((Roots ∪ (π[1] (σ[#0 = #2] ((EBase ∪ EDelta) × Reach)))) ⊆ Reach) ∧
  (Reach ⊆ RBound)
\end{Verbatim}

\Needspace{0.45\textheight}

\section{Example2012: Frontier (grey-set) marking ≡ naive mark propagation}\label{case:Example2012}

\begin{description}[leftmargin=2.6cm,style=nextline,itemsep=1pt]

\item[Category] production (derived from a production system's source or documentation); family: Production.

\item[Kind] equivalence (with a least fixpoint, via a pre-fixpoint witness).

\item[Contributors] Leo Zhang.

\item[Status] certified valid (Certificate/Valid.lean, written by the certificate emitter); expected valid.

\item[Tags] \hyperref[tag:evaluation-strategies]{\texttt{evaluation-strategies}} \hyperref[tag:fixpoint-witness]{\texttt{fixpoint-witness}} \hyperref[tag:semi-naive]{\texttt{semi-naive}}.

\item[Input] \path{Benchmark/Example2012/Input.lean} (canonical notation).

\item[Summary] equivalence\_03: Frontier (grey-set) marking ≡ naive mark propagation.

\item[Sources] {\raggedright \href{https://github.com/golang/go/blob/master/src/runtime/mgcmark.go}{golang/go src/runtime/mgcmark.go (gcDrain, scanobject, greyobject)} \par the equivalence-and-containment paper draft (the pre-fixpoint witness method for equivalence and containment) \par retired-corpus source record production\_12 of the Production family \par retired-corpus input record production\_12 of the Production family\par}

\end{description}

\paragraph{Datalog program(s).}

{\raggedright\textit{prog\_p} --- production check (the hand-written loop equals this program's naive compilation on the check instances).\par}

\begin{Verbatim}[fontsize=\scriptsize,breaklines=true,breakanywhere=true,frame=single,framesep=1.5mm]
MarkedC(x1) :- RootPtr(x1).
MarkedC(y1) :- MarkedC(x1), HeapPtr(x1, y1).
\end{Verbatim}

\paragraph{Reading.} The loop computes reachability from root pointers RootPtr through heap pointers HeapPtr using an explicit grey frontier Grey that is repeatedly advanced and deduplicated against the accumulated Marked set, the tri-color/frontier marking strategy used by a garbage collector's mark phase (e.g. a Go-runtime-style gcDrain), contrasted with naively re-propagating over the whole marked set each round. The Hoare triple is the least-fixpoint correctness statement: the precondition makes RBound a pre-fixpoint of the naive mark-propagation operator, and the postcondition shows the frontier-computed Marked set is itself a pre-fixpoint of that operator and is contained in RBound, together pinning Marked to the least fixpoint, i.e. frontier-based marking with its already-marked dedup marks exactly the same objects naive re-scanning would.

\paragraph{First-order reading of PRE/POST, translated mechanically from the relational-algebra assertion in Input.lean.} The relational-algebra assertion in \texttt{Input.lean} remains the authority.

\textit{PRE:} ∀x. (RootPtr(x) ∨ (∃y. HeapPtr(y,x) ∧ RBound(y)) → RBound(x))\par The precondition asserts RBound contains RootPtr and is closed under one more step of composing itself with HeapPtr, i.e. RBound is a pre-fixpoint of the naive mark-propagation operator.

\textit{POST:} (∀x. (RootPtr(x) ∨ (∃y. HeapPtr(y,x) ∧ Marked(y)) → Marked(x))) ∧ ∀z. (Marked(z) → RBound(z))\par The postcondition asserts the frontier-computed Marked set is itself a pre-fixpoint of the same operator (contains RootPtr and is closed under HeapPtr) and is contained in RBound, together forcing Marked to equal the least fixpoint.

\paragraph{Whiel encoding (Hoare triple).}

\begin{Verbatim}[fontsize=\scriptsize,breaklines=true,breakanywhere=true,frame=single,framesep=1.5mm]
SCHEMA
  RootPtr/1, Marked/1, Grey/1, NewG/1, RBound/1, HeapPtr/2

PRE
  ((RootPtr ∪ (π[1] (σ[#0 = #2] (HeapPtr × RBound)))) ⊆ RBound)

CMD
  Marked := RootPtr;
  Grey := RootPtr;
  WHILE Grey ≠ ∅ DO
    NewG := ((π[1] (σ[#0 = #2] (HeapPtr × Grey))) ∖ Marked);
    Marked := (Marked ∪ NewG);
    Grey := NewG
  END

POST
  ((RootPtr ∪ (π[1] (σ[#0 = #2] (HeapPtr × Marked)))) ⊆ Marked) ∧
  (Marked ⊆ RBound)
\end{Verbatim}

\Needspace{0.45\textheight}

\section{Example3001: Audited delegation alternation loses an uncertified reviewed hop}\label{case:Example3001}

\begin{description}[leftmargin=2.6cm,style=nextline,itemsep=1pt]

\item[Category] LLM-generated (GPT/Claude benchmark programs); family: two programs under TGD constraints (LLM-generated, GPT (raw\_invalid\_gpt batch); legacy raw\_new pipeline), invalid.

\item[Kind] containment.

\item[Contributors] Leo Zhang.

\item[Status] certified invalid (Certificate/Invalid.lean, written by the certificate emitter); expected invalid.

\item[Tags] \hyperref[tag:dependencies]{\texttt{dependencies}} \hyperref[tag:mutual-recursion]{\texttt{mutual-recursion}} \hyperref[tag:nonlinear-recursion]{\texttt{nonlinear-recursion}}.

\item[Input] \path{Benchmark/Example3001/Input.lean} (canonical notation).

\item[Summary] Audited delegation alternation loses an uncertified reviewed hop.

\item[Sides] naive compiler: program C1 of the retired corpus case Example3001 (output AccessOne) \par naive compiler: program C2 of the retired corpus case Example3001 (output AccessTwo)

\item[Sources] {\raggedright LLM-generated (GPT (raw\_invalid\_gpt batch)), legacy raw\_new TGD pipeline this repository the retired corpus case Example3001 (programs C1, C2 and the TGD list tgd0..) TGD pipeline run raw-new-tgd-0002-product: TGDs translated to relational algebra by RelCalc.TGD.listToAssertExpr; programs compiled with the certified naive compiler (read 2026-09-15) \par method this repository the equivalence-under-constraints paper draft (technical appendix: reduction of \path{containment/equivalence} under TGDs to a Hoare triple) (read 2026-09-15) \par legacy source record raw\_invalid\_gpt.json \par legacy source record \path{Benchmark/Example3001/Input.lean} \par Datalog source noted in the legacy inventory of Example3001\par}

\item[Obstruction note] from the metadata: invalid (legacy kernel counterexample)

\item[Note] Legacy raw\_new invalid TGD case (GPT): the second program demands a certificate on reviewed hops that the dependencies do not guarantee.

\end{description}

\paragraph{Datalog program(s).}

{\raggedright\textit{prog\_c1} --- fidelity module (kernel-checked: the compiled sides of the command are this program's naive compilation).\par}

\begin{Verbatim}[fontsize=\scriptsize,breaklines=true,breakanywhere=true,frame=single,framesep=1.5mm]
AccessOne(x1) :- Root(x1).
CertStateOne(y1) :- AccessOne(x1), Certified(x1, y1).
ReviewStateOne(y1) :- AccessOne(x1), Review(x1, y1, z1), Reviewer(z1).
AccessOne(y1) :- CertStateOne(x1), Review(x1, y1, z1), Reviewer(z1).
AccessOne(y1) :- ReviewStateOne(x1), Certified(x1, y1).
AccessOne(y1) :- CertStateOne(y1), ReviewStateOne(y1), Delegates(x1, y1).
\end{Verbatim}

{\raggedright\textit{prog\_c2} --- fidelity module (kernel-checked: the compiled sides of the command are this program's naive compilation).\par}

\begin{Verbatim}[fontsize=\scriptsize,breaklines=true,breakanywhere=true,frame=single,framesep=1.5mm]
AccessTwo(x1) :- Root(x1).
CertStateTwo(y1) :- AccessTwo(x1), Certified(x1, y1).
ReviewStateTwo(y1) :- AccessTwo(x1), Review(x1, y1, z1), Reviewer(z1).
JointStateTwo(y1) :- CertStateTwo(x1), Review(x1, y1, z1), Reviewer(z1), Certified(x1, y1).
AccessTwo(y1) :- JointStateTwo(y1), Delegates(x1, y1).
AccessTwo(y1) :- CertStateTwo(y1), ReviewStateTwo(y1), Delegates(x1, y1).
\end{Verbatim}

\paragraph{Reading.} Program C1 computes an access set AccessOne via delegation that alternates between a certified-state check and a reviewer-backed review state, allowing a step whenever either state alone combines with a background Delegates edge, and additionally whenever a certified hop and a reviewed hop chain together (certified-then-reviewed or reviewed-then-certified). Program C2 computes AccessTwo with an analogous structure but replaces that alternation with a stricter JointStateTwo requiring a reviewed hop to also be independently certified before it can combine with Delegates. This models an authorization engine comparing an alternating certified/reviewed delegation policy against a stricter one that insists a reviewed transition also be certified. The Hoare triple claims containment AccessOne ⊆ AccessTwo under five TGDs relating Certified, Review, Reviewer, and Delegates, but the claim fails: the case records an explicit witness instance on which C1's alternating chaining reaches a principal that C2's stricter jointly-certified-and-reviewed requirement cannot reach, because C2 loses a deeply authorized principal that is only reachable through an uncertified reviewed hop.

\paragraph{First-order reading of PRE/POST, translated mechanically from the relational-algebra assertion in Input.lean.} The relational-algebra assertion in \texttt{Input.lean} remains the authority.

\textit{PRE:} (∀x y. (Certified(x,y) → Delegates(x,y))) ∧ (∀z u. ((∃v. Review(z,u,v)) → Delegates(z,u))) ∧ (∀w. ((∃x1 x2. Review(x1,x2,w)) → Reviewer(w))) ∧ (∀x3 x4. (Certified(x3,x4) → ∃x5. Review(x3,x4,x5) ∧ Reviewer(x5))) ∧ ∀x6 x7. ((∃x8 x9. Certified(x6,x8) ∧ Review(x8,x7,x9)) → Delegates(x6,x7))\par The precondition is five TGDs: Certified and the first two coordinates of Review are contained in Delegates, the third coordinate of Review is a Reviewer, every Certified pair has a Review/Reviewer witness, and composing a Certified pair with a Review triple through their shared value lands in Delegates.

\textit{POST:} ∀x. (AccessOne(x) → AccessTwo(x))\par The postcondition claims AccessOne ⊆ AccessTwo, but the case is invalid: the recorded witness instance shows a principal that C1's alternating certified/reviewed delegation reaches while C2's stricter jointly-certified delegation does not.

\paragraph{Whiel encoding (Hoare triple).}

\begin{Verbatim}[fontsize=\tiny,breaklines=true,breakanywhere=true,frame=single,framesep=1.5mm]
SCHEMA
  Root/1, Reviewer/1, AccessOne/1, CertStateOne/1, ReviewStateOne/1, AccessTwo/1, CertStateTwo/1, ReviewStateTwo/1, JointStateTwo/1, AccessOne_aux/1, CertStateOne_aux/1, ReviewStateOne_aux/1, AccessTwo_aux/1, CertStateTwo_aux/1, ReviewStateTwo_aux/1, JointStateTwo_aux/1, Certified/2, Delegates/2, Review/3

PRE
  ((Certified ⊆ Delegates) ∧ ((π[0,1] (Review) ⊆ Delegates) ∧ ((π[2] (Review) ⊆ Reviewer) ∧ ((Certified ⊆ π[0,1] (σ[#2 = #3] ((Review × Reviewer)))) ∧ (π[0,3] (σ[#1 = #2] ((Certified × Review))) ⊆ Delegates)))))

CMD
  AccessOne_aux := ∅[1];
  CertStateOne_aux := ∅[1];
  ReviewStateOne_aux := ∅[1];
  AccessOne := (Root ∪ (π[2] (σ[(#0 = #1 ∧ #3 = #4)] ((CertStateOne_aux × (Review × Reviewer)))) ∪ (π[2] (σ[#0 = #1] ((ReviewStateOne_aux × Certified))) ∪ π[0] (σ[(#0 = #1 ∧ #0 = #3)] ((CertStateOne_aux × (ReviewStateOne_aux × Delegates)))))));
  CertStateOne := π[2] (σ[#0 = #1] ((AccessOne_aux × Certified)));
  ReviewStateOne := π[2] (σ[(#0 = #1 ∧ #3 = #4)] ((AccessOne_aux × (Review × Reviewer))));
  WHILE (¬(((AccessOne = AccessOne_aux) ∧ ((CertStateOne = CertStateOne_aux) ∧ (ReviewStateOne = ReviewStateOne_aux))))) DO
    AccessOne_aux := AccessOne;
    CertStateOne_aux := CertStateOne;
    ReviewStateOne_aux := ReviewStateOne;
    AccessOne := (AccessOne ∪ (Root ∪ (π[2] (σ[(#0 = #1 ∧ #3 = #4)] ((CertStateOne_aux × (Review × Reviewer)))) ∪ (π[2] (σ[#0 = #1] ((ReviewStateOne_aux × Certified))) ∪ π[0] (σ[(#0 = #1 ∧ #0 = #3)] ((CertStateOne_aux × (ReviewStateOne_aux × Delegates))))))));
    CertStateOne := (CertStateOne ∪ π[2] (σ[#0 = #1] ((AccessOne_aux × Certified))));
    ReviewStateOne := (ReviewStateOne ∪ π[2] (σ[(#0 = #1 ∧ #3 = #4)] ((AccessOne_aux × (Review × Reviewer)))))
  END;
  AccessTwo_aux := ∅[1];
  CertStateTwo_aux := ∅[1];
  ReviewStateTwo_aux := ∅[1];
  JointStateTwo_aux := ∅[1];
  AccessTwo := (Root ∪ (π[0] (σ[#0 = #2] ((JointStateTwo_aux × Delegates))) ∪ π[0] (σ[(#0 = #1 ∧ #0 = #3)] ((CertStateTwo_aux × (ReviewStateTwo_aux × Delegates))))));
  CertStateTwo := π[2] (σ[#0 = #1] ((AccessTwo_aux × Certified)));
  ReviewStateTwo := π[2] (σ[(#0 = #1 ∧ #3 = #4)] ((AccessTwo_aux × (Review × Reviewer))));
  JointStateTwo := π[2] (σ[(((#0 = #1 ∧ #3 = #4) ∧ #0 = #5) ∧ #2 = #6)] ((CertStateTwo_aux × (Review × (Reviewer × Certified)))));
  WHILE (¬(((AccessTwo = AccessTwo_aux) ∧ ((CertStateTwo = CertStateTwo_aux) ∧ ((ReviewStateTwo = ReviewStateTwo_aux) ∧ (JointStateTwo = JointStateTwo_aux)))))) DO
    AccessTwo_aux := AccessTwo;
    CertStateTwo_aux := CertStateTwo;
    ReviewStateTwo_aux := ReviewStateTwo;
    JointStateTwo_aux := JointStateTwo;
    AccessTwo := (AccessTwo ∪ (Root ∪ (π[0] (σ[#0 = #2] ((JointStateTwo_aux × Delegates))) ∪ π[0] (σ[(#0 = #1 ∧ #0 = #3)] ((CertStateTwo_aux × (ReviewStateTwo_aux × Delegates)))))));
    CertStateTwo := (CertStateTwo ∪ π[2] (σ[#0 = #1] ((AccessTwo_aux × Certified))));
    ReviewStateTwo := (ReviewStateTwo ∪ π[2] (σ[(#0 = #1 ∧ #3 = #4)] ((AccessTwo_aux × (Review × Reviewer)))));
    JointStateTwo := (JointStateTwo ∪ π[2] (σ[(((#0 = #1 ∧ #3 = #4) ∧ #0 = #5) ∧ #2 = #6)] ((CertStateTwo_aux × (Review × (Reviewer × Certified))))))
  END

POST
  (AccessOne ⊆ AccessTwo)
\end{Verbatim}

\Needspace{0.45\textheight}

\section{Example3009: Certified delegation closure exceeds audit-bound authorization}\label{case:Example3009}

\begin{description}[leftmargin=2.6cm,style=nextline,itemsep=1pt]

\item[Category] LLM-generated (GPT/Claude benchmark programs); family: two programs under TGD constraints (LLM-generated, GPT (raw\_invalid\_gpt batch); legacy raw\_new pipeline), invalid.

\item[Kind] containment.

\item[Contributors] Leo Zhang.

\item[Status] certified invalid (Certificate/Invalid.lean, written by the certificate emitter); expected invalid.

\item[Tags] \hyperref[tag:dependencies]{\texttt{dependencies}} \hyperref[tag:mutual-recursion]{\texttt{mutual-recursion}} \hyperref[tag:nonlinear-recursion]{\texttt{nonlinear-recursion}}.

\item[Input] \path{Benchmark/Example3009/Input.lean} (canonical notation).

\item[Summary] Certified delegation closure exceeds audit-bound authorization.

\item[Sides] naive compiler: program C1 of the retired corpus case Example3009 (output GrantA) \par naive compiler: program C2 of the retired corpus case Example3009 (output AllowB)

\item[Sources] {\raggedright LLM-generated (GPT (raw\_invalid\_gpt batch)), legacy raw\_new TGD pipeline this repository the retired corpus case Example3009 (programs C1, C2 and the TGD list tgd0..) TGD pipeline run raw-new-tgd-0002-product: TGDs translated to relational algebra by RelCalc.TGD.listToAssertExpr; programs compiled with the certified naive compiler (read 2026-09-15) \par method this repository the equivalence-under-constraints paper draft (technical appendix: reduction of \path{containment/equivalence} under TGDs to a Hoare triple) (read 2026-09-15) \par legacy source record raw\_invalid\_gpt\_2.json \par legacy source record \path{Benchmark/Example3009/Input.lean} \par Datalog source noted in the legacy inventory of Example3009\par}

\item[Obstruction note] from the metadata: invalid (legacy kernel counterexample)

\item[Note] Legacy raw\_new invalid TGD case (GPT): the second program demands a matching audit tuple per promotion.

\end{description}

\paragraph{Datalog program(s).}

{\raggedright\textit{prog\_c1} --- fidelity module (kernel-checked: the compiled sides of the command are this program's naive compilation).\par}

\begin{Verbatim}[fontsize=\scriptsize,breaklines=true,breakanywhere=true,frame=single,framesep=1.5mm]
GrantA(x1) :- Root(x1).
PendingA(y1) :- GrantA(x1), Delegate(x1, y1).
CheckedA(y1) :- PendingA(y1), Certified(x1, y1), Trust(x1).
GrantA(y1) :- CheckedA(y1).
PendingA(y1) :- CheckedA(x1), GrantA(z1), Review(x1, y1).
\end{Verbatim}

{\raggedright\textit{prog\_c2} --- fidelity module (kernel-checked: the compiled sides of the command are this program's naive compilation).\par}

\begin{Verbatim}[fontsize=\scriptsize,breaklines=true,breakanywhere=true,frame=single,framesep=1.5mm]
AllowB(x1) :- Root(x1).
QueueB(y1) :- AllowB(x1), Delegate(x1, y1).
VettedB(y1) :- QueueB(y1), Certified(x1, y1), Trust(x1).
AllowB(y1) :- VettedB(y1), Audit(x1, y1, z1), AllowB(z1).
QueueB(y1) :- VettedB(x1), Review(x1, y1).
\end{Verbatim}

\paragraph{Reading.} Program C1 computes a grant set GrantA in which a pending principal, once checked against only a Certified/Trust pair, is promoted, and can then reopen pending status via a further Review hop even while relying on a previously certified check. Program C2 computes AllowB with an analogous pending/vetted structure but requires a matching Audit triple, linking the certified pair to a trusted endpoint, each time a vetted principal is promoted via Delegate/Audit composition, rather than only a Review hop. This models staged authorization delegation compared against a version that insists on a matching audit record for every promotion. The Hoare triple claims containment GrantA ⊆ AllowB under four TGDs, including an existential one guaranteeing every root an audited delegate, but the claim fails: the recorded witness instance exhibits a principal that C1 grants via a reviewed hop with no matching audit tuple, which C2's audit-bound promotion rule fails to allow.

\paragraph{First-order reading of PRE/POST, translated mechanically from the relational-algebra assertion in Input.lean.} The relational-algebra assertion in \texttt{Input.lean} remains the authority.

\textit{PRE:} (∀x. (Root(x) → ∃y. Delegate(x,y) ∧ Audit(x,y,x))) ∧ (∀z u v. (Audit(z,u,v) → Certified(z,u) ∧ Trust(v))) ∧ (∀w x1. (Delegate(w,x1) ∧ Certified(w,x1) → Review(w,x1))) ∧ ∀x2 x3. ((∃x4. Review(x4,x2) ∧ Audit(x4,x2,x3)) → Delegate(x3,x2))\par The precondition is four TGDs: every Root element has some Delegate/Audit witness pairing it with itself, Audit triples are backed by a Certified pair and a Trust element, every certified Delegate pair is also a Review pair, and composing a Review pair with an Audit triple through shared coordinates yields a Delegate pair.

\textit{POST:} ∀x. (GrantA(x) → AllowB(x))\par The postcondition claims GrantA ⊆ AllowB, but the case is invalid: the recorded witness instance shows a principal in GrantA that is not in AllowB, because C2 requires an audit tuple that C1's reviewed promotion does not guarantee.

\paragraph{Whiel encoding (Hoare triple).}

\begin{Verbatim}[fontsize=\tiny,breaklines=true,breakanywhere=true,frame=single,framesep=1.5mm]
SCHEMA
  Root/1, Trust/1, GrantA/1, PendingA/1, CheckedA/1, AllowB/1, QueueB/1, VettedB/1, GrantA_aux/1, PendingA_aux/1, CheckedA_aux/1, AllowB_aux/1, QueueB_aux/1, VettedB_aux/1, Delegate/2, Certified/2, Review/2, Audit/3

PRE
  ((Root ⊆ π[0] (σ[((#0 = #2 ∧ #1 = #3) ∧ #0 = #4)] ((Delegate × Audit)))) ∧ ((Audit ⊆ (Certified × Trust)) ∧ ((π[0,1] (σ[(#0 = #2 ∧ #1 = #3)] ((Delegate × Certified))) ⊆ Review) ∧ (π[1,4] (σ[(#0 = #2 ∧ #1 = #3)] ((Review × Audit))) ⊆ π[1,0] (Delegate)))))

CMD
  GrantA_aux := ∅[1];
  PendingA_aux := ∅[1];
  CheckedA_aux := ∅[1];
  GrantA := (Root ∪ CheckedA_aux);
  PendingA := (π[2] (σ[#0 = #1] ((GrantA_aux × Delegate))) ∪ π[3] (σ[#0 = #2] ((CheckedA_aux × (GrantA_aux × Review)))));
  CheckedA := π[0] (σ[(#0 = #2 ∧ #1 = #3)] ((PendingA_aux × (Certified × Trust))));
  WHILE (¬(((GrantA = GrantA_aux) ∧ ((PendingA = PendingA_aux) ∧ (CheckedA = CheckedA_aux))))) DO
    GrantA_aux := GrantA;
    PendingA_aux := PendingA;
    CheckedA_aux := CheckedA;
    GrantA := (GrantA ∪ (Root ∪ CheckedA_aux));
    PendingA := (PendingA ∪ (π[2] (σ[#0 = #1] ((GrantA_aux × Delegate))) ∪ π[3] (σ[#0 = #2] ((CheckedA_aux × (GrantA_aux × Review))))));
    CheckedA := (CheckedA ∪ π[0] (σ[(#0 = #2 ∧ #1 = #3)] ((PendingA_aux × (Certified × Trust)))))
  END;
  AllowB_aux := ∅[1];
  QueueB_aux := ∅[1];
  VettedB_aux := ∅[1];
  AllowB := (Root ∪ π[0] (σ[(#0 = #2 ∧ #3 = #4)] ((VettedB_aux × (Audit × AllowB_aux)))));
  QueueB := (π[2] (σ[#0 = #1] ((AllowB_aux × Delegate))) ∪ π[2] (σ[#0 = #1] ((VettedB_aux × Review))));
  VettedB := π[0] (σ[(#0 = #2 ∧ #1 = #3)] ((QueueB_aux × (Certified × Trust))));
  WHILE (¬(((AllowB = AllowB_aux) ∧ ((QueueB = QueueB_aux) ∧ (VettedB = VettedB_aux))))) DO
    AllowB_aux := AllowB;
    QueueB_aux := QueueB;
    VettedB_aux := VettedB;
    AllowB := (AllowB ∪ (Root ∪ π[0] (σ[(#0 = #2 ∧ #3 = #4)] ((VettedB_aux × (Audit × AllowB_aux))))));
    QueueB := (QueueB ∪ (π[2] (σ[#0 = #1] ((AllowB_aux × Delegate))) ∪ π[2] (σ[#0 = #1] ((VettedB_aux × Review)))));
    VettedB := (VettedB ∪ π[0] (σ[(#0 = #2 ∧ #1 = #3)] ((QueueB_aux × (Certified × Trust)))))
  END

POST
  (GrantA ⊆ AllowB)
\end{Verbatim}

\Needspace{0.45\textheight}

\section{Example3014: Auditor trust diverges from resource trust under recursive elevation}\label{case:Example3014}

\begin{description}[leftmargin=2.6cm,style=nextline,itemsep=1pt]

\item[Category] LLM-generated (GPT/Claude benchmark programs); family: two programs under TGD constraints (LLM-generated, GPT (raw\_invalid\_gpt batch); legacy raw\_new pipeline), invalid.

\item[Kind] containment.

\item[Contributors] Leo Zhang.

\item[Status] certified invalid (Certificate/Invalid.lean, written by the certificate emitter); expected invalid.

\item[Tags] \hyperref[tag:dependencies]{\texttt{dependencies}} \hyperref[tag:mutual-recursion]{\texttt{mutual-recursion}} \hyperref[tag:nonlinear-recursion]{\texttt{nonlinear-recursion}}.

\item[Input] \path{Benchmark/Example3014/Input.lean} (canonical notation).

\item[Summary] Auditor trust diverges from resource trust under recursive elevation.

\item[Sides] naive compiler: program C1 of the retired corpus case Example3014 (output AccessA) \par naive compiler: program C2 of the retired corpus case Example3014 (output AccessB)

\item[Sources] {\raggedright LLM-generated (GPT (raw\_invalid\_gpt batch)), legacy raw\_new TGD pipeline this repository the retired corpus case Example3014 (programs C1, C2 and the TGD list tgd0..) TGD pipeline run raw-new-tgd-0002-product: TGDs translated to relational algebra by RelCalc.TGD.listToAssertExpr; programs compiled with the certified naive compiler (read 2026-09-15) \par method this repository the equivalence-under-constraints paper draft (technical appendix: reduction of \path{containment/equivalence} under TGDs to a Hoare triple) (read 2026-09-15) \par legacy source record raw\_invalid\_gpt\_2.json \par legacy source record \path{Benchmark/Example3014/Input.lean} \par Datalog source noted in the legacy inventory of Example3014\par}

\item[Obstruction note] from the metadata: invalid (legacy kernel counterexample)

\item[Note] Legacy raw\_new invalid TGD case (GPT): the proof is checked on the resource coordinate instead of the auditor coordinate.

\end{description}

\paragraph{Datalog program(s).}

{\raggedright\textit{prog\_c1} --- fidelity module (kernel-checked: the compiled sides of the command are this program's naive compilation).\par}

\begin{Verbatim}[fontsize=\scriptsize,breaklines=true,breakanywhere=true,frame=single,framesep=1.5mm]
AccessA(x1, y1) :- Request(x1, y1), Permit(x1, y1).
ElevatedA(x1, z1) :- AccessA(x1, y1), Grant(y1, z1), Scope(z1).
CheckedA(x1, z1) :- ElevatedA(x1, y1), Override(y1, z1), Proof(x1, z1, x2), Trusted(x2).
AccessA(x1, z1) :- CheckedA(x1, z1).
ElevatedA(x1, z1) :- AccessA(x1, y1), AccessA(x1, x2), Proof(y1, z1, x2).
\end{Verbatim}

{\raggedright\textit{prog\_c2} --- fidelity module (kernel-checked: the compiled sides of the command are this program's naive compilation).\par}

\begin{Verbatim}[fontsize=\scriptsize,breaklines=true,breakanywhere=true,frame=single,framesep=1.5mm]
AccessB(x1, y1) :- Request(x1, y1), Permit(x1, y1).
ElevatedB(x1, z1) :- AccessB(x1, y1), Grant(y1, z1), Scope(z1).
AuditB(x1, z1, x2) :- ElevatedB(x1, y1), Override(y1, z1), Proof(x1, z1, x2).
AccessB(x1, z1) :- AuditB(x1, z1, x2), Trusted(z1).
ElevatedB(x1, z1) :- AccessB(x1, y1), AccessB(x1, x2), Proof(y1, z1, x2).
\end{Verbatim}

\paragraph{Reading.} Program C1 computes an access relation AccessA that elevates a principal's access to a new resource via Grant/Scope, re-checks that elevation by requiring a Proof triple whose auditor coordinate (the proof's third argument) is Trusted, and further allows elevation whenever two already-accessible resources are linked by a Proof triple. Program C2 computes AccessB with the analogous elevation structure but tests Trusted against the proof's resource coordinate rather than its auditor coordinate. This models recursive policy elevation comparing which coordinate of a proof triple, auditor versus resource, a trust check should apply to. The Hoare triple claims containment AccessA ⊆ AccessB under five TGDs, including an existential one guaranteeing every request a trusted proof witness, but the claim fails: the recorded witness instance shows the two programs agree on the first trusted endpoint but diverge at the next audited elevation, since trusting the auditor coordinate versus the resource coordinate of the same Proof triple picks out different principals.

\paragraph{First-order reading of PRE/POST, translated mechanically from the relational-algebra assertion in Input.lean.} The relational-algebra assertion in \texttt{Input.lean} remains the authority.

\textit{PRE:} (∀x y. (Request(x,y) → ∃z. Proof(x,y,z) ∧ Trusted(z))) ∧ (∀u v. ((∃w. Proof(u,v,w)) → Permit(u,v))) ∧ (∀x1 x2. (Grant(x1,x2) ∧ Scope(x2) → Permit(x1,x2))) ∧ (∀x3. ((∃x4 x5 x6. Override(x3,x5) ∧ Proof(x4,x5,x6)) → Scope(x3))) ∧ ∀x7. ((∃x8 x9. Permit(x8,x9) ∧ Grant(x9,x7)) → Scope(x7))\par The precondition is five TGDs: every Request pair has a Proof triple whose third coordinate is Trusted, the first two coordinates of every Proof triple are a Permit pair, every Grant pair whose target is in Scope is a Permit pair, the source of an Override pair with a matching Proof witness is in Scope, and the target of a Permit pair composed with a Grant pair is in Scope.

\textit{POST:} ∀x y. (AccessA(x,y) → AccessB(x,y))\par The postcondition claims AccessA ⊆ AccessB, but the case is invalid: the recorded witness instance shows a resource in AccessA that is not in AccessB.

\paragraph{Whiel encoding (Hoare triple).}

\begin{Verbatim}[fontsize=\tiny,breaklines=true,breakanywhere=true,frame=single,framesep=1.5mm]
SCHEMA
  Scope/1, Trusted/1, Request/2, Permit/2, Grant/2, Override/2, AccessA/2, ElevatedA/2, CheckedA/2, AccessB/2, ElevatedB/2, AccessA_aux/2, ElevatedA_aux/2, CheckedA_aux/2, AccessB_aux/2, ElevatedB_aux/2, Proof/3, AuditB/3, AuditB_aux/3

PRE
  ((Request ⊆ π[0,1] (σ[#2 = #3] ((Proof × Trusted)))) ∧ ((π[0,1] (Proof) ⊆ Permit) ∧ ((π[0,1] (σ[#1 = #2] ((Grant × Scope))) ⊆ Permit) ∧ ((π[0] (σ[#1 = #3] ((Override × Proof))) ⊆ Scope) ∧ (π[3] (σ[#1 = #2] ((Permit × Grant))) ⊆ Scope)))))

CMD
  AccessA_aux := ∅[2];
  ElevatedA_aux := ∅[2];
  CheckedA_aux := ∅[2];
  AccessA := (π[0,1] (σ[(#0 = #2 ∧ #1 = #3)] ((Request × Permit))) ∪ CheckedA_aux);
  ElevatedA := (π[0,3] (σ[(#1 = #2 ∧ #3 = #4)] ((AccessA_aux × (Grant × Scope)))) ∪ π[0,5] (σ[((#0 = #2 ∧ #1 = #4) ∧ #3 = #6)] ((AccessA_aux × (AccessA_aux × Proof)))));
  CheckedA := π[0,3] (σ[(((#1 = #2 ∧ #0 = #4) ∧ #3 = #5) ∧ #6 = #7)] ((ElevatedA_aux × (Override × (Proof × Trusted)))));
  WHILE (¬(((AccessA = AccessA_aux) ∧ ((ElevatedA = ElevatedA_aux) ∧ (CheckedA = CheckedA_aux))))) DO
    AccessA_aux := AccessA;
    ElevatedA_aux := ElevatedA;
    CheckedA_aux := CheckedA;
    AccessA := (AccessA ∪ (π[0,1] (σ[(#0 = #2 ∧ #1 = #3)] ((Request × Permit))) ∪ CheckedA_aux));
    ElevatedA := (ElevatedA ∪ (π[0,3] (σ[(#1 = #2 ∧ #3 = #4)] ((AccessA_aux × (Grant × Scope)))) ∪ π[0,5] (σ[((#0 = #2 ∧ #1 = #4) ∧ #3 = #6)] ((AccessA_aux × (AccessA_aux × Proof))))));
    CheckedA := (CheckedA ∪ π[0,3] (σ[(((#1 = #2 ∧ #0 = #4) ∧ #3 = #5) ∧ #6 = #7)] ((ElevatedA_aux × (Override × (Proof × Trusted))))))
  END;
  AccessB_aux := ∅[2];
  ElevatedB_aux := ∅[2];
  AuditB_aux := ∅[3];
  AccessB := (π[0,1] (σ[(#0 = #2 ∧ #1 = #3)] ((Request × Permit))) ∪ π[0,1] (σ[#1 = #3] ((AuditB_aux × Trusted))));
  ElevatedB := (π[0,3] (σ[(#1 = #2 ∧ #3 = #4)] ((AccessB_aux × (Grant × Scope)))) ∪ π[0,5] (σ[((#0 = #2 ∧ #1 = #4) ∧ #3 = #6)] ((AccessB_aux × (AccessB_aux × Proof)))));
  AuditB := π[0,3,6] (σ[((#1 = #2 ∧ #0 = #4) ∧ #3 = #5)] ((ElevatedB_aux × (Override × Proof))));
  WHILE (¬(((AccessB = AccessB_aux) ∧ ((ElevatedB = ElevatedB_aux) ∧ (AuditB = AuditB_aux))))) DO
    AccessB_aux := AccessB;
    ElevatedB_aux := ElevatedB;
    AuditB_aux := AuditB;
    AccessB := (AccessB ∪ (π[0,1] (σ[(#0 = #2 ∧ #1 = #3)] ((Request × Permit))) ∪ π[0,1] (σ[#1 = #3] ((AuditB_aux × Trusted)))));
    ElevatedB := (ElevatedB ∪ (π[0,3] (σ[(#1 = #2 ∧ #3 = #4)] ((AccessB_aux × (Grant × Scope)))) ∪ π[0,5] (σ[((#0 = #2 ∧ #1 = #4) ∧ #3 = #6)] ((AccessB_aux × (AccessB_aux × Proof))))));
    AuditB := (AuditB ∪ π[0,3,6] (σ[((#1 = #2 ∧ #0 = #4) ∧ #3 = #5)] ((ElevatedB_aux × (Override × Proof)))))
  END

POST
  (AccessA ⊆ AccessB)
\end{Verbatim}

\Needspace{0.45\textheight}

\section{Example4001: Directed triangle query versus its path-view reformulation}\label{case:Example4001}

\begin{description}[leftmargin=2.6cm,style=nextline,itemsep=1pt]

\item[Category] verification task (VerificationTasks.pdf, section 10.1); family: Core.

\item[Kind] equivalence.

\item[Contributors] Leo Zhang, Val Tannen.

\item[Status] certified valid (Certificate/Valid.lean, written by the certificate emitter); expected valid.

\item[Tags] \hyperref[tag:views]{\texttt{views}}.

\item[Input] \path{Benchmark/Example4001/Input.lean} (canonical notation).

\item[Summary] A directed-triangle query and its reformulation over a length-two path view.

\item[Sources] {\raggedright VerificationTasks.pdf, section 10.1 (Val Tannen; \path{reports/VerificationTasks.pdf}) \par the reformulations task document \par legacy source record \path{Benchmark/Example4001/Input.lean} \par legacy source record of the reformulations task document\par}

\end{description}

\paragraph{Datalog program(s).}

{\raggedright\textit{globalProgram --- `Q x :- Exy, Eyz, Ezx`.} --- fidelity module (kernel-checked: the compiled sides of the command are this program's naive compilation).\par}

\begin{Verbatim}[fontsize=\scriptsize,breaklines=true,breakanywhere=true,frame=single,framesep=1.5mm]
Q(x1) :- E(x1, y1), E(y1, z1), E(z1, x1).
\end{Verbatim}

{\raggedright\textit{localProgram --- `Q↑ x :- V xyz, V yzx`.} --- fidelity module (kernel-checked: the compiled sides of the command are this program's naive compilation).\par}

\begin{Verbatim}[fontsize=\scriptsize,breaklines=true,breakanywhere=true,frame=single,framesep=1.5mm]
QUp(x1) :- V(x1, y1, z1), V(y1, z1, x1).
\end{Verbatim}

\paragraph{Reading.} The two Datalog programs compute the directed-triangle query Q(x) :- E(x,y), E(y,z), E(z,x) over an edge relation E, and its reformulation Q↑(x) :- V(x,y,z), V(y,z,x) over a length-two path view V. Under the view constraint that V records exactly the length-two paths of E together with their midpoint (V = E ∘ E), the Hoare triple claims the two queries compute equal results, Q = Q↑; this is an equivalence claim. The claim is certified valid: joining two copies of the view on overlapping variables reconstructs the directed-triangle detection under set semantics.

\paragraph{First-order reading of PRE/POST, translated mechanically from the relational-algebra assertion in Input.lean.} The relational-algebra assertion in \texttt{Input.lean} remains the authority.

\textit{PRE:} ∀x y z. (V(x,y,z) ↔ E(x,y) ∧ E(y,z))\par The precondition states that V equals the set of length-two paths through E: V is E composed with E, projected onto the start, midpoint and end columns with the middle columns equated.

\textit{POST:} ∀x. (Q(x) ↔ QUp(x))\par The postcondition asserts the two queries agree exactly, Q = QUp.

\paragraph{Whiel encoding (Hoare triple).}

\begin{Verbatim}[fontsize=\tiny,breaklines=true,breakanywhere=true,frame=single,framesep=1.5mm]
SCHEMA
  Q/1, QUp/1, Q_aux/1, QUp_aux/1, E/2, V/3

PRE
  (V = (π[0, 1, 3] (σ[#1 = #2] (E × E))))

CMD
  Q_aux := ∅[1];
  Q :=
    (π[0]
      (σ[((#1 = #2 ∧ #3 = #4) ∧ #0 = #5)]
        (E × (E × E))));
  WHILE (¬ (Q = Q_aux)) DO
    Q_aux := Q;
    Q :=
      (Q ∪
        (π[0]
          (σ[((#1 = #2 ∧ #3 = #4) ∧ #0 = #5)]
            (E × (E × E)))))
  END;
  QUp_aux := ∅[1];
  QUp :=
    (π[0]
      (σ[((#1 = #3 ∧ #2 = #4) ∧ #0 = #5)]
        (V × V)));
  WHILE (¬ (QUp = QUp_aux)) DO
    QUp_aux := QUp;
    QUp :=
      (QUp ∪
        (π[0]
          (σ[((#1 = #3 ∧ #2 = #4) ∧ #0 = #5)]
            (V × V))))
  END

POST
  (Q = QUp)
\end{Verbatim}

\Needspace{0.45\textheight}

\section{Example4002: Transitive closure versus its length-two path-view reformulation}\label{case:Example4002}

\begin{description}[leftmargin=2.6cm,style=nextline,itemsep=1pt]

\item[Category] verification task (VerificationTasks.pdf, section 10.1); family: Core.

\item[Kind] equivalence.

\item[Contributors] Leo Zhang, Val Tannen.

\item[Status] certified invalid (Certificate/Invalid.lean, written by the certificate emitter); expected invalid.

\item[Tags] \hyperref[tag:views]{\texttt{views}}.

\item[Input] \path{Benchmark/Example4002/Input.lean} (canonical notation).

\item[Summary] Right-linear transitive closure and a reformulation over a length-two path view, tested for equality.

\item[Sources] {\raggedright VerificationTasks.pdf, section 10.1 (Val Tannen; \path{reports/VerificationTasks.pdf}) \par the reformulations task document \par legacy source record \path{Benchmark/Example4002/Input.lean} \par legacy source record of the reformulations task document\par}

\end{description}

\paragraph{Datalog program(s).}

{\raggedright\textit{globalProgram --- `T xy :- Exy` and `T xy :- Exz, T zy`.} --- fidelity module (kernel-checked: the compiled sides of the command are this program's naive compilation).\par}

\begin{Verbatim}[fontsize=\scriptsize,breaklines=true,breakanywhere=true,frame=single,framesep=1.5mm]
T(x1, y1) :- E(x1, y1).
T(x1, y1) :- E(x1, z1), T(z1, y1).
\end{Verbatim}

{\raggedright\textit{localProgram --- `T↑ xy :- V xwy` and `T↑ xz :- V xwy, T↑ yz`.} --- fidelity module (kernel-checked: the compiled sides of the command are this program's naive compilation).\par}

\begin{Verbatim}[fontsize=\scriptsize,breaklines=true,breakanywhere=true,frame=single,framesep=1.5mm]
TUp(x1, y1) :- V(x1, z1, y1).
TUp(x1, x2) :- V(x1, z1, y1), TUp(y1, x2).
\end{Verbatim}

\paragraph{Reading.} The global program T is the standard right-linear transitive closure of edge relation E; the local reformulation TUp chains the length-two path view V (V = E ∘ E, keeping the midpoint) recursively: TUp(x,y) :- V(x,w,y), and TUp(x,z) :- V(x,w,y), TUp(y,z). The Hoare triple claims equality, T = TUp --- an equivalence claim between transitive closure and the view-chaining reformulation. This claim is invalid, certified by a kernel-checked counterexample: because every TUp derivation composes at least one V-step, TUp collects exactly the paths of even length at least two, so a single edge (e.g. E = \{(0,1)\} with everything else empty) lies in T but not in TUp.

\paragraph{First-order reading of PRE/POST, translated mechanically from the relational-algebra assertion in Input.lean.} The relational-algebra assertion in \texttt{Input.lean} remains the authority.

\textit{PRE:} ∀x y z. (V(x,y,z) ↔ E(x,y) ∧ E(y,z))\par The precondition states that V equals the set of length-two paths through E, i.e. V = E ∘ E with the midpoint kept.

\textit{POST:} ∀x y. (T(x,y) ↔ TUp(x,y))\par The postcondition asserts T = TUp, that transitive closure and the view-chaining reformulation compute the same relation --- a claim that fails.

\paragraph{Whiel encoding (Hoare triple).}

\begin{Verbatim}[fontsize=\scriptsize,breaklines=true,breakanywhere=true,frame=single,framesep=1.5mm]
SCHEMA
  E/2, T/2, TUp/2, T_aux/2, TUp_aux/2, V/3

PRE
  (V = (π[0, 1, 3] (σ[#1 = #2] (E × E))))

CMD
  T_aux := ∅[2];
  T :=
    (E ∪ (π[0, 3] (σ[#1 = #2] (E × T_aux))));
  WHILE (¬ (T = T_aux)) DO
    T_aux := T;
    T :=
      (T ∪
        (E ∪
          (π[0, 3] (σ[#1 = #2] (E × T_aux)))))
  END;
  TUp_aux := ∅[2];
  TUp :=
    ((π[0, 2] V) ∪
      (π[0, 4] (σ[#2 = #3] (V × TUp_aux))));
  WHILE (¬ (TUp = TUp_aux)) DO
    TUp_aux := TUp;
    TUp :=
      (TUp ∪
        ((π[0, 2] V) ∪
          (π[0, 4]
            (σ[#2 = #3] (V × TUp_aux)))))
  END

POST
  (T = TUp)
\end{Verbatim}

\Needspace{0.45\textheight}

\section{Example4003: Containment of the path-view reformulation in transitive closure}\label{case:Example4003}

\begin{description}[leftmargin=2.6cm,style=nextline,itemsep=1pt]

\item[Category] verification task (VerificationTasks.pdf, section 10.1); family: Core.

\item[Kind] containment.

\item[Contributors] Leo Zhang, Val Tannen.

\item[Status] certified valid (Certificate/Valid.lean, written by the certificate emitter); expected valid.

\item[Tags] \hyperref[tag:prophecy]{\texttt{prophecy}} \hyperref[tag:views]{\texttt{views}}.

\item[Input] \path{Benchmark/Example4003/Input.lean} (canonical notation).

\item[Summary] Right-linear transitive closure and a reformulation over a length-two path view, tested for containment.

\item[Sources] {\raggedright VerificationTasks.pdf, section 10.1 (Val Tannen; \path{reports/VerificationTasks.pdf}) \par the reformulations task document \par legacy source record \path{Benchmark/Example4003/Input.lean} \par legacy source record of the reformulations task document\par}

\item[Prophecy] kernel-checked certificate with a level-1 prophecy clause (Certificate/Valid.lean, Core.json).

\end{description}

\paragraph{Datalog program(s).}

{\raggedright\textit{globalProgram --- `T xy :- Exy` and `T xy :- Exz, T zy`.} --- fidelity module (kernel-checked: the compiled sides of the command are this program's naive compilation).\par}

\begin{Verbatim}[fontsize=\scriptsize,breaklines=true,breakanywhere=true,frame=single,framesep=1.5mm]
T(x1, y1) :- E(x1, y1).
T(x1, y1) :- E(x1, z1), T(z1, y1).
\end{Verbatim}

{\raggedright\textit{localProgram --- `T↑ xy :- V xwy` and `T↑ xz :- V xwy, T↑ yz`.} --- fidelity module (kernel-checked: the compiled sides of the command are this program's naive compilation).\par}

\begin{Verbatim}[fontsize=\scriptsize,breaklines=true,breakanywhere=true,frame=single,framesep=1.5mm]
TUp(x1, y1) :- V(x1, z1, y1).
TUp(x1, x2) :- V(x1, z1, y1), TUp(y1, x2).
\end{Verbatim}

{\raggedright\textit{program} --- fidelity module (kernel-checked: the compiled sides of the command are this program's naive compilation).\par}

\begin{Verbatim}[fontsize=\scriptsize,breaklines=true,breakanywhere=true,frame=single,framesep=1.5mm]
T(x1, y1) :- E(x1, y1).
T(x1, y1) :- E(x1, z1), E(z1, y1).
T(x1, y1) :- E(x1, z1), T(z1, y1).
T(x1, y1) :- E(x1, z1), E(z1, x2), T(x2, y1).
\end{Verbatim}

\paragraph{Reading.} This case uses the same global program (right-linear transitive closure T of E) and the same local reformulation (TUp, chaining the length-two path view V = E ∘ E) as Example4002, but weakens the postcondition from equality to a one-directional (forward) containment, TUp ⊆ T: every pair the view-chaining reformulation derives should also be a transitive-closure pair. The header states the expected verdict is valid, since TUp collects exactly the even-length-at-least-two paths, all of which lie in T; no certificate exists yet for this pass, so the classification stays unknown.

\paragraph{First-order reading of PRE/POST, translated mechanically from the relational-algebra assertion in Input.lean.} The relational-algebra assertion in \texttt{Input.lean} remains the authority.

\textit{PRE:} ∀x y z. (V(x,y,z) ↔ E(x,y) ∧ E(y,z))\par The precondition states that V equals the set of length-two paths through E, i.e. V = E ∘ E with the midpoint kept.

\textit{POST:} ∀x y. (TUp(x,y) → T(x,y))\par The postcondition asserts TUp ⊆ T, that every pair produced by the view-chaining reformulation is also produced by transitive closure.

\paragraph{Whiel encoding (Hoare triple).}

\begin{Verbatim}[fontsize=\scriptsize,breaklines=true,breakanywhere=true,frame=single,framesep=1.5mm]
SCHEMA
  E/2, T/2, TUp/2, T_aux/2, TUp_aux/2, V/3

PRE
  (V = (π[0, 1, 3] (σ[#1 = #2] (E × E))))

CMD
  T_aux := ∅[2];
  T :=
    (E ∪ (π[0, 3] (σ[#1 = #2] (E × T_aux))));
  WHILE (¬ (T = T_aux)) DO
    T_aux := T;
    T :=
      (T ∪
        (E ∪
          (π[0, 3] (σ[#1 = #2] (E × T_aux)))))
  END;
  TUp_aux := ∅[2];
  TUp :=
    ((π[0, 2] V) ∪
      (π[0, 4] (σ[#2 = #3] (V × TUp_aux))));
  WHILE (¬ (TUp = TUp_aux)) DO
    TUp_aux := TUp;
    TUp :=
      (TUp ∪
        ((π[0, 2] V) ∪
          (π[0, 4]
            (σ[#2 = #3] (V × TUp_aux)))))
  END

POST
  (TUp ⊆ T)
\end{Verbatim}

\Needspace{0.45\textheight}

\section{Example4004: Path-view reformulation versus the even-length path query}\label{case:Example4004}

\begin{description}[leftmargin=2.6cm,style=nextline,itemsep=1pt]

\item[Category] verification task (VerificationTasks.pdf, section 10.1); family: Core.

\item[Kind] equivalence.

\item[Contributors] Leo Zhang, Val Tannen.

\item[Status] certified valid (Certificate/Valid.lean, written by the certificate emitter); expected valid.

\item[Tags] \hyperref[tag:mutual-recursion]{\texttt{mutual-recursion}} \hyperref[tag:prophecy]{\texttt{prophecy}} \hyperref[tag:views]{\texttt{views}}.

\item[Input] \path{Benchmark/Example4004/Input.lean} (canonical notation).

\item[Summary] An odd/even path program and a reformulation over a length-two path view, tested for equality against the even half.

\item[Sources] {\raggedright VerificationTasks.pdf, section 10.1 (Val Tannen; \path{reports/VerificationTasks.pdf}) \par the reformulations task document \par legacy source record \path{Benchmark/Example4004/Input.lean} \par legacy source record of the reformulations task document\par}

\item[Prophecy] kernel-checked certificate with a level-1 prophecy clause (Certificate/Valid.lean, Core.json).

\end{description}

\paragraph{Datalog program(s).}

{\raggedright\textit{globalProgram --- The odd/even program, with output IDB `Te`.} --- fidelity module (kernel-checked: the compiled sides of the command are this program's naive compilation).\par}

\begin{Verbatim}[fontsize=\scriptsize,breaklines=true,breakanywhere=true,frame=single,framesep=1.5mm]
To(x1, y1) :- E(x1, y1).
To(x1, y1) :- E(x1, z1), Te(z1, y1).
Te(x1, y1) :- E(x1, z1), To(z1, y1).
\end{Verbatim}

{\raggedright\textit{localProgram --- `T↑ xy :- V xwy` and `T↑ xz :- V xwy, T↑ yz`.} --- fidelity module (kernel-checked: the compiled sides of the command are this program's naive compilation).\par}

\begin{Verbatim}[fontsize=\scriptsize,breaklines=true,breakanywhere=true,frame=single,framesep=1.5mm]
TUp(x1, y1) :- V(x1, z1, y1).
TUp(x1, x2) :- V(x1, z1, y1), TUp(y1, x2).
\end{Verbatim}

{\raggedright\textit{program} --- fidelity module (kernel-checked: the compiled sides of the command are this program's naive compilation).\par}

\begin{Verbatim}[fontsize=\scriptsize,breaklines=true,breakanywhere=true,frame=single,framesep=1.5mm]
To(x1, y1) :- E(x1, y1).
To(x1, y1) :- E(x1, z1), Te(z1, y1).
Te(x1, y1) :- E(x1, z1), E(z1, y1).
Te(x1, y1) :- E(x1, z1), E(z1, x2), Te(x2, y1).
\end{Verbatim}

\paragraph{Reading.} The global program computes paths of odd length (To) and even length (Te) over E by mutual recursion --- To(x,y):-E(x,y), To(x,y):-E(x,z),Te(z,y), Te(x,y):-E(x,z),To(z,y) --- with Te as the relevant output. The local reformulation TUp chains the same length-two path view V = E ∘ E used in the neighboring cases. The Hoare triple claims equality, TUp = Te: that the view-chaining reformulation exactly recovers the even-length-path relation from the odd/even decomposition of transitive closure. The header states the expected verdict is valid, since Te is exactly the paths of even length at least two, which is what TUp collects; no certificate exists yet, so the classification stays unknown.

\paragraph{First-order reading of PRE/POST, translated mechanically from the relational-algebra assertion in Input.lean.} The relational-algebra assertion in \texttt{Input.lean} remains the authority.

\textit{PRE:} ∀x y z. (V(x,y,z) ↔ E(x,y) ∧ E(y,z))\par The precondition states that V equals the set of length-two paths through E, i.e. V = E ∘ E with the midpoint kept.

\textit{POST:} ∀x y. (TUp(x,y) ↔ Te(x,y))\par The postcondition asserts TUp = Te, that the view-chaining reformulation equals the even-length-path relation from the odd/even decomposition.

\paragraph{Whiel encoding (Hoare triple).}

\begin{Verbatim}[fontsize=\tiny,breaklines=true,breakanywhere=true,frame=single,framesep=1.5mm]
SCHEMA
  E/2, To/2, Te/2, TUp/2, To_aux/2, Te_aux/2, TUp_aux/2, V/3

PRE
  (V = (π[0, 1, 3] (σ[#1 = #2] (E × E))))

CMD
  To_aux := ∅[2];
  Te_aux := ∅[2];
  To :=
    (E ∪ (π[0, 3] (σ[#1 = #2] (E × Te_aux))));
  Te := (π[0, 3] (σ[#1 = #2] (E × To_aux)));
  WHILE (¬ ((To = To_aux) ∧ (Te = Te_aux))) DO
    To_aux := To;
    Te_aux := Te;
    To :=
      (To ∪
        (E ∪
          (π[0, 3] (σ[#1 = #2] (E × Te_aux)))));
    Te :=
      (Te ∪
        (π[0, 3] (σ[#1 = #2] (E × To_aux))))
  END;
  TUp_aux := ∅[2];
  TUp :=
    ((π[0, 2] V) ∪
      (π[0, 4] (σ[#2 = #3] (V × TUp_aux))));
  WHILE (¬ (TUp = TUp_aux)) DO
    TUp_aux := TUp;
    TUp :=
      (TUp ∪
        ((π[0, 2] V) ∪
          (π[0, 4]
            (σ[#2 = #3] (V × TUp_aux)))))
  END

POST
  (TUp = Te)
\end{Verbatim}

\Needspace{0.45\textheight}

\section{Example4034: Containment of a reformulation over sound odd/even views in transitive closure}\label{case:Example4034}

\begin{description}[leftmargin=2.6cm,style=nextline,itemsep=1pt]

\item[Category] verification task (VerificationTasks.pdf, section 10.1); family: Core.

\item[Kind] conditional claim (implication).

\item[Contributors] Leo Zhang, Val Tannen.

\item[Status] certified valid (Certificate/Valid.lean, written by the certificate emitter); expected valid.

\item[Tags] \hyperref[tag:mutual-recursion]{\texttt{mutual-recursion}} \hyperref[tag:prophecy]{\texttt{prophecy}} \hyperref[tag:views]{\texttt{views}}.

\item[Input] \path{Benchmark/Example4034/Input.lean} (canonical notation).

\item[Summary] Transitive closure and a reformulation read off two sound odd/even views.

\item[Sources] {\raggedright VerificationTasks.pdf, section 10.1 (Val Tannen; \path{reports/VerificationTasks.pdf}) \par the reformulations task document\par}

\item[Prophecy] kernel-checked certificate with a level-1 prophecy clause (Certificate/Valid.lean, Core.json).

\end{description}

\paragraph{Datalog program(s).}

{\raggedright\textit{localProgram --- `T↑ xy :- Vo xy` and `T↑ xy :- Ve xy`.} --- fidelity module (kernel-checked: the compiled sides of the command are this program's naive compilation).\par}

\begin{Verbatim}[fontsize=\scriptsize,breaklines=true,breakanywhere=true,frame=single,framesep=1.5mm]
TUp(x1, y1) :- Vo(x1, y1).
TUp(x1, y1) :- Ve(x1, y1).
\end{Verbatim}

{\raggedright\textit{viewProgram --- The odd/even program defining the two views.} --- fidelity module (kernel-checked: the compiled sides of the command are this program's naive compilation).\par}

\begin{Verbatim}[fontsize=\scriptsize,breaklines=true,breakanywhere=true,frame=single,framesep=1.5mm]
To(x1, y1) :- E(x1, y1).
To(x1, y1) :- E(x1, z1), Te(z1, y1).
Te(x1, y1) :- E(x1, z1), To(z1, y1).
\end{Verbatim}

{\raggedright\textit{globalProgram --- `T xy :- Exy` and `T xy :- Exz, T zy`.} --- fidelity module (kernel-checked: the compiled sides of the command are this program's naive compilation).\par}

\begin{Verbatim}[fontsize=\scriptsize,breaklines=true,breakanywhere=true,frame=single,framesep=1.5mm]
T(x1, y1) :- E(x1, y1).
T(x1, y1) :- E(x1, z1), T(z1, y1).
\end{Verbatim}

\paragraph{Reading.} The global program T is the standard right-linear transitive closure of E. The local reformulation is loop-free: TUp is simply the union of two source relations Vo and Ve, which stand for odd- and even-length path views computed by a separate odd/even mutual-recursion program over E, but are only assumed sound (Vo ⊆ To, Ve ⊆ Te) rather than exact. Because that soundness hypothesis refers to the least fixpoints the loop itself computes, it cannot be phrased as a precondition, so it is folded into the postcondition as an implication written with ¬ and ∨: if Vo and Ve are sound approximations of the odd/even path relations, then the reformulation is contained in transitive closure, TUp ⊆ T. This is a conditional (soundness-guarded) containment claim. The claim is valid, since Vo ∪ Ve ⊆ To ∪ Te = T whenever the hypothesis holds.

\paragraph{First-order reading of PRE/POST, translated mechanically from the relational-algebra assertion in Input.lean.} The relational-algebra assertion in \texttt{Input.lean} remains the authority.

\textit{PRE:} ⊤\par The precondition is simply true; no constraint holds at entry.

\textit{POST:} ¬((∀x y. (Vo(x,y) → To(x,y))) ∧ ∀z u. (Ve(z,u) → Te(z,u))) ∨ ∀v w. (TUp(v,w) → T(v,w))\par The postcondition states that if Vo is contained in the odd-length-path relation To and Ve is contained in the even-length-path relation Te, then the reformulation TUp is contained in the transitive closure T.

\paragraph{Whiel encoding (Hoare triple).}

\begin{Verbatim}[fontsize=\tiny,breaklines=true,breakanywhere=true,frame=single,framesep=1.5mm]
SCHEMA
  E/2, Vo/2, Ve/2, T/2, To/2, Te/2, TUp/2, T_aux/2, To_aux/2, Te_aux/2, TUp_aux/2

PRE
  true 

CMD
  TUp_aux := ∅[2];
  TUp := (Vo ∪ Ve);
  WHILE (¬ (TUp = TUp_aux)) DO
    TUp_aux := TUp;
    TUp := (TUp ∪ (Vo ∪ Ve))
  END;
  To_aux := ∅[2];
  Te_aux := ∅[2];
  To :=
    (E ∪ (π[0, 3] (σ[#1 = #2] (E × Te_aux))));
  Te := (π[0, 3] (σ[#1 = #2] (E × To_aux)));
  WHILE (¬ ((To = To_aux) ∧ (Te = Te_aux))) DO
    To_aux := To;
    Te_aux := Te;
    To :=
      (To ∪
        (E ∪
          (π[0, 3] (σ[#1 = #2] (E × Te_aux)))));
    Te :=
      (Te ∪
        (π[0, 3] (σ[#1 = #2] (E × To_aux))))
  END;
  T_aux := ∅[2];
  T :=
    (E ∪ (π[0, 3] (σ[#1 = #2] (E × T_aux))));
  WHILE (¬ (T = T_aux)) DO
    T_aux := T;
    T :=
      (T ∪
        (E ∪
          (π[0, 3] (σ[#1 = #2] (E × T_aux)))))
  END

POST
  ((¬ ((Vo ⊆ To) ∧ (Ve ⊆ Te))) ∨ (TUp ⊆ T))
\end{Verbatim}

\Needspace{0.45\textheight}

\section{Example4035: Good-interior path query versus its begin-good and end-good edge-view reformulation}\label{case:Example4035}

\begin{description}[leftmargin=2.6cm,style=nextline,itemsep=1pt]

\item[Category] verification task (VerificationTasks.pdf, section 10.1); family: Core.

\item[Kind] equivalence.

\item[Contributors] Leo Zhang, Val Tannen.

\item[Status] certified valid (Certificate/Valid.lean, written by the certificate emitter); expected valid.

\item[Tags] \hyperref[tag:prophecy]{\texttt{prophecy}} \hyperref[tag:views]{\texttt{views}}.

\item[Input] \path{Benchmark/Example4035/Input.lean} (canonical notation).

\item[Summary] Walks through good vertices and a reformulation over two conjunctive views.

\item[Sources] {\raggedright VerificationTasks.pdf, section 10.1 (Val Tannen; \path{reports/VerificationTasks.pdf}) \par the reformulations task document\par}

\item[Prophecy] kernel-checked certificate with a level-1 prophecy clause (Certificate/Valid.lean, Core.json).

\end{description}

\paragraph{Datalog program(s).}

{\raggedright\textit{globalProgram --- `T xy :- Exz, Gz, Ezy` and `T xy :- Exz, Gz, T zy`.} --- fidelity module (kernel-checked: the compiled sides of the command are this program's naive compilation).\par}

\begin{Verbatim}[fontsize=\scriptsize,breaklines=true,breakanywhere=true,frame=single,framesep=1.5mm]
T(x1, y1) :- E(x1, z1), G(z1), E(z1, y1).
T(x1, y1) :- E(x1, z1), G(z1), T(z1, y1).
\end{Verbatim}

{\raggedright\textit{localProgram --- `T↑ xy :- Ve xz, Vb zy` and `T↑ xy :- Ve xz, T↑ zy`.} --- fidelity module (kernel-checked: the compiled sides of the command are this program's naive compilation).\par}

\begin{Verbatim}[fontsize=\scriptsize,breaklines=true,breakanywhere=true,frame=single,framesep=1.5mm]
TUp(x1, y1) :- Ve(x1, z1), Vb(z1, y1).
TUp(x1, y1) :- Ve(x1, z1), TUp(z1, y1).
\end{Verbatim}

\paragraph{Reading.} The global program T computes walks of length at least two through digraph E whose interior vertices are all "good" (marked by unary relation G): T(x,y):-E(x,z),G(z),E(z,y), and T(x,y):-E(x,z),G(z),T(z,y) --- a good-interior-vertex reachability query. The local reformulation TUp is defined over two exact conjunctive views, Vb (begin-good edges, G(x) ∧ E(x,y)) and Ve (end-good edges, E(x,y) ∧ G(y)): TUp(x,y):-Ve(x,z),Vb(z,y), and TUp(x,y):-Ve(x,z),TUp(z,y). The Hoare triple claims the two programs compute the same relation, T = TUp --- an equivalence claim that follows from substituting the view definitions into each TUp rule to recover the matching T rule. Expected valid; no certificate exists yet, so the classification stays unknown.

\paragraph{First-order reading of PRE/POST, translated mechanically from the relational-algebra assertion in Input.lean.} The relational-algebra assertion in \texttt{Input.lean} remains the authority.

\textit{PRE:} (∀x y. (Vb(x,y) ↔ G(x) ∧ E(x,y))) ∧ ∀z u. (Ve(z,u) ↔ E(z,u) ∧ G(u))\par The precondition states the two view equations exactly: Vb is the join of G and E giving good-start edges, and Ve is the join of E and G giving good-end edges.

\textit{POST:} ∀x y. (T(x,y) ↔ TUp(x,y))\par The postcondition asserts T = TUp, that the good-interior path query and its two-view reformulation agree.

\paragraph{Whiel encoding (Hoare triple).}

\begin{Verbatim}[fontsize=\tiny,breaklines=true,breakanywhere=true,frame=single,framesep=1.5mm]
SCHEMA
  E/2, T/2, Vb/2, Ve/2, TUp/2, T_aux/2, TUp_aux/2, G/1

PRE
  ((Vb = (π[1, 2] (σ[#0 = #1] (G × E)))) ∧
    (Ve = (π[0, 1] (σ[#1 = #2] (E × G)))))

CMD
  T_aux := ∅[2];
  T :=
    ((π[0, 4]
        (σ[(#1 = #2 ∧ #1 = #3)] (E × (G × E)))) ∪
      (π[0, 4]
        (σ[(#1 = #2 ∧ #1 = #3)]
          (E × (G × T_aux)))));
  WHILE (¬ (T = T_aux)) DO
    T_aux := T;
    T :=
      (T ∪
        ((π[0, 4]
            (σ[(#1 = #2 ∧ #1 = #3)]
              (E × (G × E)))) ∪
          (π[0, 4]
            (σ[(#1 = #2 ∧ #1 = #3)]
              (E × (G × T_aux))))))
  END;
  TUp_aux := ∅[2];
  TUp :=
    ((π[0, 3] (σ[#1 = #2] (Ve × Vb))) ∪
      (π[0, 3] (σ[#1 = #2] (Ve × TUp_aux))));
  WHILE (¬ (TUp = TUp_aux)) DO
    TUp_aux := TUp;
    TUp :=
      (TUp ∪
        ((π[0, 3] (σ[#1 = #2] (Ve × Vb))) ∪
          (π[0, 3]
            (σ[#1 = #2] (Ve × TUp_aux)))))
  END

POST
  (T = TUp)
\end{Verbatim}

\Needspace{0.45\textheight}

\section{Example4036: Airline reachability query versus its three-view reformulation, forward containment}\label{case:Example4036}

\begin{description}[leftmargin=2.6cm,style=nextline,itemsep=1pt]

\item[Category] verification task (VerificationTasks.pdf, section 10.1); family: Core.

\item[Kind] containment.

\item[Contributors] Leo Zhang, Val Tannen.

\item[Status] certified valid (Certificate/Valid.lean, written by the certificate emitter); expected valid.

\item[Tags] \hyperref[tag:prophecy]{\texttt{prophecy}} \hyperref[tag:views]{\texttt{views}}.

\item[Input] \path{Benchmark/Example4036/Input.lean} (canonical notation).

\item[Summary] Single-airline reachability and a reformulation over three flight views, tested for containment.

\item[Sources] {\raggedright VerificationTasks.pdf, section 10.1 (Val Tannen; \path{reports/VerificationTasks.pdf}) \par the reformulations task document\par}

\item[Prophecy] kernel-checked certificate with a level-1 prophecy clause (Certificate/Valid.lean, Core.json).

\item[Note] In an earlier (rate-matched) encoding this containment had a quantifier-free certificate; whether the compiled-sequence encoding needs a prophecy level has not been checked.

\end{description}

\paragraph{Datalog program(s).}

{\raggedright\textit{globalProgram --- `T xyu :- Fxyu`, `T xyu :- Fxzu, T zyu`, and `Qy :- T byu`.} --- fidelity module (kernel-checked: the compiled sides of the command are this program's naive compilation).\par}

\begin{Verbatim}[fontsize=\scriptsize,breaklines=true,breakanywhere=true,frame=single,framesep=1.5mm]
T(x1, y1, x2) :- F(x1, y1, x2).
T(x1, y1, x2) :- F(x1, z1, x2), T(z1, y1, x2).
Q(y1) :- T("b", y1, x2).
\end{Verbatim}

{\raggedright\textit{localProgram --- `Q↑ y :- Va by`, `Q↑ z :- Vb yc, Vc yz`, and `Q↑ y :- Vb yu`.} --- fidelity module (kernel-checked: the compiled sides of the command are this program's naive compilation).\par}

\begin{Verbatim}[fontsize=\scriptsize,breaklines=true,breakanywhere=true,frame=single,framesep=1.5mm]
QUp(y1) :- Va("b", y1).
QUp(z1) :- Vb(y1, "c"), Vc(y1, z1).
QUp(y1) :- Vb(y1, x2).
\end{Verbatim}

\paragraph{Reading.} The global program computes per-airline reachability T over a ternary flight relation F(x,y,u) (non-stop flight from x to y on airline u), then restricts to Q(y):-T("b",y,u), the airports reachable from "b" while staying on a single airline. The local reformulation QUp is built, without recursion, from three conjunctive views: Va (direct flights on airline "a"), Vb (flights out of "b"), and Vc (two-hop connections on airline "c"), combined by three rules. The Hoare triple claims a one-directional (forward) containment, QUp ⊆ Q: every airport the reformulation derives is genuinely reachable from "b" on one airline. Expected valid, since each QUp rule unfolds to a genuine same-airline itinerary out of "b"; no certificate exists yet, so the classification stays unknown. This case is the twin of Example4037, which tests the reverse direction of the same containment.

\paragraph{First-order reading of PRE/POST, translated mechanically from the relational-algebra assertion in Input.lean.} The relational-algebra assertion in \texttt{Input.lean} remains the authority.

\textit{PRE:} (∀x y. (Va(x,y) ↔ F(x,y,"a"))) ∧ (∀z u. (Vb(z,u) ↔ F("b",z,u))) ∧ ∀v w. (Vc(v,w) ↔ ∃x1. F(v,x1,"c") ∧ F(x1,w,"c"))\par The precondition states the three view equations: Va is F restricted to airline "a", Vb is F restricted to flights departing "b", and Vc is the two-hop composition of two airline-"c" flights.

\textit{POST:} ∀x. (QUp(x) → Q(x))\par The postcondition asserts QUp ⊆ Q, that every airport reached by the three-view reformulation is reached by the global per-airline query.

\paragraph{Whiel encoding (Hoare triple).}

\begin{Verbatim}[fontsize=\tiny,breaklines=true,breakanywhere=true,frame=single,framesep=1.5mm]
SCHEMA
  Q/1, QUp/1, Q_aux/1, QUp_aux/1, Va/2, Vb/2, Vc/2, F/3, T/3, T_aux/3

PRE
  ((Va = (π[0, 1] (σ[#2 = "a"] F))) ∧
    ((Vb = (π[1, 2] (σ[#0 = "b"] F))) ∧
      (Vc =
        (π[0, 4]
          (σ[((#2 = "c" ∧ #1 = #3) ∧ #5 = "c")]
            (F × F))))))

CMD
  T_aux := ∅[3];
  Q_aux := ∅[1];
  T :=
    (F ∪
      (π[0, 4, 2]
        (σ[(#1 = #3 ∧ #2 = #5)] (F × T_aux))));
  Q := (π[1] (σ[#0 = "b"] T_aux));
  WHILE (¬ ((T = T_aux) ∧ (Q = Q_aux))) DO
    T_aux := T;
    Q_aux := Q;
    T :=
      (T ∪
        (F ∪
          (π[0, 4, 2]
            (σ[(#1 = #3 ∧ #2 = #5)]
              (F × T_aux)))));
    Q := (Q ∪ (π[1] (σ[#0 = "b"] T_aux)))
  END;
  QUp_aux := ∅[1];
  QUp :=
    ((π[1] (σ[#0 = "b"] Va)) ∪
      ((π[3]
          (σ[(#1 = "c" ∧ #0 = #2)] (Vb × Vc))) ∪
        (π[0] Vb)));
  WHILE (¬ (QUp = QUp_aux)) DO
    QUp_aux := QUp;
    QUp :=
      (QUp ∪
        ((π[1] (σ[#0 = "b"] Va)) ∪
          ((π[3]
              (σ[(#1 = "c" ∧ #0 = #2)]
                (Vb × Vc))) ∪
            (π[0] Vb))))
  END

POST
  (QUp ⊆ Q)
\end{Verbatim}

\Needspace{0.45\textheight}

\section{Example4037: Airline reachability query versus its three-view reformulation, reverse containment}\label{case:Example4037}

\begin{description}[leftmargin=2.6cm,style=nextline,itemsep=1pt]

\item[Category] verification task (VerificationTasks.pdf, section 10.1); family: Core.

\item[Kind] containment.

\item[Contributors] Leo Zhang, Val Tannen.

\item[Status] certified invalid (Certificate/Invalid.lean, written by the certificate emitter); expected invalid.

\item[Tags] \hyperref[tag:views]{\texttt{views}}.

\item[Input] \path{Benchmark/Example4037/Input.lean} (canonical notation).

\item[Summary] Single-airline reachability and a reformulation over three flight views, tested for the reverse containment.

\item[Sources] {\raggedright VerificationTasks.pdf, section 10.1 (Val Tannen; \path{reports/VerificationTasks.pdf}) \par the reformulations task document\par}

\end{description}

\paragraph{Datalog program(s).}

{\raggedright\textit{globalProgram --- `T xyu :- Fxyu`, `T xyu :- Fxzu, T zyu`, and `Qy :- T byu`.} --- fidelity module (kernel-checked: the compiled sides of the command are this program's naive compilation).\par}

\begin{Verbatim}[fontsize=\scriptsize,breaklines=true,breakanywhere=true,frame=single,framesep=1.5mm]
T(x1, y1, x2) :- F(x1, y1, x2).
T(x1, y1, x2) :- F(x1, z1, x2), T(z1, y1, x2).
Q(y1) :- T("b", y1, x2).
\end{Verbatim}

{\raggedright\textit{localProgram --- `Q↑ y :- Va by`, `Q↑ z :- Vb yc, Vc yz`, and `Q↑ y :- Vb yu`.} --- fidelity module (kernel-checked: the compiled sides of the command are this program's naive compilation).\par}

\begin{Verbatim}[fontsize=\scriptsize,breaklines=true,breakanywhere=true,frame=single,framesep=1.5mm]
QUp(y1) :- Va("b", y1).
QUp(z1) :- Vb(y1, "c"), Vc(y1, z1).
QUp(y1) :- Vb(y1, x2).
\end{Verbatim}

\paragraph{Reading.} This case uses the same global and local programs, views, and precondition as Example4036 (per-airline reachability T restricted to departures from "b" for Q, versus the three-view reformulation QUp), but tests the opposite direction of containment, Q ⊆ QUp: that every airport genuinely reachable from "b" on a single airline is also captured by the reformulation. This claim is invalid, certified by a kernel-checked counterexample: the reformulation is only maximally (forward-) contained in Q, not surjective onto it, so a two-hop route out of "b" on an airline other than "a" or "c" is invisible to all three views. The witness is F = \{("b","x","d"), ("x","y","d")\}, where at exit Q = \{"x","y"\} but QUp = \{"x"\}.

\paragraph{First-order reading of PRE/POST, translated mechanically from the relational-algebra assertion in Input.lean.} The relational-algebra assertion in \texttt{Input.lean} remains the authority.

\textit{PRE:} (∀x y. (Va(x,y) ↔ F(x,y,"a"))) ∧ (∀z u. (Vb(z,u) ↔ F("b",z,u))) ∧ ∀v w. (Vc(v,w) ↔ ∃x1. F(v,x1,"c") ∧ F(x1,w,"c"))\par The precondition states the same three view equations as Example4036: Va is F restricted to airline "a", Vb is F restricted to flights departing "b", and Vc is the two-hop composition of two airline-"c" flights.

\textit{POST:} ∀x. (Q(x) → QUp(x))\par The postcondition asserts Q ⊆ QUp, that every airport genuinely reachable from "b" on one airline is captured by the reformulation --- a claim that fails.

\paragraph{Whiel encoding (Hoare triple).}

\begin{Verbatim}[fontsize=\tiny,breaklines=true,breakanywhere=true,frame=single,framesep=1.5mm]
SCHEMA
  Q/1, QUp/1, Q_aux/1, QUp_aux/1, Va/2, Vb/2, Vc/2, F/3, T/3, T_aux/3

PRE
  ((Va = (π[0, 1] (σ[#2 = "a"] F))) ∧
    ((Vb = (π[1, 2] (σ[#0 = "b"] F))) ∧
      (Vc =
        (π[0, 4]
          (σ[((#2 = "c" ∧ #1 = #3) ∧ #5 = "c")]
            (F × F))))))

CMD
  T_aux := ∅[3];
  Q_aux := ∅[1];
  T :=
    (F ∪
      (π[0, 4, 2]
        (σ[(#1 = #3 ∧ #2 = #5)] (F × T_aux))));
  Q := (π[1] (σ[#0 = "b"] T_aux));
  WHILE (¬ ((T = T_aux) ∧ (Q = Q_aux))) DO
    T_aux := T;
    Q_aux := Q;
    T :=
      (T ∪
        (F ∪
          (π[0, 4, 2]
            (σ[(#1 = #3 ∧ #2 = #5)]
              (F × T_aux)))));
    Q := (Q ∪ (π[1] (σ[#0 = "b"] T_aux)))
  END;
  QUp_aux := ∅[1];
  QUp :=
    ((π[1] (σ[#0 = "b"] Va)) ∪
      ((π[3]
          (σ[(#1 = "c" ∧ #0 = #2)] (Vb × Vc))) ∪
        (π[0] Vb)));
  WHILE (¬ (QUp = QUp_aux)) DO
    QUp_aux := QUp;
    QUp :=
      (QUp ∪
        ((π[1] (σ[#0 = "b"] Va)) ∪
          ((π[3]
              (σ[(#1 = "c" ∧ #0 = #2)]
                (Vb × Vc))) ∪
            (π[0] Vb))))
  END

POST
  (Q ⊆ QUp)
\end{Verbatim}

\Needspace{0.45\textheight}

\section{Example4040: Closure of red and blue edges versus its reformulation over Datalog path views}\label{case:Example4040}

\begin{description}[leftmargin=2.6cm,style=nextline,itemsep=1pt]

\item[Category] verification task (VerificationTasks.pdf, section 10.1); family: Core.

\item[Kind] equivalence.

\item[Contributors] Leo Zhang, Val Tannen.

\item[Status] certified valid (certificate checked outside the library build) (Certificate/Valid.lean, written by the certificate emitter); expected valid. Placement note: a proof module imported by \path{Certificate/Valid.lean} does not build under the library build's 4194304 KB limit; it builds under 12582912 KB (peak 1651 MB)

\item[Tags] \hyperref[tag:linear-vs-nonlinear]{\texttt{linear-vs-nonlinear}} \hyperref[tag:nonlinear-recursion]{\texttt{nonlinear-recursion}} \hyperref[tag:prophecy]{\texttt{prophecy}} \hyperref[tag:views]{\texttt{views}}.

\item[Input] \path{Benchmark/Example4040/Input.lean} (canonical notation).

\item[Summary] Two-colour reachability and a reformulation read off two recursive path views.

\item[Sources] {\raggedright VerificationTasks.pdf, section 10.1 (Val Tannen; \path{reports/VerificationTasks.pdf}) \par the reformulations task document\par}

\item[Prophecy] kernel-checked certificate with a level-1 prophecy clause (Certificate/Valid.lean, Core.json).

\end{description}

\paragraph{Datalog program(s).}

{\raggedright\textit{globalProgram --- The transitive closure of `R ∪ B`, left-linearly.} --- fidelity module (kernel-checked: the compiled sides of the command are this program's naive compilation).\par}

\begin{Verbatim}[fontsize=\scriptsize,breaklines=true,breakanywhere=true,frame=single,framesep=1.5mm]
T(x1, y1) :- R(x1, y1).
T(x1, y1) :- B(x1, y1).
T(x1, y1) :- R(x1, z1), T(z1, y1).
T(x1, y1) :- B(x1, z1), T(z1, y1).
\end{Verbatim}

{\raggedright\textit{viewProgram --- The red and the blue monochrome path views.} --- fidelity module (kernel-checked: the compiled sides of the command are this program's naive compilation).\par}

\begin{Verbatim}[fontsize=\scriptsize,breaklines=true,breakanywhere=true,frame=single,framesep=1.5mm]
Vr(x1, y1) :- R(x1, y1).
Vr(x1, y1) :- R(x1, z1), Vr(z1, y1).
Vb(x1, y1) :- B(x1, y1).
Vb(x1, y1) :- B(x1, z1), Vb(z1, y1).
\end{Verbatim}

{\raggedright\textit{reformProgram --- The reformulation over the two path views.} --- fidelity module (kernel-checked: the compiled sides of the command are this program's naive compilation).\par}

\begin{Verbatim}[fontsize=\scriptsize,breaklines=true,breakanywhere=true,frame=single,framesep=1.5mm]
TUp(x1, y1) :- Vr(x1, y1).
TUp(x1, y1) :- Vb(x1, y1).
TUp(x1, y1) :- TUp(x1, z1), TUp(z1, y1).
\end{Verbatim}

{\raggedright\textit{program} --- fidelity module (kernel-checked: the compiled sides of the command are this program's naive compilation).\par}

\begin{Verbatim}[fontsize=\scriptsize,breaklines=true,breakanywhere=true,frame=single,framesep=1.5mm]
T(x1, y1) :- R(x1, y1).
T(x1, y1) :- B(x1, y1).
T(x1, y1) :- R(x1, z1), T(z1, y1).
T(x1, y1) :- B(x1, z1), T(z1, y1).
T(x1, y1) :- T(x1, z1), T(z1, y1).
\end{Verbatim}

\paragraph{Reading.} The global program T computes the transitive closure of R ∪ B (red and blue edges of a digraph combined), advancing one edge at a time. The local reformulation is itself recursive Datalog: it first computes two monochrome path views, Vr (red-only paths) and Vb (blue-only paths), each by its own right-linear recursion over R or B, then defines TUp(x,y):-Vr(x,y), TUp(x,y):-Vb(x,y), and the nonlinear rule TUp(x,y):-TUp(x,z),TUp(z,y), composing whole monochrome walks together. The Hoare triple claims equality, T = TUp: both compute the same transitive closure of R ∪ B, one by single-edge steps and the other by self-composing entire same-colour walks. Certified valid.

\paragraph{First-order reading of PRE/POST, translated mechanically from the relational-algebra assertion in Input.lean.} The relational-algebra assertion in \texttt{Input.lean} remains the authority.

\textit{PRE:} ⊤\par The precondition is simply true; there is no separate view hypothesis, since the views Vr and Vb are themselves computed by Datalog programs within the command.

\textit{POST:} ∀x y. (T(x,y) ↔ TUp(x,y))\par The postcondition asserts T = TUp, that the closure of the combined red/blue edges equals the path-view-composition reformulation.

\paragraph{Whiel encoding (Hoare triple).}

\begin{Verbatim}[fontsize=\tiny,breaklines=true,breakanywhere=true,frame=single,framesep=1.5mm]
SCHEMA
  R/2, B/2, T/2, Vr/2, Vb/2, TUp/2, T_aux/2, Vr_aux/2, Vb_aux/2, TUp_aux/2

PRE
  true

CMD
  T_aux := ∅[2];
  T :=
    (R ∪
      (B ∪
        ((π[0, 3] (σ[#1 = #2] (R × T_aux))) ∪
          (π[0, 3] (σ[#1 = #2] (B × T_aux))))));
  WHILE (¬ (T = T_aux)) DO
    T_aux := T;
    T :=
      (T ∪
        (R ∪
          (B ∪
            ((π[0, 3]
                (σ[#1 = #2] (R × T_aux))) ∪
              (π[0, 3]
                (σ[#1 = #2] (B × T_aux)))))))
  END;
  Vr_aux := ∅[2];
  Vb_aux := ∅[2];
  Vr :=
    (R ∪ (π[0, 3] (σ[#1 = #2] (R × Vr_aux))));
  Vb :=
    (B ∪ (π[0, 3] (σ[#1 = #2] (B × Vb_aux))));
  WHILE (¬ ((Vr = Vr_aux) ∧ (Vb = Vb_aux))) DO
    Vr_aux := Vr;
    Vb_aux := Vb;
    Vr :=
      (Vr ∪
        (R ∪
          (π[0, 3] (σ[#1 = #2] (R × Vr_aux)))));
    Vb :=
      (Vb ∪
        (B ∪
          (π[0, 3] (σ[#1 = #2] (B × Vb_aux)))))
  END;
  TUp_aux := ∅[2];
  TUp :=
    (Vr ∪
      (Vb ∪
        (π[0, 3]
          (σ[#1 = #2] (TUp_aux × TUp_aux)))));
  WHILE (¬ (TUp = TUp_aux)) DO
    TUp_aux := TUp;
    TUp :=
      (TUp ∪
        (Vr ∪
          (Vb ∪
            (π[0, 3]
              (σ[#1 = #2]
                (TUp_aux × TUp_aux))))))
  END

POST
  (T = TUp)
\end{Verbatim}

\Needspace{0.45\textheight}

\section{Example4041: Transitive closure versus its reformulation over exact odd/even views}\label{case:Example4041}

\begin{description}[leftmargin=2.6cm,style=nextline,itemsep=1pt]

\item[Category] verification task (VerificationTasks.pdf, section 10.1); family: Core.

\item[Kind] equivalence.

\item[Contributors] Jesse Comer, Val Tannen.

\item[Status] certified valid (Certificate/Valid.lean, written by the certificate emitter); expected valid.

\item[Tags] \hyperref[tag:mutual-recursion]{\texttt{mutual-recursion}} \hyperref[tag:prophecy]{\texttt{prophecy}} \hyperref[tag:views]{\texttt{views}}.

\item[Input] \path{Benchmark/Example4041/Input.lean} (canonical notation).

\item[Summary] Transitive closure and a reformulation read off two exact odd/even views.

\item[Sources] {\raggedright VerificationTasks.pdf, section 10.1 (Val Tannen; \path{reports/VerificationTasks.pdf}) \par the reformulations task document\par}

\item[Prophecy] kernel-checked certificate with a level-1 prophecy clause (Certificate/Valid.lean, Core.json).

\end{description}

\paragraph{Datalog program(s).}

{\raggedright\textit{globalProgram --- `T xy :- Exy` and `T xy :- Exz, T zy`.} --- fidelity module (kernel-checked: the compiled sides of the command are this program's naive compilation).\par}

\begin{Verbatim}[fontsize=\scriptsize,breaklines=true,breakanywhere=true,frame=single,framesep=1.5mm]
T(x1, y1) :- E(x1, y1).
T(x1, y1) :- E(x1, z1), T(z1, y1).
\end{Verbatim}

{\raggedright\textit{viewProgram --- The odd/even program that computes the two views.} --- fidelity module (kernel-checked: the compiled sides of the command are this program's naive compilation).\par}

\begin{Verbatim}[fontsize=\scriptsize,breaklines=true,breakanywhere=true,frame=single,framesep=1.5mm]
To(x1, y1) :- E(x1, y1).
To(x1, y1) :- E(x1, z1), Te(z1, y1).
Te(x1, y1) :- E(x1, z1), To(z1, y1).
\end{Verbatim}

{\raggedright\textit{reformProgram --- The reformulation `T↑ :- To`, `T↑ :- Te`.} --- fidelity module (kernel-checked: the compiled sides of the command are this program's naive compilation).\par}

\begin{Verbatim}[fontsize=\scriptsize,breaklines=true,breakanywhere=true,frame=single,framesep=1.5mm]
TUp(x1, y1) :- To(x1, y1).
TUp(x1, y1) :- Te(x1, y1).
\end{Verbatim}

\paragraph{Reading.} The global program T is the standard right-linear transitive closure of E. The local reformulation computes the odd/even path decomposition of E by mutual recursion (To for odd-length paths, Te for even-length paths) and then unions them, TUp(x,y):-To(x,y), TUp(x,y):-Te(x,y). Because the views are exact --- computed within the same command rather than assumed as free source tables --- the precondition is trivially true and correctness is carried entirely by the computation itself. The Hoare triple claims equality, TUp = T: the union of the odd- and even-length-path relations is exactly the transitive closure. Certified valid, via an eleven-clause loop invariant (eight clauses at level 0, three at level 1) closing all 23 verification obligations.

\paragraph{First-order reading of PRE/POST, translated mechanically from the relational-algebra assertion in Input.lean.} The relational-algebra assertion in \texttt{Input.lean} remains the authority.

\textit{PRE:} ⊤\par The precondition is simply true.

\textit{POST:} ∀x y. (TUp(x,y) ↔ T(x,y))\par The postcondition asserts TUp = T, that the union of the odd- and even-length path views equals the transitive closure.

\paragraph{Whiel encoding (Hoare triple).}

\begin{Verbatim}[fontsize=\tiny,breaklines=true,breakanywhere=true,frame=single,framesep=1.5mm]
SCHEMA
  E/2, T/2, To/2, Te/2, TUp/2, T_aux/2, To_aux/2, Te_aux/2, TUp_aux/2

PRE
  true 

CMD
  T_aux := ∅[2];
  T :=
    (E ∪ (π[0, 3] (σ[#1 = #2] (E × T_aux))));
  WHILE (¬ (T = T_aux)) DO
    T_aux := T;
    T :=
      (T ∪
        (E ∪
          (π[0, 3] (σ[#1 = #2] (E × T_aux)))))
  END;
  To_aux := ∅[2];
  Te_aux := ∅[2];
  To :=
    (E ∪ (π[0, 3] (σ[#1 = #2] (E × Te_aux))));
  Te := (π[0, 3] (σ[#1 = #2] (E × To_aux)));
  WHILE (¬ ((To = To_aux) ∧ (Te = Te_aux))) DO
    To_aux := To;
    Te_aux := Te;
    To :=
      (To ∪
        (E ∪
          (π[0, 3] (σ[#1 = #2] (E × Te_aux)))));
    Te :=
      (Te ∪
        (π[0, 3] (σ[#1 = #2] (E × To_aux))))
  END;
  TUp_aux := ∅[2];
  TUp := (To ∪ Te);
  WHILE (¬ (TUp = TUp_aux)) DO
    TUp_aux := TUp;
    TUp := (TUp ∪ (To ∪ Te))
  END

POST
  (TUp = T)
\end{Verbatim}

\Needspace{0.45\textheight}

\section{Example5001: PostgreSQL recursive UNION versus naive seeded closure}\label{case:Example5001}

\begin{description}[leftmargin=2.6cm,style=nextline,itemsep=1pt]

\item[Category] production pair (two algorithms from production sources or legacy production cases); family: direct-equality pair.

\item[Kind] equivalence.

\item[Contributors] Leo Zhang.

\item[Status] certified valid (Certificate/Valid.lean, written by the certificate emitter); expected valid.

\item[Tags] \hyperref[tag:database-engine]{\texttt{database-engine}} \hyperref[tag:evaluation-strategies]{\texttt{evaluation-strategies}} \hyperref[tag:prophecy]{\texttt{prophecy}} \hyperref[tag:semi-naive]{\texttt{semi-naive}}.

\item[Input] \path{Benchmark/Example5001/Input.lean} (canonical notation).

\item[Summary] PostgreSQL's WITH RECURSIVE executor against the naive evaluation of the same recursive query.

\item[Sides] legacy Example2037 (PostgreSQL nodeRecursiveunion.c model) \par naive compiler on the seeded closure program

\item[Sources] {\raggedright \href{https://github.com/postgres/postgres/blob/master/src/backend/executor/nodeRecursiveunion.c}{postgres/postgres src/backend/executor/nodeRecursiveunion.c (ExecRecursiveUnion)} \par \href{https://www.postgresql.org/docs/current/queries-with.html}{PostgreSQL documentation, WITH queries (recursive)}\par}

\item[Prophecy] kernel-checked certificate with a level-1 prophecy clause (Certificate/Valid.lean, Core.json).

\end{description}

\paragraph{Datalog program(s).}

{\raggedright\textit{prog\_p} --- fidelity module (kernel-checked: the compiled sides of the command are this program's naive compilation).\par}

\begin{Verbatim}[fontsize=\scriptsize,breaklines=true,breakanywhere=true,frame=single,framesep=1.5mm]
Tcl(x1, y1) :- Base(x1, y1).
Tcl(x1, y1) :- Tcl(x1, z1), Edge(z1, y1).
\end{Verbatim}

\paragraph{Reading.} Both programs compute Base ∘ Edge*, the rows reachable from a base relation by following zero or more Edge steps --- a standard recursive-transitive-closure computation of the kind PostgreSQL's WITH RECURSIVE executor performs. Side 1 models that executor: it maintains a working table WorkT holding only the previous round's newly derived rows, joins WorkT with Edge, subtracts rows already in the accumulated result Res, and halts once nothing new is produced --- a frontier/delta-based fixpoint evaluation. Side 2 is the standard naive evaluation of Tcl(x,y):-Base(x,y), Tcl(x,y):-Tcl(x,z),Edge(z,y), recomputing the whole relation each round. Run one after the other as independent programs, the Hoare triple claims the two agree, Res = Tcl --- the standard correctness statement that a delta-restricted (semi-naive-style) executor computes the same least fixpoint as full naive evaluation.

\paragraph{First-order reading of PRE/POST, translated mechanically from the relational-algebra assertion in Input.lean.} The relational-algebra assertion in \texttt{Input.lean} remains the authority.

\textit{PRE:} ⊤\par The precondition is simply true.

\textit{POST:} ∀x y. (Res(x,y) ↔ Tcl(x,y))\par The postcondition asserts Res = Tcl, that the PostgreSQL-style recursive-union executor and naive evaluation compute the same closure.

\paragraph{Whiel encoding (Hoare triple).}

\begin{Verbatim}[fontsize=\scriptsize,breaklines=true,breakanywhere=true,frame=single,framesep=1.5mm]
SCHEMA
  Base/2, Edge/2, Res/2, WorkT/2, Inter/2, Tcl/2, Tcl_aux/2

PRE
  true 

CMD
  Res := Base;
  WorkT := Base;
  WHILE (WorkT ≠ ∅) DO
    Inter := ((π[0, 3] (σ[#1 = #2] (WorkT × Edge))) ∖ Res);
    Res := (Res ∪ Inter);
    WorkT := Inter
  END;
  Tcl_aux := ∅[2];
  Tcl := (Base ∪ π[0,3] (σ[#1 = #2] ((Tcl_aux × Edge))));
  WHILE (¬((Tcl = Tcl_aux))) DO
    Tcl_aux := Tcl;
    Tcl := (Tcl ∪ (Base ∪ π[0,3] (σ[#1 = #2] ((Tcl_aux × Edge)))))
  END

POST
  (Res = Tcl)
\end{Verbatim}

\Needspace{0.45\textheight}

\section{Example5002: Datafrog semi-naive delta evaluation versus naive evaluation}\label{case:Example5002}

\begin{description}[leftmargin=2.6cm,style=nextline,itemsep=1pt]

\item[Category] production pair (two algorithms from production sources or legacy production cases); family: direct-equality pair.

\item[Kind] equivalence.

\item[Contributors] Leo Zhang.

\item[Status] certified valid (Certificate/Valid.lean, written by the certificate emitter); expected valid.

\item[Tags] \hyperref[tag:database-engine]{\texttt{database-engine}} \hyperref[tag:evaluation-strategies]{\texttt{evaluation-strategies}} \hyperref[tag:prophecy]{\texttt{prophecy}} \hyperref[tag:semi-naive]{\texttt{semi-naive}}.

\item[Input] \path{Benchmark/Example5002/Input.lean} (canonical notation).

\item[Summary] Datafrog's semi-naive delta evaluation of a reachability query against the naive evaluation of the same query.

\item[Sides] legacy Example2039 (Datafrog / Soufflé semi-naive delta loop) \par naive compiler on the transposed reachability program

\item[Sources] {\raggedright \href{https://github.com/rust-lang/datafrog}{rust-lang/datafrog (README example, src/lib.rs Variable/from\_join/changed)}\par}

\item[Prophecy] kernel-checked certificate with a level-1 prophecy clause (Certificate/Valid.lean, Core.json).

\item[Note] Kept after the duplicate review: the loop coincides with 5001 up to transposition and seeding, but it is the corpus's only Datafrog case (2039 was retired).

\end{description}

\paragraph{Datalog program(s).}

{\raggedright\textit{prog\_p} --- fidelity module (kernel-checked: the compiled sides of the command are this program's naive compilation).\par}

\begin{Verbatim}[fontsize=\scriptsize,breaklines=true,breakanywhere=true,frame=single,framesep=1.5mm]
Reach(y1, x1) :- EdgeF(x1, y1).
Reach(z1, x1) :- Reach(y1, x1), EdgeF(y1, z1).
\end{Verbatim}

\paragraph{Reading.} Both sides compute the transpose of the transitive closure of EdgeF: Nodes/Reach(y,x) holds when y is reachable from x by one or more EdgeF steps. Side 1 models Datafrog's (and similarly Soufflé's) semi-naive delta evaluation: it maintains only the newly derived facts DeltaN from the previous round, joins them with EdgeF, subtracts what is already known, and halts when the delta becomes empty. Side 2 is naive evaluation of Reach(y,x):-EdgeF(x,y), Reach(z,x):-Reach(y,x),EdgeF(y,z), recomputing the full relation every round. The Hoare triple, running the two independent programs side by side, claims equality, Nodes = Reach --- the standard semi-naive-evaluation correctness statement, that restricting joins to newly derived facts each round yields the same least fixpoint as full naive evaluation.

\paragraph{First-order reading of PRE/POST, translated mechanically from the relational-algebra assertion in Input.lean.} The relational-algebra assertion in \texttt{Input.lean} remains the authority.

\textit{PRE:} ⊤\par The precondition is simply true.

\textit{POST:} ∀x y. (Nodes(x,y) ↔ Reach(x,y))\par The postcondition asserts Nodes = Reach, that the semi-naive delta evaluation and naive evaluation compute the same reachability relation.

\paragraph{Whiel encoding (Hoare triple).}

\begin{Verbatim}[fontsize=\scriptsize,breaklines=true,breakanywhere=true,frame=single,framesep=1.5mm]
SCHEMA
  EdgeF/2, Nodes/2, DeltaN/2, NewN/2, Reach/2, Reach_aux/2

PRE
  true 

CMD
  Nodes := (π[1, 0] EdgeF);
  DeltaN := (π[1, 0] EdgeF);
  WHILE (DeltaN ≠ ∅) DO
    NewN := ((π[3, 1] (σ[#0 = #2] (DeltaN × EdgeF))) ∖ Nodes);
    Nodes := (Nodes ∪ NewN);
    DeltaN := NewN
  END;
  Reach_aux := ∅[2];
  Reach := (π[1,0] (EdgeF) ∪ π[3,1] (σ[#0 = #2] ((Reach_aux × EdgeF))));
  WHILE (¬((Reach = Reach_aux))) DO
    Reach_aux := Reach;
    Reach := (Reach ∪ (π[1,0] (EdgeF) ∪ π[3,1] (σ[#0 = #2] ((Reach_aux × EdgeF)))))
  END

POST
  (Nodes = Reach)
\end{Verbatim}

\Needspace{0.45\textheight}

\section{Example5003: Worklist breadth-first reachability versus naive evaluation}\label{case:Example5003}

\begin{description}[leftmargin=2.6cm,style=nextline,itemsep=1pt]

\item[Category] production pair (two algorithms from production sources or legacy production cases); family: direct-equality pair.

\item[Kind] equivalence.

\item[Contributors] Leo Zhang.

\item[Status] certified valid (Certificate/Valid.lean, written by the certificate emitter); expected valid.

\item[Tags] \hyperref[tag:evaluation-strategies]{\texttt{evaluation-strategies}} \hyperref[tag:prophecy]{\texttt{prophecy}} \hyperref[tag:semi-naive]{\texttt{semi-naive}}.

\item[Input] \path{Benchmark/Example5003/Input.lean} (canonical notation).

\item[Summary] A worklist breadth-first search against the naive evaluation of the reachability query it implements.

\item[Sides] hard case Example0161 (worklist BFS with explicit frontier) \par naive compiler on the reachability program

\item[Sources] {\raggedright hard-case source records of the retired classic-source collection (hard-*.json)\par}

\item[Prophecy] kernel-checked certificate with a level-1 prophecy clause (Certificate/Valid.lean, Core.json).

\end{description}

\paragraph{Datalog program(s).}

{\raggedright\textit{prog\_p} --- fidelity module (kernel-checked: the compiled sides of the command are this program's naive compilation).\par}

\begin{Verbatim}[fontsize=\scriptsize,breaklines=true,breakanywhere=true,frame=single,framesep=1.5mm]
Reach(x1) :- Src(x1).
Reach(y1) :- Reach(x1), E(x1, y1).
\end{Verbatim}

\paragraph{Reading.} Both sides compute the set of vertices reachable from a source set Src over edge relation E. Side 1 is an explicit worklist/frontier breadth-first search: R is the visited set and W the current frontier; each round computes the newly discovered vertices N as E-successors of W not already in R, adds them to R, and advances the frontier W := N --- the classic BFS-as-fixpoint-iteration pattern. Side 2 is naive Datalog evaluation of Reach(x):-Src(x), Reach(y):-Reach(x),E(x,y). The Hoare triple, running the two independent programs side by side, claims equality, R = Reach --- that the explicit worklist BFS and full naive evaluation compute the same reachable set.

\paragraph{First-order reading of PRE/POST, translated mechanically from the relational-algebra assertion in Input.lean.} The relational-algebra assertion in \texttt{Input.lean} remains the authority.

\textit{PRE:} ⊤\par The precondition is simply true.

\textit{POST:} ∀x. (R(x) ↔ Reach(x))\par The postcondition asserts R = Reach, that worklist BFS and naive Datalog evaluation compute the same reachable set.

\paragraph{Whiel encoding (Hoare triple).}

\begin{Verbatim}[fontsize=\scriptsize,breaklines=true,breakanywhere=true,frame=single,framesep=1.5mm]
SCHEMA
  Src/1, R/1, W/1, N/1, Reach/1, Reach_aux/1, E/2

PRE
  true 

CMD
  R := Src;
  W := Src;
  N := Src;
  WHILE (W ≠ ∅) DO
    N := ((π[2] (σ[#0 = #1] (W × E))) ∖ R);
    R := (R ∪ N);
    W := N
  END;
  Reach_aux := ∅[1];
  Reach := (Src ∪ π[2] (σ[#0 = #1] ((Reach_aux × E))));
  WHILE (¬((Reach = Reach_aux))) DO
    Reach_aux := Reach;
    Reach := (Reach ∪ (Src ∪ π[2] (σ[#0 = #1] ((Reach_aux × E)))))
  END

POST
  (R = Reach)
\end{Verbatim}

\Needspace{0.45\textheight}

\section{Example5004: Materialized role-hierarchy closure with a join versus a recursive has-role query}\label{case:Example5004}

\begin{description}[leftmargin=2.6cm,style=nextline,itemsep=1pt]

\item[Category] production pair (two algorithms from production sources or legacy production cases); family: direct-equality pair.

\item[Kind] equivalence.

\item[Contributors] Leo Zhang.

\item[Status] certified valid (Certificate/Valid.lean, written by the certificate emitter); expected valid.

\item[Tags] \hyperref[tag:access-control]{\texttt{access-control}} \hyperref[tag:evaluation-strategies]{\texttt{evaluation-strategies}} \hyperref[tag:prophecy]{\texttt{prophecy}}.

\item[Input] \path{Benchmark/Example5004/Input.lean} (canonical notation).

\item[Summary] Two ways access-control engines answer "does user u hold role r": materialize the role hierarchy's closure and join once (Casbin's role manager), or recurse over the hierarchy from the user's direct roles (OPA/Zanzibar style, legacy Example2035 without its violation query).

\item[Sides] naive compiler: role closure then join (Casbin RoleManager model) \par naive compiler: recursive has-role (legacy Example2035 program)

\item[Sources] {\raggedright \href{https://github.com/casbin/casbin/blob/master/rbac/default-role-manager/role_manager.go}{casbin/casbin rbac/default-role-manager/role\_manager.go (HasLink, hasLinkHelper)} \par \href{https://github.com/casbin/casbin/blob/master/enforcer.go}{casbin/casbin enforcer.go (maxHierarchyLevel = 10)}\par}

\item[Prophecy] kernel-checked certificate with a level-1 prophecy clause (Certificate/Valid.lean, Core.json).

\end{description}

\paragraph{Datalog program(s).}

{\raggedright\textit{prog\_pa} --- fidelity module (kernel-checked: the compiled sides of the command are this program's naive compilation).\par}

\begin{Verbatim}[fontsize=\scriptsize,breaklines=true,breakanywhere=true,frame=single,framesep=1.5mm]
Inh(x1, y1) :- RoleInh(x1, y1).
Inh(x1, z1) :- Inh(x1, y1), RoleInh(y1, z1).
\end{Verbatim}

{\raggedright\textit{prog\_pb} --- fidelity module (kernel-checked: the compiled sides of the command are this program's naive compilation).\par}

\begin{Verbatim}[fontsize=\scriptsize,breaklines=true,breakanywhere=true,frame=single,framesep=1.5mm]
HasM(x1, y1) :- UserRole(x1, y1).
HasM(x1, z1) :- UserRole(x1, y1), Inh(y1, z1).
\end{Verbatim}

{\raggedright\textit{prog\_pc} --- fidelity module (kernel-checked: the compiled sides of the command are this program's naive compilation).\par}

\begin{Verbatim}[fontsize=\scriptsize,breaklines=true,breakanywhere=true,frame=single,framesep=1.5mm]
Has(x1, y1) :- UserRole(x1, y1).
Has(x1, z1) :- Has(x1, y1), RoleInh(y1, z1).
\end{Verbatim}

\paragraph{Reading.} Both sides answer "does user u hold role r" via role-inheritance closure over UserRole and RoleInh. Side 1 models a materialize-then-join strategy (as in Casbin's default role manager): it first computes the full, unbounded role-inheritance closure Inh by recursion over RoleInh, then joins it once against UserRole to get HasM --- with the join program reading the closure program's output, making the sequence dependent. Side 2 models an OPA/Zanzibar-style engine that instead recurses directly from a user's own roles without ever materializing the hierarchy: Has(u,r):-UserRole(u,r), Has(u,t):-Has(u,r),RoleInh(r,t). The Hoare triple claims the two access-control strategies agree, HasM = Has --- both compute UserRole ∘ RoleInh*, the standard role-hierarchy-closure correctness statement, even though the two loops converge at different rates since the join side must wait for the closure computation to finish first.

\paragraph{First-order reading of PRE/POST, translated mechanically from the relational-algebra assertion in Input.lean.} The relational-algebra assertion in \texttt{Input.lean} remains the authority.

\textit{PRE:} ⊤\par The precondition is simply true.

\textit{POST:} ∀x y. (HasM(x,y) ↔ Has(x,y))\par The postcondition asserts HasM = Has, that materialize-then-join role resolution and direct recursive role resolution compute the same has-role relation.

\paragraph{Whiel encoding (Hoare triple).}

\begin{Verbatim}[fontsize=\scriptsize,breaklines=true,breakanywhere=true,frame=single,framesep=1.5mm]
SCHEMA
  UserRole/2, RoleInh/2, Inh/2, Inh_aux/2, HasM/2, HasM_aux/2, Has/2, Has_aux/2

PRE
  true 

CMD
  Inh_aux := ∅[2];
  Inh := (RoleInh ∪ π[0,3] (σ[#1 = #2] ((Inh_aux × RoleInh))));
  WHILE (¬((Inh = Inh_aux))) DO
    Inh_aux := Inh;
    Inh := (Inh ∪ (RoleInh ∪ π[0,3] (σ[#1 = #2] ((Inh_aux × RoleInh)))))
  END;
  HasM_aux := ∅[2];
  HasM := (UserRole ∪ π[0,3] (σ[#1 = #2] ((UserRole × Inh))));
  WHILE (¬((HasM = HasM_aux))) DO
    HasM_aux := HasM;
    HasM := (HasM ∪ (UserRole ∪ π[0,3] (σ[#1 = #2] ((UserRole × Inh)))))
  END;
  Has_aux := ∅[2];
  Has := (UserRole ∪ π[0,3] (σ[#1 = #2] ((Has_aux × RoleInh))));
  WHILE (¬((Has = Has_aux))) DO
    Has_aux := Has;
    Has := (Has ∪ (UserRole ∪ π[0,3] (σ[#1 = #2] ((Has_aux × RoleInh)))))
  END

POST
  (HasM = Has)
\end{Verbatim}

\Needspace{0.45\textheight}

\section{Example5005: Incremental maintenance of a reachability view versus recomputation from scratch}\label{case:Example5005}

\begin{description}[leftmargin=2.6cm,style=nextline,itemsep=1pt]

\item[Category] production pair (two algorithms from production sources or legacy production cases); family: direct-equality pair.

\item[Kind] equivalence.

\item[Contributors] Leo Zhang.

\item[Status] certified valid (Certificate/Valid.lean, written by the certificate emitter); expected valid.

\item[Tags] \hyperref[tag:evaluation-strategies]{\texttt{evaluation-strategies}} \hyperref[tag:prophecy]{\texttt{prophecy}} \hyperref[tag:view-maintenance]{\texttt{view-maintenance}}.

\item[Input] \path{Benchmark/Example5005/Input.lean} (canonical notation).

\item[Summary] Incremental view maintenance against recomputation, the direct form of legacy Example2011.

\item[Sides] naive compiler: old view \par legacy Example2025 (incremental maintainer) \par naive compiler: recomputation on EBase ∪ EDelta

\item[Sources] {\raggedright \href{https://github.com/TimelyDataflow/differential-dataflow}{TimelyDataflow/differential-dataflow} \par \href{https://materialize.com/docs/get-started/arrangements/}{Materialize arrangements / incremental dataflows (documentation)}\par}

\item[Prophecy] kernel-checked certificate with a level-1 prophecy clause (Certificate/Valid.lean, Core.json).

\end{description}

\paragraph{Datalog program(s).}

{\raggedright\textit{prog\_pa} --- fidelity module (kernel-checked: the compiled sides of the command are this program's naive compilation).\par}

\begin{Verbatim}[fontsize=\scriptsize,breaklines=true,breakanywhere=true,frame=single,framesep=1.5mm]
ReachOld(x1) :- Roots(x1).
ReachOld(y1) :- ReachOld(x1), EBase(x1, y1).
\end{Verbatim}

{\raggedright\textit{prog\_pc} --- fidelity module (kernel-checked: the compiled sides of the command are this program's naive compilation).\par}

\begin{Verbatim}[fontsize=\scriptsize,breaklines=true,breakanywhere=true,frame=single,framesep=1.5mm]
ReachNew(x1) :- Roots(x1).
ReachNew(y1) :- ReachNew(x1), EBase(x1, y1).
ReachNew(y1) :- ReachNew(x1), EDelta(x1, y1).
\end{Verbatim}

\paragraph{Reading.} All three programs concern reachability from Roots over an edge relation that is updated by an insertion EDelta. Program 1 (compiled) computes the old view ReachOld using only the base edges EBase. Program 2 is an incremental-view-maintenance algorithm, in the style used by systems like Differential Dataflow or Materialize: it restarts from ReachOld ∪ Roots and iterates over the combined edges EBase ∪ EDelta, an insert-only re-evaluation seeded from the previously materialized view rather than computed from scratch; it reads program 1's output, so the sequence is dependent. Program 3 (compiled) recomputes ReachNew from scratch directly over the updated graph EBase ∪ EDelta. The Hoare triple claims the incrementally maintained result equals full recomputation, Reach = ReachNew --- the standard correctness statement for incremental view maintenance under an insertion-only update.

\paragraph{First-order reading of PRE/POST, translated mechanically from the relational-algebra assertion in Input.lean.} The relational-algebra assertion in \texttt{Input.lean} remains the authority.

\textit{PRE:} ⊤\par The precondition is simply true.

\textit{POST:} ∀x. (Reach(x) ↔ ReachNew(x))\par The postcondition asserts Reach = ReachNew, that incrementally maintaining the reachability view from the old result equals recomputing it from scratch on the updated graph.

\paragraph{Whiel encoding (Hoare triple).}

\begin{Verbatim}[fontsize=\scriptsize,breaklines=true,breakanywhere=true,frame=single,framesep=1.5mm]
SCHEMA
  Roots/1, ReachOld/1, ReachOld_aux/1, Reach/1, Reach_2/1, ReachNew/1, ReachNew_aux/1, EBase/2, EDelta/2

PRE
  true 

CMD
  ReachOld_aux := ∅[1];
  ReachOld := (Roots ∪ π[2] (σ[#0 = #1] ((ReachOld_aux × EBase))));
  WHILE (¬((ReachOld = ReachOld_aux))) DO
    ReachOld_aux := ReachOld;
    ReachOld := (ReachOld ∪ (Roots ∪ π[2] (σ[#0 = #1] ((ReachOld_aux × EBase)))))
  END;
  Reach_2 := ∅[1];
  Reach := (ReachOld ∪ Roots);
  WHILE (Reach ≠ Reach_2) DO
    Reach_2 := Reach;
    Reach := ((Roots ∪ ReachOld) ∪ (π[1] (σ[#0 = #2] ((EBase ∪ EDelta) × Reach_2))))
  END;
  ReachNew_aux := ∅[1];
  ReachNew := (Roots ∪ (π[2] (σ[#0 = #1] ((ReachNew_aux × EBase))) ∪ π[2] (σ[#0 = #1] ((ReachNew_aux × EDelta)))));
  WHILE (¬((ReachNew = ReachNew_aux))) DO
    ReachNew_aux := ReachNew;
    ReachNew := (ReachNew ∪ (Roots ∪ (π[2] (σ[#0 = #1] ((ReachNew_aux × EBase))) ∪ π[2] (σ[#0 = #1] ((ReachNew_aux × EDelta))))))
  END

POST
  (Reach = ReachNew)
\end{Verbatim}

\Needspace{0.45\textheight}

\section{Example5006: Semi-naive closure with an appended delta versus a prepended delta}\label{case:Example5006}

\begin{description}[leftmargin=2.6cm,style=nextline,itemsep=1pt]

\item[Category] production pair (two algorithms from production sources or legacy production cases); family: direct-equality pair.

\item[Kind] equivalence.

\item[Contributors] Leo Zhang.

\item[Status] certified valid (Certificate/Valid.lean, written by the certificate emitter); expected valid.

\item[Tags] \hyperref[tag:evaluation-strategies]{\texttt{evaluation-strategies}} \hyperref[tag:linear-vs-nonlinear]{\texttt{linear-vs-nonlinear}} \hyperref[tag:prophecy]{\texttt{prophecy}} \hyperref[tag:semi-naive]{\texttt{semi-naive}}.

\item[Input] \path{Benchmark/Example5006/Input.lean} (canonical notation).

\item[Summary] Two semi-naive evaluations of transitive closure that differ in where the delta is joined: appended (`Delta ∘ E`, legacy Example1042) or prepended (`E ∘ Delta`, legacy Example1058).

\item[Sides] legacy Example1042 (append delta) \par legacy Example1058 (prepend delta)

\item[Sources] {\raggedright LLM benchmark source records of the retired corpus (GPT and Claude generated; no external source)\par}

\item[Prophecy] kernel-checked certificate with a level-1 prophecy clause (Certificate/Valid.lean, Core.json).

\end{description}

\paragraph{Datalog program(s).}

{\raggedright\textit{Side 1, appended delta (legacy Example1042); side 2, prepended delta (legacy Example1058); written for this report} --- hand-written reading (this report).\par}

\begin{Verbatim}[fontsize=\scriptsize,breaklines=true,breakanywhere=true,frame=single,framesep=1.5mm]
AccA(x, y) :- E(x, y).
AccA(x, z) :- AccA(x, y), E(y, z).

AccP(x, y) :- E(x, y).
AccP(x, z) :- E(x, y), AccP(y, z).
\end{Verbatim}

\paragraph{Reading.} Both sides compute the transitive closure of E via semi-naive delta evaluation, differing only in which side of the join the growing frontier is placed on: side A appends the delta (joining the frontier on the left of E, Delta ∘ E) while side P prepends it (joining E onto the frontier from the left, E ∘ Delta), each round adding the frontier of walks one edge longer than the previous round. The Hoare triple, running the two independent programs side by side, claims the two evaluation orders produce equal results, AccA = AccP --- since both frontiers hold the same set of walks of length n+1 after round n regardless of which side of the join the new edge is added on, the fixpoints coincide.

\paragraph{First-order reading of PRE/POST, translated mechanically from the relational-algebra assertion in Input.lean.} The relational-algebra assertion in \texttt{Input.lean} remains the authority.

\textit{PRE:} ⊤\par The precondition is simply true.

\textit{POST:} ∀x y. (AccA(x,y) ↔ AccP(x,y))\par The postcondition asserts AccA = AccP, that the append-delta and prepend-delta semi-naive evaluations of transitive closure compute the same relation.

\paragraph{Whiel encoding (Hoare triple).}

\begin{Verbatim}[fontsize=\scriptsize,breaklines=true,breakanywhere=true,frame=single,framesep=1.5mm]
SCHEMA
  E/2, AccA/2, DeltaA/2, AccP/2, DeltaP/2

PRE
  true 

CMD
  AccA := E;
  DeltaA := E;
  WHILE (DeltaA ≠ ∅) DO
    DeltaA := ((π[0, 3] (σ[#1 = #2] (DeltaA × E))) ∖ AccA);
    AccA := (AccA ∪ DeltaA)
  END;
  AccP := E;
  DeltaP := E;
  WHILE (DeltaP ≠ ∅) DO
    DeltaP := ((π[0, 3] (σ[#1 = #2] (E × DeltaP))) ∖ AccP);
    AccP := (AccP ∪ DeltaP)
  END

POST
  (AccA = AccP)
\end{Verbatim}

\Needspace{0.45\textheight}

\section{Example5007: Transitive closure by repeated squaring versus linear closure}\label{case:Example5007}

\begin{description}[leftmargin=2.6cm,style=nextline,itemsep=1pt]

\item[Category] production pair (two algorithms from production sources or legacy production cases); family: direct-equality pair.

\item[Kind] equivalence.

\item[Contributors] Leo Zhang.

\item[Status] certified valid (Certificate/Valid.lean, written by the certificate emitter); expected valid.

\item[Tags] \hyperref[tag:evaluation-strategies]{\texttt{evaluation-strategies}} \hyperref[tag:linear-vs-nonlinear]{\texttt{linear-vs-nonlinear}} \hyperref[tag:nonlinear-recursion]{\texttt{nonlinear-recursion}} \hyperref[tag:prophecy]{\texttt{prophecy}}.

\item[Input] \path{Benchmark/Example5007/Input.lean} (canonical notation).

\item[Summary] Closure by repeated squaring (legacy Example1056) against the naive evaluation of the linear closure program.

\item[Sides] legacy Example1056 (closure by squaring) \par naive compiler on the linear closure program

\item[Sources] {\raggedright LLM benchmark source records of the retired corpus (GPT and Claude generated; no external source)\par}

\item[Prophecy] kernel-checked certificate with a level-1 prophecy clause (Certificate/Valid.lean, Core.json).

\end{description}

\paragraph{Datalog program(s).}

{\raggedright\textit{prog\_p} --- fidelity module (kernel-checked: the compiled sides of the command are this program's naive compilation).\par}

\begin{Verbatim}[fontsize=\scriptsize,breaklines=true,breakanywhere=true,frame=single,framesep=1.5mm]
T(x1, y1) :- E(x1, y1).
T(x1, y1) :- E(x1, z1), T(z1, y1).
\end{Verbatim}

\paragraph{Reading.} Both sides compute the transitive closure of E. Side 1 computes it by repeated squaring: starting from E, each round unions in TSq ∘ TSq (the self-composition of the accumulated relation), doubling the walk length covered every iteration --- a logarithmic-depth closure algorithm. Side 2 is standard linear naive evaluation, T(x,y):-E(x,y), T(x,y):-E(x,z),T(z,y), which extends the covered walk length by one edge per round. The Hoare triple claims the two algorithms converge to the same relation, TSq = T --- the same least fixpoint reached at very different rates, walks up to length 2\textasciicircum{}(n+1) after round n for the squaring side versus length n+2 for the linear side.

\paragraph{First-order reading of PRE/POST, translated mechanically from the relational-algebra assertion in Input.lean.} The relational-algebra assertion in \texttt{Input.lean} remains the authority.

\textit{PRE:} ⊤\par The precondition is simply true.

\textit{POST:} ∀x y. (TSq(x,y) ↔ T(x,y))\par The postcondition asserts TSq = T, that closure by repeated squaring and linear naive closure compute the same transitive-closure relation.

\paragraph{Whiel encoding (Hoare triple).}

\begin{Verbatim}[fontsize=\scriptsize,breaklines=true,breakanywhere=true,frame=single,framesep=1.5mm]
SCHEMA
  E/2, TSq/2, TSq_2/2, T/2, T_aux/2

PRE
  true 

CMD
  TSq_2 := ∅[2];
  TSq := E;
  WHILE (TSq ≠ TSq_2) DO
    TSq_2 := TSq;
    TSq := (TSq ∪ (π[0, 3] (σ[#1 = #2] (TSq_2 × TSq_2))))
  END;
  T_aux := ∅[2];
  T := (E ∪ π[0,3] (σ[#1 = #2] ((E × T_aux))));
  WHILE (¬((T = T_aux))) DO
    T_aux := T;
    T := (T ∪ (E ∪ π[0,3] (σ[#1 = #2] ((E × T_aux)))))
  END

POST
  (TSq = T)
\end{Verbatim}

\Needspace{0.45\textheight}

\section{Example5008: Soufflé tc benchmark: naive versus semi-naive evaluation}\label{case:Example5008}

\begin{description}[leftmargin=2.6cm,style=nextline,itemsep=1pt]

\item[Category] Soufflé benchmark program (naive versus semi-naive translation); family: Soufflé naive vs semi-naive.

\item[Kind] equivalence.

\item[Contributors] Leo Zhang.

\item[Status] certified valid (certificate checked outside the library build) (Certificate/Valid.lean, written by the certificate emitter); expected valid. Placement note: a proof module imported by \path{Certificate/Valid.lean} does not build under the library build's 4194304 KB limit; it builds under 12582912 KB (peak 1686 MB)

\item[Tags] \hyperref[tag:database-engine]{\texttt{database-engine}} \hyperref[tag:evaluation-strategies]{\texttt{evaluation-strategies}} \hyperref[tag:nonlinear-recursion]{\texttt{nonlinear-recursion}} \hyperref[tag:prophecy]{\texttt{prophecy}} \hyperref[tag:semi-naive]{\texttt{semi-naive}}.

\item[Input] \path{Benchmark/Example5008/Input.lean} (canonical notation).

\item[Summary] Soufflé's `tc` benchmark program evaluated naively and semi-naively, tested for equality of every output.

\item[Sides] naive translation of tc.dl (Soufflé benchmark suite) \par semi-naive translation of tc.dl (Soufflé benchmark suite)

\item[Sources] {\raggedright \href{https://github.com/souffle-lang/benchmarks}{souffle-lang/benchmarks (the Soufflé benchmark programs)} \par \href{https://souffle-lang.github.io/translate}{Soufflé translation to C++ / semi-naive loop (documentation)} \par program source file tc.dl of the Soufflé benchmark suite (naive translation)\par}

\item[Prophecy] accepted invariant with a prophecy clause above level zero (Core.json); no certificate built from it yet.

\end{description}

\paragraph{Datalog program(s).}

{\raggedright\textit{normalized Soufflé source (tc.dl of the Soufflé benchmark suite (naive translation))} --- Soufflé benchmark program, normalized (both sides are script translations of it).\par}

\begin{Verbatim}[fontsize=\scriptsize,breaklines=true,breakanywhere=true,frame=single,framesep=1.5mm]
.decl tcl(v0:symbol, v1:symbol) output
.decl tcr(v0:symbol, v1:symbol) output
.decl base(v0:symbol, v1:symbol) input
.decl tc(v0:symbol, v1:symbol) output

tcl(X, Y) :- base(X, Y).
tcl(X, Y) :- tcl(X, Z), base(Z, Y).
tcr(X, Y) :- base(X, Y).
tcr(X, Y) :- base(X, Z), tcr(Z, Y).
tc(X, Y) :- base(X, Y).
tc(X, Y) :- tc(X, Z), tc(Z, Y).
\end{Verbatim}

\paragraph{Reading.} Soufflé's tc.dl benchmark computes graph transitive closure via three intertwined recursive relations over an edge relation base: tc (the ordinary both-sides closure), tcl (built by left-extending tcl with base), and tcr (built by right-extending base with tcr) --- a standard reachability computation. The Hoare triple compares Soufflé's naive fixpoint evaluation of these rules against classical semi-naive (delta-based) evaluation of the same rules, claiming the two strategies compute identical output relations, the standard naive/semi-naive evaluation equivalence for Datalog.

\paragraph{First-order reading of PRE/POST, translated mechanically from the relational-algebra assertion in Input.lean.} The relational-algebra assertion in \texttt{Input.lean} remains the authority.

\textit{PRE:} ⊤\par True; no precondition constrains the inputs.

\textit{POST:} (∀x y. (tc(x,y) ↔ tcS(x,y))) ∧ (∀z u. (tcl(z,u) ↔ tclS(z,u))) ∧ ∀v w. (tcr(v,w) ↔ tcrS(v,w))\par The naively evaluated tc, tcl, and tcr each equal the corresponding semi-naively evaluated relation (tcS, tclS, tcrS).

\paragraph{Whiel encoding (Hoare triple).}

\begin{Verbatim}[fontsize=\tiny,breaklines=true,breakanywhere=true,frame=single,framesep=1.5mm]
SCHEMA
  base/2, tc/2, tcl/2, tcr/2, tcNew/2, tcOld/2, tclNew/2, tclOld/2, tcrNew/2, tcrOld/2, tcS/2, tclS/2, tcrS/2, deltaTcS/2, deltaTcOldS/2, deltaTclS/2, deltaTclOldS/2, deltaTcrS/2, deltaTcrOldS/2

PRE
  true 

CMD
  tcNew := (π[0, 1] base);
  tcOld := ∅[2];
  tclNew := (π[0, 1] base);
  tclOld := ∅[2];
  tcrNew := (π[0, 1] base);
  tcrOld := ∅[2];
  WHILE (((tcNew ≠ tcOld) ∨ (tclNew ≠ tclOld)) ∨ (tcrNew ≠ tcrOld)) DO
  tcOld := tcNew;
  tclOld := tclNew;
  tcrOld := tcrNew;
  tcNew := (tcNew ∪ (π[0, 3] (σ[#2 = #1] (tcOld × tcOld))));
  tclNew := (tclNew ∪ (π[0, 3] (σ[#2 = #1] (tclOld × base))));
  tcrNew := (tcrNew ∪ (π[0, 3] (σ[#2 = #1] (base × tcrOld))))
  END;
  tc := tcNew;
  tcl := tclNew;
  tcr := tcrNew;
  tcS := ∅[2];
  deltaTcS := (π[0, 1] base);
  deltaTcOldS := ∅[2];
  tclS := ∅[2];
  deltaTclS := (π[0, 1] base);
  deltaTclOldS := ∅[2];
  tcrS := ∅[2];
  deltaTcrS := (π[0, 1] base);
  deltaTcrOldS := ∅[2];
  WHILE (((deltaTcS ≠ ∅[2]) ∨ (deltaTclS ≠ ∅[2])) ∨ (deltaTcrS ≠ ∅[2])) DO
  deltaTcOldS := deltaTcS;
  deltaTclOldS := deltaTclS;
  deltaTcrOldS := deltaTcrS;
  tcS := (tcS ∪ deltaTcOldS);
  tclS := (tclS ∪ deltaTclOldS);
  tcrS := (tcrS ∪ deltaTcrOldS);
  deltaTcS := ((((π[0, 3] (σ[#2 = #1] (deltaTcOldS × tcS))) ∪ (π[0, 3] (σ[#2 = #1] (tcS × deltaTcOldS)))) ∪ (π[0, 3] (σ[#2 = #1] (deltaTcOldS × deltaTcOldS)))) ∖ tcS);
  deltaTclS := ((π[0, 3] (σ[#2 = #1] (deltaTclOldS × base))) ∖ tclS);
  deltaTcrS := ((π[0, 3] (σ[#2 = #1] (base × deltaTcrOldS))) ∖ tcrS)
  END

POST
  ((tc = tcS) ∧ (tcl = tclS) ∧ (tcr = tcrS))
\end{Verbatim}

\Needspace{0.45\textheight}

\section{Example5009: Soufflé family benchmark: naive versus semi-naive evaluation}\label{case:Example5009}

\begin{description}[leftmargin=2.6cm,style=nextline,itemsep=1pt]

\item[Category] Soufflé benchmark program (naive versus semi-naive translation); family: Soufflé naive vs semi-naive.

\item[Kind] equivalence.

\item[Contributors] Leo Zhang.

\item[Status] certified valid (certificate checked outside the library build) (Certificate/Valid.lean, written by the certificate emitter); expected valid. Placement note: a proof module imported by \path{Certificate/Valid.lean} does not build under the library build's 4194304 KB limit; it builds under 12582912 KB (peak 1668 MB)

\item[Tags] \hyperref[tag:database-engine]{\texttt{database-engine}} \hyperref[tag:evaluation-strategies]{\texttt{evaluation-strategies}} \hyperref[tag:negation]{\texttt{negation}} \hyperref[tag:prophecy]{\texttt{prophecy}} \hyperref[tag:semi-naive]{\texttt{semi-naive}}.

\item[Input] \path{Benchmark/Example5009/Input.lean} (canonical notation).

\item[Summary] Soufflé's `family` benchmark program evaluated naively and semi-naively, tested for equality of every output.

\item[Sides] naive translation of family.dl (Soufflé benchmark suite) \par semi-naive translation of family.dl (Soufflé benchmark suite)

\item[Sources] {\raggedright \href{https://github.com/souffle-lang/benchmarks}{souffle-lang/benchmarks (the Soufflé benchmark programs)} \par \href{https://souffle-lang.github.io/translate}{Soufflé translation to C++ / semi-naive loop (documentation)} \par program source file family.dl of the Soufflé benchmark suite (naive translation)\par}

\item[Prophecy] accepted invariant with a prophecy clause above level zero (Core.json); no certificate built from it yet.

\end{description}

\paragraph{Datalog program(s).}

{\raggedright\textit{normalized Soufflé source (family.dl of the Soufflé benchmark suite (naive translation))} --- Soufflé benchmark program, normalized (both sides are script translations of it).\par}

\begin{Verbatim}[fontsize=\scriptsize,breaklines=true,breakanywhere=true,frame=single,framesep=1.5mm]
.decl father(v0:symbol, v1:symbol) input
.decl mother(v0:symbol, v1:symbol) input
.decl parent(v0:symbol, v1:symbol) output
.decl ancestor(v0:symbol, v1:symbol) output
.decl grandmother(v0:symbol, v1:symbol) output
.decl sibling(v0:symbol, v1:symbol) output
.decl cousin(v0:symbol, v1:symbol) output
.decl relative(v0:symbol, v1:symbol) output

parent(X, Y) :- father(X, Y).
parent(X, Y) :- mother(X, Y).
ancestor(X, Y) :- parent(X, Y).
ancestor(X, Y) :- parent(X, Z), ancestor(Z, Y).
grandmother(X, Y) :- mother(X, Z), ancestor(Z, Y).
sibling(X, Y) :- parent(Z, X), parent(Z, Y), X != Y.
cousin(X, Y) :- ancestor(Z, X), ancestor(Z, Y), X != Y, !sibling(X, Y), !parent(X, Y).
relative(X, Y) :- sibling(X, Z), parent(Z, Y).
\end{Verbatim}

\paragraph{Reading.} Soufflé's family.dl benchmark derives standard kinship relations --- parent, the transitive-closure relation ancestor, sibling, cousin, grandmother, and relative --- from father/mother facts, a typical family-relationships Datalog program with one recursive rule (ancestor) and several non-recursive relations built on top of it. The triple compares naive fixpoint evaluation of the whole program against semi-naive evaluation, claiming every one of the six output relations is computed identically by both strategies.

\paragraph{First-order reading of PRE/POST, translated mechanically from the relational-algebra assertion in Input.lean.} The relational-algebra assertion in \texttt{Input.lean} remains the authority.

\textit{PRE:} ⊤\par True; no precondition constrains the inputs.

\textit{POST:} (∀x y. (ancestor(x,y) ↔ ancestorS(x,y))) ∧ (∀z u. (cousin(z,u) ↔ cousinS(z,u))) ∧ (∀v w. (grandmother(v,w) ↔ grandmotherS(v,w))) ∧ (∀x1 x2. (parent(x1,x2) ↔ parentS(x1,x2))) ∧ (∀x3 x4. (relative(x3,x4) ↔ relativeS(x3,x4))) ∧ ∀x5 x6. (sibling(x5,x6) ↔ siblingS(x5,x6))\par Each of the six output relations (ancestor, cousin, grandmother, parent, relative, sibling) computed by naive evaluation equals its semi-naive counterpart.

\paragraph{Whiel encoding (Hoare triple).}

\begin{Verbatim}[fontsize=\scriptsize,breaklines=true,breakanywhere=true,frame=single,framesep=1.5mm]
SCHEMA
  father/2, mother/2, ancestor/2, cousin/2, grandmother/2, parent/2, relative/2, sibling/2, ancestorNew/2, ancestorOld/2, ancestorS/2, cousinS/2, grandmotherS/2, parentS/2, relativeS/2, siblingS/2, deltaAncestorS/2

PRE
  true 

CMD
  parent := ((π[0, 1] father) ∪ (π[0, 1] mother));
  sibling := (π[1, 3] (σ[(#2 = #0 ∧ ¬(#1 = #3))] (parent × parent)));
  relative := (π[0, 3] (σ[#2 = #1] (sibling × parent)));
  ancestorNew := (π[0, 1] parent);
  ancestorOld := ∅[2];
  WHILE (ancestorNew ≠ ancestorOld) DO
  ancestorOld := ancestorNew;
  ancestorNew := (ancestorNew ∪ (π[0, 3] (σ[#2 = #1] (parent × ancestorOld))))
  END;
  ancestor := ancestorNew;
  cousin := (π[1, 3] (((σ[(#2 = #0 ∧ ¬(#1 = #3))] (ancestor × ancestor)) ∖ (π[0, 1, 2, 3] (σ[(#4 = #1 ∧ #5 = #3)] ((σ[(#2 = #0 ∧ ¬(#1 = #3))] (ancestor × ancestor)) × sibling)))) ∖ (π[0, 1, 2, 3] (σ[(#4 = #1 ∧ #5 = #3)] (((σ[(#2 = #0 ∧ ¬(#1 = #3))] (ancestor × ancestor)) ∖ (π[0, 1, 2, 3] (σ[(#4 = #1 ∧ #5 = #3)] ((σ[(#2 = #0 ∧ ¬(#1 = #3))] (ancestor × ancestor)) × sibling)))) × parent)))));
  grandmother := (π[0, 3] (σ[#2 = #1] (mother × ancestor)));
  parentS := ((π[0, 1] father) ∪ (π[0, 1] mother));
  siblingS := (π[1, 3] (σ[(#2 = #0 ∧ ¬(#1 = #3))] (parentS × parentS)));
  relativeS := (π[0, 3] (σ[#2 = #1] (siblingS × parentS)));
  ancestorS := ∅[2];
  deltaAncestorS := (π[0, 1] parentS);
  WHILE (deltaAncestorS ≠ ∅[2]) DO
  ancestorS := (ancestorS ∪ deltaAncestorS);
  deltaAncestorS := ((π[0, 3] (σ[#2 = #1] (parentS × deltaAncestorS))) ∖ ancestorS)
  END;
  cousinS := (π[1, 3] (((σ[(#2 = #0 ∧ ¬(#1 = #3))] (ancestorS × ancestorS)) ∖ (π[0, 1, 2, 3] (σ[(#4 = #1 ∧ #5 = #3)] ((σ[(#2 = #0 ∧ ¬(#1 = #3))] (ancestorS × ancestorS)) × siblingS)))) ∖ (π[0, 1, 2, 3] (σ[(#4 = #1 ∧ #5 = #3)] (((σ[(#2 = #0 ∧ ¬(#1 = #3))] (ancestorS × ancestorS)) ∖ (π[0, 1, 2, 3] (σ[(#4 = #1 ∧ #5 = #3)] ((σ[(#2 = #0 ∧ ¬(#1 = #3))] (ancestorS × ancestorS)) × siblingS)))) × parentS)))));
  grandmotherS := (π[0, 3] (σ[#2 = #1] (mother × ancestorS)))

POST
  ((ancestor = ancestorS) ∧ (cousin = cousinS) ∧ (grandmother = grandmotherS) ∧ (parent = parentS) ∧ (relative = relativeS) ∧ (sibling = siblingS))
\end{Verbatim}

\Needspace{0.45\textheight}

\section{Example5010: Soufflé orbits benchmark: naive versus semi-naive evaluation}\label{case:Example5010}

\begin{description}[leftmargin=2.6cm,style=nextline,itemsep=1pt]

\item[Category] Soufflé benchmark program (naive versus semi-naive translation); family: Soufflé naive vs semi-naive.

\item[Kind] equivalence.

\item[Contributors] Leo Zhang.

\item[Status] certified valid (Certificate/Valid.lean, written by the certificate emitter); expected valid.

\item[Tags] \hyperref[tag:database-engine]{\texttt{database-engine}} \hyperref[tag:evaluation-strategies]{\texttt{evaluation-strategies}} \hyperref[tag:negation]{\texttt{negation}} \hyperref[tag:nonlinear-recursion]{\texttt{nonlinear-recursion}} \hyperref[tag:prophecy]{\texttt{prophecy}} \hyperref[tag:semi-naive]{\texttt{semi-naive}}.

\item[Input] \path{Benchmark/Example5010/Input.lean} (canonical notation).

\item[Summary] Soufflé's `orbits` benchmark program evaluated naively and semi-naively, tested for equality of every output.

\item[Sides] naive translation of orbits.dl (Soufflé benchmark suite) \par semi-naive translation of orbits.dl (Soufflé benchmark suite)

\item[Sources] {\raggedright \href{https://github.com/souffle-lang/benchmarks}{souffle-lang/benchmarks (the Soufflé benchmark programs)} \par \href{https://souffle-lang.github.io/translate}{Soufflé translation to C++ / semi-naive loop (documentation)} \par program source file orbits.dl of the Soufflé benchmark suite (naive translation)\par}

\item[Prophecy] kernel-checked certificate with a level-1 prophecy clause (Certificate/Valid.lean, Core.json).

\end{description}

\paragraph{Datalog program(s).}

{\raggedright\textit{normalized Soufflé source (orbits.dl of the Soufflé benchmark suite (naive translation))} --- Soufflé benchmark program, normalized (both sides are script translations of it).\par}

\begin{Verbatim}[fontsize=\scriptsize,breaklines=true,breakanywhere=true,frame=single,framesep=1.5mm]
.decl orbits_seed(v0:symbol, v1:symbol) input
.decl orbits(v0:symbol, v1:symbol) output
.decl intermediate(v0:symbol, v1:symbol) output
.decl satellite(v0:symbol, v1:symbol) output

orbits(V0, V1) :- orbits_seed(V0, V1).
orbits(X, Y) :- orbits(X, Z), orbits(Z, Y).
satellite(X, Y) :- orbits(X, Y), !intermediate(X, Y).
intermediate(X, Y) :- orbits(X, Z), orbits(Z, Y).
\end{Verbatim}

\paragraph{Reading.} Soufflé's orbits.dl benchmark computes the transitive closure of a seed relation, orbits (the orbits of the seed relation under composition), together with intermediate (a two-step composition of orbits with itself) and satellite (orbit pairs not reachable by such an intermediate two-step path) built on top of it --- a graph-composition closure computation. The triple compares naive versus semi-naive evaluation of the whole program, claiming all three output relations agree between the two strategies.

\paragraph{First-order reading of PRE/POST, translated mechanically from the relational-algebra assertion in Input.lean.} The relational-algebra assertion in \texttt{Input.lean} remains the authority.

\textit{PRE:} ⊤\par True; no precondition constrains the inputs.

\textit{POST:} (∀x y. (intermediate(x,y) ↔ intermediateS(x,y))) ∧ (∀z u. (orbits(z,u) ↔ orbitsS(z,u))) ∧ ∀v w. (satellite(v,w) ↔ satelliteS(v,w))\par The naively computed intermediate, orbits, and satellite relations each equal their semi-naive counterparts.

\paragraph{Whiel encoding (Hoare triple).}

\begin{Verbatim}[fontsize=\scriptsize,breaklines=true,breakanywhere=true,frame=single,framesep=1.5mm]
SCHEMA
  orbitsSeed/2, intermediate/2, orbits/2, satellite/2, orbitsNew/2, orbitsOld/2, intermediateS/2, orbitsS/2, satelliteS/2, deltaOrbitsS/2

PRE
  true 

CMD
  orbitsNew := (π[0, 1] orbitsSeed);
  orbitsOld := ∅[2];
  WHILE (orbitsNew ≠ orbitsOld) DO
  orbitsOld := orbitsNew;
  orbitsNew := (orbitsNew ∪ (π[0, 3] (σ[#2 = #1] (orbitsOld × orbitsOld))))
  END;
  orbits := orbitsNew;
  intermediate := (π[0, 3] (σ[#2 = #1] (orbits × orbits)));
  satellite := (π[0, 1] (orbits ∖ (π[0, 1] (σ[(#2 = #0 ∧ #3 = #1)] (orbits × intermediate)))));
  orbitsS := ∅[2];
  deltaOrbitsS := (π[0, 1] orbitsSeed);
  WHILE (deltaOrbitsS ≠ ∅[2]) DO
  orbitsS := (orbitsS ∪ deltaOrbitsS);
  deltaOrbitsS := ((((π[0, 3] (σ[#2 = #1] (deltaOrbitsS × orbitsS))) ∪ (π[0, 3] (σ[#2 = #1] (orbitsS × deltaOrbitsS)))) ∪ (π[0, 3] (σ[#2 = #1] (deltaOrbitsS × deltaOrbitsS)))) ∖ orbitsS)
  END;
  intermediateS := (π[0, 3] (σ[#2 = #1] (orbitsS × orbitsS)));
  satelliteS := (π[0, 1] (orbitsS ∖ (π[0, 1] (σ[(#2 = #0 ∧ #3 = #1)] (orbitsS × intermediateS)))))

POST
  ((intermediate = intermediateS) ∧ (orbits = orbitsS) ∧ (satellite = satelliteS))
\end{Verbatim}

\Needspace{0.45\textheight}

\section{Example5011: Soufflé topological\_ordering benchmark: naive versus semi-naive evaluation}\label{case:Example5011}

\begin{description}[leftmargin=2.6cm,style=nextline,itemsep=1pt]

\item[Category] Soufflé benchmark program (naive versus semi-naive translation); family: Soufflé naive vs semi-naive.

\item[Kind] equivalence.

\item[Contributors] Leo Zhang.

\item[Status] certified invalid (Certificate/Invalid.lean, written by the certificate emitter); expected invalid.

\item[Tags] \hyperref[tag:database-engine]{\texttt{database-engine}} \hyperref[tag:evaluation-strategies]{\texttt{evaluation-strategies}} \hyperref[tag:negation]{\texttt{negation}} \hyperref[tag:nonlinear-recursion]{\texttt{nonlinear-recursion}} \hyperref[tag:semi-naive]{\texttt{semi-naive}}.

\item[Input] \path{Benchmark/Example5011/Input.lean} (canonical notation).

\item[Summary] Soufflé's `topological\_ordering` benchmark program evaluated naively and semi-naively, tested for equality of every output.

\item[Sides] naive translation of topological\_ordering.dl (Soufflé benchmark suite) \par semi-naive translation of topological\_ordering.dl (Soufflé benchmark suite)

\item[Sources] {\raggedright \href{https://github.com/souffle-lang/benchmarks}{souffle-lang/benchmarks (the Soufflé benchmark programs)} \par \href{https://souffle-lang.github.io/translate}{Soufflé translation to C++ / semi-naive loop (documentation)} \par program source file topological\_ordering.dl of the Soufflé benchmark suite (naive translation)\par}

\end{description}

\paragraph{Datalog program(s).}

{\raggedright\textit{normalized Soufflé source (topological\_ordering.dl of the Soufflé benchmark suite (naive translation))} --- Soufflé benchmark program, normalized (both sides are script translations of it).\par}

\begin{Verbatim}[fontsize=\scriptsize,breaklines=true,breakanywhere=true,frame=single,framesep=1.5mm]
.decl edge(v0:symbol, v1:symbol) input
.decl RankZero(v0:symbol) input
.decl RankOne(v0:symbol) input
.decl Succ(v0:symbol, v1:symbol) input
.decl vertex(v0:symbol) output
.decl is_before(v0:symbol, v1:symbol) output
.decl is_after(v0:symbol, v1:symbol) output
.decl isolated(v0:symbol) output
.decl source(v0:symbol) output
.decl indices(v0:symbol, v1:symbol) output
.decl has_successor_rank(v0:symbol, v1:symbol) output
.decl index(v0:symbol, v1:symbol) output

vertex(v) :- edge(v, __anon_1).
vertex(v) :- edge(__anon_2, v).
is_before(x, y) :- edge(x, y).
is_before(x, y) :- is_before(x, z), is_before(z, y).
is_after(x, y) :- edge(y, x).
is_after(x, y) :- is_after(z, x), is_after(y, z).
isolated(x) :- vertex(x), !edge(__anon_3, x), !edge(x, __anon_4).
source(x) :- vertex(x), !edge(__anon_5, x), edge(x, __anon_6).
indices(x, i) :- isolated(x), RankZero(i).
indices(x, i) :- source(x), RankOne(i).
indices(x, i1) :- is_after(x, y), indices(y, i), Succ(i, i1).
has_successor_rank(x, i) :- indices(x, i), Succ(i, i1), indices(x, i1).
index(x, i) :- indices(x, i), !has_successor_rank(x, i).
\end{Verbatim}

\paragraph{Reading.} Soufflé's topological\_ordering.dl benchmark computes a topological layering of a directed graph. is\_before and is\_after are the transitive closure of the edge relation and of its reverse; isolated vertices receive rank zero and source vertices rank one, and a vertex that comes after a vertex of rank i also receives the next rank (ranks are symbols ordered by the input relation Succ). index keeps, for each vertex, the rank that has no successor rank assigned to the same vertex, that is, its maximal rank: on an acyclic graph, one more than the length of the longest path reaching it from a source. The triple compares naive and semi-naive evaluation of this whole stratified program and claims that all eight output relations agree.

\paragraph{First-order reading of PRE/POST, translated mechanically from the relational-algebra assertion in Input.lean.} The relational-algebra assertion in \texttt{Input.lean} remains the authority.

\textit{PRE:} ⊤\par True; no precondition constrains the inputs.

\textit{POST:} (∀x y. (hasSuccessorRank(x,y) ↔ hasSuccessorRankS(x,y))) ∧ (∀z u. (index(z,u) ↔ indexS(z,u))) ∧ (∀v w. (indices(v,w) ↔ indicesS(v,w))) ∧ (∀x1 x2. (isAfter(x1,x2) ↔ isAfterS(x1,x2))) ∧ (∀x3 x4. (isBefore(x3,x4) ↔ isBeforeS(x3,x4))) ∧ (∀x5. (isolated(x5) ↔ isolatedS(x5))) ∧ (∀x6. (source(x6) ↔ sourceS(x6))) ∧ ∀x7. (vertex(x7) ↔ vertexS(x7))\par Each of the eight output relations (hasSuccessorRank, index, indices, isAfter, isBefore, isolated, source, vertex) computed naively equals its semi-naive counterpart.

\paragraph{Whiel encoding (Hoare triple).}

\begin{Verbatim}[fontsize=\tiny,breaklines=true,breakanywhere=true,frame=single,framesep=1.5mm]
SCHEMA
  RankOne/1, RankZero/1, isolated/1, source/1, vertex/1, isolatedS/1, sourceS/1, vertexS/1, Succ/2, edge/2, hasSuccessorRank/2, index/2, indices/2, isAfter/2, isBefore/2, indicesNew/2, indicesOld/2, isAfterNew/2, isAfterOld/2, isBeforeNew/2, isBeforeOld/2, hasSuccessorRankS/2, indexS/2, indicesS/2, isAfterS/2, isBeforeS/2, deltaIndicesS/2, deltaIndicesOldS/2, deltaIsAfterS/2, deltaIsAfterOldS/2, deltaIsBeforeS/2, deltaIsBeforeOldS/2

PRE
  true 

CMD
  vertex := ((π[0] edge) ∪ (π[1] edge));
  isolated := (π[0] ((vertex ∖ (π[0] (σ[#2 = #0] (vertex × edge)))) ∖ (π[0] (σ[#1 = #0] ((vertex ∖ (π[0] (σ[#2 = #0] (vertex × edge)))) × edge)))));
  source := (π[0] ((σ[#1 = #0] (vertex × edge)) ∖ (π[0, 1, 2] (σ[#4 = #0] ((σ[#1 = #0] (vertex × edge)) × edge)))));
  indicesNew := ((π[0, 1] (isolated × RankZero)) ∪ (π[0, 1] (source × RankOne)));
  indicesOld := ∅[2];
  isAfterNew := (π[1, 0] edge);
  isAfterOld := ∅[2];
  isBeforeNew := (π[0, 1] edge);
  isBeforeOld := ∅[2];
  WHILE (((indicesNew ≠ indicesOld) ∨ (isAfterNew ≠ isAfterOld)) ∨ (isBeforeNew ≠ isBeforeOld)) DO
  indicesOld := indicesNew;
  isAfterOld := isAfterNew;
  isBeforeOld := isBeforeNew;
  indicesNew := (indicesNew ∪ (π[0, 5] (σ[(#2 = #1 ∧ #4 = #3)] ((isAfterOld × indicesOld) × Succ))));
  isAfterNew := (isAfterNew ∪ (π[1, 2] (σ[#3 = #0] (isAfterOld × isAfterOld))));
  isBeforeNew := (isBeforeNew ∪ (π[0, 3] (σ[#2 = #1] (isBeforeOld × isBeforeOld))))
  END;
  indices := indicesNew;
  isAfter := isAfterNew;
  isBefore := isBeforeNew;
  hasSuccessorRank := (π[0, 1] (σ[((#2 = #1 ∧ #4 = #0) ∧ #5 = #3)] ((indices × Succ) × indices)));
  index := (π[0, 1] (indices ∖ (π[0, 1] (σ[(#2 = #0 ∧ #3 = #1)] (indices × hasSuccessorRank)))));
  vertexS := ((π[0] edge) ∪ (π[1] edge));
  isolatedS := (π[0] ((vertexS ∖ (π[0] (σ[#2 = #0] (vertexS × edge)))) ∖ (π[0] (σ[#1 = #0] ((vertexS ∖ (π[0] (σ[#2 = #0] (vertexS × edge)))) × edge)))));
  sourceS := (π[0] ((σ[#1 = #0] (vertexS × edge)) ∖ (π[0, 1, 2] (σ[#4 = #0] ((σ[#1 = #0] (vertexS × edge)) × edge)))));
  indicesS := ∅[2];
  deltaIndicesS := ((π[0, 1] (isolatedS × RankZero)) ∪ (π[0, 1] (sourceS × RankOne)));
  deltaIndicesOldS := ∅[2];
  isAfterS := ∅[2];
  deltaIsAfterS := (π[1, 0] edge);
  deltaIsAfterOldS := ∅[2];
  isBeforeS := ∅[2];
  deltaIsBeforeS := (π[0, 1] edge);
  deltaIsBeforeOldS := ∅[2];
  WHILE (((deltaIndicesS ≠ ∅[2]) ∨ (deltaIsAfterS ≠ ∅[2])) ∨ (deltaIsBeforeS ≠ ∅[2])) DO
  deltaIndicesOldS := deltaIndicesS;
  deltaIsAfterOldS := deltaIsAfterS;
  deltaIsBeforeOldS := deltaIsBeforeS;
  indicesS := (indicesS ∪ deltaIndicesOldS);
  isAfterS := (isAfterS ∪ deltaIsAfterOldS);
  isBeforeS := (isBeforeS ∪ deltaIsBeforeOldS);
  deltaIndicesS := ((π[0, 5] (σ[(#2 = #1 ∧ #4 = #3)] ((isAfterS × deltaIndicesOldS) × Succ))) ∖ indicesS);
  deltaIsAfterS := ((((π[1, 2] (σ[#3 = #0] (deltaIsAfterOldS × isAfterS))) ∪ (π[1, 2] (σ[#3 = #0] (isAfterS × deltaIsAfterOldS)))) ∪ (π[1, 2] (σ[#3 = #0] (deltaIsAfterOldS × deltaIsAfterOldS)))) ∖ isAfterS);
  deltaIsBeforeS := ((((π[0, 3] (σ[#2 = #1] (deltaIsBeforeOldS × isBeforeS))) ∪ (π[0, 3] (σ[#2 = #1] (isBeforeS × deltaIsBeforeOldS)))) ∪ (π[0, 3] (σ[#2 = #1] (deltaIsBeforeOldS × deltaIsBeforeOldS)))) ∖ isBeforeS)
  END;
  hasSuccessorRankS := (π[0, 1] (σ[((#2 = #1 ∧ #4 = #0) ∧ #5 = #3)] ((indicesS × Succ) × indicesS)));
  indexS := (π[0, 1] (indicesS ∖ (π[0, 1] (σ[(#2 = #0 ∧ #3 = #1)] (indicesS × hasSuccessorRankS)))))

POST
  ((hasSuccessorRank = hasSuccessorRankS) ∧ (index = indexS) ∧ (indices = indicesS) ∧ (isAfter = isAfterS) ∧ (isBefore = isBeforeS) ∧ (isolated = isolatedS) ∧ (source = sourceS) ∧ (vertex = vertexS))
\end{Verbatim}

\Needspace{0.45\textheight}

\section{Example5012: Casbin depth-bounded role resolution contained in the unbounded role closure}\label{case:Example5012}

\begin{description}[leftmargin=2.6cm,style=nextline,itemsep=1pt]

\item[Category] production pair (two algorithms from production sources or legacy production cases); family: production discrepancy pair (Casbin depth bound).

\item[Kind] containment.

\item[Contributors] Leo Zhang.

\item[Status] certified valid (Certificate/Valid.lean, written by the certificate emitter); expected valid.

\item[Tags] \hyperref[tag:access-control]{\texttt{access-control}} \hyperref[tag:prophecy]{\texttt{prophecy}} \hyperref[tag:recursion-limit]{\texttt{recursion-limit}}.

\item[Input] \path{Benchmark/Example5012/Input.lean} (canonical notation).

\item[Summary] Casbin's depth-bounded role resolution against the unbounded role-hierarchy closure: the containment direction.

\item[Sides] Casbin HasLink with maxHierarchyLevel = 10, encoded as StepA = UserRole and ten distinct step relations StepB..StepK (role\_manager.go, enforcer.go) \par naive compiler on the recursive has-role program (Example2035 / 5004)

\item[Sources] {\raggedright \href{https://github.com/casbin/casbin/blob/master/rbac/default-role-manager/role_manager.go}{casbin/casbin rbac/default-role-manager/role\_manager.go (HasLink, hasLinkHelper)} \par \href{https://github.com/casbin/casbin/blob/master/enforcer.go}{casbin/casbin enforcer.go (maxHierarchyLevel = 10)}\par}

\item[Prophecy] kernel-checked certificate with a level-1 prophecy clause (Certificate/Valid.lean, Core.json).

\end{description}

\paragraph{Datalog program(s).}

{\raggedright\textit{prog\_p} --- fidelity module (kernel-checked: the compiled sides of the command are this program's naive compilation).\par}

\begin{Verbatim}[fontsize=\scriptsize,breaklines=true,breakanywhere=true,frame=single,framesep=1.5mm]
Has(x1, y1) :- UserRole(x1, y1).
Has(x1, z1) :- Has(x1, y1), RoleInh(y1, z1).
\end{Verbatim}

\paragraph{Reading.} Side 1 encodes Casbin's default role manager, which resolves role inheritance only up to a fixed depth bound (maxHierarchyLevel = 10) by unrolling ten loop-free steps of composition with RoleInh into HasB; side 2 is the naive-compiler translation of the unbounded recursive has-role Datalog program (Has derived from UserRole directly, and from Has composed with RoleInh) evaluated to its least fixpoint. The Hoare triple states that Casbin's depth-bounded role resolution is contained in the true unbounded transitive closure of role inheritance: every role grant the bounded lookup produces is also justified by the full closure.

\paragraph{First-order reading of PRE/POST, translated mechanically from the relational-algebra assertion in Input.lean.} The relational-algebra assertion in \texttt{Input.lean} remains the authority.

\textit{PRE:} ⊤\par True; no precondition constrains the inputs.

\textit{POST:} ∀x y. (HasB(x,y) → Has(x,y))\par The depth-bounded role relation HasB is a subset of the unbounded role-closure relation Has.

\paragraph{Whiel encoding (Hoare triple).}

\begin{Verbatim}[fontsize=\scriptsize,breaklines=true,breakanywhere=true,frame=single,framesep=1.5mm]
SCHEMA
  UserRole/2, RoleInh/2, StepA/2, StepB/2, StepC/2, StepD/2, StepE/2, StepF/2, StepG/2, StepH/2, StepI/2, StepJ/2, StepK/2, HasB/2, Has/2, Has_aux/2

PRE
  true 

CMD
  StepA := UserRole;
  StepB := (StepA ∪ (π[0, 3] (σ[#1 = #2] (StepA × RoleInh))));
  StepC := (StepB ∪ (π[0, 3] (σ[#1 = #2] (StepB × RoleInh))));
  StepD := (StepC ∪ (π[0, 3] (σ[#1 = #2] (StepC × RoleInh))));
  StepE := (StepD ∪ (π[0, 3] (σ[#1 = #2] (StepD × RoleInh))));
  StepF := (StepE ∪ (π[0, 3] (σ[#1 = #2] (StepE × RoleInh))));
  StepG := (StepF ∪ (π[0, 3] (σ[#1 = #2] (StepF × RoleInh))));
  StepH := (StepG ∪ (π[0, 3] (σ[#1 = #2] (StepG × RoleInh))));
  StepI := (StepH ∪ (π[0, 3] (σ[#1 = #2] (StepH × RoleInh))));
  StepJ := (StepI ∪ (π[0, 3] (σ[#1 = #2] (StepI × RoleInh))));
  StepK := (StepJ ∪ (π[0, 3] (σ[#1 = #2] (StepJ × RoleInh))));
  HasB := StepK;
  Has_aux := ∅[2];
  Has := (UserRole ∪ π[0,3] (σ[#1 = #2] ((Has_aux × RoleInh))));
  WHILE (¬((Has = Has_aux))) DO
    Has_aux := Has;
    Has := (Has ∪ (UserRole ∪ π[0,3] (σ[#1 = #2] ((Has_aux × RoleInh)))))
  END

POST
  (HasB ⊆ Has)
\end{Verbatim}

\Needspace{0.45\textheight}

\section{Example5013: Casbin depth-bounded role resolution versus the unbounded role closure}\label{case:Example5013}

\begin{description}[leftmargin=2.6cm,style=nextline,itemsep=1pt]

\item[Category] production pair (two algorithms from production sources or legacy production cases); family: production discrepancy pair (Casbin depth bound).

\item[Kind] equivalence.

\item[Contributors] Leo Zhang.

\item[Status] certified invalid (certificate checked outside the library build) (Certificate/Invalid.lean, written by the certificate emitter); expected invalid. Placement note: peak resident memory 6935552 KB while checking \path{Certificate/Invalid.lean} exceeds the library build's 4194304 KB limit (peak 6773 MB)

\item[Tags] \hyperref[tag:access-control]{\texttt{access-control}} \hyperref[tag:recursion-limit]{\texttt{recursion-limit}}.

\item[Input] \path{Benchmark/Example5013/Input.lean} (canonical notation).

\item[Summary] Casbin's depth-bounded role resolution against the unbounded role-hierarchy closure: the equality, which fails.

\item[Sides] Casbin HasLink with maxHierarchyLevel = 10, encoded as StepA = UserRole and ten distinct step relations StepB..StepK (role\_manager.go, enforcer.go) \par naive compiler on the recursive has-role program (Example2035 / 5004)

\item[Sources] {\raggedright \href{https://github.com/casbin/casbin/blob/master/rbac/default-role-manager/role_manager.go}{casbin/casbin rbac/default-role-manager/role\_manager.go (HasLink, hasLinkHelper)} \par \href{https://github.com/casbin/casbin/blob/master/enforcer.go}{casbin/casbin enforcer.go (maxHierarchyLevel = 10)}\par}

\end{description}

\paragraph{Datalog program(s).}

{\raggedright\textit{prog\_p} --- fidelity module (kernel-checked: the compiled sides of the command are this program's naive compilation).\par}

\begin{Verbatim}[fontsize=\scriptsize,breaklines=true,breakanywhere=true,frame=single,framesep=1.5mm]
Has(x1, y1) :- UserRole(x1, y1).
Has(x1, z1) :- Has(x1, y1), RoleInh(y1, z1).
\end{Verbatim}

\paragraph{Reading.} Same two sides as Example5012 (Casbin's depth-10-bounded role resolution HasB versus the unbounded recursive-closure relation Has), but here the postcondition asserts full equality rather than containment. The expected verdict is invalid: because Casbin stops resolving inheritance after ten steps while the true closure has no such bound, an eleven-edge inheritance chain lets the unbounded closure grant a role that the bounded resolution denies, so HasB is a strict subset of Has, not equal to it.

\paragraph{First-order reading of PRE/POST, translated mechanically from the relational-algebra assertion in Input.lean.} The relational-algebra assertion in \texttt{Input.lean} remains the authority.

\textit{PRE:} ⊤\par True; no precondition constrains the inputs.

\textit{POST:} ∀x y. (HasB(x,y) ↔ Has(x,y))\par The claim (which fails) is that the depth-bounded role relation HasB equals the unbounded role-closure relation Has.

\paragraph{Whiel encoding (Hoare triple).}

\begin{Verbatim}[fontsize=\scriptsize,breaklines=true,breakanywhere=true,frame=single,framesep=1.5mm]
SCHEMA
  UserRole/2, RoleInh/2, StepA/2, StepB/2, StepC/2, StepD/2, StepE/2, StepF/2, StepG/2, StepH/2, StepI/2, StepJ/2, StepK/2, HasB/2, Has/2, Has_aux/2

PRE
  true 

CMD
  StepA := UserRole;
  StepB := (StepA ∪ (π[0, 3] (σ[#1 = #2] (StepA × RoleInh))));
  StepC := (StepB ∪ (π[0, 3] (σ[#1 = #2] (StepB × RoleInh))));
  StepD := (StepC ∪ (π[0, 3] (σ[#1 = #2] (StepC × RoleInh))));
  StepE := (StepD ∪ (π[0, 3] (σ[#1 = #2] (StepD × RoleInh))));
  StepF := (StepE ∪ (π[0, 3] (σ[#1 = #2] (StepE × RoleInh))));
  StepG := (StepF ∪ (π[0, 3] (σ[#1 = #2] (StepF × RoleInh))));
  StepH := (StepG ∪ (π[0, 3] (σ[#1 = #2] (StepG × RoleInh))));
  StepI := (StepH ∪ (π[0, 3] (σ[#1 = #2] (StepH × RoleInh))));
  StepJ := (StepI ∪ (π[0, 3] (σ[#1 = #2] (StepI × RoleInh))));
  StepK := (StepJ ∪ (π[0, 3] (σ[#1 = #2] (StepJ × RoleInh))));
  HasB := StepK;
  Has_aux := ∅[2];
  Has := (UserRole ∪ π[0,3] (σ[#1 = #2] ((Has_aux × RoleInh))));
  WHILE (¬((Has = Has_aux))) DO
    Has_aux := Has;
    Has := (Has ∪ (UserRole ∪ π[0,3] (σ[#1 = #2] ((Has_aux × RoleInh)))))
  END

POST
  (HasB = Has)
\end{Verbatim}

\Needspace{0.45\textheight}

\section{Example5014: W3C SPARQL 1.1 ALP evaluation of p* and p+ versus the Datalog closure with reflexive pairs over nodes(G)}\label{case:Example5014}

\begin{description}[leftmargin=2.6cm,style=nextline,itemsep=1pt]

\item[Category] specification (a standard or a paper's procedure versus a declarative program); family: specification versus declarative program.

\item[Kind] equivalence.

\item[Contributors] Leo Zhang.

\item[Status] certified valid (Certificate/Valid.lean, written by the certificate emitter); expected valid.

\item[Tags] \hyperref[tag:evaluation-strategies]{\texttt{evaluation-strategies}} \hyperref[tag:semi-naive]{\texttt{semi-naive}}.

\item[Input] \path{Benchmark/Example5014/Input.lean} (canonical notation).

\item[Summary] SPARQL 1.1 property paths: the W3C evaluation procedure for `?x p* ?y` and `?x p+ ?y` (the ALP function with a visited set) against the Datalog transitive closure, with the reflexive pairs that `p*` adds over every term of the graph.

\item[Sides] hand-written: W3C ALP procedure set-at-a-time, visited set indexed by start, seeded with the reflexive pairs over nodes(G); OneOrMorePath = P ∘ Alp stated in the postcondition \par naive compiler: Tc = closure of P; Star = Tc plus reflexive pairs over nodes(G)

\item[Sources] {\raggedright \href{https://www.w3.org/TR/sparql11-query/}{W3C SPARQL 1.1 Query Language, Recommendation 21 March 2013} (sections 9.1 (ZeroOrMorePath, OneOrMorePath), 9.4 (Arbitrary Length Path Matching, set semantics), 18.4 (Definition: Node set of a graph; Definition: Function ALP; Evaluation of ZeroOrMorePath; Evaluation of OneOrMorePath); read 2026-09-15)\par}

\item[Twin] Example5015 (invalid: reflexive pairs over the endpoints of P only).

\item[Obstruction note] from the metadata: anchored (independent loops merged side by side, lockstep frontiers: Front = Star ∖ Star\_aux; the p+ conjunct needs the closure inclusions P ⊆ Tc, P ∘ Tc ⊆ Tc, Tc ⊆ P ∪ P ∘ Tc)

\end{description}

\paragraph{Datalog program(s).}

{\raggedright\textit{prog\_p --- `p*` over nodes(G) and the closure `p+`, as Datalog.} --- fidelity module (kernel-checked: the compiled sides of the command are this program's naive compilation).\par}

\begin{Verbatim}[fontsize=\scriptsize,breaklines=true,breakanywhere=true,frame=single,framesep=1.5mm]
Tc(x1, y1) :- P(x1, y1).
Tc(x1, z1) :- Tc(x1, y1), P(y1, z1).
Star(x1, x1) :- G(x1, y1).
Star(y1, y1) :- G(x1, y1).
Star(x1, y1) :- Tc(x1, y1).
\end{Verbatim}

\paragraph{Reading.} Side 1 implements the W3C SPARQL 1.1 ALP (arbitrary-length-path) evaluation procedure for property paths ?x p* ?y and ?x p+ ?y, run set-at-a-time from every node of the graph via a frontier/visited-set loop; side 2 is the naive-compiler translation of the corresponding Datalog transitive-closure program (Tc, the closure of the predicate-specific relation P) together with a reflexive relation Star seeded over every node of the whole graph G. The triple states the specification's ALP procedure computes exactly the same result as this declarative closure formulation: the visited set Alp equals Star, and one P-step composed with Alp equals the closure Tc, the equivalence that licenses implementing p*/p+ by a recursive query instead of explicit graph traversal.

\paragraph{First-order reading of PRE/POST, translated mechanically from the relational-algebra assertion in Input.lean.} The relational-algebra assertion in \texttt{Input.lean} remains the authority.

\textit{PRE:} ∀x y. (P(x,y) → G(x,y))\par The predicate-specific relation P (triples with the path's predicate) is a subset of the full graph relation G.

\textit{POST:} (∀x y. (Alp(x,y) ↔ Star(x,y))) ∧ ∀z u. ((∃v. P(z,v) ∧ Alp(v,u)) ↔ Tc(z,u))\par The ALP visited-set relation Alp equals Star (the closure of P plus reflexive pairs over every node of G), and composing one P-step with Alp equals the transitive-closure relation Tc.

\paragraph{Whiel encoding (Hoare triple).}

\begin{Verbatim}[fontsize=\scriptsize,breaklines=true,breakanywhere=true,frame=single,framesep=1.5mm]
SCHEMA
  G/2, P/2, Alp/2, Front/2, Next/2, Tc/2, Tc_aux/2, Star/2, Star_aux/2

PRE
  (P ⊆ G) 

CMD
  Alp := (π[0, 0] G ∪ π[1, 1] G);
  Front := Alp;
  WHILE (Front ≠ ∅) DO
    Next := (π[0, 3] (σ[#1 = #2] (Front × P)) ∖ Alp);
    Alp := (Alp ∪ Next);
    Front := Next
  END;
  Tc_aux := ∅[2];
  Star_aux := ∅[2];
  Tc := (P ∪ π[0,3] (σ[#1 = #2] ((Tc_aux × P))));
  Star := (π[0,0] (G) ∪ (π[1,1] (G) ∪ Tc_aux));
  WHILE (¬(((Tc = Tc_aux) ∧ (Star = Star_aux)))) DO
    Tc_aux := Tc;
    Star_aux := Star;
    Tc := (Tc ∪ (P ∪ π[0,3] (σ[#1 = #2] ((Tc_aux × P)))));
    Star := (Star ∪ (π[0,0] (G) ∪ (π[1,1] (G) ∪ Tc_aux)))
  END

POST
  ((Alp = Star) ∧ (π[0, 3] (σ[#1 = #2] (P × Alp)) = Tc))
\end{Verbatim}

\Needspace{0.45\textheight}

\section{Example5015: W3C SPARQL 1.1 ALP evaluation of p* versus closure plus identity over the p endpoints (invalid twin of Example5014)}\label{case:Example5015}

\begin{description}[leftmargin=2.6cm,style=nextline,itemsep=1pt]

\item[Category] specification (a standard or a paper's procedure versus a declarative program); family: specification versus declarative program (invalid twin).

\item[Kind] equivalence.

\item[Contributors] Leo Zhang.

\item[Status] certified invalid (Certificate/Invalid.lean, written by the certificate emitter); expected invalid.

\item[Tags] \hyperref[tag:evaluation-strategies]{\texttt{evaluation-strategies}} \hyperref[tag:semi-naive]{\texttt{semi-naive}}.

\item[Input] \path{Benchmark/Example5015/Input.lean} (canonical notation).

\item[Summary] SPARQL 1.1 property paths, the invalid twin of Example5014: the W3C procedure for `?x p* ?y` against the closure of the `p` triples plus the identity on their own endpoints, the common misreading of the zero-length case.

\item[Sides] hand-written: W3C ALP procedure as in Example5014 \par naive compiler: Tc = closure of P; Star = Tc plus reflexive pairs over the endpoints of P (the misreading)

\item[Sources] {\raggedright \href{https://www.w3.org/TR/sparql11-query/}{W3C SPARQL 1.1 Query Language, Recommendation 21 March 2013} (sections 9.1 (ZeroOrMorePath, OneOrMorePath), 9.4 (Arbitrary Length Path Matching, set semantics), 18.4 (Definition: Node set of a graph; Definition: Function ALP; Evaluation of ZeroOrMorePath; Evaluation of OneOrMorePath); read 2026-09-15)\par}

\item[Twin] Example5014 (valid).

\item[Obstruction note] from the metadata: n/a (invalid; kernel counterexample)

\end{description}

\paragraph{Datalog program(s).}

{\raggedright\textit{prog\_p --- The closure `p+` and "closure plus identity" over the endpoints of `p` (the misreading of `p*`).} --- fidelity module (kernel-checked: the compiled sides of the command are this program's naive compilation).\par}

\begin{Verbatim}[fontsize=\scriptsize,breaklines=true,breakanywhere=true,frame=single,framesep=1.5mm]
Tc(x1, y1) :- P(x1, y1).
Tc(x1, z1) :- Tc(x1, y1), P(y1, z1).
Star(x1, x1) :- P(x1, y1).
Star(y1, y1) :- P(x1, y1).
Star(x1, y1) :- Tc(x1, y1).
\end{Verbatim}

\paragraph{Reading.} Same ALP procedure (side 1) as Example5014, but side 2's reflexive pairs are seeded only over the endpoints of P (the p-labeled triples) rather than over every node of the whole graph G --- what one gets by naively translating p* as closure-plus-identity restricted to the p-subgraph. The expected verdict is invalid: the W3C specification's node set ranges over all of G, so a node occurring only in triples with other predicates still satisfies ?x p* ?x under the specification but is missing from this narrower Star relation, breaking the claimed equality.

\paragraph{First-order reading of PRE/POST, translated mechanically from the relational-algebra assertion in Input.lean.} The relational-algebra assertion in \texttt{Input.lean} remains the authority.

\textit{PRE:} ∀x y. (P(x,y) → G(x,y))\par The predicate-specific relation P is a subset of the full graph relation G.

\textit{POST:} (∀x y. (Alp(x,y) ↔ Star(x,y))) ∧ ∀z u. ((∃v. P(z,v) ∧ Alp(v,u)) ↔ Tc(z,u))\par The claim (which fails) is that the ALP visited-set relation Alp equals Star, here seeded reflexively only over the endpoints of P, and that one P-step composed with Alp equals Tc.

\paragraph{Whiel encoding (Hoare triple).}

\begin{Verbatim}[fontsize=\scriptsize,breaklines=true,breakanywhere=true,frame=single,framesep=1.5mm]
SCHEMA
  G/2, P/2, Alp/2, Front/2, Next/2, Tc/2, Tc_aux/2, Star/2, Star_aux/2

PRE
  (P ⊆ G) 

CMD
  Alp := (π[0, 0] G ∪ π[1, 1] G);
  Front := Alp;
  WHILE (Front ≠ ∅) DO
    Next := (π[0, 3] (σ[#1 = #2] (Front × P)) ∖ Alp);
    Alp := (Alp ∪ Next);
    Front := Next
  END;
  Tc_aux := ∅[2];
  Star_aux := ∅[2];
  Tc := (P ∪ π[0,3] (σ[#1 = #2] ((Tc_aux × P))));
  Star := (π[0,0] (P) ∪ (π[1,1] (P) ∪ Tc_aux));
  WHILE (¬(((Tc = Tc_aux) ∧ (Star = Star_aux)))) DO
    Tc_aux := Tc;
    Star_aux := Star;
    Tc := (Tc ∪ (P ∪ π[0,3] (σ[#1 = #2] ((Tc_aux × P)))));
    Star := (Star ∪ (π[0,0] (P) ∪ (π[1,1] (P) ∪ Tc_aux)))
  END

POST
  ((Alp = Star) ∧ (π[0, 3] (σ[#1 = #2] (P × Alp)) = Tc))
\end{Verbatim}

\Needspace{0.45\textheight}

\section{Example5016: DRed deletion maintenance of a materialised transitive closure versus recomputation}\label{case:Example5016}

\begin{description}[leftmargin=2.6cm,style=nextline,itemsep=1pt]

\item[Category] specification (a standard or a paper's procedure versus a declarative program); family: maintenance algorithm versus recomputation.

\item[Kind] equivalence.

\item[Contributors] Leo Zhang.

\item[Status] certified valid (certificate checked outside the library build) (Certificate/Valid.lean, written by the certificate emitter); expected valid. Placement note: a proof module imported by \path{Certificate/Valid.lean} does not build under the library build's 4194304 KB limit; it builds under 12582912 KB (peak 1673 MB)

\item[Tags] \hyperref[tag:evaluation-strategies]{\texttt{evaluation-strategies}} \hyperref[tag:negation]{\texttt{negation}} \hyperref[tag:prophecy]{\texttt{prophecy}} \hyperref[tag:semi-naive]{\texttt{semi-naive}} \hyperref[tag:view-maintenance]{\texttt{view-maintenance}}.

\item[Input] \path{Benchmark/Example5016/Input.lean} (canonical notation).

\item[Summary] DRed (Delete/Rederive) maintenance of a materialised transitive closure under deletions against recomputation of the closure from the updated edges.

\item[Sides] naive compiler: T = closure of E (old materialisation) \par hand-written DRed after Motik et al. DR1-DR5 / GMS section 7 steps 1-2 for the closure program, deletion-only update \par naive compiler: Tn = closure of EN = E ∖ Del

\item[Sources] {\raggedright \href{https://doi.org/10.1145/170035.170066}{A. Gupta, I. S. Mumick, V. S. Subrahmanian, Maintaining Views Incrementally, SIGMOD 1993} (sections section 7 (Incremental Maintenance of Recursive Views: steps 1-3, Theorem 7.1); read 2026-09-15) \par \href{https://doi.org/10.1609/aaai.v29i1.9409}{B. Motik, Y. Nenov, R. Piro, I. Horrocks, Incremental Update of Datalog Materialisation: the Backward/Forward Algorithm, AAAI 2015} (sections section 3 (The Delete/Rederive Algorithm, steps DR1-DR5, rules (5) and (6)); Algorithm 1 and Theorem 1 (B/F); read 2026-09-15)\par}

\item[Twin] Example5017 (invalid: overdeletion without rederivation).

\item[Prophecy] accepted invariant with a prophecy clause above level zero (Core.json); no certificate built from it yet.

\end{description}

\paragraph{Datalog program(s).}

{\raggedright\textit{prog\_pa --- The maintained program: the closure of the old edges.} --- fidelity module (kernel-checked: the compiled sides of the command are this program's naive compilation).\par}

\begin{Verbatim}[fontsize=\scriptsize,breaklines=true,breakanywhere=true,frame=single,framesep=1.5mm]
T(x1, y1) :- E(x1, y1).
T(x1, z1) :- T(x1, y1), E(y1, z1).
\end{Verbatim}

{\raggedright\textit{prog\_pc --- Recomputation: the closure of the updated edges.} --- fidelity module (kernel-checked: the compiled sides of the command are this program's naive compilation).\par}

\begin{Verbatim}[fontsize=\scriptsize,breaklines=true,breakanywhere=true,frame=single,framesep=1.5mm]
Tn(x1, y1) :- EN(x1, y1).
Tn(x1, z1) :- Tn(x1, y1), EN(y1, z1).
\end{Verbatim}

{\raggedright\textit{rules as stated in the header comment of Input.lean} --- header comment.\par}

\begin{Verbatim}[fontsize=\scriptsize,breaklines=true,breakanywhere=true,frame=single,framesep=1.5mm]
T(x, y) :- E(x, y)
T(x, z) :- T(x, y), E(y, z)
T⁻(x, y) :- E⁻(x, y)
T⁻(x, z) :- T(x, y), E⁻(y, z)
T⁻(x, z) :- T⁻(x, y), E(y, z)
T⁺(x, y) :- T⁻(x, y), EN(x, y)
T⁺(x, z) :- T⁻(x, z), T(x, y), EN(y, z)
T⁺(x, z) :- T⁺(x, y), EN(y, z)
Tn(x, y) :- EN(x, y)
Tn(x, z) :- Tn(x, y), EN(y, z)
\end{Verbatim}

\paragraph{Reading.} The maintained program is the transitive closure T of an edge relation E; side 1 implements DRed (Delete/Rederive), the standard incremental-view-maintenance algorithm of Gupta, Mumick and Subrahmanian: it overdeletes closure facts that could depend on deleted edges, removes them, then rederives any of those facts that still have a surviving derivation on the updated edge set, propagating rederived facts back through the closure rule; side 2 recomputes the closure Tn from scratch on the updated edge set EN = E minus Del. The triple states DRed's maintained closure equals the recomputed closure, the standard correctness result for incremental deletion maintenance of a recursive Datalog view.

\paragraph{First-order reading of PRE/POST, translated mechanically from the relational-algebra assertion in Input.lean.} The relational-algebra assertion in \texttt{Input.lean} remains the authority.

\textit{PRE:} ∀x y. (Del(x,y) → E(x,y))\par The set of edges to delete, Del, is a subset of the edge relation E.

\textit{POST:} ∀x y. (T(x,y) ↔ Tn(x,y))\par The DRed-maintained closure relation T equals the freshly recomputed closure Tn on the edges remaining after deletion.

\paragraph{Whiel encoding (Hoare triple).}

\begin{Verbatim}[fontsize=\tiny,breaklines=true,breakanywhere=true,frame=single,framesep=1.5mm]
SCHEMA
  E/2, Del/2, EN/2, T/2, T_aux/2, Over/2, Cand/2, New/2, Tn/2, Tn_aux/2

PRE
  (Del ⊆ E) 

CMD
  T_aux := ∅[2];
  T := (E ∪ π[0,3] (σ[#1 = #2] ((T_aux × E))));
  WHILE (¬((T = T_aux))) DO
    T_aux := T;
    T := (T ∪ (E ∪ π[0,3] (σ[#1 = #2] ((T_aux × E)))))
  END;
  EN := (E ∖ Del);
  Over := ∅[2];
  Cand := (Del ∪ π[0, 3] (σ[#1 = #2] (T × Del)));
  WHILE ((Cand ∖ Over) ≠ ∅) DO
    New := (Cand ∖ Over);
    Over := (Over ∪ New);
    Cand := π[0, 3] (σ[#1 = #2] (New × E))
  END;
  T := (T ∖ Over);
  Cand := (Over ∖ (Over ∖ (EN ∪ π[0, 3] (σ[#1 = #2] (T × EN)))));
  WHILE ((Cand ∖ T) ≠ ∅) DO
    New := (Cand ∖ T);
    T := (T ∪ New);
    Cand := π[0, 3] (σ[#1 = #2] (New × EN))
  END;
  Tn_aux := ∅[2];
  Tn := (EN ∪ π[0,3] (σ[#1 = #2] ((Tn_aux × EN))));
  WHILE (¬((Tn = Tn_aux))) DO
    Tn_aux := Tn;
    Tn := (Tn ∪ (EN ∪ π[0,3] (σ[#1 = #2] ((Tn_aux × EN)))))
  END

POST
  (T = Tn)
\end{Verbatim}

\Needspace{0.45\textheight}

\section{Example5017: Overdeletion without rederivation versus recomputation (invalid twin of Example5016)}\label{case:Example5017}

\begin{description}[leftmargin=2.6cm,style=nextline,itemsep=1pt]

\item[Category] specification (a standard or a paper's procedure versus a declarative program); family: maintenance algorithm versus recomputation (invalid twin).

\item[Kind] equivalence.

\item[Contributors] Leo Zhang.

\item[Status] certified invalid (Certificate/Invalid.lean, written by the certificate emitter); expected invalid.

\item[Tags] \hyperref[tag:evaluation-strategies]{\texttt{evaluation-strategies}} \hyperref[tag:negation]{\texttt{negation}} \hyperref[tag:semi-naive]{\texttt{semi-naive}} \hyperref[tag:view-maintenance]{\texttt{view-maintenance}}.

\item[Input] \path{Benchmark/Example5017/Input.lean} (canonical notation).

\item[Summary] Deletion by overdeletion alone (semi-naive propagation of the deleted edges through the closure, without DRed's rederivation phase) against recomputation: the invalid twin of Example5016.

\item[Sides] naive compiler: T = closure of E \par hand-written: DRed step 1 only (DR1-DR3: overdeletion, removal), no rederivation \par naive compiler: Tn = closure of EN = E ∖ Del

\item[Sources] {\raggedright \href{https://doi.org/10.1145/170035.170066}{A. Gupta, I. S. Mumick, V. S. Subrahmanian, Maintaining Views Incrementally, SIGMOD 1993} (sections section 7 (Incremental Maintenance of Recursive Views: steps 1-3, Theorem 7.1); read 2026-09-15) \par \href{https://doi.org/10.1609/aaai.v29i1.9409}{B. Motik, Y. Nenov, R. Piro, I. Horrocks, Incremental Update of Datalog Materialisation: the Backward/Forward Algorithm, AAAI 2015} (sections section 3 (The Delete/Rederive Algorithm, steps DR1-DR5, rules (5) and (6)); Algorithm 1 and Theorem 1 (B/F); read 2026-09-15)\par}

\item[Twin] Example5016 (valid).

\item[Obstruction note] from the metadata: n/a (invalid; kernel counterexample)

\end{description}

\paragraph{Datalog program(s).}

{\raggedright\textit{prog\_pa --- The maintained program: the closure of the old edges.} --- fidelity module (kernel-checked: the compiled sides of the command are this program's naive compilation).\par}

\begin{Verbatim}[fontsize=\scriptsize,breaklines=true,breakanywhere=true,frame=single,framesep=1.5mm]
T(x1, y1) :- E(x1, y1).
T(x1, z1) :- T(x1, y1), E(y1, z1).
\end{Verbatim}

{\raggedright\textit{prog\_pc --- Recomputation: the closure of the updated edges.} --- fidelity module (kernel-checked: the compiled sides of the command are this program's naive compilation).\par}

\begin{Verbatim}[fontsize=\scriptsize,breaklines=true,breakanywhere=true,frame=single,framesep=1.5mm]
Tn(x1, y1) :- EN(x1, y1).
Tn(x1, z1) :- Tn(x1, y1), EN(y1, z1).
\end{Verbatim}

{\raggedright\textit{rules as stated in the header comment of Input.lean} --- header comment.\par}

\begin{Verbatim}[fontsize=\scriptsize,breaklines=true,breakanywhere=true,frame=single,framesep=1.5mm]
T(x, y) :- E(x, y)
T(x, z) :- T(x, y), E(y, z)
T⁻(x, y) :- E⁻(x, y)
T⁻(x, z) :- T(x, y), E⁻(y, z)
T⁻(x, z) :- T⁻(x, y), E(y, z)
Tn(x, y) :- EN(x, y)
Tn(x, z) :- Tn(x, y), EN(y, z)
\end{Verbatim}

\paragraph{Reading.} Same maintained closure and recomputation side as Example5016, but side 1 performs only the overdeletion phase of DRed (propagating deleted edges through the closure and removing every fact that could depend on them) without the rederivation phase that restores facts with alternative surviving derivations. The expected verdict is invalid: overdeletion is a genuine overestimate, so a closure fact reachable via a deleted edge but also via another surviving path is deleted and never restored, leaving the maintained closure T missing facts that the recomputed closure Tn still contains.

\paragraph{First-order reading of PRE/POST, translated mechanically from the relational-algebra assertion in Input.lean.} The relational-algebra assertion in \texttt{Input.lean} remains the authority.

\textit{PRE:} ∀x y. (Del(x,y) → E(x,y))\par The set of edges to delete, Del, is a subset of the edge relation E.

\textit{POST:} ∀x y. (T(x,y) ↔ Tn(x,y))\par The claim (which fails) is that the overdeletion-only maintained closure T equals the recomputed closure Tn on the edges remaining after deletion.

\paragraph{Whiel encoding (Hoare triple).}

\begin{Verbatim}[fontsize=\scriptsize,breaklines=true,breakanywhere=true,frame=single,framesep=1.5mm]
SCHEMA
  E/2, Del/2, EN/2, T/2, T_aux/2, Over/2, Cand/2, New/2, Tn/2, Tn_aux/2

PRE
  (Del ⊆ E) 

CMD
  T_aux := ∅[2];
  T := (E ∪ π[0,3] (σ[#1 = #2] ((T_aux × E))));
  WHILE (¬((T = T_aux))) DO
    T_aux := T;
    T := (T ∪ (E ∪ π[0,3] (σ[#1 = #2] ((T_aux × E)))))
  END;
  EN := (E ∖ Del);
  Over := ∅[2];
  Cand := (Del ∪ π[0, 3] (σ[#1 = #2] (T × Del)));
  WHILE ((Cand ∖ Over) ≠ ∅) DO
    New := (Cand ∖ Over);
    Over := (Over ∪ New);
    Cand := π[0, 3] (σ[#1 = #2] (New × E))
  END;
  T := (T ∖ Over);
  Tn_aux := ∅[2];
  Tn := (EN ∪ π[0,3] (σ[#1 = #2] ((Tn_aux × EN))));
  WHILE (¬((Tn = Tn_aux))) DO
    Tn_aux := Tn;
    Tn := (Tn ∪ (EN ∪ π[0,3] (σ[#1 = #2] ((Tn_aux × EN)))))
  END

POST
  (T = Tn)
\end{Verbatim}

\Needspace{0.45\textheight}

\section{Example5018: Chase order-independence on the full-TGD subset of ChaseBench correctness/tgds: dependency-at-a-time versus parallel chase}\label{case:Example5018}

\begin{description}[leftmargin=2.6cm,style=nextline,itemsep=1pt]

\item[Category] specification (a standard or a paper's procedure versus a declarative program); family: two evaluation orders of one chase.

\item[Kind] equivalence.

\item[Contributors] Leo Zhang.

\item[Status] certified valid (certificate checked outside the library build) (Certificate/Valid.lean, written by the certificate emitter); expected valid. Placement note: a proof module imported by \path{Certificate/Valid.lean} does not build under the library build's 4194304 KB limit; it builds under 12582912 KB (peak 1714 MB)

\item[Tags] \hyperref[tag:chase]{\texttt{chase}} \hyperref[tag:evaluation-strategies]{\texttt{evaluation-strategies}} \hyperref[tag:mutual-recursion]{\texttt{mutual-recursion}} \hyperref[tag:prophecy]{\texttt{prophecy}}.

\item[Input] \path{Benchmark/Example5018/Input.lean} (canonical notation).

\item[Summary] Chase order-independence on a full-TGD scenario of ChaseBench: the parallel chase (every dependency applied to every active trigger in each round) against the dependency-at-a-time chase (one dependency chased to its own fixpoint, then the next, the sweep repeated until no dependency has an active trigger).

\item[Sides] hand-written: s-t TGDs once, then target TGDs in file order each to its fixpoint (inner loops), sweep repeated (outer loop) \par naive compiler: the six full TGDs as a Datalog program (parallel chase)

\item[Sources] {\raggedright \href{https://github.com/dbunibas/chasebench}{ChaseBench (dbunibas/chasebench), scenario scenarios/correctness/tgds} (commit 7427e1c1e3196769ee4ea04557465965ec26d02b (master, 2017-09-13); read 2026-09-15) \par \href{https://doi.org/10.1145/3034786.3034796}{M. Benedikt, G. Konstantinidis, G. Mecca, B. Motik, P. Papotti, D. Santoro, E. Tsamoura, Benchmarking the Chase, PODS 2017} (sections Section 2 (chase sequences, fairness), Section 3 (Algorithm 1, Chase variants), Section 8 (Restricted vs. unrestricted vs. parallel chase); read 2026-09-15) \par program source file ChaseBench \path{scenarios/correctness/tgds} (full TGDs; the existential t2 -\textgreater{} t3 dropped, t3 an input)\par}

\item[Twin] Example5019 (invalid: one pass in file order).

\item[Prophecy] kernel-checked certificate with a level-1 prophecy clause (Certificate/Valid.lean, Core.json).

\end{description}

\paragraph{Datalog program(s).}

{\raggedright\textit{prog\_p --- The full TGDs of ChaseBench scenarios/correctness/tgds as a Datalog program (S = s, TA = t1, TB = t2, TC = t3, WA = w1, WB = w2).} --- fidelity module (kernel-checked: the compiled sides of the command are this program's naive compilation).\par}

\begin{Verbatim}[fontsize=\scriptsize,breaklines=true,breakanywhere=true,frame=single,framesep=1.5mm]
TA(x1, y1, z1) :- S(x1, y1, z1).
WA(x1, y1) :- S(x1, y1, z1).
TB(x1, y1) :- TA(x1, y1, z1).
TB(y1, y1) :- TC(x1, y1, z1).
WB(x1, y1) :- WA(x1, y1).
WA(y1, y1) :- WB(x1, y1).
\end{Verbatim}

\paragraph{Reading.} The program is the full-TGD subset of ChaseBench's tgds correctness scenario, a small chain of six dependencies (renamed S/TA/TB/TC/WA/WB) relating source and target relations in a data-exchange setting; side 1 hand-encodes the dependency-at-a-time chase (each dependency chased to its own fixpoint in file order, the sweep over all dependencies repeated until nothing new is derived), while side 2 is the naive-compiler translation applying all dependencies in parallel each round. The triple states the two chase orders converge to the same universal solution, illustrating the chase's order-independence for full (non-existential) dependencies.

\paragraph{First-order reading of PRE/POST, translated mechanically from the relational-algebra assertion in Input.lean.} The relational-algebra assertion in \texttt{Input.lean} remains the authority.

\textit{PRE:} ⊤\par True; no precondition constrains the inputs.

\textit{POST:} (∀x y z. (TAQ(x,y,z) ↔ TA(x,y,z))) ∧ (∀u v. (TBQ(u,v) ↔ TB(u,v))) ∧ (∀w x1. (WAQ(w,x1) ↔ WA(w,x1))) ∧ ∀x2 x3. (WBQ(x2,x3) ↔ WB(x2,x3))\par Each target relation computed by the dependency-at-a-time chase (TAQ, TBQ, WAQ, WBQ) equals the corresponding relation computed by the parallel chase (TA, TB, WA, WB).

\paragraph{Whiel encoding (Hoare triple).}

\begin{Verbatim}[fontsize=\tiny,breaklines=true,breakanywhere=true,frame=single,framesep=1.5mm]
SCHEMA
  S/3, TA/3, TC/3, TA_aux/3, TAQ/3, TB/2, WA/2, WB/2, TB_aux/2, WA_aux/2, WB_aux/2, TBQ/2, WAQ/2, WBQ/2

PRE
  true 

CMD
  TAQ := S;
  WAQ := π[0, 1] S;
  TBQ := ∅[2];
  WBQ := ∅[2];
  WHILE (¬((((π[0, 1] TAQ) ⊆ TBQ) ∧ ((π[1, 1] TC) ⊆ TBQ)) ∧ ((WAQ ⊆ WBQ) ∧ ((π[1, 1] WBQ) ⊆ WAQ)))) DO
    WHILE (¬((π[0, 1] TAQ) ⊆ TBQ)) DO
      TBQ := (TBQ ∪ π[0, 1] TAQ)
    END;
    WHILE (¬((π[1, 1] TC) ⊆ TBQ)) DO
      TBQ := (TBQ ∪ π[1, 1] TC)
    END;
    WHILE (¬(WAQ ⊆ WBQ)) DO
      WBQ := (WBQ ∪ WAQ)
    END;
    WHILE (¬((π[1, 1] WBQ) ⊆ WAQ)) DO
      WAQ := (WAQ ∪ π[1, 1] WBQ)
    END
  END;
  TA_aux := ∅[3];
  WA_aux := ∅[2];
  TB_aux := ∅[2];
  WB_aux := ∅[2];
  TA := S;
  WA := (π[0,1] (S) ∪ π[1,1] (WB_aux));
  TB := (π[0,1] (TA_aux) ∪ π[1,1] (TC));
  WB := WA_aux;
  WHILE (¬(((TA = TA_aux) ∧ ((WA = WA_aux) ∧ ((TB = TB_aux) ∧ (WB = WB_aux)))))) DO
  TA_aux := TA;
  WA_aux := WA;
  TB_aux := TB;
  WB_aux := WB;
  TA := (TA ∪ S);
  WA := (WA ∪ (π[0,1] (S) ∪ π[1,1] (WB_aux)));
  TB := (TB ∪ (π[0,1] (TA_aux) ∪ π[1,1] (TC)));
  WB := (WB ∪ WA_aux)
  END

POST
  (((TAQ = TA) ∧ (TBQ = TB)) ∧ ((WAQ = WA) ∧ (WBQ = WB)))
\end{Verbatim}

\Needspace{0.45\textheight}

\section{Example5019: One pass over the dependencies versus the parallel chase on ChaseBench correctness/tgds (invalid twin of Example5018)}\label{case:Example5019}

\begin{description}[leftmargin=2.6cm,style=nextline,itemsep=1pt]

\item[Category] specification (a standard or a paper's procedure versus a declarative program); family: two evaluation orders of one chase (invalid twin).

\item[Kind] equivalence.

\item[Contributors] Leo Zhang.

\item[Status] certified invalid (Certificate/Invalid.lean, written by the certificate emitter); expected invalid.

\item[Tags] \hyperref[tag:chase]{\texttt{chase}} \hyperref[tag:evaluation-strategies]{\texttt{evaluation-strategies}} \hyperref[tag:mutual-recursion]{\texttt{mutual-recursion}}.

\item[Input] \path{Benchmark/Example5019/Input.lean} (canonical notation).

\item[Summary] One pass over the dependencies of a full-TGD ChaseBench scenario against the parallel chase: the invalid twin of Example5018.

\item[Sides] hand-written: s-t TGDs once, then each target TGD to its fixpoint once in file order, no repetition (not a fair chase sequence) \par naive compiler: the six full TGDs as a Datalog program (parallel chase)

\item[Sources] {\raggedright \href{https://github.com/dbunibas/chasebench}{ChaseBench (dbunibas/chasebench), scenario scenarios/correctness/tgds} (commit 7427e1c1e3196769ee4ea04557465965ec26d02b (master, 2017-09-13); read 2026-09-15) \par \href{https://doi.org/10.1145/3034786.3034796}{M. Benedikt, G. Konstantinidis, G. Mecca, B. Motik, P. Papotti, D. Santoro, E. Tsamoura, Benchmarking the Chase, PODS 2017} (sections Section 2 (chase sequences, fairness), Section 3 (Algorithm 1, Chase variants), Section 8 (Restricted vs. unrestricted vs. parallel chase); read 2026-09-15) \par program source file ChaseBench \path{scenarios/correctness/tgds} (full TGDs)\par}

\item[Twin] Example5018 (valid).

\item[Obstruction note] from the metadata: n/a (invalid; kernel counterexample)

\end{description}

\paragraph{Datalog program(s).}

{\raggedright\textit{prog\_p --- The full TGDs of ChaseBench scenarios/correctness/tgds as a Datalog program (S = s, TA = t1, TB = t2, TC = t3, WA = w1, WB = w2).} --- fidelity module (kernel-checked: the compiled sides of the command are this program's naive compilation).\par}

\begin{Verbatim}[fontsize=\scriptsize,breaklines=true,breakanywhere=true,frame=single,framesep=1.5mm]
TA(x1, y1, z1) :- S(x1, y1, z1).
WA(x1, y1) :- S(x1, y1, z1).
TB(x1, y1) :- TA(x1, y1, z1).
TB(y1, y1) :- TC(x1, y1, z1).
WB(x1, y1) :- WA(x1, y1).
WA(y1, y1) :- WB(x1, y1).
\end{Verbatim}

\paragraph{Reading.} Same scenario and parallel-chase side as Example5018, but side 1 applies the s-t dependencies once and then chases each target dependency to its own fixpoint exactly once in file order, without repeating the sweep over all dependencies. The expected verdict is invalid: that single pass is not a fair chase sequence, since a dependency whose trigger becomes newly active only after a later dependency fires is never re-applied, so the one-pass result falls short of the true least fixpoint the parallel chase computes.

\paragraph{First-order reading of PRE/POST, translated mechanically from the relational-algebra assertion in Input.lean.} The relational-algebra assertion in \texttt{Input.lean} remains the authority.

\textit{PRE:} ⊤\par True; no precondition constrains the inputs.

\textit{POST:} (∀x y z. (TAQ(x,y,z) ↔ TA(x,y,z))) ∧ (∀u v. (TBQ(u,v) ↔ TB(u,v))) ∧ (∀w x1. (WAQ(w,x1) ↔ WA(w,x1))) ∧ ∀x2 x3. (WBQ(x2,x3) ↔ WB(x2,x3))\par The claim (which fails) is that the single-pass target relations (TAQ, TBQ, WAQ, WBQ) equal the parallel chase's fixpoint relations (TA, TB, WA, WB).

\paragraph{Whiel encoding (Hoare triple).}

\begin{Verbatim}[fontsize=\tiny,breaklines=true,breakanywhere=true,frame=single,framesep=1.5mm]
SCHEMA
  S/3, TA/3, TC/3, TA_aux/3, TAQ/3, TB/2, WA/2, WB/2, TB_aux/2, WA_aux/2, WB_aux/2, TBQ/2, WAQ/2, WBQ/2

PRE
  true 

CMD
  TAQ := S;
  WAQ := π[0, 1] S;
  TBQ := ∅[2];
  WBQ := ∅[2];
  WHILE (¬((π[0, 1] TAQ) ⊆ TBQ)) DO
    TBQ := (TBQ ∪ π[0, 1] TAQ)
  END;
  WHILE (¬((π[1, 1] TC) ⊆ TBQ)) DO
    TBQ := (TBQ ∪ π[1, 1] TC)
  END;
  WHILE (¬(WAQ ⊆ WBQ)) DO
    WBQ := (WBQ ∪ WAQ)
  END;
  WHILE (¬((π[1, 1] WBQ) ⊆ WAQ)) DO
    WAQ := (WAQ ∪ π[1, 1] WBQ)
  END;
  TA_aux := ∅[3];
  WA_aux := ∅[2];
  TB_aux := ∅[2];
  WB_aux := ∅[2];
  TA := S;
  WA := (π[0,1] (S) ∪ π[1,1] (WB_aux));
  TB := (π[0,1] (TA_aux) ∪ π[1,1] (TC));
  WB := WA_aux;
  WHILE (¬(((TA = TA_aux) ∧ ((WA = WA_aux) ∧ ((TB = TB_aux) ∧ (WB = WB_aux)))))) DO
  TA_aux := TA;
  WA_aux := WA;
  TB_aux := TB;
  WB_aux := WB;
  TA := (TA ∪ S);
  WA := (WA ∪ (π[0,1] (S) ∪ π[1,1] (WB_aux)));
  TB := (TB ∪ (π[0,1] (TA_aux) ∪ π[1,1] (TC)));
  WB := (WB ∪ WA_aux)
  END

POST
  (((TAQ = TA) ∧ (TBQ = TB)) ∧ ((WAQ = WA) ∧ (WBQ = WB)))
\end{Verbatim}

\Needspace{0.45\textheight}

\section{Example5020: RDFS closure by the rho-df rules versus the full W3C RDFS rule set: equal up to the reflexive subPropertyOf/subClassOf triples}\label{case:Example5020}

\begin{description}[leftmargin=2.6cm,style=nextline,itemsep=1pt]

\item[Category] specification (a standard or a paper's procedure versus a declarative program); family: two rule sets for one semantics.

\item[Kind] equivalence.

\item[Contributors] Leo Zhang.

\item[Status] certified valid (certificate checked outside the library build) (Certificate/Valid.lean, written by the certificate emitter); expected valid. Placement note: a proof module imported by \path{Certificate/Valid.lean} does not build under the library build's 4194304 KB limit; it builds under 12582912 KB (peak 1753 MB)

\item[Tags] \hyperref[tag:mutual-recursion]{\texttt{mutual-recursion}} \hyperref[tag:nonlinear-recursion]{\texttt{nonlinear-recursion}}.

\item[Input] \path{Benchmark/Example5020/Input.lean} (canonical notation).

\item[Summary] RDFS entailment computed by two rule sets: the minimal rho-df system of Munoz, Perez and Gutierrez against the full RDFS entailment patterns of the W3C, on the same input. The two closures agree on triples and on typing, and differ exactly by the reflexive subPropertyOf and subClassOf triples that RDFS derives for every property and class.

\item[Sides] naive compiler: rho-df system ⊢mrdf, rules (2a),(2b),(3a),(3b),(4a),(4b) = rdfs5, rdfs7, rdfs11, rdfs9, rdfs2, rdfs3 \par naive compiler: W3C RDFS entailment patterns on mrdf-graphs (the six shared patterns, rdfs6, rdfs10, rdfD2, and the typing forced by the RDFS axiomatic triples of the rho-df vocabulary, in IsProp/IsClass)

\item[Sources] {\raggedright \href{https://doi.org/10.1016/j.websem.2009.07.003}{S. Munoz, J. Perez, C. Gutierrez, Simple and Efficient Minimal RDFS, Journal of Web Semantics 7(3), 2009} (sections Definition 3 (rho-df vocabulary), Table 1 (rules (1)-(7)), Definition 15 (mrdf graphs, ⊢mrdf), Proposition 17, Corollary 18; read 2026-09-15) \par \href{https://www.w3.org/TR/rdf11-mt/}{W3C RDF 1.1 Semantics, Recommendation 25 February 2014} (sections 8.1.1 (rdfD2), 9 (RDFS axiomatic triples), 9.2.1 (rdfs1-rdfs13); read 2026-09-15) \par \href{https://github.com/apache/jena/blob/main/jena-core/src/main/resources/etc/rdfs.rules}{Apache Jena, jena-core/src/main/resources/etc/rdfs.rules} (commit f9eed353614c178b80b1e099ea2598e648014074 (main, 2026-09-14); read 2026-09-15)\par}

\item[Twin] Example5021 (invalid: reflexive subClassOf triples forgotten).

\item[Obstruction note] from the metadata: anchored (independent compiled loops merged side by side, one round of the same patterns per iteration; reflexive pairs lag one round: SubPropF = SubPropM ∪ π[0,0] IsProp\_aux, SubClassF = SubClassM ∪ π[0,0] IsClass\_aux)

\end{description}

\paragraph{Datalog program(s).}

{\raggedright\textit{prog\_pm --- rho-df, the system ⊢mrdf: rules (2a), (2b), (3a), (3b), (4a), (4b) of Munoz, Perez, Gutierrez (= rdfs5, rdfs7, rdfs11, rdfs9, rdfs2, rdfs3).} --- fidelity module (kernel-checked: the compiled sides of the command are this program's naive compilation).\par}

\begin{Verbatim}[fontsize=\scriptsize,breaklines=true,breakanywhere=true,frame=single,framesep=1.5mm]
SubPropM(x1, y1) :- SubProp(x1, y1).
SubPropM(x1, z1) :- SubPropM(x1, y1), SubPropM(y1, z1).
TripleM(x1, y1, z1) :- Triple(x1, y1, z1).
TripleM(x1, z1, y1) :- SubPropM(x2, z1), TripleM(x1, x2, y1).
SubClassM(x1, y1) :- SubClass(x1, y1).
SubClassM(x1, z1) :- SubClassM(x1, y1), SubClassM(y1, z1).
TypeM(x1, y1) :- TypeOf(x1, y1).
TypeM(x1, z1) :- SubClassM(y1, z1), TypeM(x1, y1).
TypeM(x1, z1) :- Dom(y1, z1), TripleM(x1, y1, x2).
TypeM(y1, z1) :- Range(x2, z1), TripleM(x1, x2, y1).
\end{Verbatim}

{\raggedright\textit{prog\_pf --- The W3C RDFS entailment patterns on mrdf-graphs: the six shared patterns, rdfs6, rdfs10, rdfD2 and the typing forced by the RDFS axiomatic triples of the rho-df vocabulary (IsProp = rdf:type rdf:Property, IsClass = rdf:type rdfs:Class).} --- fidelity module (kernel-checked: the compiled sides of the command are this program's naive compilation).\par}

\begin{Verbatim}[fontsize=\scriptsize,breaklines=true,breakanywhere=true,frame=single,framesep=1.5mm]
SubPropF(x1, y1) :- SubProp(x1, y1).
SubPropF(x1, z1) :- SubPropF(x1, y1), SubPropF(y1, z1).
SubPropF(x1, x1) :- IsProp(x1).
TripleF(x1, y1, z1) :- Triple(x1, y1, z1).
TripleF(x1, z1, y1) :- SubPropF(x2, z1), TripleF(x1, x2, y1).
SubClassF(x1, y1) :- SubClass(x1, y1).
SubClassF(x1, z1) :- SubClassF(x1, y1), SubClassF(y1, z1).
SubClassF(x1, x1) :- IsClass(x1).
TypeF(x1, y1) :- TypeOf(x1, y1).
TypeF(x1, z1) :- SubClassF(y1, z1), TypeF(x1, y1).
TypeF(x1, z1) :- Dom(y1, z1), TripleF(x1, y1, x2).
TypeF(y1, z1) :- Range(x2, z1), TripleF(x1, x2, y1).
IsProp(y1) :- TripleF(x1, y1, z1).
IsProp(x1) :- SubPropF(x1, y1).
IsProp(y1) :- SubPropF(x1, y1).
IsProp(x1) :- Dom(x1, y1).
IsProp(x1) :- Range(x1, y1).
IsClass(y1) :- TypeF(x1, y1).
IsClass(x1) :- SubClassF(x1, y1).
IsClass(y1) :- SubClassF(x1, y1).
IsClass(y1) :- Dom(x1, y1).
IsClass(y1) :- Range(x1, y1).
\end{Verbatim}

\paragraph{Reading.} Both sides compute RDFS entailment closures over ground triples using rho-df vocabulary relations (subclass, subproperty, domain, range, type, generic triples); side 1 implements the minimal rho-df rule set of Munoz, Perez and Gutierrez (subproperty/subclass transitivity plus type and triple propagation), while side 2 implements the full W3C RDFS rule set restricted to the same vocabulary, additionally deriving that every recognized property or class is reflexively a subproperty/subclass of itself. The triple states the two closures agree exactly on triples and typing, and that the full closure's subproperty/subclass relations equal the minimal closure's relations unioned with the extra reflexive pairs forced by which terms are recognized as properties or classes --- a precise account of what the dropped RDFS rules contribute.

\paragraph{First-order reading of PRE/POST, translated mechanically from the relational-algebra assertion in Input.lean.} The relational-algebra assertion in \texttt{Input.lean} remains the authority.

\textit{PRE:} ⊤\par True; no precondition constrains the inputs.

\textit{POST:} (∀x y z. (TripleM(x,y,z) ↔ TripleF(x,y,z))) ∧ (∀u v. (TypeM(u,v) ↔ TypeF(u,v))) ∧ (∀w x1. (SubPropM(w,x1) ∨ IsProp(w) ∧ w = x1 ↔ SubPropF(w,x1))) ∧ ∀x2 x3. (SubClassM(x2,x3) ∨ IsClass(x2) ∧ x2 = x3 ↔ SubClassF(x2,x3))\par The two rule sets derive identical Triple and Type relations, and the full rule set's subproperty/subclass relations equal the minimal rule set's subproperty/subclass relations unioned with the reflexive pairs over the terms recognized as properties/classes.

\paragraph{Whiel encoding (Hoare triple).}

\begin{Verbatim}[fontsize=\tiny,breaklines=true,breakanywhere=true,frame=single,framesep=1.5mm]
SCHEMA
  TypeOf/2, SubClass/2, SubProp/2, Dom/2, Range/2, TypeM/2, SubClassM/2, SubPropM/2, TypeM_aux/2, SubClassM_aux/2, SubPropM_aux/2, TypeF/2, SubClassF/2, SubPropF/2, TypeF_aux/2, SubClassF_aux/2, SubPropF_aux/2, Triple/3, TripleM/3, TripleM_aux/3, TripleF/3, TripleF_aux/3, IsProp/1, IsClass/1, IsProp_aux/1, IsClass_aux/1

PRE
  true 

CMD
  SubPropM_aux := ∅[2];
  TripleM_aux := ∅[3];
  SubClassM_aux := ∅[2];
  TypeM_aux := ∅[2];
  SubPropM := (SubProp ∪ π[0,3] (σ[#1 = #2] ((SubPropM_aux × SubPropM_aux))));
  TripleM := (Triple ∪ π[2,1,4] (σ[#0 = #3] ((SubPropM_aux × TripleM_aux))));
  SubClassM := (SubClass ∪ π[0,3] (σ[#1 = #2] ((SubClassM_aux × SubClassM_aux))));
  TypeM := (TypeOf ∪ (π[2,1] (σ[#0 = #3] ((SubClassM_aux × TypeM_aux))) ∪ (π[2,1] (σ[#0 = #3] ((Dom × TripleM_aux))) ∪ π[4,1] (σ[#0 = #3] ((Range × TripleM_aux))))));
  WHILE (¬(((SubPropM = SubPropM_aux) ∧ ((TripleM = TripleM_aux) ∧ ((SubClassM = SubClassM_aux) ∧ (TypeM = TypeM_aux)))))) DO
  SubPropM_aux := SubPropM;
  TripleM_aux := TripleM;
  SubClassM_aux := SubClassM;
  TypeM_aux := TypeM;
  SubPropM := (SubPropM ∪ (SubProp ∪ π[0,3] (σ[#1 = #2] ((SubPropM_aux × SubPropM_aux)))));
  TripleM := (TripleM ∪ (Triple ∪ π[2,1,4] (σ[#0 = #3] ((SubPropM_aux × TripleM_aux)))));
  SubClassM := (SubClassM ∪ (SubClass ∪ π[0,3] (σ[#1 = #2] ((SubClassM_aux × SubClassM_aux)))));
  TypeM := (TypeM ∪ (TypeOf ∪ (π[2,1] (σ[#0 = #3] ((SubClassM_aux × TypeM_aux))) ∪ (π[2,1] (σ[#0 = #3] ((Dom × TripleM_aux))) ∪ π[4,1] (σ[#0 = #3] ((Range × TripleM_aux)))))))
  END;
  SubPropF_aux := ∅[2];
  TripleF_aux := ∅[3];
  SubClassF_aux := ∅[2];
  TypeF_aux := ∅[2];
  IsProp_aux := ∅[1];
  IsClass_aux := ∅[1];
  SubPropF := (SubProp ∪ (π[0,3] (σ[#1 = #2] ((SubPropF_aux × SubPropF_aux))) ∪ π[0,0] (IsProp_aux)));
  TripleF := (Triple ∪ π[2,1,4] (σ[#0 = #3] ((SubPropF_aux × TripleF_aux))));
  SubClassF := (SubClass ∪ (π[0,3] (σ[#1 = #2] ((SubClassF_aux × SubClassF_aux))) ∪ π[0,0] (IsClass_aux)));
  TypeF := (TypeOf ∪ (π[2,1] (σ[#0 = #3] ((SubClassF_aux × TypeF_aux))) ∪ (π[2,1] (σ[#0 = #3] ((Dom × TripleF_aux))) ∪ π[4,1] (σ[#0 = #3] ((Range × TripleF_aux))))));
  IsProp := (π[1] (TripleF_aux) ∪ (π[0] (SubPropF_aux) ∪ (π[1] (SubPropF_aux) ∪ (π[0] (Dom) ∪ π[0] (Range)))));
  IsClass := (π[1] (TypeF_aux) ∪ (π[0] (SubClassF_aux) ∪ (π[1] (SubClassF_aux) ∪ (π[1] (Dom) ∪ π[1] (Range)))));
  WHILE (¬(((SubPropF = SubPropF_aux) ∧ ((TripleF = TripleF_aux) ∧ ((SubClassF = SubClassF_aux) ∧ ((TypeF = TypeF_aux) ∧ ((IsProp = IsProp_aux) ∧ (IsClass = IsClass_aux)))))))) DO
  SubPropF_aux := SubPropF;
  TripleF_aux := TripleF;
  SubClassF_aux := SubClassF;
  TypeF_aux := TypeF;
  IsProp_aux := IsProp;
  IsClass_aux := IsClass;
  SubPropF := (SubPropF ∪ (SubProp ∪ (π[0,3] (σ[#1 = #2] ((SubPropF_aux × SubPropF_aux))) ∪ π[0,0] (IsProp_aux))));
  TripleF := (TripleF ∪ (Triple ∪ π[2,1,4] (σ[#0 = #3] ((SubPropF_aux × TripleF_aux)))));
  SubClassF := (SubClassF ∪ (SubClass ∪ (π[0,3] (σ[#1 = #2] ((SubClassF_aux × SubClassF_aux))) ∪ π[0,0] (IsClass_aux))));
  TypeF := (TypeF ∪ (TypeOf ∪ (π[2,1] (σ[#0 = #3] ((SubClassF_aux × TypeF_aux))) ∪ (π[2,1] (σ[#0 = #3] ((Dom × TripleF_aux))) ∪ π[4,1] (σ[#0 = #3] ((Range × TripleF_aux)))))));
  IsProp := (IsProp ∪ (π[1] (TripleF_aux) ∪ (π[0] (SubPropF_aux) ∪ (π[1] (SubPropF_aux) ∪ (π[0] (Dom) ∪ π[0] (Range))))));
  IsClass := (IsClass ∪ (π[1] (TypeF_aux) ∪ (π[0] (SubClassF_aux) ∪ (π[1] (SubClassF_aux) ∪ (π[1] (Dom) ∪ π[1] (Range))))))
  END

POST
  (((TripleM = TripleF) ∧ (TypeM = TypeF)) ∧
    (((SubPropM ∪ π[0, 0] IsProp) = SubPropF) ∧ ((SubClassM ∪ π[0, 0] IsClass) = SubClassF)))
\end{Verbatim}

\Needspace{0.45\textheight}

\section{Example5021: rho-df versus full RDFS with the reflexive subClassOf triples forgotten (invalid twin of Example5020)}\label{case:Example5021}

\begin{description}[leftmargin=2.6cm,style=nextline,itemsep=1pt]

\item[Category] specification (a standard or a paper's procedure versus a declarative program); family: two rule sets for one semantics (invalid twin).

\item[Kind] equivalence.

\item[Contributors] Leo Zhang.

\item[Status] certified invalid (Certificate/Invalid.lean, written by the certificate emitter); expected invalid.

\item[Tags] \hyperref[tag:mutual-recursion]{\texttt{mutual-recursion}} \hyperref[tag:nonlinear-recursion]{\texttt{nonlinear-recursion}}.

\item[Input] \path{Benchmark/Example5021/Input.lean} (canonical notation).

\item[Summary] RDFS entailment by the rho-df rules against the full W3C rule set, with the reflexive subClassOf triples forgotten: the invalid twin of Example5020.

\item[Sides] naive compiler: rho-df system ⊢mrdf (as Example5020) \par naive compiler: W3C RDFS patterns on mrdf-graphs (as Example5020)

\item[Sources] {\raggedright \href{https://doi.org/10.1016/j.websem.2009.07.003}{S. Munoz, J. Perez, C. Gutierrez, Simple and Efficient Minimal RDFS, Journal of Web Semantics 7(3), 2009} (sections Definition 3 (rho-df vocabulary), Table 1 (rules (1)-(7)), Definition 15 (mrdf graphs, ⊢mrdf), Proposition 17, Corollary 18; read 2026-09-15) \par \href{https://www.w3.org/TR/rdf11-mt/}{W3C RDF 1.1 Semantics, Recommendation 25 February 2014} (sections 8.1.1 (rdfD2), 9 (RDFS axiomatic triples), 9.2.1 (rdfs1-rdfs13); read 2026-09-15) \par \href{https://github.com/apache/jena/blob/main/jena-core/src/main/resources/etc/rdfs.rules}{Apache Jena, jena-core/src/main/resources/etc/rdfs.rules} (commit f9eed353614c178b80b1e099ea2598e648014074 (main, 2026-09-14); read 2026-09-15)\par}

\item[Twin] Example5020 (valid).

\item[Obstruction note] from the metadata: n/a (invalid; kernel counterexample)

\end{description}

\paragraph{Datalog program(s).}

{\raggedright\textit{prog\_pm --- rho-df, the system ⊢mrdf: rules (2a), (2b), (3a), (3b), (4a), (4b) of Munoz, Perez, Gutierrez (= rdfs5, rdfs7, rdfs11, rdfs9, rdfs2, rdfs3).} --- fidelity module (kernel-checked: the compiled sides of the command are this program's naive compilation).\par}

\begin{Verbatim}[fontsize=\scriptsize,breaklines=true,breakanywhere=true,frame=single,framesep=1.5mm]
SubPropM(x1, y1) :- SubProp(x1, y1).
SubPropM(x1, z1) :- SubPropM(x1, y1), SubPropM(y1, z1).
TripleM(x1, y1, z1) :- Triple(x1, y1, z1).
TripleM(x1, z1, y1) :- SubPropM(x2, z1), TripleM(x1, x2, y1).
SubClassM(x1, y1) :- SubClass(x1, y1).
SubClassM(x1, z1) :- SubClassM(x1, y1), SubClassM(y1, z1).
TypeM(x1, y1) :- TypeOf(x1, y1).
TypeM(x1, z1) :- SubClassM(y1, z1), TypeM(x1, y1).
TypeM(x1, z1) :- Dom(y1, z1), TripleM(x1, y1, x2).
TypeM(y1, z1) :- Range(x2, z1), TripleM(x1, x2, y1).
\end{Verbatim}

{\raggedright\textit{prog\_pf --- The W3C RDFS entailment patterns on mrdf-graphs: the six shared patterns, rdfs6, rdfs10, rdfD2 and the typing forced by the RDFS axiomatic triples of the rho-df vocabulary (IsProp = rdf:type rdf:Property, IsClass = rdf:type rdfs:Class).} --- fidelity module (kernel-checked: the compiled sides of the command are this program's naive compilation).\par}

\begin{Verbatim}[fontsize=\scriptsize,breaklines=true,breakanywhere=true,frame=single,framesep=1.5mm]
SubPropF(x1, y1) :- SubProp(x1, y1).
SubPropF(x1, z1) :- SubPropF(x1, y1), SubPropF(y1, z1).
SubPropF(x1, x1) :- IsProp(x1).
TripleF(x1, y1, z1) :- Triple(x1, y1, z1).
TripleF(x1, z1, y1) :- SubPropF(x2, z1), TripleF(x1, x2, y1).
SubClassF(x1, y1) :- SubClass(x1, y1).
SubClassF(x1, z1) :- SubClassF(x1, y1), SubClassF(y1, z1).
SubClassF(x1, x1) :- IsClass(x1).
TypeF(x1, y1) :- TypeOf(x1, y1).
TypeF(x1, z1) :- SubClassF(y1, z1), TypeF(x1, y1).
TypeF(x1, z1) :- Dom(y1, z1), TripleF(x1, y1, x2).
TypeF(y1, z1) :- Range(x2, z1), TripleF(x1, x2, y1).
IsProp(y1) :- TripleF(x1, y1, z1).
IsProp(x1) :- SubPropF(x1, y1).
IsProp(y1) :- SubPropF(x1, y1).
IsProp(x1) :- Dom(x1, y1).
IsProp(x1) :- Range(x1, y1).
IsClass(y1) :- TypeF(x1, y1).
IsClass(x1) :- SubClassF(x1, y1).
IsClass(y1) :- SubClassF(x1, y1).
IsClass(y1) :- Dom(x1, y1).
IsClass(y1) :- Range(x1, y1).
\end{Verbatim}

\paragraph{Reading.} Same rho-df versus full-RDFS comparison as Example5020, but the claimed postcondition omits the reflexive-subclass correction: it accounts for the reflexive subproperty pairs (rdfs6) but asserts the subclass relations equal outright, with no correction for the reflexive subclass pairs (rdfs10). The expected verdict is invalid: rdfs10 makes every recognized class a subclass of itself under the full RDFS rules, a pair the minimal rho-df closure never derives, so the uncorrected equality fails whenever some class is recognized in the data.

\paragraph{First-order reading of PRE/POST, translated mechanically from the relational-algebra assertion in Input.lean.} The relational-algebra assertion in \texttt{Input.lean} remains the authority.

\textit{PRE:} ⊤\par True; no precondition constrains the inputs.

\textit{POST:} (∀x y z. (TripleM(x,y,z) ↔ TripleF(x,y,z))) ∧ (∀u v. (TypeM(u,v) ↔ TypeF(u,v))) ∧ (∀w x1. (SubPropM(w,x1) ∨ IsProp(w) ∧ w = x1 ↔ SubPropF(w,x1))) ∧ ∀x2 x3. (SubClassM(x2,x3) ↔ SubClassF(x2,x3))\par The claim (which fails) is that Triple and Type agree between the two rule sets, the subproperty relations agree once the reflexive property pairs are added back, but the subclass relations are asserted equal outright, with no reflexive-subclass correction.

\paragraph{Whiel encoding (Hoare triple).}

\begin{Verbatim}[fontsize=\tiny,breaklines=true,breakanywhere=true,frame=single,framesep=1.5mm]
SCHEMA
  TypeOf/2, SubClass/2, SubProp/2, Dom/2, Range/2, TypeM/2, SubClassM/2, SubPropM/2, TypeM_aux/2, SubClassM_aux/2, SubPropM_aux/2, TypeF/2, SubClassF/2, SubPropF/2, TypeF_aux/2, SubClassF_aux/2, SubPropF_aux/2, Triple/3, TripleM/3, TripleM_aux/3, TripleF/3, TripleF_aux/3, IsProp/1, IsClass/1, IsProp_aux/1, IsClass_aux/1

PRE
  true 

CMD
  SubPropM_aux := ∅[2];
  TripleM_aux := ∅[3];
  SubClassM_aux := ∅[2];
  TypeM_aux := ∅[2];
  SubPropM := (SubProp ∪ π[0,3] (σ[#1 = #2] ((SubPropM_aux × SubPropM_aux))));
  TripleM := (Triple ∪ π[2,1,4] (σ[#0 = #3] ((SubPropM_aux × TripleM_aux))));
  SubClassM := (SubClass ∪ π[0,3] (σ[#1 = #2] ((SubClassM_aux × SubClassM_aux))));
  TypeM := (TypeOf ∪ (π[2,1] (σ[#0 = #3] ((SubClassM_aux × TypeM_aux))) ∪ (π[2,1] (σ[#0 = #3] ((Dom × TripleM_aux))) ∪ π[4,1] (σ[#0 = #3] ((Range × TripleM_aux))))));
  WHILE (¬(((SubPropM = SubPropM_aux) ∧ ((TripleM = TripleM_aux) ∧ ((SubClassM = SubClassM_aux) ∧ (TypeM = TypeM_aux)))))) DO
  SubPropM_aux := SubPropM;
  TripleM_aux := TripleM;
  SubClassM_aux := SubClassM;
  TypeM_aux := TypeM;
  SubPropM := (SubPropM ∪ (SubProp ∪ π[0,3] (σ[#1 = #2] ((SubPropM_aux × SubPropM_aux)))));
  TripleM := (TripleM ∪ (Triple ∪ π[2,1,4] (σ[#0 = #3] ((SubPropM_aux × TripleM_aux)))));
  SubClassM := (SubClassM ∪ (SubClass ∪ π[0,3] (σ[#1 = #2] ((SubClassM_aux × SubClassM_aux)))));
  TypeM := (TypeM ∪ (TypeOf ∪ (π[2,1] (σ[#0 = #3] ((SubClassM_aux × TypeM_aux))) ∪ (π[2,1] (σ[#0 = #3] ((Dom × TripleM_aux))) ∪ π[4,1] (σ[#0 = #3] ((Range × TripleM_aux)))))))
  END;
  SubPropF_aux := ∅[2];
  TripleF_aux := ∅[3];
  SubClassF_aux := ∅[2];
  TypeF_aux := ∅[2];
  IsProp_aux := ∅[1];
  IsClass_aux := ∅[1];
  SubPropF := (SubProp ∪ (π[0,3] (σ[#1 = #2] ((SubPropF_aux × SubPropF_aux))) ∪ π[0,0] (IsProp_aux)));
  TripleF := (Triple ∪ π[2,1,4] (σ[#0 = #3] ((SubPropF_aux × TripleF_aux))));
  SubClassF := (SubClass ∪ (π[0,3] (σ[#1 = #2] ((SubClassF_aux × SubClassF_aux))) ∪ π[0,0] (IsClass_aux)));
  TypeF := (TypeOf ∪ (π[2,1] (σ[#0 = #3] ((SubClassF_aux × TypeF_aux))) ∪ (π[2,1] (σ[#0 = #3] ((Dom × TripleF_aux))) ∪ π[4,1] (σ[#0 = #3] ((Range × TripleF_aux))))));
  IsProp := (π[1] (TripleF_aux) ∪ (π[0] (SubPropF_aux) ∪ (π[1] (SubPropF_aux) ∪ (π[0] (Dom) ∪ π[0] (Range)))));
  IsClass := (π[1] (TypeF_aux) ∪ (π[0] (SubClassF_aux) ∪ (π[1] (SubClassF_aux) ∪ (π[1] (Dom) ∪ π[1] (Range)))));
  WHILE (¬(((SubPropF = SubPropF_aux) ∧ ((TripleF = TripleF_aux) ∧ ((SubClassF = SubClassF_aux) ∧ ((TypeF = TypeF_aux) ∧ ((IsProp = IsProp_aux) ∧ (IsClass = IsClass_aux)))))))) DO
  SubPropF_aux := SubPropF;
  TripleF_aux := TripleF;
  SubClassF_aux := SubClassF;
  TypeF_aux := TypeF;
  IsProp_aux := IsProp;
  IsClass_aux := IsClass;
  SubPropF := (SubPropF ∪ (SubProp ∪ (π[0,3] (σ[#1 = #2] ((SubPropF_aux × SubPropF_aux))) ∪ π[0,0] (IsProp_aux))));
  TripleF := (TripleF ∪ (Triple ∪ π[2,1,4] (σ[#0 = #3] ((SubPropF_aux × TripleF_aux)))));
  SubClassF := (SubClassF ∪ (SubClass ∪ (π[0,3] (σ[#1 = #2] ((SubClassF_aux × SubClassF_aux))) ∪ π[0,0] (IsClass_aux))));
  TypeF := (TypeF ∪ (TypeOf ∪ (π[2,1] (σ[#0 = #3] ((SubClassF_aux × TypeF_aux))) ∪ (π[2,1] (σ[#0 = #3] ((Dom × TripleF_aux))) ∪ π[4,1] (σ[#0 = #3] ((Range × TripleF_aux)))))));
  IsProp := (IsProp ∪ (π[1] (TripleF_aux) ∪ (π[0] (SubPropF_aux) ∪ (π[1] (SubPropF_aux) ∪ (π[0] (Dom) ∪ π[0] (Range))))));
  IsClass := (IsClass ∪ (π[1] (TypeF_aux) ∪ (π[0] (SubClassF_aux) ∪ (π[1] (SubClassF_aux) ∪ (π[1] (Dom) ∪ π[1] (Range))))))
  END

POST
  (((TripleM = TripleF) ∧ (TypeM = TypeF)) ∧
    (((SubPropM ∪ π[0, 0] IsProp) = SubPropF) ∧ (SubClassM = SubClassF)))
\end{Verbatim}

\Needspace{0.45\textheight}

\section{Example5022: SQL recursion limits (MySQL cte\_max\_recursion\_depth, SQL Server MAXRECURSION): completion within the limit implies the closure}\label{case:Example5022}

\begin{description}[leftmargin=2.6cm,style=nextline,itemsep=1pt]

\item[Category] specification (a standard or a paper's procedure versus a declarative program); family: engine semantics versus declarative program (conditional claim).

\item[Kind] implication (disjunction in the postcondition).

\item[Contributors] Leo Zhang.

\item[Status] certified valid (Certificate/Valid.lean, written by the certificate emitter); expected valid.

\item[Tags] \hyperref[tag:database-engine]{\texttt{database-engine}} \hyperref[tag:prophecy]{\texttt{prophecy}} \hyperref[tag:recursion-limit]{\texttt{recursion-limit}}.

\item[Input] \path{Benchmark/Example5022/Input.lean} (canonical notation).

\item[Summary] Recursion limits of SQL engines: a recursive CTE run under MySQL's cte\_max\_recursion\_depth or SQL Server's MAXRECURSION against the Datalog closure it is meant to compute. Claim: when the engine does not abort, its result is the closure.

\item[Sides] hand-written, loop-free: four semi-naive iterations of the recursive CTE under limit k = 4 (ResA = anchor; NewB..NewE new rows per iteration; ResB..ResE cumulative), MySQL UNION DISTINCT semantics; SQL Server's UNION ALL termination implies the modelled condition \par naive compiler: Cl = closure (the CTE's Datalog meaning)

\item[Sources] {\raggedright \href{https://dev.mysql.com/doc/refman/8.4/en/with.html}{MySQL 8.4 Reference Manual, WITH (Common Table Expressions): Recursive Common Table Expressions; Limiting Common Table Expression Recursion} (read 2026-09-15) \par \href{https://github.com/mysql/mysql-server/blob/8.4/share/messages_to_clients.txt}{MySQL server source, share/messages\_to\_clients.txt (ER\_CTE\_MAX\_RECURSION\_DEPTH = 3636)} (commit 99960bf74fa919347e4f4e3ca47672f333d6e91f (branch 8.4); read 2026-09-15) \par \href{https://learn.microsoft.com/en-us/sql/t-sql/queries/with-common-table-expression-transact-sql}{Microsoft SQL Server, WITH common\_table\_expression (Transact-SQL)} (read 2026-09-15) \par \href{https://learn.microsoft.com/en-us/sql/t-sql/queries/hints-transact-sql-query}{Microsoft SQL Server, Hints (Transact-SQL) - Query (MAXRECURSION)} (read 2026-09-15) \par \href{https://learn.microsoft.com/en-us/sql/relational-databases/errors-events/database-engine-events-and-errors-0-to-999}{Microsoft SQL Server, Database Engine events and errors (0 to 999) (error 530)} (read 2026-09-15)\par}

\item[Prophecy] accepted invariant with a prophecy clause above level zero (Core.json); no certificate built from it yet.

\end{description}

\paragraph{Datalog program(s).}

{\raggedright\textit{prog\_p --- The recursive CTE as Datalog.} --- fidelity module (kernel-checked: the compiled sides of the command are this program's naive compilation).\par}

\begin{Verbatim}[fontsize=\scriptsize,breaklines=true,breakanywhere=true,frame=single,framesep=1.5mm]
Cl(x1, y1) :- Base(x1, y1).
Cl(x1, z1) :- Cl(x1, y1), E(y1, z1).
\end{Verbatim}

\paragraph{Reading.} Side 1 encodes, loop-free, four semi-naive iterations of a recursive common table expression under a fixed recursion-depth limit (k = 4), as implemented by MySQL's cte\_max\_recursion\_depth or SQL Server's MAXRECURSION, tracking new rows and cumulative results per iteration; side 2 is the naive-compiler translation of the closure the CTE is meant to compute. The triple states a conditional correctness property for bounded-recursion SQL engines: if the engine's iteration limit is not exceeded (the final iteration produces no new rows), the truncated result actually equals the full transitive closure --- a limited engine only ever silently under-approximates by aborting, never by returning a wrong-but-apparently-complete answer.

\paragraph{First-order reading of PRE/POST, translated mechanically from the relational-algebra assertion in Input.lean.} The relational-algebra assertion in \texttt{Input.lean} remains the authority.

\textit{PRE:} ⊤\par True; no precondition constrains the inputs.

\textit{POST:} (∃x y. NewE(x,y)) ∨ ∀z u. (Cl(z,u) ↔ ResE(z,u))\par Either the final iteration still produced new rows (the recursion-depth limit was exhausted), or the limited engine's result equals the true closure.

\paragraph{Whiel encoding (Hoare triple).}

\begin{Verbatim}[fontsize=\scriptsize,breaklines=true,breakanywhere=true,frame=single,framesep=1.5mm]
SCHEMA
  Base/2, E/2, ResA/2, NewB/2, ResB/2, NewC/2, ResC/2, NewD/2, ResD/2, NewE/2, ResE/2, Cl/2, Cl_aux/2

PRE
  true 

CMD
  ResA := Base;
  NewB := (π[0, 3] (σ[#1 = #2] (ResA × E)) ∖ ResA);
  ResB := (ResA ∪ NewB);
  NewC := (π[0, 3] (σ[#1 = #2] (NewB × E)) ∖ ResB);
  ResC := (ResB ∪ NewC);
  NewD := (π[0, 3] (σ[#1 = #2] (NewC × E)) ∖ ResC);
  ResD := (ResC ∪ NewD);
  NewE := (π[0, 3] (σ[#1 = #2] (NewD × E)) ∖ ResD);
  ResE := (ResD ∪ NewE);
  Cl_aux := ∅[2];
  Cl := (Base ∪ π[0,3] (σ[#1 = #2] ((Cl_aux × E))));
  WHILE (¬((Cl = Cl_aux))) DO
  Cl_aux := Cl;
  Cl := (Cl ∪ (Base ∪ π[0,3] (σ[#1 = #2] ((Cl_aux × E)))))
  END

POST
  ((NewE ≠ ∅) ∨ (Cl = ResE))
\end{Verbatim}

\Needspace{0.45\textheight}

\section{Example5023: Win-move: the alternating fixpoint of the well-founded semantics (nested loops) versus the attractor computation}\label{case:Example5023}

\begin{description}[leftmargin=2.6cm,style=nextline,itemsep=1pt]

\item[Category] specification (a standard or a paper's procedure versus a declarative program); family: two computations of one semantics.

\item[Kind] equivalence.

\item[Contributors] Leo Zhang.

\item[Status] certified valid (certificate checked outside the library build) (Certificate/Valid.lean, written by the certificate emitter); expected valid. Placement note: a proof module imported by \path{Certificate/Valid.lean} does not build under the library build's 4194304 KB limit; it builds under 12582912 KB (peak 1672 MB)

\item[Tags] \hyperref[tag:evaluation-strategies]{\texttt{evaluation-strategies}} \hyperref[tag:greatest-fixpoint]{\texttt{greatest-fixpoint}} \hyperref[tag:negation]{\texttt{negation}} \hyperref[tag:prophecy]{\texttt{prophecy}}.

\item[Input] \path{Benchmark/Example5023/Input.lean} (canonical notation).

\item[Summary] Win-move under the well-founded semantics: the alternating fixpoint (Van Gelder) computed with nested loops against the attractor computation of Example0162 and legacy Example2019 (positions with a move to a losing position win; positions all of whose moves lead to winning positions lose).

\item[Sides] hand-written: alternating fixpoint A\_0 = ∅, A\_\{i+1\} = S\_P(A\_i) with Under (even) / Over (odd), each S\_P an inner least-fixpoint loop in the naive compiler's form for OverI(x) :- Move(x,y), NotUnder(y) and UnderI(x) :- Move(x,y), NotOver(y) \par hand-written: attractor (Example0162 / legacy 2019 style, simultaneous update of Win and Lose)

\item[Sources] {\raggedright \href{https://doi.org/10.1145/116825.116838}{A. Van Gelder, K. A. Ross, J. S. Schlipf, The Well-Founded Semantics for General Logic Programs, JACM 38(3), 1991} (read not retrievable 2026-09-15 (ACM DL 403 / bot check)) \par \href{https://doi.org/10.1145/73721.73722}{A. Van Gelder, The Alternating Fixpoint of Logic Programs with Negation, PODS 1989} (read not retrievable 2026-09-15 (ACM DL 403 / bot check)) \par \href{https://doi.org/10.1016/0022-0000(93)90024-Q}{A. Van Gelder, The Alternating Fixpoint of Logic Programs with Negation, JCSS 47(1), 1993} (read not retrievable 2026-09-15 (ScienceDirect 403, also in the in-app browser)) \par \href{http://webdam.inria.fr/Alice/pdfs/Chapter-15.pdf}{S. Abiteboul, R. Hull, V. Vianu, Foundations of Databases, chapter 15 (Example 15.3.1; 15.3 A Fixpoint Definition; Example 15.3.8(b); Theorem 15.3.9)} (read 2026-09-15) \par \href{local}{Example0162 and the retired corpus case Example2019 (attractor side)} (read 2026-09-15) \par program source file win(X) :- move(X,Y), not win(Y)\par}

\item[Twin] Example5024 (invalid: Over = Win, i.e. no draws).

\item[Prophecy] accepted invariant with a prophecy clause above level zero (Core.json); no certificate built from it yet.

\end{description}

\paragraph{Datalog program(s).}

{\raggedright\textit{prog\_pO --- The positive programs P\textasciicircum{}I of the two half-steps of the alternation (`not win(Y)` read as `NotUnder(Y)`, resp. `NotOver(Y)`); their naive compilation is the inner loop of each half-step.} --- fidelity module (kernel-checked: the compiled sides of the command are this program's naive compilation).\par}

\begin{Verbatim}[fontsize=\scriptsize,breaklines=true,breakanywhere=true,frame=single,framesep=1.5mm]
OverI(x1) :- Move(x1, y1), NotUnder(y1).
\end{Verbatim}

{\raggedright\textit{prog\_pU} --- fidelity module (kernel-checked: the compiled sides of the command are this program's naive compilation).\par}

\begin{Verbatim}[fontsize=\scriptsize,breaklines=true,breakanywhere=true,frame=single,framesep=1.5mm]
UnderI(x1) :- Move(x1, y1), NotOver(y1).
\end{Verbatim}

{\raggedright\textit{rules as stated in the header comment of Input.lean} --- header comment.\par}

\begin{Verbatim}[fontsize=\scriptsize,breaklines=true,breakanywhere=true,frame=single,framesep=1.5mm]
win(X) :- move(X, Y), not win(Y)
OverI(x) :- Move(x, y), NotUnder(y)
UnderI(x) :- Move(x, y), NotOver(y)
\end{Verbatim}

\paragraph{Reading.} Both sides compute the well-founded semantics of the classic win-move game program (a position wins if it has a move to a position that is not a win), whose model is three-valued (winning/losing/drawn); side 1 computes Van Gelder's alternating fixpoint directly, alternating an outer loop between an even sequence (Under) and odd sequence (Over), each obtained via an inner positive-program least-fixpoint loop --- the first nested-loop case in this corpus --- while side 2 computes the same sets via the game-theoretic attractor construction (simultaneously updating Win/Lose sets by one half-step each iteration), as in Example0162. The triple states the alternating fixpoint's even limit and the complement of its odd limit coincide exactly with the attractor's Win and Lose sets, the standard correspondence between well-founded semantics and attractor-based game solving.

\paragraph{First-order reading of PRE/POST, translated mechanically from the relational-algebra assertion in Input.lean.} The relational-algebra assertion in \texttt{Input.lean} remains the authority.

\textit{PRE:} ⊤\par True; no precondition constrains the inputs.

\textit{POST:} (∀x. (Under(x) ↔ Win(x))) ∧ ∀y. (Nodes(y) ∧ ¬Over(y) ↔ Lose(y))\par The alternating fixpoint's even limit Under equals the attractor's Win set, and the complement of the odd limit Over (relative to all positions) equals the attractor's Lose set.

\paragraph{Whiel encoding (Hoare triple).}

\begin{Verbatim}[fontsize=\tiny,breaklines=true,breakanywhere=true,frame=single,framesep=1.5mm]
SCHEMA
  Move/2, Nodes/1, Under/1, Over/1, Prev/1, NotUnder/1, NotOver/1, OverI/1, OverI_aux/1, UnderI/1, UnderI_aux/1, Win/1, Lose/1, WinP/1, LoseP/1

PRE
  true 

CMD
  Nodes := (π[0] Move ∪ π[1] Move);
  Under := ∅[1];
  Over := ∅[1];
  Prev := Nodes;
  WHILE (¬((Under = Prev))) DO
    Prev := Under;
    NotUnder := (Nodes ∖ Under);
    OverI_aux := ∅[1];
    OverI := π[0] (σ[#1 = #2] ((Move × NotUnder)));
    WHILE (¬((OverI = OverI_aux))) DO
      OverI_aux := OverI;
      OverI := (OverI ∪ π[0] (σ[#1 = #2] ((Move × NotUnder))))
    END;
    Over := OverI;
    NotOver := (Nodes ∖ Over);
    UnderI_aux := ∅[1];
    UnderI := π[0] (σ[#1 = #2] ((Move × NotOver)));
    WHILE (¬((UnderI = UnderI_aux))) DO
      UnderI_aux := UnderI;
      UnderI := (UnderI ∪ π[0] (σ[#1 = #2] ((Move × NotOver))))
    END;
    Under := UnderI
  END;
  WinP := ∅[1];
  LoseP := ∅[1];
  Win := π[0] (σ[#1 = #2] (Move × LoseP));
  Lose := (Nodes ∖ π[0] (Move ∖ π[0, 1] (σ[#1 = #2] (Move × WinP))));
  WHILE (¬((Win = WinP) ∧ (Lose = LoseP))) DO
    WinP := Win;
    LoseP := Lose;
    Win := π[0] (σ[#1 = #2] (Move × LoseP));
    Lose := (Nodes ∖ π[0] (Move ∖ π[0, 1] (σ[#1 = #2] (Move × WinP))))
  END

POST
  ((Under = Win) ∧ ((Nodes ∖ Over) = Lose))
\end{Verbatim}

\Needspace{0.45\textheight}

\section{Example5024: Win-move: alternating fixpoint versus attractor, claiming the odd limit is the win set (invalid twin of Example5023)}\label{case:Example5024}

\begin{description}[leftmargin=2.6cm,style=nextline,itemsep=1pt]

\item[Category] specification (a standard or a paper's procedure versus a declarative program); family: two computations of one semantics (invalid twin).

\item[Kind] equivalence.

\item[Contributors] Leo Zhang.

\item[Status] certified invalid (Certificate/Invalid.lean, written by the certificate emitter); expected invalid.

\item[Tags] \hyperref[tag:evaluation-strategies]{\texttt{evaluation-strategies}} \hyperref[tag:greatest-fixpoint]{\texttt{greatest-fixpoint}} \hyperref[tag:negation]{\texttt{negation}}.

\item[Input] \path{Benchmark/Example5024/Input.lean} (canonical notation).

\item[Summary] Win-move under the well-founded semantics, the invalid twin of Example5023: the alternating fixpoint against the attractor, claiming in addition that the odd limit (the positions that are not false) is the attractor's win set, that is, that the well-founded model of win-move is total.

\item[Sides] hand-written: alternating fixpoint as in Example5023 \par hand-written: attractor as in Example5023

\item[Sources] {\raggedright \href{https://doi.org/10.1145/116825.116838}{A. Van Gelder, K. A. Ross, J. S. Schlipf, The Well-Founded Semantics for General Logic Programs, JACM 38(3), 1991} (read not retrievable 2026-09-15 (ACM DL 403 / bot check)) \par \href{https://doi.org/10.1145/73721.73722}{A. Van Gelder, The Alternating Fixpoint of Logic Programs with Negation, PODS 1989} (read not retrievable 2026-09-15 (ACM DL 403 / bot check)) \par \href{https://doi.org/10.1016/0022-0000(93)90024-Q}{A. Van Gelder, The Alternating Fixpoint of Logic Programs with Negation, JCSS 47(1), 1993} (read not retrievable 2026-09-15 (ScienceDirect 403, also in the in-app browser)) \par \href{http://webdam.inria.fr/Alice/pdfs/Chapter-15.pdf}{S. Abiteboul, R. Hull, V. Vianu, Foundations of Databases, chapter 15 (Example 15.3.1; 15.3 A Fixpoint Definition; Example 15.3.8(b); Theorem 15.3.9)} (read 2026-09-15) \par \href{local}{Example0162 and the retired corpus case Example2019 (attractor side)} (read 2026-09-15) \par program source file win(X) :- move(X,Y), not win(Y)\par}

\item[Twin] Example5023 (valid).

\item[Obstruction note] from the metadata: n/a (invalid; kernel counterexample)

\end{description}

\paragraph{Datalog program(s).}

{\raggedright\textit{prog\_pO --- The positive programs P\textasciicircum{}I of the two half-steps of the alternation (`not win(Y)` read as `NotUnder(Y)`, resp. `NotOver(Y)`); their naive compilation is the inner loop of each half-step.} --- fidelity module (kernel-checked: the compiled sides of the command are this program's naive compilation).\par}

\begin{Verbatim}[fontsize=\scriptsize,breaklines=true,breakanywhere=true,frame=single,framesep=1.5mm]
OverI(x1) :- Move(x1, y1), NotUnder(y1).
\end{Verbatim}

{\raggedright\textit{prog\_pU} --- fidelity module (kernel-checked: the compiled sides of the command are this program's naive compilation).\par}

\begin{Verbatim}[fontsize=\scriptsize,breaklines=true,breakanywhere=true,frame=single,framesep=1.5mm]
UnderI(x1) :- Move(x1, y1), NotOver(y1).
\end{Verbatim}

{\raggedright\textit{rules as stated in the header comment of Input.lean} --- header comment.\par}

\begin{Verbatim}[fontsize=\scriptsize,breaklines=true,breakanywhere=true,frame=single,framesep=1.5mm]
win(X) :- move(X, Y), not win(Y)
OverI(x) :- Move(x, y), NotUnder(y)
UnderI(x) :- Move(x, y), NotOver(y)
\end{Verbatim}

\paragraph{Reading.} Same two computations as Example5023 (Van Gelder's alternating fixpoint versus the game attractor for well-founded win-move), but the postcondition additionally claims Over equals Win directly --- that the odd limit of the alternation coincides with the win set, equivalently that the well-founded model has no undefined (drawn) positions. The expected verdict is invalid: on a two-cycle where neither player can force a win, both the true win set and the attractor's Win set are empty, but Over stabilizes at both positions (each has a move to a position not yet known to win), so those positions are genuinely undefined rather than losing, and Over does not equal Win.

\paragraph{First-order reading of PRE/POST, translated mechanically from the relational-algebra assertion in Input.lean.} The relational-algebra assertion in \texttt{Input.lean} remains the authority.

\textit{PRE:} ⊤\par True; no precondition constrains the inputs.

\textit{POST:} (∀x. (Under(x) ↔ Win(x))) ∧ ∀y. (Over(y) ↔ Win(y))\par The claim (which fails) is that both the even limit Under and the odd limit Over of the alternating fixpoint equal the attractor's Win set.

\paragraph{Whiel encoding (Hoare triple).}

\begin{Verbatim}[fontsize=\tiny,breaklines=true,breakanywhere=true,frame=single,framesep=1.5mm]
SCHEMA
  Move/2, Nodes/1, Under/1, Over/1, Prev/1, NotUnder/1, NotOver/1, OverI/1, OverI_aux/1, UnderI/1, UnderI_aux/1, Win/1, Lose/1, WinP/1, LoseP/1

PRE
  true 

CMD
  Nodes := (π[0] Move ∪ π[1] Move);
  Under := ∅[1];
  Over := ∅[1];
  Prev := Nodes;
  WHILE (¬((Under = Prev))) DO
    Prev := Under;
    NotUnder := (Nodes ∖ Under);
    OverI_aux := ∅[1];
    OverI := π[0] (σ[#1 = #2] ((Move × NotUnder)));
    WHILE (¬((OverI = OverI_aux))) DO
      OverI_aux := OverI;
      OverI := (OverI ∪ π[0] (σ[#1 = #2] ((Move × NotUnder))))
    END;
    Over := OverI;
    NotOver := (Nodes ∖ Over);
    UnderI_aux := ∅[1];
    UnderI := π[0] (σ[#1 = #2] ((Move × NotOver)));
    WHILE (¬((UnderI = UnderI_aux))) DO
      UnderI_aux := UnderI;
      UnderI := (UnderI ∪ π[0] (σ[#1 = #2] ((Move × NotOver))))
    END;
    Under := UnderI
  END;
  WinP := ∅[1];
  LoseP := ∅[1];
  Win := π[0] (σ[#1 = #2] (Move × LoseP));
  Lose := (Nodes ∖ π[0] (Move ∖ π[0, 1] (σ[#1 = #2] (Move × WinP))));
  WHILE (¬((Win = WinP) ∧ (Lose = LoseP))) DO
    WinP := Win;
    LoseP := Lose;
    Win := π[0] (σ[#1 = #2] (Move × LoseP));
    Lose := (Nodes ∖ π[0] (Move ∖ π[0, 1] (σ[#1 = #2] (Move × WinP))))
  END

POST
  ((Under = Win) ∧ (Over = Win))
\end{Verbatim}

\Needspace{0.45\textheight}

\section{Example5030: NetworkX per-source transitive closure versus PostgreSQL recursive-union rounds}\label{case:Example5030}

\begin{description}[leftmargin=2.6cm,style=nextline,itemsep=1pt]

\item[Category] production pair (two algorithms from production sources or legacy production cases); family: production algorithm pair (rate mismatch).

\item[Kind] equivalence.

\item[Contributors] Leo Zhang.

\item[Status] certified valid (certificate checked outside the library build) (Certificate/Valid.lean, written by the certificate emitter); expected valid. Placement note: a proof module imported by \path{Certificate/Valid.lean} does not build under the library build's 4194304 KB limit; it builds under 12582912 KB (peak 1679 MB)

\item[Tags] \hyperref[tag:database-engine]{\texttt{database-engine}} \hyperref[tag:evaluation-strategies]{\texttt{evaluation-strategies}} \hyperref[tag:prophecy]{\texttt{prophecy}} \hyperref[tag:semi-naive]{\texttt{semi-naive}}.

\item[Input] \path{Benchmark/Example5030/Input.lean} (canonical notation).

\item[Summary] NetworkX's per-source transitive closure against PostgreSQL's recursive-union rounds.

\item[Sides] hand-written: NetworkX transitive\_closure(G, reflexive=False) = for each node in Lt order, the edge BFS of edge\_bfs set-at-a-time (nested loop) \par hand-written: PostgreSQL ExecRecursiveUnion working/intermediate-table rounds (legacy Example2037) seeded with Edge

\item[Sources] {\raggedright \href{https://github.com/networkx/networkx/blob/main/networkx/algorithms/dag.py}{NetworkX networkx/networkx networkx/algorithms/dag.py transitive\_closure (reflexive=False branch: TC.add\_edges\_from((v, e[1]) for e in nx.edge\_bfs(G, v) if e[1] not in TC[v])); descendants} (branch main; commit 4e74880b0da0 (2026-09-12); read 2026-09-15) \par \href{https://github.com/networkx/networkx/blob/main/networkx/algorithms/traversal/edgebfs.py}{NetworkX networkx/networkx networkx/algorithms/traversal/edgebfs.py edge\_bfs (visited\_nodes, visited\_edges, deque of (node, edges\_from(node)))} (branch main; commit 4e74880b0da0 (2026-09-12); read 2026-09-15) \par \href{https://github.com/postgres/postgres/blob/master/src/backend/executor/nodeRecursiveunion.c}{PostgreSQL postgres/postgres src/backend/executor/nodeRecursiveunion.c ExecRecursiveUnion (header comment steps 2.1-2.6; 'Ignore tuple if already seen'; working/intermediate table swap)} (branch master; commit 6e58d6356fcb (2026-09-15); read 2026-09-15; encoded as legacy Example2037 / Example5001 hand side with the base term equal to the edge relation)\par}

\item[Twin] Example5031 (invalid: per-source search with cutoff 2).

\item[Prophecy] accepted invariant with a prophecy clause above level zero (Core.json); no certificate built from it yet.

\end{description}

\paragraph{Datalog program(s).}

{\raggedright\textit{prog\_T --- The Datalog program(s) of the header: compiled sides are their naive compilation; reference programs are compared with the hand-written sides on the check instances.} --- fidelity module (kernel-checked: the compiled sides of the command are this program's naive compilation).\par}

\begin{Verbatim}[fontsize=\scriptsize,breaklines=true,breakanywhere=true,frame=single,framesep=1.5mm]
T(x1, y1) :- Edge(x1, y1).
T(x1, y1) :- T(x1, z1), Edge(z1, y1).
\end{Verbatim}

\paragraph{Reading.} The Datalog program is ordinary graph transitive closure, T(x,y) :- Edge(x,y) and T(x,y) :- T(x,z), Edge(z,y). Side 1 encodes NetworkX's transitive\_closure(reflexive=False): for each node taken in a fixed order it runs an edge-wise breadth-first search (edge\_bfs) set-at-a-time, accumulating into Tca every head reached by one or more edges from that node. Side 2 encodes PostgreSQL's recursive-union executor (ExecRecursiveUnion): a working table of newly derived edges is repeatedly joined with Edge and folded into a result Res, already-seen rows dropped, until no new rows appear. The Hoare triple claims the two independently produced relations coincide, Tca = Res, i.e. the per-source edge-BFS search and PostgreSQL's round-based recursive union compute exactly the same transitive closure, for every strict total order used to sequence the per-source search.

\paragraph{First-order reading of PRE/POST, translated mechanically from the relational-algebra assertion in Input.lean.} The relational-algebra assertion in \texttt{Input.lean} remains the authority.

\textit{PRE:} (¬∃x. Lt(x,x)) ∧ (∀y z. ((∃u. Lt(y,u) ∧ Lt(u,z)) → Lt(y,z))) ∧ (∀v w. (Lt(v,w) → ((∃x1. Edge(v,x1)) ∨ ∃x2. Edge(x2,v)) ∧ ((∃x3. Edge(w,x3)) ∨ ∃x4. Edge(x4,w)))) ∧ ∀x5 x6. (((∃x7. Edge(x5,x7)) ∨ ∃x8. Edge(x8,x5)) ∧ ((∃x9. Edge(x6,x9)) ∨ ∃x10. Edge(x10,x6)) ∧ ¬x5 = x6 → Lt(x5,x6) ∨ Lt(x6,x5))\par Lt is a strict total order (irreflexive, transitive, total, and confined to the nodes that occur in Edge) over the nodes of the graph.

\textit{POST:} ∀x y. (Tca(x,y) ↔ Res(x,y))\par The relation Tca, built by NetworkX's per-source edge-BFS search taken node by node in Lt order, equals the relation Res, built by PostgreSQL's round-based recursive union.

\paragraph{Whiel encoding (Hoare triple).}

\begin{Verbatim}[fontsize=\tiny,breaklines=true,breakanywhere=true,frame=single,framesep=1.5mm]
SCHEMA
  Rest/1, Cur/1, Edge/2, Lt/2, Vis/2, Fr/2, Step/2, New/2, Tca/2, Res/2, WorkT/2, Inter/2

PRE
  ((((σ[#0 = #1] Lt) = ∅) ∧
      ((π[0, 3] (σ[#1 = #2] (Lt × Lt))) ⊆ Lt)) ∧
    ((Lt ⊆ ((π[0] Edge ∪ π[1] Edge) × (π[0] Edge ∪ π[1] Edge))) ∧
      ((σ[¬(#0 = #1)] ((π[0] Edge ∪ π[1] Edge) × (π[0] Edge ∪ π[1] Edge))) ⊆ (Lt ∪ π[1, 0] Lt))))

CMD
  Tca := ∅[2];
  Rest := (π[0] Edge ∪ π[1] Edge);
  WHILE (Rest ≠ ∅) DO
    Cur := (Rest ∖ (π[1] (σ[#0 = #2] (Lt × Rest))));
    Vis := (π[0, 0] Cur);
    Fr := Vis;
    WHILE (Fr ≠ ∅) DO
      Step := (π[0, 3] (σ[#1 = #2] (Fr × Edge)));
      Tca := (Tca ∪ Step);
      New := (Step ∖ Vis);
      Vis := (Vis ∪ New);
      Fr := New
    END;
    Rest := (Rest ∖ Cur)
  END;
  Res := Edge;
  WorkT := Edge;
  WHILE (WorkT ≠ ∅) DO
    Inter := ((π[0, 3] (σ[#1 = #2] (WorkT × Edge))) ∖ Res);
    Res := (Res ∪ Inter);
    WorkT := Inter
  END

POST
  (Tca = Res)
\end{Verbatim}

\Needspace{0.45\textheight}

\section{Example5031: NetworkX per-source search with cutoff 2 versus PostgreSQL recursive-union rounds (invalid twin of Example5030)}\label{case:Example5031}

\begin{description}[leftmargin=2.6cm,style=nextline,itemsep=1pt]

\item[Category] production pair (two algorithms from production sources or legacy production cases); family: production algorithm pair (rate mismatch), invalid twin.

\item[Kind] equivalence.

\item[Contributors] Leo Zhang.

\item[Status] certified invalid (Certificate/Invalid.lean, written by the certificate emitter); expected invalid.

\item[Tags] \hyperref[tag:database-engine]{\texttt{database-engine}} \hyperref[tag:evaluation-strategies]{\texttt{evaluation-strategies}} \hyperref[tag:recursion-limit]{\texttt{recursion-limit}} \hyperref[tag:semi-naive]{\texttt{semi-naive}}.

\item[Input] \path{Benchmark/Example5031/Input.lean} (canonical notation).

\item[Summary] The invalid twin of Example5030: a per-source search with a depth cutoff against PostgreSQL's recursive-union rounds.

\item[Sides] hand-written: for each node in Lt order, the levels 1 and 2 of single\_source\_shortest\_path\_length(G, v, cutoff=2) (loop-free per node) \par hand-written: PostgreSQL ExecRecursiveUnion rounds (legacy Example2037) seeded with Edge

\item[Sources] {\raggedright \href{https://github.com/networkx/networkx/blob/main/networkx/algorithms/shortest_paths/unweighted.py}{NetworkX networkx/networkx networkx/algorithms/shortest\_paths/unweighted.py single\_source\_shortest\_path\_length, \_single\_shortest\_path\_length (seen = set(firstlevel); while nextlevel and cutoff \textgreater{} level)} (branch main; commit 4e74880b0da0 (2026-09-12); read 2026-09-15) \par \href{https://github.com/postgres/postgres/blob/master/src/backend/executor/nodeRecursiveunion.c}{PostgreSQL postgres/postgres src/backend/executor/nodeRecursiveunion.c ExecRecursiveUnion (header comment steps 2.1-2.6; 'Ignore tuple if already seen'; working/intermediate table swap)} (branch master; commit 6e58d6356fcb (2026-09-15); read 2026-09-15; encoded as legacy Example2037 / Example5001 hand side with the base term equal to the edge relation)\par}

\item[Twin] Example5030 (valid: unbounded per-source closure).

\item[Obstruction note] from the metadata: invalid (depth cutoff); rate mismatch as in Example5030

\end{description}

\paragraph{Datalog program(s).}

{\raggedright\textit{prog\_T --- The Datalog program(s) of the header: compiled sides are their naive compilation; reference programs are compared with the hand-written sides on the check instances.} --- fidelity module (kernel-checked: the compiled sides of the command are this program's naive compilation).\par}

\begin{Verbatim}[fontsize=\scriptsize,breaklines=true,breakanywhere=true,frame=single,framesep=1.5mm]
T(x1, y1) :- Edge(x1, y1).
T(x1, y1) :- T(x1, z1), Edge(z1, y1).
\end{Verbatim}

\paragraph{Reading.} This is the invalid twin of Example5030: side 2 is the same PostgreSQL recursive-union closure, but side 1 replaces the unbounded per-source search with NetworkX's single\_source\_shortest\_path\_length run with cutoff=2, so for each source it only collects the nodes reachable within two edges (levels 1 and 2), accumulated into TcB. The Hoare triple claims TcB = Res, i.e. that this depth-2-bounded per-source search still equals the full transitive closure computed by PostgreSQL's rounds. This fails because a depth cutoff silently drops longer paths: the case gives an explicit witness (a length-4 path under a compatible order) where the executor derives a pair the bounded search never reaches, so the claim is refuted by a kernel-checked counterexample.

\paragraph{First-order reading of PRE/POST, translated mechanically from the relational-algebra assertion in Input.lean.} The relational-algebra assertion in \texttt{Input.lean} remains the authority.

\textit{PRE:} (¬∃x. Lt(x,x)) ∧ (∀y z. ((∃u. Lt(y,u) ∧ Lt(u,z)) → Lt(y,z))) ∧ (∀v w. (Lt(v,w) → ((∃x1. Edge(v,x1)) ∨ ∃x2. Edge(x2,v)) ∧ ((∃x3. Edge(w,x3)) ∨ ∃x4. Edge(x4,w)))) ∧ ∀x5 x6. (((∃x7. Edge(x5,x7)) ∨ ∃x8. Edge(x8,x5)) ∧ ((∃x9. Edge(x6,x9)) ∨ ∃x10. Edge(x10,x6)) ∧ ¬x5 = x6 → Lt(x5,x6) ∨ Lt(x6,x5))\par Lt is a strict total order over the nodes of Edge, exactly as in Example5030.

\textit{POST:} ∀x y. (TcB(x,y) ↔ Res(x,y))\par It is claimed that TcB, the union over sources of the nodes reachable within 2 edges (per NetworkX's single\_source\_shortest\_path\_length with cutoff=2), equals Res, the full transitive closure computed by PostgreSQL's recursive union --- a claim the case shows to be false.

\paragraph{Whiel encoding (Hoare triple).}

\begin{Verbatim}[fontsize=\scriptsize,breaklines=true,breakanywhere=true,frame=single,framesep=1.5mm]
SCHEMA
  Rest/1, Cur/1, Edge/2, Lt/2, SeenA/2, LvlA/2, SeenB/2, LvlB/2, TcB/2, Res/2, WorkT/2, Inter/2

PRE
  ((((σ[#0 = #1] Lt) = ∅) ∧
      ((π[0, 3] (σ[#1 = #2] (Lt × Lt))) ⊆ Lt)) ∧
    ((Lt ⊆ ((π[0] Edge ∪ π[1] Edge) × (π[0] Edge ∪ π[1] Edge))) ∧
      ((σ[¬(#0 = #1)] ((π[0] Edge ∪ π[1] Edge) × (π[0] Edge ∪ π[1] Edge))) ⊆ (Lt ∪ π[1, 0] Lt))))

CMD
  TcB := ∅[2];
  Rest := (π[0] Edge ∪ π[1] Edge);
  WHILE (Rest ≠ ∅) DO
    Cur := (Rest ∖ (π[1] (σ[#0 = #2] (Lt × Rest))));
    SeenA := (π[0, 0] Cur);
    LvlA := ((π[0, 3] (σ[#1 = #2] (SeenA × Edge))) ∖ SeenA);
    SeenB := (SeenA ∪ LvlA);
    LvlB := ((π[0, 3] (σ[#1 = #2] (LvlA × Edge))) ∖ SeenB);
    TcB := ((TcB ∪ LvlA) ∪ LvlB);
    Rest := (Rest ∖ Cur)
  END;
  Res := Edge;
  WorkT := Edge;
  WHILE (WorkT ≠ ∅) DO
    Inter := ((π[0, 3] (σ[#1 = #2] (WorkT × Edge))) ∖ Res);
    Res := (Res ∪ Inter);
    WorkT := Inter
  END

POST
  (TcB = Res)
\end{Verbatim}

\Needspace{0.45\textheight}

\section{Example5032: Delta-driven incremental maintenance of a transitive closure under insertions versus PostgreSQL recomputation on the updated edges}\label{case:Example5032}

\begin{description}[leftmargin=2.6cm,style=nextline,itemsep=1pt]

\item[Category] production pair (two algorithms from production sources or legacy production cases); family: production algorithm pair (dependent phases, rate mismatch).

\item[Kind] equivalence.

\item[Contributors] Leo Zhang.

\item[Status] certified valid (Certificate/Valid.lean, written by the certificate emitter); expected valid.

\item[Tags] \hyperref[tag:database-engine]{\texttt{database-engine}} \hyperref[tag:evaluation-strategies]{\texttt{evaluation-strategies}} \hyperref[tag:prophecy]{\texttt{prophecy}} \hyperref[tag:semi-naive]{\texttt{semi-naive}} \hyperref[tag:view-maintenance]{\texttt{view-maintenance}}.

\item[Input] \path{Benchmark/Example5032/Input.lean} (canonical notation).

\item[Summary] Delta-driven incremental maintenance of a transitive closure under insertions against PostgreSQL's recomputation on the updated edge set.

\item[Sides] naive compiler: TOld = closure of EBase (phase 1) \par hand-written: delta-driven maintainer (seed EDelta ∪ TOld ∘ EDelta minus TOld, propagate Delta ∘ (EBase ∪ EDelta) minus T until empty), the set-level delta rule of DBSP/Feldera and Materialize for an insert-only update \par hand-written: PostgreSQL ExecRecursiveUnion rounds (legacy Example2037) on EBase ∪ EDelta

\item[Sources] {\raggedright \href{https://materialize.com/docs/sql/select/recursive-ctes/}{Materialize MaterializeInc/materialize doc/user/content/sql/select/recursive-ctes.md WITH MUTUALLY RECURSIVE semantics (initialise empty, update bindings in sequence, stop when no binding changes; 'work resulting only from the input changes for this iteration ... feeds back the resulting output changes')} (branch main; commit 4e1db20efcc4 (2026-09-15); read 2026-09-15) \par \href{https://arxiv.org/abs/2203.16684}{DBSP} (sections 5.1 (recursion circuit, D yields the new facts of each iteration, semi-naive form 5.1), 6 (incremental recursive programs), 6.1 (transitive closure example R(x,y) :- E(x,y); R(x,y) :- E(x,z), R(z,y)), 6.2 (complexity); read 2026-09-15) \par \href{https://github.com/feldera/feldera/blob/main/crates/dbsp/src/operator/recursive.rs}{Feldera feldera/feldera crates/dbsp/src/operator/recursive.rs; docs.feldera.com/docs/sql/recursion.mdx Circuit::recursive (doc: Δx = y - x with y = f(i+Δi, y); implicit distinct); SQL recursion semantics (least fixed point, DISTINCT, deletions recomputed incrementally)} (branch main; commit 16c3f5324fa0 (2026-09-15); read 2026-09-15) \par \href{https://github.com/postgres/postgres/blob/master/src/backend/executor/nodeRecursiveunion.c}{PostgreSQL postgres/postgres src/backend/executor/nodeRecursiveunion.c ExecRecursiveUnion (header comment steps 2.1-2.6; 'Ignore tuple if already seen'; working/intermediate table swap)} (branch master; commit 6e58d6356fcb (2026-09-15); read 2026-09-15; encoded as legacy Example2037 / Example5001 hand side with the base term equal to the edge relation)\par}

\item[Prophecy] accepted invariant with a prophecy clause above level zero (Core.json); no certificate built from it yet.

\end{description}

\paragraph{Datalog program(s).}

{\raggedright\textit{prog\_TOld --- The Datalog program(s) of the header: compiled sides are their naive compilation; reference programs are compared with the hand-written sides on the check instances.} --- fidelity module (kernel-checked: the compiled sides of the command are this program's naive compilation).\par}

\begin{Verbatim}[fontsize=\scriptsize,breaklines=true,breakanywhere=true,frame=single,framesep=1.5mm]
TOld(x1, y1) :- EBase(x1, y1).
TOld(x1, y1) :- TOld(x1, z1), EBase(z1, y1).
\end{Verbatim}

{\raggedright\textit{prog\_TNew} --- fidelity module (kernel-checked: the compiled sides of the command are this program's naive compilation).\par}

\begin{Verbatim}[fontsize=\scriptsize,breaklines=true,breakanywhere=true,frame=single,framesep=1.5mm]
TNew(x1, y1) :- EBase(x1, y1).
TNew(x1, y1) :- EDelta(x1, y1).
TNew(x1, y1) :- TNew(x1, z1), EBase(z1, y1).
TNew(x1, y1) :- TNew(x1, z1), EDelta(z1, y1).
\end{Verbatim}

{\raggedright\textit{rules as stated in the header comment of Input.lean} --- header comment.\par}

\begin{Verbatim}[fontsize=\scriptsize,breaklines=true,breakanywhere=true,frame=single,framesep=1.5mm]
T(x, y) :- E(x, y)
T(x, y) :- T(x, z), E(z, y)
TOld(x, y) :- EBase(x, y)
TOld(x, y) :- TOld(x, z), EBase(z, y)
\end{Verbatim}

\paragraph{Reading.} The case has three phases over an edge set that is split into a base part EBase and an inserted delta EDelta. Phase 1 is the naive fixpoint computation of TOld, the transitive closure of EBase alone. Phase 2 is a hand-written delta-driven maintainer modelled on DBSP/Feldera/Materialize-style incremental view maintenance: it seeds Delta with the consequences of the inserted edges against the already-known old closure (EDelta together with TOld composed with EDelta, minus what TOld already has) and then repeatedly extends only that delta by joining it with the whole new edge set EBase ∪ EDelta, folding the result into T, until the delta is empty --- the old view TOld is never recomputed. Phase 3 recomputes the closure Res of EBase ∪ EDelta from scratch via PostgreSQL's recursive-union rounds. The Hoare triple's claim T = Res is the correctness statement for incremental view maintenance under an insertion: propagating only the change from the old fixpoint yields the same relation as recomputing the closure of the updated edge set from scratch.

\paragraph{First-order reading of PRE/POST, translated mechanically from the relational-algebra assertion in Input.lean.} The relational-algebra assertion in \texttt{Input.lean} remains the authority.

\textit{PRE:} ⊤\par The precondition is simply true; no constraints are placed on EBase, EDelta, or their relationship.

\textit{POST:} ∀x y. (T(x,y) ↔ Res(x,y))\par T, the closure obtained by seeding from the old closure TOld and propagating only the delta caused by the inserted edges EDelta, equals Res, the closure of EBase ∪ EDelta recomputed from scratch by PostgreSQL's recursive-union rounds.

\paragraph{Whiel encoding (Hoare triple).}

\begin{Verbatim}[fontsize=\scriptsize,breaklines=true,breakanywhere=true,frame=single,framesep=1.5mm]
SCHEMA
  EBase/2, EDelta/2, TOld/2, TOld_aux/2, Delta/2, T/2, Res/2, WorkT/2, Inter/2

PRE
  true

CMD
  TOld_aux := ∅[2];
  TOld := (EBase ∪ π[0,3] (σ[#1 = #2] ((TOld_aux × EBase))));
  WHILE (¬((TOld = TOld_aux))) DO
    TOld_aux := TOld;
    TOld := (TOld ∪ (EBase ∪ π[0,3] (σ[#1 = #2] ((TOld_aux × EBase)))))
  END;
  Delta := ((EDelta ∪ (π[0, 3] (σ[#1 = #2] (TOld × EDelta)))) ∖ TOld);
  T := (TOld ∪ Delta);
  WHILE (Delta ≠ ∅) DO
    Delta := ((π[0, 3] (σ[#1 = #2] (Delta × (EBase ∪ EDelta)))) ∖ T);
    T := (T ∪ Delta)
  END;
  Res := (EBase ∪ EDelta);
  WorkT := (EBase ∪ EDelta);
  WHILE (WorkT ≠ ∅) DO
    Inter := ((π[0, 3] (σ[#1 = #2] (WorkT × (EBase ∪ EDelta)))) ∖ Res);
    Res := (Res ∪ Inter);
    WorkT := Inter
  END

POST
  (T = Res)
\end{Verbatim}

\Needspace{0.45\textheight}

\section{Example5033: Soufflé magic-set demand-driven evaluation versus PostgreSQL bulk evaluation restricted to the demanded sources}\label{case:Example5033}

\begin{description}[leftmargin=2.6cm,style=nextline,itemsep=1pt]

\item[Category] production pair (two algorithms from production sources or legacy production cases); family: production algorithm pair (rate mismatch).

\item[Kind] equivalence.

\item[Contributors] Leo Zhang.

\item[Status] certified valid (Certificate/Valid.lean, written by the certificate emitter); expected valid.

\item[Tags] \hyperref[tag:database-engine]{\texttt{database-engine}} \hyperref[tag:evaluation-strategies]{\texttt{evaluation-strategies}} \hyperref[tag:magic-sets]{\texttt{magic-sets}} \hyperref[tag:prophecy]{\texttt{prophecy}} \hyperref[tag:semi-naive]{\texttt{semi-naive}}.

\item[Input] \path{Benchmark/Example5033/Input.lean} (canonical notation).

\item[Summary] Soufflé's magic-set (demand-driven) evaluation of a right-linear path query with a bound source against PostgreSQL's bulk evaluation of the original query restricted to the demanded sources.

\item[Sides] naive compiler: Soufflé MagicSetCoreTransformer output for Path with adornment bf (the program of Example2003/2004) \par hand-written: PostgreSQL ExecRecursiveUnion rounds (legacy Example2037) for the original right-linear query, joined with the demand set in the postcondition

\item[Sources] {\raggedright \href{https://github.com/souffle-lang/souffle/blob/master/src/ast/transform/MagicSet.cpp}{Soufflé souffle-lang/souffle src/ast/transform/MagicSet.cpp MagicSetCoreTransformer::transform, createMagicClause, createMagicAtom, getMagicName} (branch master; commit a1303be3c016 (2026-07-13); read 2026-09-15) \par \href{https://souffle-lang.github.io/magicset}{Soufflé documentation souffle-lang/souffle-lang.github.io pages/docs/magicset.md The Algorithm: constraint normalisation, adornment (SIPS), the magic set transformation ('For all adorned predicates ... Add the clause magic(Ai\_ai) :- magic(A\_a), A1\_a1, ..., A(i-1)\_a(i-1) ... Replace C with A\_a :- magic(A\_a), T')} (read 2026-09-15) \par \href{https://github.com/postgres/postgres/blob/master/src/backend/executor/nodeRecursiveunion.c}{PostgreSQL postgres/postgres src/backend/executor/nodeRecursiveunion.c ExecRecursiveUnion (header comment steps 2.1-2.6; 'Ignore tuple if already seen'; working/intermediate table swap)} (branch master; commit 6e58d6356fcb (2026-09-15); read 2026-09-15; encoded as legacy Example2037 / Example5001 hand side with the base term equal to the edge relation)\par}

\item[Prophecy] accepted invariant with a prophecy clause above level zero (Core.json); no certificate built from it yet.

\end{description}

\paragraph{Datalog program(s).}

{\raggedright\textit{prog\_magic} --- fidelity module (kernel-checked: the compiled sides of the command are this program's naive compilation).\par}

\begin{Verbatim}[fontsize=\scriptsize,breaklines=true,breakanywhere=true,frame=single,framesep=1.5mm]
MPath(x1) :- Root(x1).
MPath(y1) :- MPath(x1), Edge(x1, y1).
PathBF(x1, y1) :- MPath(x1), Edge(x1, y1).
PathBF(x1, z1) :- MPath(x1), Edge(x1, y1), PathBF(y1, z1).
\end{Verbatim}

{\raggedright\textit{prog\_PathR} --- fidelity module (kernel-checked: the compiled sides of the command are this program's naive compilation).\par}

\begin{Verbatim}[fontsize=\scriptsize,breaklines=true,breakanywhere=true,frame=single,framesep=1.5mm]
PathR(x1, y1) :- Edge(x1, y1).
PathR(x1, z1) :- Edge(x1, y1), PathR(y1, z1).
\end{Verbatim}

{\raggedright\textit{rules as stated in the header comment of Input.lean} --- header comment.\par}

\begin{Verbatim}[fontsize=\scriptsize,breaklines=true,breakanywhere=true,frame=single,framesep=1.5mm]
Path(x, y) :- Edge(x, y)
Path(x, z) :- Edge(x, y), Path(y, z)
A_a :- A1_a1, ..., An_an
MPath(x) :- Root(x)
MPath(y) :- MPath(x), Edge(x, y)
PathBF(x, y) :- MPath(x), Edge(x, y)
PathBF(x, z) :- MPath(x), Edge(x, y), PathBF(y, z)
\end{Verbatim}

\paragraph{Reading.} The underlying query is the right-linear path program Path(x,y) :- Edge(x,y); Path(x,z) :- Edge(x,y), Path(y,z), asked with the first argument bound to a single constant. Side 1 encodes Soufflé's magic-set transformation for this adornment (bf): it introduces a demand relation MPath seeded from the bound source Root and propagated forward along edges, and an answer relation PathBF built only over pairs whose source is currently demanded --- the demand-driven, top-down-restricted evaluation Soufflé's --magic-transform would produce. Side 2 encodes PostgreSQL's bulk recursive-union evaluation of the unrestricted query, producing the full path closure Res for every source, independent of any demand. The Hoare triple claims that the demand-driven answers PathBF equal the bulk answers Res restricted, by a join with MPath, to the pairs whose source is in the demand set --- the standard correctness obligation of the magic-set transformation, stated directly against PostgreSQL's actual round-based executor rather than against an idealized reference program.

\paragraph{First-order reading of PRE/POST, translated mechanically from the relational-algebra assertion in Input.lean.} The relational-algebra assertion in \texttt{Input.lean} remains the authority.

\textit{PRE:} ∀x. (Root(x) ↔ x = "root")\par Root is fixed to the singleton set containing exactly the constant "root", the single bound source of the query.

\textit{POST:} ∀x y. (PathBF(x,y) ↔ MPath(x) ∧ Res(x,y))\par PathBF, the magic-set/demand-driven answer relation, equals the projection of MPath × Res onto pairs whose source column matches a demanded node --- i.e. the demand-driven answers equal the bulk-evaluated answers restricted to the demanded sources.

\paragraph{Whiel encoding (Hoare triple).}

\begin{Verbatim}[fontsize=\scriptsize,breaklines=true,breakanywhere=true,frame=single,framesep=1.5mm]
SCHEMA
  Root/1, MPath/1, MPath_aux/1, Edge/2, PathBF/2, PathBF_aux/2, Res/2, WorkT/2, Inter/2

PRE
  (Root = { "root" })

CMD
  MPath_aux := ∅[1];
  PathBF_aux := ∅[2];
  MPath := (Root ∪ π[2] (σ[#0 = #1] ((MPath_aux × Edge))));
  PathBF := (π[0,2] (σ[#0 = #1] ((MPath_aux × Edge))) ∪ π[0,4] (σ[(#0 = #1 ∧ #2 = #3)] ((MPath_aux × (Edge × PathBF_aux)))));
  WHILE (¬(((MPath = MPath_aux) ∧ (PathBF = PathBF_aux)))) DO
    MPath_aux := MPath;
    PathBF_aux := PathBF;
    MPath := (MPath ∪ (Root ∪ π[2] (σ[#0 = #1] ((MPath_aux × Edge)))));
    PathBF := (PathBF ∪ (π[0,2] (σ[#0 = #1] ((MPath_aux × Edge))) ∪ π[0,4] (σ[(#0 = #1 ∧ #2 = #3)] ((MPath_aux × (Edge × PathBF_aux))))))
  END;
  Res := Edge;
  WorkT := Edge;
  WHILE (WorkT ≠ ∅) DO
    Inter := ((π[0, 3] (σ[#1 = #2] (Edge × WorkT))) ∖ Res);
    Res := (Res ∪ Inter);
    WorkT := Inter
  END

POST
  (PathBF = (π[1, 2] (σ[#0 = #1] (MPath × Res))))
\end{Verbatim}

\Needspace{0.45\textheight}

\section{Example5034: Warshall's intermediate-vertex sweep (boolean Floyd-Warshall) versus PostgreSQL recursive-union rounds}\label{case:Example5034}

\begin{description}[leftmargin=2.6cm,style=nextline,itemsep=1pt]

\item[Category] production pair (two algorithms from production sources or legacy production cases); family: production algorithm pair (rate mismatch).

\item[Kind] equivalence.

\item[Contributors] Leo Zhang.

\item[Status] certified valid (Certificate/Valid.lean, written by the certificate emitter); expected valid.

\item[Tags] \hyperref[tag:database-engine]{\texttt{database-engine}} \hyperref[tag:evaluation-strategies]{\texttt{evaluation-strategies}} \hyperref[tag:prophecy]{\texttt{prophecy}} \hyperref[tag:semi-naive]{\texttt{semi-naive}}.

\item[Input] \path{Benchmark/Example5034/Input.lean} (canonical notation).

\item[Summary] Warshall's intermediate-vertex sweep, the boolean core of Floyd-Warshall, against PostgreSQL's recursive-union rounds.

\item[Sides] hand-written: boolean Floyd-Warshall sweep over the vertices in Lt order (SciPy \_floyd\_warshall / NetworkX floyd\_warshall\_predecessor\_and\_distance, weights not modelled) \par hand-written: PostgreSQL ExecRecursiveUnion rounds (legacy Example2037) seeded with Edge

\item[Sources] {\raggedright \href{https://github.com/scipy/scipy/blob/main/scipy/sparse/csgraph/_shortest_path.pyx}{SciPy scipy/scipy scipy/sparse/csgraph/\_shortest\_path.pyx floyd\_warshall (wrapper, unweighted option sets finite entries to 1), \_floyd\_warshall (for k in range(N): for i in range(N): if dist\_matrix[i, k] == INFINITY: continue; for j in range(N): d\_ijk = dist\_matrix[i, k] + dist\_matrix[k, j]; if d\_ijk \textless{} dist\_matrix[i, j]: dist\_matrix[i, j] = d\_ijk; comment 'In each loop, this finds the shortest path from point i to point j using intermediate nodes 0 ... k')} (branch main; commit 884ac0b6fc53 (2026-09-15); read 2026-09-15) \par \href{https://github.com/networkx/networkx/blob/main/networkx/algorithms/shortest_paths/dense.py}{NetworkX networkx/networkx networkx/algorithms/shortest\_paths/dense.py floyd\_warshall\_predecessor\_and\_distance (for w in G: for u in G: for v in G: d = dist\_u[w] + dist\_w[v]; if dist\_u[v] \textgreater{} d: ...), floyd\_warshall, floyd\_warshall\_numpy (for i in range(n): A = np.minimum(A, A[i, :][np.newaxis, :] + A[:, i][:, np.newaxis]))} (branch main; commit 4e74880b0da0 (2026-09-12); read 2026-09-15) \par \href{https://github.com/postgres/postgres/blob/master/src/backend/executor/nodeRecursiveunion.c}{PostgreSQL postgres/postgres src/backend/executor/nodeRecursiveunion.c ExecRecursiveUnion (header comment steps 2.1-2.6; 'Ignore tuple if already seen'; working/intermediate table swap)} (branch master; commit 6e58d6356fcb (2026-09-15); read 2026-09-15; encoded as legacy Example2037 / Example5001 hand side with the base term equal to the edge relation)\par}

\item[Twin] Example5035 (invalid: only the first two vertices of the order are swept).

\item[Prophecy] accepted invariant with a prophecy clause above level zero (Core.json); no certificate built from it yet.

\end{description}

\paragraph{Datalog program(s).}

{\raggedright\textit{prog\_T --- The Datalog program(s) of the header: compiled sides are their naive compilation; reference programs are compared with the hand-written sides on the check instances.} --- fidelity module (kernel-checked: the compiled sides of the command are this program's naive compilation).\par}

\begin{Verbatim}[fontsize=\scriptsize,breaklines=true,breakanywhere=true,frame=single,framesep=1.5mm]
T(x1, y1) :- Edge(x1, y1).
T(x1, y1) :- T(x1, z1), Edge(z1, y1).
\end{Verbatim}

\paragraph{Reading.} The Datalog program is again transitive closure, T(x,y) :- Edge(x,y); T(x,y) :- T(x,z), Edge(z,y). Side 1 encodes the boolean core of Floyd-Warshall (SciPy's \_floyd\_warshall / NetworkX's floyd\_warshall\_predecessor\_and\_distance with weights ignored, i.e. Warshall's algorithm): for each intermediate vertex taken in a fixed order, it composes all pairs currently entering that vertex with all pairs currently leaving it and folds the result into TW, one vertex-sweep at a time. Side 2 is PostgreSQL's recursive-union rounds on Edge, as in the other cases in this batch. The Hoare triple claims TW = Res: Warshall's intermediate-vertex sweep and PostgreSQL's round-based recursive union compute the same transitive closure, for every strict total order used to sequence the sweep.

\paragraph{First-order reading of PRE/POST, translated mechanically from the relational-algebra assertion in Input.lean.} The relational-algebra assertion in \texttt{Input.lean} remains the authority.

\textit{PRE:} (¬∃x. Lt(x,x)) ∧ (∀y z. ((∃u. Lt(y,u) ∧ Lt(u,z)) → Lt(y,z))) ∧ (∀v w. (Lt(v,w) → ((∃x1. Edge(v,x1)) ∨ ∃x2. Edge(x2,v)) ∧ ((∃x3. Edge(w,x3)) ∨ ∃x4. Edge(x4,w)))) ∧ ∀x5 x6. (((∃x7. Edge(x5,x7)) ∨ ∃x8. Edge(x8,x5)) ∧ ((∃x9. Edge(x6,x9)) ∨ ∃x10. Edge(x10,x6)) ∧ ¬x5 = x6 → Lt(x5,x6) ∨ Lt(x6,x5))\par Lt is a strict total order over the nodes that occur in Edge.

\textit{POST:} ∀x y. (TW(x,y) ↔ Res(x,y))\par TW, the relation built by sweeping intermediate vertices in Lt order and composing incoming with outgoing pairs at each vertex, equals Res, built by PostgreSQL's recursive-union rounds.

\paragraph{Whiel encoding (Hoare triple).}

\begin{Verbatim}[fontsize=\scriptsize,breaklines=true,breakanywhere=true,frame=single,framesep=1.5mm]
SCHEMA
  Rest/1, Cur/1, Edge/2, Lt/2, TW/2, InK/2, OutK/2, Res/2, WorkT/2, Inter/2

PRE
  ((((σ[#0 = #1] Lt) = ∅) ∧
      ((π[0, 3] (σ[#1 = #2] (Lt × Lt))) ⊆ Lt)) ∧
    ((Lt ⊆ ((π[0] Edge ∪ π[1] Edge) × (π[0] Edge ∪ π[1] Edge))) ∧
      ((σ[¬(#0 = #1)] ((π[0] Edge ∪ π[1] Edge) × (π[0] Edge ∪ π[1] Edge))) ⊆ (Lt ∪ π[1, 0] Lt))))

CMD
  TW := Edge;
  Rest := (π[0] Edge ∪ π[1] Edge);
  WHILE (Rest ≠ ∅) DO
    Cur := (Rest ∖ (π[1] (σ[#0 = #2] (Lt × Rest))));
    InK := (π[0, 1] (σ[#1 = #2] (TW × Cur)));
    OutK := (π[1, 2] (σ[#0 = #1] (Cur × TW)));
    TW := (TW ∪ (π[0, 3] (σ[#1 = #2] (InK × OutK))));
    Rest := (Rest ∖ Cur)
  END;
  Res := Edge;
  WorkT := Edge;
  WHILE (WorkT ≠ ∅) DO
    Inter := ((π[0, 3] (σ[#1 = #2] (WorkT × Edge))) ∖ Res);
    Res := (Res ∪ Inter);
    WorkT := Inter
  END

POST
  (TW = Res)
\end{Verbatim}

\Needspace{0.45\textheight}

\section{Example5035: Warshall's sweep truncated to the first two vertices versus PostgreSQL recursive-union rounds (invalid twin of Example5034)}\label{case:Example5035}

\begin{description}[leftmargin=2.6cm,style=nextline,itemsep=1pt]

\item[Category] production pair (two algorithms from production sources or legacy production cases); family: production algorithm pair (rate mismatch), invalid twin.

\item[Kind] equivalence.

\item[Contributors] Leo Zhang.

\item[Status] certified invalid (Certificate/Invalid.lean, written by the certificate emitter); expected invalid.

\item[Tags] \hyperref[tag:database-engine]{\texttt{database-engine}} \hyperref[tag:evaluation-strategies]{\texttt{evaluation-strategies}} \hyperref[tag:recursion-limit]{\texttt{recursion-limit}} \hyperref[tag:semi-naive]{\texttt{semi-naive}}.

\item[Input] \path{Benchmark/Example5035/Input.lean} (canonical notation).

\item[Summary] The invalid twin of Example5034: Warshall's sweep stopped after the first two vertices of the order, against PostgreSQL's recursive-union rounds.

\item[Sides] hand-written: boolean Floyd-Warshall sweeps of the first two vertices of the order only (loop-free) \par hand-written: PostgreSQL ExecRecursiveUnion rounds (legacy Example2037) seeded with Edge

\item[Sources] {\raggedright \href{https://github.com/scipy/scipy/blob/main/scipy/sparse/csgraph/_shortest_path.pyx}{SciPy scipy/scipy scipy/sparse/csgraph/\_shortest\_path.pyx floyd\_warshall (wrapper, unweighted option sets finite entries to 1), \_floyd\_warshall (for k in range(N): for i in range(N): if dist\_matrix[i, k] == INFINITY: continue; for j in range(N): d\_ijk = dist\_matrix[i, k] + dist\_matrix[k, j]; if d\_ijk \textless{} dist\_matrix[i, j]: dist\_matrix[i, j] = d\_ijk; comment 'In each loop, this finds the shortest path from point i to point j using intermediate nodes 0 ... k')} (branch main; commit 884ac0b6fc53 (2026-09-15); read 2026-09-15) \par \href{https://github.com/networkx/networkx/blob/main/networkx/algorithms/shortest_paths/dense.py}{NetworkX networkx/networkx networkx/algorithms/shortest\_paths/dense.py floyd\_warshall\_predecessor\_and\_distance (for w in G: for u in G: for v in G: d = dist\_u[w] + dist\_w[v]; if dist\_u[v] \textgreater{} d: ...), floyd\_warshall, floyd\_warshall\_numpy (for i in range(n): A = np.minimum(A, A[i, :][np.newaxis, :] + A[:, i][:, np.newaxis]))} (branch main; commit 4e74880b0da0 (2026-09-12); read 2026-09-15) \par \href{https://github.com/postgres/postgres/blob/master/src/backend/executor/nodeRecursiveunion.c}{PostgreSQL postgres/postgres src/backend/executor/nodeRecursiveunion.c ExecRecursiveUnion (header comment steps 2.1-2.6; 'Ignore tuple if already seen'; working/intermediate table swap)} (branch master; commit 6e58d6356fcb (2026-09-15); read 2026-09-15; encoded as legacy Example2037 / Example5001 hand side with the base term equal to the edge relation)\par}

\item[Twin] Example5034 (valid: all vertices swept).

\item[Obstruction note] from the metadata: invalid (bounded number of sweeps); the bounded side is loop-free and enters the precondition, the verified loop is the executor's

\end{description}

\paragraph{Datalog program(s).}

{\raggedright\textit{prog\_T --- The Datalog program(s) of the header: compiled sides are their naive compilation; reference programs are compared with the hand-written sides on the check instances.} --- fidelity module (kernel-checked: the compiled sides of the command are this program's naive compilation).\par}

\begin{Verbatim}[fontsize=\scriptsize,breaklines=true,breakanywhere=true,frame=single,framesep=1.5mm]
T(x1, y1) :- Edge(x1, y1).
T(x1, y1) :- T(x1, z1), Edge(z1, y1).
\end{Verbatim}

\paragraph{Reading.} This is the invalid twin of Example5034: side 2 is again PostgreSQL's recursive-union closure, but side 1 truncates Warshall's sweep to only the first two vertices of the order (a loop-free computation of TWa then TWb, each adding the pairs composed through one more intermediate vertex). The Hoare triple claims TWb = Res, i.e. that sweeping only two intermediate vertices already produces the full transitive closure. This is false in general: cutting a sweep short is not the same as bounding path length, since it misses every path whose intermediate vertices lie outside the swept prefix regardless of the path's length, and the case gives an explicit witness graph (a 4-node path, and separately a 3-cycle) where the two relations differ, backed by a kernel-checked counterexample.

\paragraph{First-order reading of PRE/POST, translated mechanically from the relational-algebra assertion in Input.lean.} The relational-algebra assertion in \texttt{Input.lean} remains the authority.

\textit{PRE:} (¬∃x. Lt(x,x)) ∧ (∀y z. ((∃u. Lt(y,u) ∧ Lt(u,z)) → Lt(y,z))) ∧ (∀v w. (Lt(v,w) → ((∃x1. Edge(v,x1)) ∨ ∃x2. Edge(x2,v)) ∧ ((∃x3. Edge(w,x3)) ∨ ∃x4. Edge(x4,w)))) ∧ ∀x5 x6. (((∃x7. Edge(x5,x7)) ∨ ∃x8. Edge(x8,x5)) ∧ ((∃x9. Edge(x6,x9)) ∨ ∃x10. Edge(x10,x6)) ∧ ¬x5 = x6 → Lt(x5,x6) ∨ Lt(x6,x5))\par Lt is a strict total order over the nodes of Edge, exactly as in Example5034.

\textit{POST:} ∀x y. (TWb(x,y) ↔ Res(x,y))\par It is claimed that TWb, the result of sweeping only the first two vertices of Lt as intermediate vertices, equals Res, the full transitive closure from PostgreSQL's recursive union --- a claim the case shows to be false.

\paragraph{Whiel encoding (Hoare triple).}

\begin{Verbatim}[fontsize=\scriptsize,breaklines=true,breakanywhere=true,frame=single,framesep=1.5mm]
SCHEMA
  KA/1, RestB/1, KB/1, Edge/2, Lt/2, InA/2, OutA/2, TWa/2, InB/2, OutB/2, TWb/2, Res/2, WorkT/2, Inter/2

PRE
  ((((σ[#0 = #1] Lt) = ∅) ∧
      ((π[0, 3] (σ[#1 = #2] (Lt × Lt))) ⊆ Lt)) ∧
    ((Lt ⊆ ((π[0] Edge ∪ π[1] Edge) × (π[0] Edge ∪ π[1] Edge))) ∧
      ((σ[¬(#0 = #1)] ((π[0] Edge ∪ π[1] Edge) × (π[0] Edge ∪ π[1] Edge))) ⊆ (Lt ∪ π[1, 0] Lt))))

CMD
  KA := ((π[0] Edge ∪ π[1] Edge) ∖ (π[1] Lt));
  RestB := ((π[0] Edge ∪ π[1] Edge) ∖ KA);
  KB := (RestB ∖ (π[1] (σ[#0 = #2] (Lt × RestB))));
  InA := (π[0, 1] (σ[#1 = #2] (Edge × KA)));
  OutA := (π[1, 2] (σ[#0 = #1] (KA × Edge)));
  TWa := (Edge ∪ (π[0, 3] (σ[#1 = #2] (InA × OutA))));
  InB := (π[0, 1] (σ[#1 = #2] (TWa × KB)));
  OutB := (π[1, 2] (σ[#0 = #1] (KB × TWa)));
  TWb := (TWa ∪ (π[0, 3] (σ[#1 = #2] (InB × OutB))));
  Res := Edge;
  WorkT := Edge;
  WHILE (WorkT ≠ ∅) DO
    Inter := ((π[0, 3] (σ[#1 = #2] (WorkT × Edge))) ∖ Res);
    Res := (Res ∪ Inter);
    WorkT := Inter
  END

POST
  (TWb = Res)
\end{Verbatim}

\Needspace{0.45\textheight}

\section{Example5036: NetworkX DAG transitive closure by a reverse topological sweep versus PostgreSQL recursive-union rounds}\label{case:Example5036}

\begin{description}[leftmargin=2.6cm,style=nextline,itemsep=1pt]

\item[Category] production pair (two algorithms from production sources or legacy production cases); family: production algorithm pair (rate mismatch).

\item[Kind] equivalence.

\item[Contributors] Leo Zhang.

\item[Status] certified valid (Certificate/Valid.lean, written by the certificate emitter); expected valid.

\item[Tags] \hyperref[tag:database-engine]{\texttt{database-engine}} \hyperref[tag:evaluation-strategies]{\texttt{evaluation-strategies}} \hyperref[tag:prophecy]{\texttt{prophecy}} \hyperref[tag:semi-naive]{\texttt{semi-naive}}.

\item[Input] \path{Benchmark/Example5036/Input.lean} (canonical notation).

\item[Summary] NetworkX's transitive closure of a directed acyclic graph by a reverse topological sweep against PostgreSQL's recursive-union rounds.

\item[Sides] hand-written: NetworkX transitive\_closure\_dag = for each node in reverse Lt order, add its descendants at distance 2 in the current TC (set-at-a-time: Row ∘ TC) \par hand-written: PostgreSQL ExecRecursiveUnion rounds (legacy Example2037) seeded with Edge

\item[Sources] {\raggedright \href{https://github.com/networkx/networkx/blob/main/networkx/algorithms/dag.py}{NetworkX networkx/networkx networkx/algorithms/dag.py transitive\_closure\_dag (TC = G.copy(); for v in reversed(topo\_order): TC.add\_edges\_from((v, u) for u in nx.descendants\_at\_distance(TC, v, 2)))} (branch main; commit 4e74880b0da0 (2026-09-12); read 2026-09-15) \par \href{https://github.com/networkx/networkx/blob/main/networkx/algorithms/traversal/breadth_first_search.py}{NetworkX networkx/networkx networkx/algorithms/traversal/breadth\_first\_search.py descendants\_at\_distance (layer `distance` of bfs\_layers(G, source)), bfs\_layers} (branch main; commit 4e74880b0da0 (2026-09-12); read 2026-09-15) \par \href{https://github.com/postgres/postgres/blob/master/src/backend/executor/nodeRecursiveunion.c}{PostgreSQL postgres/postgres src/backend/executor/nodeRecursiveunion.c ExecRecursiveUnion (header comment steps 2.1-2.6; 'Ignore tuple if already seen'; working/intermediate table swap)} (branch master; commit 6e58d6356fcb (2026-09-15); read 2026-09-15; encoded as legacy Example2037 / Example5001 hand side with the base term equal to the edge relation)\par}

\item[Prophecy] accepted invariant with a prophecy clause above level zero (Core.json); no certificate built from it yet.

\end{description}

\paragraph{Datalog program(s).}

{\raggedright\textit{prog\_T --- The Datalog program(s) of the header: compiled sides are their naive compilation; reference programs are compared with the hand-written sides on the check instances.} --- fidelity module (kernel-checked: the compiled sides of the command are this program's naive compilation).\par}

\begin{Verbatim}[fontsize=\scriptsize,breaklines=true,breakanywhere=true,frame=single,framesep=1.5mm]
T(x1, y1) :- Edge(x1, y1).
T(x1, y1) :- T(x1, z1), Edge(z1, y1).
\end{Verbatim}

\paragraph{Reading.} The Datalog program is transitive closure, T(x,y) :- Edge(x,y); T(x,y) :- T(x,z), Edge(z,y), over a graph that the precondition forces to be a DAG (Lt is a topological order containing Edge). Side 1 encodes NetworkX's transitive\_closure\_dag: processing nodes in reverse topological order, it computes each node's full row of descendants in one step by unioning its current row with that row composed with the whole closure built so far, relying on the fact that every successor of the current node, coming earlier in the reverse sweep, already has a complete row. Side 2 is PostgreSQL's recursive-union rounds on Edge. The Hoare triple claims TC = Res: the reverse-topological-sweep algorithm, which is only correct because the order respects the DAG's edges, computes the same transitive closure as PostgreSQL's generic recursive union, which needs no such assumption.

\paragraph{First-order reading of PRE/POST, translated mechanically from the relational-algebra assertion in Input.lean.} The relational-algebra assertion in \texttt{Input.lean} remains the authority.

\textit{PRE:} (¬∃x. Lt(x,x)) ∧ (∀y z. ((∃u. Lt(y,u) ∧ Lt(u,z)) → Lt(y,z))) ∧ (∀v w. (Lt(v,w) → ((∃x1. Edge(v,x1)) ∨ ∃x2. Edge(x2,v)) ∧ ((∃x3. Edge(w,x3)) ∨ ∃x4. Edge(x4,w)))) ∧ (∀x5 x6. (((∃x7. Edge(x5,x7)) ∨ ∃x8. Edge(x8,x5)) ∧ ((∃x9. Edge(x6,x9)) ∨ ∃x10. Edge(x10,x6)) ∧ ¬x5 = x6 → Lt(x5,x6) ∨ Lt(x6,x5))) ∧ ∀x11 x12. (Edge(x11,x12) → Lt(x11,x12))\par Lt is a strict total order over the nodes of Edge and additionally contains all of Edge, i.e. Lt is a topological order of the (necessarily acyclic) graph.

\textit{POST:} ∀x y. (TC(x,y) ↔ Res(x,y))\par TC, built by sweeping nodes in reverse topological (Lt) order and merging each node's row with the closure computed so far, equals Res, built by PostgreSQL's recursive-union rounds.

\paragraph{Whiel encoding (Hoare triple).}

\begin{Verbatim}[fontsize=\scriptsize,breaklines=true,breakanywhere=true,frame=single,framesep=1.5mm]
SCHEMA
  Rest/1, Cur/1, Edge/2, Lt/2, TC/2, Row/2, Res/2, WorkT/2, Inter/2

PRE
  (((((σ[#0 = #1] Lt) = ∅) ∧
      ((π[0, 3] (σ[#1 = #2] (Lt × Lt))) ⊆ Lt)) ∧
    ((Lt ⊆ ((π[0] Edge ∪ π[1] Edge) × (π[0] Edge ∪ π[1] Edge))) ∧
      ((σ[¬(#0 = #1)] ((π[0] Edge ∪ π[1] Edge) × (π[0] Edge ∪ π[1] Edge))) ⊆ (Lt ∪ π[1, 0] Lt)))) ∧
    (Edge ⊆ Lt))

CMD
  TC := Edge;
  Rest := (π[0] Edge ∪ π[1] Edge);
  WHILE (Rest ≠ ∅) DO
    Cur := (Rest ∖ (π[0] (σ[#1 = #2] (Lt × Rest))));
    Row := (π[1, 2] (σ[#0 = #1] (Cur × TC)));
    TC := (TC ∪ (π[0, 3] (σ[#1 = #2] (Row × TC))));
    Rest := (Rest ∖ Cur)
  END;
  Res := Edge;
  WorkT := Edge;
  WHILE (WorkT ≠ ∅) DO
    Inter := ((π[0, 3] (σ[#1 = #2] (WorkT × Edge))) ∖ Res);
    Res := (Res ∪ Inter);
    WorkT := Inter
  END

POST
  (TC = Res)
\end{Verbatim}

\Needspace{0.45\textheight}

\section{Example5037: Go runtime frontier marking versus the materialised Soufflé closure of the pointer graph joined with the roots}\label{case:Example5037}

\begin{description}[leftmargin=2.6cm,style=nextline,itemsep=1pt]

\item[Category] production pair (two algorithms from production sources or legacy production cases); family: production algorithm pair (dependent phases, schema link by view).

\item[Kind] equivalence.

\item[Contributors] Leo Zhang.

\item[Status] certified valid (Certificate/Valid.lean, written by the certificate emitter); expected valid.

\item[Tags] \hyperref[tag:evaluation-strategies]{\texttt{evaluation-strategies}} \hyperref[tag:nonlinear-recursion]{\texttt{nonlinear-recursion}} \hyperref[tag:prophecy]{\texttt{prophecy}} \hyperref[tag:semi-naive]{\texttt{semi-naive}}.

\item[Input] \path{Benchmark/Example5037/Input.lean} (canonical notation).

\item[Summary] Go's garbage-collector marking, a frontier traversal with a marked-bit check, against the materialised transitive closure of the pointer graph joined with the roots.

\item[Sides] hand-written: Go gcDrain/scanObject/greyobject frontier marking (legacy Example2012 program) \par naive compiler: Soufflé tc (non-linear closure of HeapPtr), joined with RootPtr in the postcondition

\item[Sources] {\raggedright \href{https://github.com/golang/go/blob/master/src/runtime/mgcmark.go}{Go runtime golang/go src/runtime/mgcmark.go gcDrain (root jobs via markroot, then the heap marking loop popping objects/spans and calling scanObject/scanSpan until no work), greyobject ('If marked we have nothing to do'; setMarked; noscan fast path; putObj)} (branch master; commit 8f5d82065574 (2026-09-15); read 2026-09-15) \par \href{https://github.com/golang/go/blob/master/src/runtime/mgcmark_greenteagc.go}{Go runtime golang/go src/runtime/mgcmark\_greenteagc.go; src/runtime/mgcmark\_nogreenteagc.go scanObject (for each pointer slot: findObject then greyobject)} (branch master; commit 8f5d82065574 (2026-09-15); read 2026-09-15)\par}

\item[Prophecy] kernel-checked certificate with a level-1 prophecy clause (Certificate/Valid.lean, Core.json).

\end{description}

\paragraph{Datalog program(s).}

{\raggedright\textit{prog\_Tc --- The Datalog program(s) of the header: compiled sides are their naive compilation; reference programs are compared with the hand-written sides on the check instances.} --- fidelity module (kernel-checked: the compiled sides of the command are this program's naive compilation).\par}

\begin{Verbatim}[fontsize=\scriptsize,breaklines=true,breakanywhere=true,frame=single,framesep=1.5mm]
Tc(x1, y1) :- HeapPtr(x1, y1).
Tc(x1, z1) :- Tc(x1, y1), Tc(y1, z1).
\end{Verbatim}

{\raggedright\textit{prog\_MarkedR} --- fidelity module (kernel-checked: the compiled sides of the command are this program's naive compilation).\par}

\begin{Verbatim}[fontsize=\scriptsize,breaklines=true,breakanywhere=true,frame=single,framesep=1.5mm]
MarkedR(x1) :- RootPtr(x1).
MarkedR(y1) :- MarkedR(x1), HeapPtr(x1, y1).
\end{Verbatim}

{\raggedright\textit{rules as stated in the header comment of Input.lean} --- header comment.\par}

\begin{Verbatim}[fontsize=\scriptsize,breaklines=true,breakanywhere=true,frame=single,framesep=1.5mm]
Marked(x) :- RootPtr(x)
Marked(y) :- Marked(x), HeapPtr(x, y)
Tc(x, y) :- HeapPtr(x, y)
Tc(x, z) :- Tc(x, y), Tc(y, z)
\end{Verbatim}

\paragraph{Reading.} Two distinct Datalog programs are involved: a unary marking program Marked(x) :- RootPtr(x); Marked(y) :- Marked(x), HeapPtr(x,y), and Soufflé's benchmark transitive-closure program Tc(x,y) :- HeapPtr(x,y); Tc(x,z) :- Tc(x,y), Tc(y,z) over the same pointer-graph edges HeapPtr. Side 1 encodes the Go runtime's garbage-collector marking (gcDrain/scanObject/greyobject): a worklist ('grey set') traversal that starts at the roots and, level by level, marks and enqueues every newly reached object, stopping when the queue is empty. Side 2 is the non-linear (doubling) naive fixpoint computation of Tc, the full pointer-graph reachability closure, independent of any roots. The Hoare triple links the two different schemas by a view stated in the postcondition: it claims the GC's marked set equals the roots together with everything the materialized closure Tc says is reachable from some root --- i.e. that frontier-by-frontier marking from the roots computes exactly the root-restricted transitive closure of the pointer graph.

\paragraph{First-order reading of PRE/POST, translated mechanically from the relational-algebra assertion in Input.lean.} The relational-algebra assertion in \texttt{Input.lean} remains the authority.

\textit{PRE:} ⊤\par The precondition is simply true; no constraints are placed on RootPtr or HeapPtr.

\textit{POST:} ∀x. (Marked(x) ↔ RootPtr(x) ∨ ∃y. RootPtr(y) ∧ Tc(y,x))\par Marked, the set produced by the GC's root-seeded frontier marking traversal, equals RootPtr unioned with the objects reachable from some root according to the separately computed pointer-graph transitive closure Tc.

\paragraph{Whiel encoding (Hoare triple).}

\begin{Verbatim}[fontsize=\scriptsize,breaklines=true,breakanywhere=true,frame=single,framesep=1.5mm]
SCHEMA
  RootPtr/1, Marked/1, Grey/1, NewG/1, HeapPtr/2, Tc/2, Tc_aux/2

PRE
  true

CMD
  Marked := RootPtr;
  Grey := RootPtr;
  WHILE (Grey ≠ ∅) DO
    NewG := ((π[1] (σ[#0 = #2] (HeapPtr × Grey))) ∖ Marked);
    Marked := (Marked ∪ NewG);
    Grey := NewG
  END;
  Tc_aux := ∅[2];
  Tc := (HeapPtr ∪ π[0,3] (σ[#1 = #2] ((Tc_aux × Tc_aux))));
  WHILE (¬((Tc = Tc_aux))) DO
    Tc_aux := Tc;
    Tc := (Tc ∪ (HeapPtr ∪ π[0,3] (σ[#1 = #2] ((Tc_aux × Tc_aux)))))
  END

POST
  (Marked = (RootPtr ∪ (π[2] (σ[#0 = #1] (RootPtr × Tc)))))
\end{Verbatim}

\Needspace{0.45\textheight}

\section{Example5038: OpenFGA Check by top-down userset dispatch versus OpenFGA ListObjects by reverse expansion from the user}\label{case:Example5038}

\begin{description}[leftmargin=2.6cm,style=nextline,itemsep=1pt]

\item[Category] production pair (two algorithms from production sources or legacy production cases); family: production algorithm pair (rate mismatch; two evaluation directions in one codebase; replaces legacy correctness case Example2032).

\item[Kind] equivalence.

\item[Contributors] Leo Zhang.

\item[Status] certified valid (Certificate/Valid.lean, written by the certificate emitter); expected valid.

\item[Tags] \hyperref[tag:access-control]{\texttt{access-control}} \hyperref[tag:evaluation-strategies]{\texttt{evaluation-strategies}} \hyperref[tag:magic-sets]{\texttt{magic-sets}} \hyperref[tag:prophecy]{\texttt{prophecy}}.

\item[Input] \path{Benchmark/Example5038/Input.lean} (canonical notation).

\item[Summary] OpenFGA's two evaluation directions for one question, membership of users in groups through nested usersets: Check, which descends from the object, and ListObjects, which ascends from the user by reverse expansion.

\item[Sides] naive compiler: OpenFGA Check, default top-down userset dispatch (checkDirect / checkDirectUsersetTuples), as the demand-driven form of the membership program restricted to the queried users \par naive compiler: OpenFGA ListObjects, reverse expansion (readTuplesAndExecute / ReadStartingWithUser) as the membership closure seeded at the queried users

\item[Sources] {\raggedright \href{https://github.com/openfga/openfga/blob/main/internal/graph/check.go}{OpenFGA openfga/openfga internal/graph/check.go LocalChecker.checkDirect, checkDirectUserTuple, checkDirectUsersetTuples (default strategy defaultUserset dispatches Check(user, h\#member) for every userset tuple of the object), ResolveCheck (maxResolutionDepth, hasCycle)} (branch main; commit 73591ef16ce5 (2026-09-08); read 2026-09-15 (the derivation of legacy Example2032, verified against the source)) \par \href{https://github.com/openfga/openfga/blob/main/internal/graph/recursive_resolver.go}{OpenFGA openfga/openfga internal/graph/recursive\_resolver.go recursiveFastPath, recursiveMatchUserUserset, breadthFirstRecursiveMatch (the planner's alternative Check strategy for recursive usersets: usersets of the user loaded once, breadth-first from the object until a level meets them; not modelled)} (branch main; commit 73591ef16ce5 (2026-09-08); read 2026-09-15) \par \href{https://github.com/openfga/openfga/blob/main/pkg/server/commands/reverseexpand/reverse_expand.go}{OpenFGA openfga/openfga pkg/server/commands/reverseexpand/reverse\_expand.go ReverseExpandQuery.execute (edges from the type system, candidate emission by trySendCandidate with candidateObjectsMap, visitedUsersetsMap, resolveNodeLimit), reverseExpandDirect, readTuplesAndExecute (ReadStartingWithUser, each found object\#relation becomes a new source)} (branch main; commit 73591ef16ce5 (2026-09-08); read 2026-09-15) \par \href{https://github.com/openfga/openfga/blob/main/pkg/server/commands/list_objects.go}{OpenFGA openfga/openfga pkg/server/commands/list\_objects.go ListObjectsQuery.evaluate (delegates to ReverseExpand.Execute; candidates with NoFurtherEvalStatus are results, RequiresFurtherEvalStatus goes to Check)} (branch main; commit 73591ef16ce5 (2026-09-08); read 2026-09-15)\par}

\item[Prophecy] accepted invariant with a prophecy clause above level zero (Core.json); no certificate built from it yet.

\item[Note] Replaces the legacy correctness case Example2032, whose bottom-up loop is the reference program of the check; both sides are OpenFGA algorithms (Check by top-down dispatch, ListObjects by reverse expansion).

\end{description}

\paragraph{Datalog program(s).}

{\raggedright\textit{prog\_demand} --- fidelity module (kernel-checked: the compiled sides of the command are this program's naive compilation).\par}

\begin{Verbatim}[fontsize=\scriptsize,breaklines=true,breakanywhere=true,frame=single,framesep=1.5mm]
Demand(x1) :- Query(x1).
Demand(y1) :- Demand(x1), SubGroup(y1, x1).
Ans(x1, y1) :- Demand(y1), QUser(x1), DirectMember(x1, y1).
Ans(x1, y1) :- Demand(y1), SubGroup(z1, y1), Ans(x1, z1).
\end{Verbatim}

{\raggedright\textit{prog\_reach} --- fidelity module (kernel-checked: the compiled sides of the command are this program's naive compilation).\par}

\begin{Verbatim}[fontsize=\scriptsize,breaklines=true,breakanywhere=true,frame=single,framesep=1.5mm]
Reach(x1, y1) :- QUser(x1), DirectMember(x1, y1).
Reach(x1, z1) :- Reach(x1, y1), SubGroup(y1, z1).
\end{Verbatim}

{\raggedright\textit{prog\_MemberR} --- fidelity module (kernel-checked: the compiled sides of the command are this program's naive compilation).\par}

\begin{Verbatim}[fontsize=\scriptsize,breaklines=true,breakanywhere=true,frame=single,framesep=1.5mm]
MemberR(x1, y1) :- DirectMember(x1, y1).
MemberR(x1, z1) :- MemberR(x1, y1), SubGroup(y1, z1).
\end{Verbatim}

{\raggedright\textit{rules as stated in the header comment of Input.lean} --- header comment.\par}

\begin{Verbatim}[fontsize=\scriptsize,breaklines=true,breakanywhere=true,frame=single,framesep=1.5mm]
Member(u, g) :- DirectMember(u, g)
Member(u, g) :- Member(u, h), SubGroup(h, g)
Demand(g) :- Query(g)
Demand(h) :- Demand(g), SubGroup(h, g)
Ans(u, g) :- Demand(g), QUser(u), DirectMember(u, g)
Ans(u, g) :- Demand(g), SubGroup(h, g), Ans(u, h)
Reach(u, g) :- QUser(u), DirectMember(u, g)
Reach(u, g) :- Reach(u, h), SubGroup(h, g)
\end{Verbatim}

\paragraph{Reading.} The underlying query is group membership through nested usersets, Member(u,g) :- DirectMember(u,g); Member(u,g) :- Member(u,h), SubGroup(h,g), the object-relationship-model membership relation used by systems like OpenFGA. Side 1 encodes OpenFGA's Check API in its default top-down strategy (checkDirect/checkDirectUsersetTuples): it first pushes demand for a queried object down through subgroups (Demand), then builds membership answers Ans back up from the direct memberships of demanded groups. Side 2 encodes OpenFGA's ListObjects API via reverse expansion (readTuplesAndExecute/ReadStartingWithUser): it grows a reachability relation Reach upward from the queried users' direct group memberships through subgroups, with no notion of a queried object. The Hoare triple claims that, restricted to the queried users and queried objects, the two APIs' answers agree --- Check's top-down, demand-restricted evaluation and ListObjects' bottom-up reverse expansion compute the same membership relation on the questions actually asked.

\paragraph{First-order reading of PRE/POST, translated mechanically from the relational-algebra assertion in Input.lean.} The relational-algebra assertion in \texttt{Input.lean} remains the authority.

\textit{PRE:} ⊤\par The precondition is simply true; no constraints are placed on DirectMember, SubGroup, QUser, or Query.

\textit{POST:} ∀x y. (Ans(x,y) ∧ Query(y) ↔ Reach(x,y) ∧ Query(y))\par Restricted to pairs whose object is in Query, the answer relation Ans produced by OpenFGA's top-down Check equals the answer relation Reach produced by OpenFGA's reverse-expansion ListObjects.

\paragraph{Whiel encoding (Hoare triple).}

\begin{Verbatim}[fontsize=\scriptsize,breaklines=true,breakanywhere=true,frame=single,framesep=1.5mm]
SCHEMA
  Query/1, QUser/1, Demand/1, Demand_aux/1, DirectMember/2, SubGroup/2, Ans/2, Ans_aux/2, Reach/2, Reach_aux/2

PRE
  true

CMD
  Demand_aux := ∅[1];
  Ans_aux := ∅[2];
  Demand := (Query ∪ π[1] (σ[#0 = #2] ((Demand_aux × SubGroup))));
  Ans := (π[1,0] (σ[(#1 = #2 ∧ #0 = #3)] ((Demand_aux × (QUser × DirectMember)))) ∪ π[3,0] (σ[(#0 = #2 ∧ #1 = #4)] ((Demand_aux × (SubGroup × Ans_aux)))));
  WHILE (¬(((Demand = Demand_aux) ∧ (Ans = Ans_aux)))) DO
    Demand_aux := Demand;
    Ans_aux := Ans;
    Demand := (Demand ∪ (Query ∪ π[1] (σ[#0 = #2] ((Demand_aux × SubGroup)))));
    Ans := (Ans ∪ (π[1,0] (σ[(#1 = #2 ∧ #0 = #3)] ((Demand_aux × (QUser × DirectMember)))) ∪ π[3,0] (σ[(#0 = #2 ∧ #1 = #4)] ((Demand_aux × (SubGroup × Ans_aux))))))
  END;
  Reach_aux := ∅[2];
  Reach := (π[0,2] (σ[#0 = #1] ((QUser × DirectMember))) ∪ π[0,3] (σ[#1 = #2] ((Reach_aux × SubGroup))));
  WHILE (¬((Reach = Reach_aux))) DO
    Reach_aux := Reach;
    Reach := (Reach ∪ (π[0,2] (σ[#0 = #1] ((QUser × DirectMember))) ∪ π[0,3] (σ[#1 = #2] ((Reach_aux × SubGroup)))))
  END

POST
  ((π[0, 1] (σ[#1 = #2] (Ans × Query))) = (π[0, 1] (σ[#1 = #2] (Reach × Query))))
\end{Verbatim}

\Needspace{0.45\textheight}

\section{Example5039: Permission exclusion applied at every object (SpiceDB/OpenFGA) is contained in the exclusion applied once to the finished closure}\label{case:Example5039}

\begin{description}[leftmargin=2.6cm,style=nextline,itemsep=1pt]

\item[Category] production pair (two algorithms from production sources or legacy production cases); family: engine semantics versus a rewriting (replaces legacy correctness case Example2036).

\item[Kind] containment.

\item[Contributors] Leo Zhang.

\item[Status] certified valid (Certificate/Valid.lean, written by the certificate emitter); expected valid.

\item[Tags] \hyperref[tag:access-control]{\texttt{access-control}} \hyperref[tag:negation]{\texttt{negation}}.

\item[Input] \path{Benchmark/Example5039/Input.lean} (canonical notation).

\item[Summary] Where the exclusion of a permission is applied: at every object, as SpiceDB and OpenFGA do, or once to the finished closure.

\item[Sides] hand-written: the permission closure with the exclusion applied at every object (legacy Example2036: SpiceDB (reader + parent-\textgreater{}view) - banned, OpenFGA but-not) \par naive compiler: the unrestricted closure reader + parent-\textgreater{}view, then the exclusion applied once to the finished closure (loop-free suffix)

\item[Sources] {\raggedright \href{https://authzed.com/docs/spicedb/concepts/schema}{SpiceDB authzed/spicedb (documentation) Schema language: permissions, union, arrow, exclusion permission view = (reader + parent-\textgreater{}view) - banned: the exclusion applies to the permission computed at each object, and the arrow reads the parent's permission} (read 2026-09-15) \par \href{https://github.com/openfga/openfga/blob/main/internal/graph/check.go}{OpenFGA openfga/openfga internal/graph/check.go checkSetOperation with the exclusion handler (base ∧ ¬subtract, evaluated per object) and checkTTU (the arrow parent-\textgreater{}view: the check is dispatched to the parent object)} (branch main; read 2026-09-15) \par \href{https://authzed.com/docs/spicedb/concepts/schema}{SpiceDB schema language: union, intersection, exclusion, arrows (documentation)} \par \href{https://github.com/openfga/openfga/blob/main/internal/graph/check.go}{openfga/openfga internal/graph/check.go (checkDirect, checkDirectUsersetTuples, checkTTU, checkSetOperation)} \par \href{https://openfga.dev/docs/modeling/building-blocks/usersets}{OpenFGA usersets (documentation)}\par}

\item[Twin] Example5040 (invalid: the two claimed equal).

\item[Obstruction note] from the metadata: anchored: View ⊆ ViewC and View ∩ Banned = ∅ hold round by round and the exclusion of side 2 is a loop-free suffix; no prophecy level expected

\item[Note] Replaces the legacy correctness case Example2036 (its loop is side 1).

\end{description}

\paragraph{Datalog program(s).}

{\raggedright\textit{prog\_ViewC} --- fidelity module (kernel-checked: the compiled sides of the command are this program's naive compilation).\par}

\begin{Verbatim}[fontsize=\scriptsize,breaklines=true,breakanywhere=true,frame=single,framesep=1.5mm]
ViewC(x1, y1) :- Reader(x1, y1).
ViewC(x1, y1) :- ParentRel(y1, z1), ViewC(x1, z1).
\end{Verbatim}

{\raggedright\textit{prog\_ViewR} --- fidelity module (kernel-checked: the compiled sides of the command are this program's naive compilation).\par}

\begin{Verbatim}[fontsize=\scriptsize,breaklines=true,breakanywhere=true,frame=single,framesep=1.5mm]
ViewR(x1, y1) :- Reader(x1, y1), Allowed(x1, y1).
ViewR(x1, y1) :- ParentRel(y1, z1), ViewR(x1, z1), Allowed(x1, y1).
\end{Verbatim}

{\raggedright\textit{rules as stated in the header comment of Input.lean} --- header comment.\par}

\begin{Verbatim}[fontsize=\scriptsize,breaklines=true,breakanywhere=true,frame=single,framesep=1.5mm]
View(u, d) :- Reader(u, d), not Banned(u, d)
View(u, d) :- ParentRel(d, p), View(u, p), not Banned(u, d)
ViewC(u, d) :- Reader(u, d)
ViewC(u, d) :- ParentRel(d, p), ViewC(u, p)
\end{Verbatim}

\paragraph{Reading.} The permission modelled is SpiceDB/OpenFGA-style exclusion over an inherited relation, permission view = (reader + parent-\textgreater{}view) - banned: a user views a document if they read it directly or view its parent, except when banned on that document. Side 1 encodes how SpiceDB and OpenFGA's checkSetOperation/checkTTU actually evaluate this, applying the exclusion at every object in the recursive closure, so a ban on a document also blocks the permission from propagating further down through its descendants. Side 2 computes the unrestricted inheritance closure ViewC by the naive compiler and then subtracts the banned pairs Banned only once, after the closure is complete, giving Filtered. The Hoare triple claims View ⊆ Filtered: applying the exclusion at every recursive step only ever removes pairs relative to the unrestricted closure and never produces a banned pair, so the per-step-filtered relation is necessarily contained in the once-filtered relation --- though, as Example5040 shows, the two are not equal.

\paragraph{First-order reading of PRE/POST, translated mechanically from the relational-algebra assertion in Input.lean.} The relational-algebra assertion in \texttt{Input.lean} remains the authority.

\textit{PRE:} ⊤\par The precondition is simply true; no constraints are placed on Reader, Banned, or ParentRel.

\textit{POST:} ∀x y. (View(x,y) → Filtered(x,y))\par View, computed with the ban exclusion applied at every object of the recursive permission definition, is contained in Filtered, computed as the unrestricted inheritance closure with the ban exclusion subtracted only once at the end.

\paragraph{Whiel encoding (Hoare triple).}

\begin{Verbatim}[fontsize=\scriptsize,breaklines=true,breakanywhere=true,frame=single,framesep=1.5mm]
SCHEMA
  Reader/2, Banned/2, ParentRel/2, View/2, View_aux/2, ViewC/2, ViewC_aux/2, Filtered/2

PRE
  true

CMD
  View_aux := ∅[2];
  View := (Reader ∖ Banned);
  WHILE (View ≠ View_aux) DO
    View_aux := View;
    View := ((Reader ∪ (π[2, 0] (σ[#1 = #3] (ParentRel × View_aux)))) ∖ Banned)
  END;
  ViewC_aux := ∅[2];
  ViewC := (Reader ∪ π[2,0] (σ[#1 = #3] ((ParentRel × ViewC_aux))));
  WHILE (¬((ViewC = ViewC_aux))) DO
    ViewC_aux := ViewC;
    ViewC := (ViewC ∪ (Reader ∪ π[2,0] (σ[#1 = #3] ((ParentRel × ViewC_aux)))))
  END;
  Filtered := (ViewC ∖ Banned)

POST
  (View ⊆ Filtered)
\end{Verbatim}

\Needspace{0.45\textheight}

\section{Example5040: Permission exclusion applied at every object versus once to the finished closure, claimed equal (invalid twin of Example5039)}\label{case:Example5040}

\begin{description}[leftmargin=2.6cm,style=nextline,itemsep=1pt]

\item[Category] production pair (two algorithms from production sources or legacy production cases); family: engine semantics versus a rewriting (invalid twin).

\item[Kind] equivalence.

\item[Contributors] Leo Zhang.

\item[Status] certified invalid (Certificate/Invalid.lean, written by the certificate emitter); expected invalid.

\item[Tags] \hyperref[tag:access-control]{\texttt{access-control}} \hyperref[tag:negation]{\texttt{negation}}.

\item[Input] \path{Benchmark/Example5040/Input.lean} (canonical notation).

\item[Summary] Where the exclusion of a permission is applied: at every object, as SpiceDB and OpenFGA do, or once to the finished closure.

\item[Sides] hand-written: the permission closure with the exclusion applied at every object (legacy Example2036: SpiceDB (reader + parent-\textgreater{}view) - banned, OpenFGA but-not) \par naive compiler: the unrestricted closure reader + parent-\textgreater{}view, then the exclusion applied once to the finished closure (loop-free suffix)

\item[Sources] {\raggedright \href{https://authzed.com/docs/spicedb/concepts/schema}{SpiceDB authzed/spicedb (documentation) Schema language: permissions, union, arrow, exclusion permission view = (reader + parent-\textgreater{}view) - banned: the exclusion applies to the permission computed at each object, and the arrow reads the parent's permission} (read 2026-09-15) \par \href{https://github.com/openfga/openfga/blob/main/internal/graph/check.go}{OpenFGA openfga/openfga internal/graph/check.go checkSetOperation with the exclusion handler (base ∧ ¬subtract, evaluated per object) and checkTTU (the arrow parent-\textgreater{}view: the check is dispatched to the parent object)} (branch main; read 2026-09-15) \par \href{https://authzed.com/docs/spicedb/concepts/schema}{SpiceDB schema language: union, intersection, exclusion, arrows (documentation)} \par \href{https://github.com/openfga/openfga/blob/main/internal/graph/check.go}{openfga/openfga internal/graph/check.go (checkDirect, checkDirectUsersetTuples, checkTTU, checkSetOperation)} \par \href{https://openfga.dev/docs/modeling/building-blocks/usersets}{OpenFGA usersets (documentation)}\par}

\item[Twin] Example5039 (valid: the containment).

\item[Obstruction note] from the metadata: invalid (the exclusion inside the recursion blocks inheritance through a banned parent; filtering the closure once does not)

\end{description}

\paragraph{Datalog program(s).}

{\raggedright\textit{prog\_ViewC} --- fidelity module (kernel-checked: the compiled sides of the command are this program's naive compilation).\par}

\begin{Verbatim}[fontsize=\scriptsize,breaklines=true,breakanywhere=true,frame=single,framesep=1.5mm]
ViewC(x1, y1) :- Reader(x1, y1).
ViewC(x1, y1) :- ParentRel(y1, z1), ViewC(x1, z1).
\end{Verbatim}

{\raggedright\textit{prog\_ViewR} --- fidelity module (kernel-checked: the compiled sides of the command are this program's naive compilation).\par}

\begin{Verbatim}[fontsize=\scriptsize,breaklines=true,breakanywhere=true,frame=single,framesep=1.5mm]
ViewR(x1, y1) :- Reader(x1, y1), Allowed(x1, y1).
ViewR(x1, y1) :- ParentRel(y1, z1), ViewR(x1, z1), Allowed(x1, y1).
\end{Verbatim}

{\raggedright\textit{rules as stated in the header comment of Input.lean} --- header comment.\par}

\begin{Verbatim}[fontsize=\scriptsize,breaklines=true,breakanywhere=true,frame=single,framesep=1.5mm]
View(u, d) :- Reader(u, d), not Banned(u, d)
View(u, d) :- ParentRel(d, p), View(u, p), not Banned(u, d)
ViewC(u, d) :- Reader(u, d)
ViewC(u, d) :- ParentRel(d, p), ViewC(u, p)
\end{Verbatim}

\paragraph{Reading.} This is the invalid twin of Example5039, with identical inputs and identical two sides (the per-object-excluded View versus the once-filtered Filtered), but the Hoare triple here claims the stronger equality View = Filtered instead of the containment. This is false: the case gives an explicit witness where a user reads a parent document but is banned on that parent, and a child document inherits from the parent. Applying the exclusion at every step blocks the permission from ever reaching the child through the banned parent, whereas filtering the unrestricted closure once at the end only removes the directly banned pair on the parent and leaves the inherited permission on the child intact, so the two relations differ --- a kernel-checked counterexample confirms this, illustrating that pushing a negation out of a recursive definition is unsound.

\paragraph{First-order reading of PRE/POST, translated mechanically from the relational-algebra assertion in Input.lean.} The relational-algebra assertion in \texttt{Input.lean} remains the authority.

\textit{PRE:} ⊤\par The precondition is simply true; no constraints are placed on Reader, Banned, or ParentRel.

\textit{POST:} ∀x y. (View(x,y) ↔ Filtered(x,y))\par It is claimed that View (exclusion applied at every recursive step) equals Filtered (exclusion applied once, after the unrestricted closure is complete) --- a claim the case shows to be false via a banned-parent witness.

\paragraph{Whiel encoding (Hoare triple).}

\begin{Verbatim}[fontsize=\scriptsize,breaklines=true,breakanywhere=true,frame=single,framesep=1.5mm]
SCHEMA
  Reader/2, Banned/2, ParentRel/2, View/2, View_aux/2, ViewC/2, ViewC_aux/2, Filtered/2

PRE
  true

CMD
  View_aux := ∅[2];
  View := (Reader ∖ Banned);
  WHILE (View ≠ View_aux) DO
    View_aux := View;
    View := ((Reader ∪ (π[2, 0] (σ[#1 = #3] (ParentRel × View_aux)))) ∖ Banned)
  END;
  ViewC_aux := ∅[2];
  ViewC := (Reader ∪ π[2,0] (σ[#1 = #3] ((ParentRel × ViewC_aux))));
  WHILE (¬((ViewC = ViewC_aux))) DO
    ViewC_aux := ViewC;
    ViewC := (ViewC ∪ (Reader ∪ π[2,0] (σ[#1 = #3] ((ParentRel × ViewC_aux)))))
  END;
  Filtered := (ViewC ∖ Banned)

POST
  (View = Filtered)
\end{Verbatim}

\Needspace{0.45\textheight}

\section{Example5041: PostgreSQL foreign-key cascade as a trigger worklist with the surviving-reference RESTRICT check versus the declarative delete closure}\label{case:Example5041}

\begin{description}[leftmargin=2.6cm,style=nextline,itemsep=1pt]

\item[Category] production pair (two algorithms from production sources or legacy production cases); family: production algorithm pair (replaces legacy correctness case Example2024).

\item[Kind] equivalence.

\item[Contributors] Leo Zhang.

\item[Status] certified valid (Certificate/Valid.lean, written by the certificate emitter); expected valid.

\item[Tags] \hyperref[tag:database-engine]{\texttt{database-engine}} \hyperref[tag:evaluation-strategies]{\texttt{evaluation-strategies}} \hyperref[tag:integrity-constraints]{\texttt{integrity-constraints}} \hyperref[tag:negation]{\texttt{negation}} \hyperref[tag:prophecy]{\texttt{prophecy}} \hyperref[tag:semi-naive]{\texttt{semi-naive}}.

\item[Input] \path{Benchmark/Example5041/Input.lean} (canonical notation).

\item[Summary] PostgreSQL's foreign-key cascade, as the trigger queue runs it, against the declarative delete closure, each followed by the RESTRICT check.

\item[Sides] hand-written: PostgreSQL RI\_FKey\_cascade\_del as a worklist of newly deleted rows, then ri\_restrict as the surviving-reference check \par naive compiler: the delete closure, then the same check as loop-free code

\item[Sources] {\raggedright \href{https://github.com/postgres/postgres/blob/master/src/backend/utils/adt/ri_triggers.c}{PostgreSQL postgres/postgres src/backend/utils/adt/ri\_triggers.c RI\_FKey\_cascade\_del (DELETE FROM ONLY fktable WHERE \$1 = fkatt1, queued as an AFTER trigger per deleted row and re-fired transitively), ri\_restrict via RI\_FKey\_restrict\_del / RI\_FKey\_noaction\_del (SELECT 1 FROM ONLY fktable x WHERE \$1 = fkatt1 FOR KEY SHARE OF x: a violation only if a referencing row still exists when the check runs)} (branch master; read 2026-09-15) \par \href{https://www.postgresql.org/docs/current/ddl-constraints.html}{PostgreSQL documentation postgresql.org Chapter 5, Constraints, Foreign Keys (ON DELETE CASCADE, RESTRICT, NO ACTION)} (read 2026-09-15) \par \href{https://github.com/postgres/postgres/blob/master/src/backend/utils/adt/ri_triggers.c}{postgres/postgres src/backend/utils/adt/ri\_triggers.c (RI\_FKey\_cascade\_del, ri\_restrict)}\par}

\item[Twin] Example5042 (invalid: the over-approximating RESTRICT check of legacy Example2024).

\item[Prophecy] kernel-checked certificate with a level-1 prophecy clause (Certificate/Valid.lean, Core.json).

\item[Note] Replaces the legacy correctness case Example2024; its over-approximating RESTRICT check is the invalid twin Example5042.

\end{description}

\paragraph{Datalog program(s).}

{\raggedright\textit{prog\_DelC} --- fidelity module (kernel-checked: the compiled sides of the command are this program's naive compilation).\par}

\begin{Verbatim}[fontsize=\scriptsize,breaklines=true,breakanywhere=true,frame=single,framesep=1.5mm]
DelC(x1) :- SeedDel(x1).
DelC(x1) :- RefBy(x1, y1), DelC(y1).
\end{Verbatim}

{\raggedright\textit{prog\_DelR} --- fidelity module (kernel-checked: the compiled sides of the command are this program's naive compilation).\par}

\begin{Verbatim}[fontsize=\scriptsize,breaklines=true,breakanywhere=true,frame=single,framesep=1.5mm]
DelR(x1) :- SeedDel(x1).
DelR(x1) :- RefBy(x1, y1), DelR(y1).
\end{Verbatim}

{\raggedright\textit{rules as stated in the header comment of Input.lean} --- header comment.\par}

\begin{Verbatim}[fontsize=\scriptsize,breaklines=true,breakanywhere=true,frame=single,framesep=1.5mm]
Del(k) :- SeedDel(k)
Del(r) :- RefBy(r, k), Del(k)
\end{Verbatim}

\paragraph{Reading.} The Datalog program is a delete-reachability closure, Del(k) :- SeedDel(k); Del(r) :- RefBy(r,k), Del(k), modelling ON DELETE CASCADE foreign keys, together with a RESTRICT check for rows referenced through non-cascading (RESTRICT/NO ACTION) foreign keys. Side 1 encodes PostgreSQL's actual trigger-based implementation (ri\_triggers.c): RI\_FKey\_cascade\_del deletes rows via a worklist that re-fires for every newly deleted row (Del/Front/New), and ri\_restrict then reports a violation only for rows that still exist and reference a deleted key, i.e. rows not themselves cascaded away in the same statement. Side 2 computes the same delete closure DelC by the naive compiler and applies the identical surviving-reference RESTRICT check to it (ViolC). The Hoare triple claims Del = DelC and Viol = ViolC: PostgreSQL's row-level trigger worklist computes exactly the declarative delete closure, and, since both sides use the same surviving-reference check on equal relations, the reported RESTRICT violations agree too.

\paragraph{First-order reading of PRE/POST, translated mechanically from the relational-algebra assertion in Input.lean.} The relational-algebra assertion in \texttt{Input.lean} remains the authority.

\textit{PRE:} ⊤\par The precondition is simply true; no constraints are placed on SeedDel, RefBy, or RestrictRef.

\textit{POST:} (∀x. (Del(x) ↔ DelC(x))) ∧ ∀y. (Viol(y) ↔ ViolC(y))\par Del, the set of keys deleted by PostgreSQL's cascading trigger worklist, equals DelC, the same set computed as a declarative closure, and Viol, the RESTRICT violations reported by the surviving-reference check on Del, equals ViolC, the same check applied to DelC.

\paragraph{Whiel encoding (Hoare triple).}

\begin{Verbatim}[fontsize=\scriptsize,breaklines=true,breakanywhere=true,frame=single,framesep=1.5mm]
SCHEMA
  SeedDel/1, Del/1, Front/1, New/1, Viol/1, DelC/1, DelC_aux/1, ViolC/1, RefBy/2, RestrictRef/2

PRE
  true

CMD
  Del := SeedDel;
  Front := SeedDel;
  WHILE (Front ≠ ∅) DO
    New := ((π[0] (σ[#1 = #2] (RefBy × Front))) ∖ Del);
    Del := (Del ∪ New);
    Front := New
  END;
  Viol := ((π[0] (σ[#1 = #2] (RestrictRef × Del))) ∖ Del);
  DelC_aux := ∅[1];
  DelC := (SeedDel ∪ π[0] (σ[#1 = #2] ((RefBy × DelC_aux))));
  WHILE (¬((DelC = DelC_aux))) DO
    DelC_aux := DelC;
    DelC := (DelC ∪ (SeedDel ∪ π[0] (σ[#1 = #2] ((RefBy × DelC_aux)))))
  END;
  ViolC := ((π[0] (σ[#1 = #2] (RestrictRef × DelC))) ∖ DelC)

POST
  ((Del = DelC) ∧ (Viol = ViolC))
\end{Verbatim}

\Needspace{0.45\textheight}

\section{Example5042: The over-approximating RESTRICT check of legacy Example2024 versus PostgreSQL's surviving-reference check, claimed equal (invalid twin of Example5041)}\label{case:Example5042}

\begin{description}[leftmargin=2.6cm,style=nextline,itemsep=1pt]

\item[Category] production pair (two algorithms from production sources or legacy production cases); family: production discrepancy pair (invalid twin).

\item[Kind] equivalence.

\item[Contributors] Leo Zhang.

\item[Status] certified invalid (Certificate/Invalid.lean, written by the certificate emitter); expected invalid.

\item[Tags] \hyperref[tag:database-engine]{\texttt{database-engine}} \hyperref[tag:evaluation-strategies]{\texttt{evaluation-strategies}} \hyperref[tag:integrity-constraints]{\texttt{integrity-constraints}} \hyperref[tag:negation]{\texttt{negation}} \hyperref[tag:semi-naive]{\texttt{semi-naive}}.

\item[Input] \path{Benchmark/Example5042/Input.lean} (canonical notation).

\item[Summary] PostgreSQL's foreign-key cascade, as the trigger queue runs it, against the declarative delete closure, each followed by the RESTRICT check.

\item[Sides] hand-written: the cascade worklist of Example5041 followed by the check of legacy Example2024, every row referencing a deleted key under RESTRICT (no ∖ Del) \par naive compiler: the delete closure, then PostgreSQL's surviving-reference check

\item[Sources] {\raggedright \href{https://github.com/postgres/postgres/blob/master/src/backend/utils/adt/ri_triggers.c}{PostgreSQL postgres/postgres src/backend/utils/adt/ri\_triggers.c RI\_FKey\_cascade\_del (DELETE FROM ONLY fktable WHERE \$1 = fkatt1, queued as an AFTER trigger per deleted row and re-fired transitively), ri\_restrict via RI\_FKey\_restrict\_del / RI\_FKey\_noaction\_del (SELECT 1 FROM ONLY fktable x WHERE \$1 = fkatt1 FOR KEY SHARE OF x: a violation only if a referencing row still exists when the check runs)} (branch master; read 2026-09-15) \par \href{https://www.postgresql.org/docs/current/ddl-constraints.html}{PostgreSQL documentation postgresql.org Chapter 5, Constraints, Foreign Keys (ON DELETE CASCADE, RESTRICT, NO ACTION)} (read 2026-09-15) \par \href{https://github.com/postgres/postgres/blob/master/src/backend/utils/adt/ri_triggers.c}{postgres/postgres src/backend/utils/adt/ri\_triggers.c (RI\_FKey\_cascade\_del, ri\_restrict)}\par}

\item[Twin] Example5041 (valid: both sides use the surviving-reference check).

\item[Obstruction note] from the metadata: invalid (the legacy check counts a row that the same statement cascades away; PostgreSQL's ri\_restrict does not)

\end{description}

\paragraph{Datalog program(s).}

{\raggedright\textit{prog\_DelC} --- fidelity module (kernel-checked: the compiled sides of the command are this program's naive compilation).\par}

\begin{Verbatim}[fontsize=\scriptsize,breaklines=true,breakanywhere=true,frame=single,framesep=1.5mm]
DelC(x1) :- SeedDel(x1).
DelC(x1) :- RefBy(x1, y1), DelC(y1).
\end{Verbatim}

{\raggedright\textit{prog\_DelR} --- fidelity module (kernel-checked: the compiled sides of the command are this program's naive compilation).\par}

\begin{Verbatim}[fontsize=\scriptsize,breaklines=true,breakanywhere=true,frame=single,framesep=1.5mm]
DelR(x1) :- SeedDel(x1).
DelR(x1) :- RefBy(x1, y1), DelR(y1).
\end{Verbatim}

{\raggedright\textit{rules as stated in the header comment of Input.lean} --- header comment.\par}

\begin{Verbatim}[fontsize=\scriptsize,breaklines=true,breakanywhere=true,frame=single,framesep=1.5mm]
Del(k) :- SeedDel(k)
Del(r) :- RefBy(r, k), Del(k)
\end{Verbatim}

\paragraph{Reading.} This is the invalid twin of Example5041: both sides compute the same cascade-delete closure (Del via the trigger worklist, DelC via the declarative closure), but side 1's RESTRICT check is changed to the over-approximating check inherited from the retired legacy Example2024, Viol := the rows referencing any deleted key, without excluding the rows that the same delete statement already cascaded away, whereas side 2 keeps PostgreSQL's faithful surviving-reference check (which subtracts Del). The Hoare triple still claims Del = DelC and Viol = ViolC, but this is now false for the Viol conjunct: the case gives an explicit witness where a row is deleted under CASCADE via one key and simultaneously referenced under RESTRICT via another deleted key, so PostgreSQL's real check (which only flags rows still present when the check runs) reports no violation, while the legacy over-approximating check flags the already-deleted row as a violation, refuted by a kernel-checked counterexample.

\paragraph{First-order reading of PRE/POST, translated mechanically from the relational-algebra assertion in Input.lean.} The relational-algebra assertion in \texttt{Input.lean} remains the authority.

\textit{PRE:} ⊤\par The precondition is simply true; no constraints are placed on SeedDel, RefBy, or RestrictRef.

\textit{POST:} (∀x. (Del(x) ↔ DelC(x))) ∧ ∀y. (Viol(y) ↔ ViolC(y))\par It is claimed that Del equals DelC (true, as in Example5041) and that Viol, computed here without excluding rows already cascaded away, equals ViolC, computed with PostgreSQL's surviving-reference exclusion --- the Viol equality is false, refuted by a cascade-and-restrict witness.

\paragraph{Whiel encoding (Hoare triple).}

\begin{Verbatim}[fontsize=\scriptsize,breaklines=true,breakanywhere=true,frame=single,framesep=1.5mm]
SCHEMA
  SeedDel/1, Del/1, Front/1, New/1, Viol/1, DelC/1, DelC_aux/1, ViolC/1, RefBy/2, RestrictRef/2

PRE
  true

CMD
  Del := SeedDel;
  Front := SeedDel;
  WHILE (Front ≠ ∅) DO
    New := ((π[0] (σ[#1 = #2] (RefBy × Front))) ∖ Del);
    Del := (Del ∪ New);
    Front := New
  END;
  Viol := (π[0] (σ[#1 = #2] (RestrictRef × Del)));
  DelC_aux := ∅[1];
  DelC := (SeedDel ∪ π[0] (σ[#1 = #2] ((RefBy × DelC_aux))));
  WHILE (¬((DelC = DelC_aux))) DO
    DelC_aux := DelC;
    DelC := (DelC ∪ (SeedDel ∪ π[0] (σ[#1 = #2] ((RefBy × DelC_aux)))))
  END;
  ViolC := ((π[0] (σ[#1 = #2] (RestrictRef × DelC))) ∖ DelC)

POST
  ((Del = DelC) ∧ (Viol = ViolC))
\end{Verbatim}
\end{document}
